% Matemáticas I para Ingeniería Industrial - Apuntes y Ejercicios
% Basado en la plantilla Apuntes_Estad_Modelos (classicthesis).

\RequirePackage{silence}
\WarningFilter{scrreprt}{Usage of package `titlesec'}
\WarningFilter{titlesec}{Non standard sectioning command detected}

\documentclass[
  twoside,openright,titlepage,numbers=noenddot,
  headinclude,footinclude,cleardoublepage=empty,abstract=on,
  BCOR=5mm,paper=a4,fontsize=11pt
]{scrreprt}

% Compatibilidad TeX Live 2025/2026: babel-spanish intenta usar el hook
% genérico de \@startsection, mientras que KOMA-Script redefine internamente
% el mecanismo de seccionado. Declarar el hook antes de cargar babel evita
% el error "Generic hooks cannot be added to '\@startsection'".
\NewHook{cmd/@startsection/before}

% ****************************************************************************************************
% classicthesis-config.tex
% formerly known as loadpackages.sty, classicthesis-ldpkg.sty, and classicthesis-preamble.sty
% Use it at the beginning of your ClassicThesis.tex, or as a LaTeX Preamble
% in your ClassicThesis.{tex,lyx} with % ****************************************************************************************************
% classicthesis-config.tex
% formerly known as loadpackages.sty, classicthesis-ldpkg.sty, and classicthesis-preamble.sty
% Use it at the beginning of your ClassicThesis.tex, or as a LaTeX Preamble
% in your ClassicThesis.{tex,lyx} with % ****************************************************************************************************
% classicthesis-config.tex
% formerly known as loadpackages.sty, classicthesis-ldpkg.sty, and classicthesis-preamble.sty
% Use it at the beginning of your ClassicThesis.tex, or as a LaTeX Preamble
% in your ClassicThesis.{tex,lyx} with \input{classicthesis-config}
% ****************************************************************************************************
% If you like the classicthesis, then I would appreciate a postcard.
% My address can be found in the file ClassicThesis.pdf. A collection
% of the postcards I received so far is available online at
% http://postcards.miede.de
% ****************************************************************************************************
\usepackage{float}
\usepackage{multirow}
% ****************************************************************************************************
% 0. Set the encoding of your files. UTF-8 is the only sensible encoding nowadays. If you can't read
% äöüßáéçèê∂åëæƒÏ€ then change the encoding setting in your editor, not the line below. If your editor
% does not support utf8 use another editor!
% ****************************************************************************************************
\PassOptionsToPackage{utf8}{inputenc}
  \usepackage{inputenc}

\PassOptionsToPackage{T1}{fontenc} % T2A for cyrillics
  \usepackage{fontenc}

% ****************************************************************************************************
% 1. Configure classicthesis for your needs here, e.g., remove "drafting" below
% in order to deactivate the time-stamp on the pages
% (see ClassicThesis.pdf for more information):
% ****************************************************************************************************
\PassOptionsToPackage{
  drafting=true,    % print version information on the bottom of the pages
  tocaligned=false, % the left column of the toc will be aligned (no indentation)
  dottedtoc=false,  % page numbers in ToC flushed right
  eulerchapternumbers=true, % use AMS Euler for chapter font (otherwise Palatino)
  linedheaders=false,       % chaper headers will have line above and beneath
  floatperchapter=true,     % numbering per chapter for all floats (i.e., Figure 1.1)
  eulermath=false,  % use awesome Euler fonts for mathematical formulae (only with pdfLaTeX)
  beramono=true,    % toggle a nice monospaced font (w/ bold)
  palatino=true,    % deactivate standard font for loading another one, see the last section at the end of this file for suggestions
  style=classicthesis % classicthesis, arsclassica
}{classicthesis}


% ****************************************************************************************************
% 2. Personal data and user ad-hoc commands (insert your own data here)
% ****************************************************************************************************
\newcommand{\myTitle}{Matemáticas I\xspace}
\newcommand{\myCourse}{Ingenierías de la rama industrial}
\newcommand{\mySubtitle}{para Ingeniería Industrial\xspace}
%\newcommand{\myDegree}{Doktor-Ingenieur (Dr.-Ing.)\xspace}
\newcommand{\myName}{Diego García Zamora y Nuria V. Castillo González\xspace}
\newcommand{\myDepartment}{Departamento de Matemáticas\xspace}
\newcommand{\myUni}{Universidad de Jaén\xspace}
\newcommand{\myLocation}{Jaén\xspace}
\newcommand{\myTime}{Septiembre de 2026\xspace}
\newcommand{\myYear}{2026--2027\xspace}
\newcommand{\myPosition}{Profesor Ayudante Doctor\xspace}



%\newcommand{\myVersion}{\classicthesis}

% ********************************************************************
% Setup, finetuning, and useful commands
% ********************************************************************
\providecommand{\mLyX}{L\kern-.1667em\lower.25em\hbox{Y}\kern-.125emX\@}
\newcommand{\ie}{i.\,e.}
\newcommand{\Ie}{I.\,e.}
\newcommand{\eg}{e.\,g.}
\newcommand{\Eg}{E.\,g.}
% ****************************************************************************************************


% ****************************************************************************************************
% 3. Loading some handy packages
% ****************************************************************************************************

%********************************************************************
\PassOptionsToPackage{american}{babel} % change this to your language(s), main language last
% Spanish languages need extra options in order to work with this template
\PassOptionsToPackage{spanish,es-lcroman}{babel}
\usepackage{babel}
\addto\captionsamerican{\renewcommand{\contentsname}{Índice General}}
\addto\captionsamerican{\renewcommand{\listfigurename}{Lista de figuras}}
\addto\captionsamerican{\renewcommand{\figurename}{Fig.}}
\addto\captionsamerican{\renewcommand{\listtablename}{Lista de tablas}}
\addto\captionsamerican{\renewcommand{\tablename}{Tabla}}
\addto\captionsamerican{\renewcommand{\partname}{Parte}}
%\addto\captionsamerican{\renewcommand{\bibname}{Referencias}}
%\renewcommand{\bibname}{Referencias}
%\addto\captionsspanish{  \renewcommand{\contentsname}{Whatever}}

\usepackage{csquotes}
\PassOptionsToPackage{%
  %backend=biber,bibencoding=utf8, %instead of bibtex
  backend=bibtex8,bibencoding=ascii,%
  language=spanish,%
  %style=numeric-comp,%
  style=authoryear-comp, % Author 1999, 2010
  %bibstyle=authoryear,dashed=false, % dashed: substitute rep. author with ---
  sorting=nyt, % name, year, title
  maxbibnames=10, % default: 3, et al.
  %backref=true,%
  natbib=true % natbib compatibility mode (\citep and \citet still work)
}{biblatex}
    \usepackage{biblatex}

\PassOptionsToPackage{fleqn}{amsmath}       % math environments and more by the AMS
  \usepackage{amsmath}

% ********************************************************************
% General useful packages
% ********************************************************************
\usepackage{graphicx} %
\usepackage{scrhack} % fix warnings when using KOMA with listings package
\usepackage{xspace} % to get the spacing after macros right
\PassOptionsToPackage{printonlyused,smaller}{acronym}
  \usepackage{acronym} % nice macros for handling all acronyms in the thesis
  %\renewcommand{\bflabel}[1]{{#1}\hfill} % fix the list of acronyms --> no longer working
  %\renewcommand*{\acsfont}[1]{\textsc{#1}}
  %\renewcommand*{\aclabelfont}[1]{\acsfont{#1}}
  %\def\bflabel#1{{#1\hfill}}
  \def\bflabel#1{{\acsfont{#1}\hfill}}
  \def\aclabelfont#1{\acsfont{#1}}
% ****************************************************************************************************
%\usepackage{pgfplots} % External TikZ/PGF support (thanks to Andreas Nautsch)
%\usetikzlibrary{external}
%\tikzexternalize[mode=list and make, prefix=ext-tikz/]
% ****************************************************************************************************


% ****************************************************************************************************
% 4. Setup floats: tables, (sub)figures, and captions
% ****************************************************************************************************
\usepackage{tabularx} % better tables
  \setlength{\extrarowheight}{3pt} % increase table row height
\newcommand{\tableheadline}[1]{\multicolumn{1}{l}{\spacedlowsmallcaps{#1}}}
\newcommand{\myfloatalign}{\centering} % to be used with each float for alignment
\usepackage{subfig}
% ****************************************************************************************************


% ****************************************************************************************************
% 5. Setup code listings
% ****************************************************************************************************
\usepackage{listings}
%\lstset{emph={trueIndex,root},emphstyle=\color{BlueViolet}}%\underbar} % for special keywords
\lstset{language=[LaTeX]Tex,%C++,
  morekeywords={PassOptionsToPackage,selectlanguage},
  keywordstyle=\color{RoyalBlue},%\bfseries,
  basicstyle=\small\ttfamily,
  %identifierstyle=\color{NavyBlue},
  commentstyle=\color{Green}\ttfamily,
  stringstyle=\rmfamily,
  numbers=none,%left,%
  numberstyle=\scriptsize,%\tiny
  stepnumber=5,
  numbersep=8pt,
  showstringspaces=false,
  breaklines=true,
  %frameround=ftff,
  %frame=single,
  belowcaptionskip=.75\baselineskip
  %frame=L
}
% ****************************************************************************************************




% ****************************************************************************************************
% 6. Last calls before the bar closes
% ****************************************************************************************************
% ********************************************************************
% Her Majesty herself
% ********************************************************************
\usepackage{classicthesis}


% ********************************************************************
% Fine-tune hyperreferences (hyperref should be called last)
% ********************************************************************
\hypersetup{%
  %draft, % hyperref's draft mode, for printing see below
  colorlinks=true, linktocpage=true, pdfstartpage=3, pdfstartview=FitV,%
  % uncomment the following line if you want to have black links (e.g., for printing)
  %colorlinks=false, linktocpage=false, pdfstartpage=3, pdfstartview=FitV, pdfborder={0 0 0},%
  breaklinks=true, pageanchor=true,%
  pdfpagemode=UseNone, %
  % pdfpagemode=UseOutlines,%
  plainpages=false, bookmarksnumbered, bookmarksopen=true, bookmarksopenlevel=1,%
  hypertexnames=true, pdfhighlight=/O,%nesting=true,%frenchlinks,%
  urlcolor=CTurl, linkcolor=CTlink, citecolor=CTcitation, %pagecolor=RoyalBlue,%
  %urlcolor=Black, linkcolor=Black, citecolor=Black, %pagecolor=Black,%
  pdftitle={Matemáticas I para Ingeniería Industrial - Apuntes y Ejercicios},%
  pdfauthor={\myName, \myUni, \myDepartment	},%\textcopyright
  pdfsubject={},%
  pdfkeywords={},%
  pdfcreator={pdfLaTeX},%
  pdfproducer={LaTeX with hyperref and classicthesis}%
}


% ********************************************************************
% Setup autoreferences (hyperref and babel)
% ********************************************************************
% There are some issues regarding autorefnames
% http://www.tex.ac.uk/cgi-bin/texfaq2html?label=latexwords
% you have to redefine the macros for the
% language you use, e.g., american, ngerman
% (as chosen when loading babel/AtBeginDocument)
% ********************************************************************
\makeatletter
\@ifpackageloaded{babel}%
  {%
    \addto\extrasamerican{%
      \renewcommand*{\figureautorefname}{Figure}%
      \renewcommand*{\tableautorefname}{Table}%
      \renewcommand*{\partautorefname}{Part}%
      \renewcommand*{\chapterautorefname}{Chapter}%
      \renewcommand*{\sectionautorefname}{Section}%
      \renewcommand*{\subsectionautorefname}{Section}%
      \renewcommand*{\subsubsectionautorefname}{Section}%
    }%
    \addto\extrasngerman{%
      \renewcommand*{\paragraphautorefname}{Absatz}%
      \renewcommand*{\subparagraphautorefname}{Unterabsatz}%
      \renewcommand*{\footnoteautorefname}{Fu\"snote}%
      \renewcommand*{\FancyVerbLineautorefname}{Zeile}%
      \renewcommand*{\theoremautorefname}{Theorem}%
      \renewcommand*{\appendixautorefname}{Anhang}%
      \renewcommand*{\equationautorefname}{Gleichung}%
      \renewcommand*{\itemautorefname}{Punkt}%
    }%
      % Fix to getting autorefs for subfigures right (thanks to Belinda Vogt for changing the definition)
      \providecommand{\subfigureautorefname}{\figureautorefname}%
    }{\relax}
\makeatother


% ********************************************************************
% Development Stuff
% ********************************************************************
\listfiles
%\PassOptionsToPackage{l2tabu,orthodox,abort}{nag}
%  \usepackage{nag}
%\PassOptionsToPackage{warning, all}{onlyamsmath}
%  \usepackage{onlyamsmath}


% ****************************************************************************************************
% 7. Further adjustments (experimental)
% ****************************************************************************************************
% ********************************************************************
% Changing the text area
% ********************************************************************
%\areaset[current]{312pt}{761pt} % 686 (factor 2.2) + 33 head + 42 head \the\footskip
%\setlength{\marginparwidth}{7em}%
%\setlength{\marginparsep}{2em}%

% ********************************************************************
% Using different fonts
% ********************************************************************
%\usepackage[oldstylenums]{kpfonts} % oldstyle notextcomp
% \usepackage[osf]{libertine}
%\usepackage[light,condensed,math]{iwona}
%\renewcommand{\sfdefault}{iwona}
%\usepackage{lmodern} % <-- no osf support :-(
%\usepackage{cfr-lm} %
%\usepackage[urw-garamond]{mathdesign} <-- no osf support :-(
%\usepackage[default,osfigures]{opensans} % scale=0.95
%\usepackage[sfdefault]{FiraSans}
% \usepackage[opticals,mathlf]{MinionPro} % onlytext
% ********************************************************************
%\usepackage[largesc,osf]{newpxtext}
%\linespread{1.05} % a bit more for Palatino
% Used to fix these:
% https://bitbucket.org/amiede/classicthesis/issues/139/italics-in-pallatino-capitals-chapter
% https://bitbucket.org/amiede/classicthesis/issues/45/problema-testatine-su-classicthesis-style
% ********************************************************************
% ****************************************************************************************************

% ****************************************************************************************************
% If you like the classicthesis, then I would appreciate a postcard.
% My address can be found in the file ClassicThesis.pdf. A collection
% of the postcards I received so far is available online at
% http://postcards.miede.de
% ****************************************************************************************************
\usepackage{float}
\usepackage{multirow}
% ****************************************************************************************************
% 0. Set the encoding of your files. UTF-8 is the only sensible encoding nowadays. If you can't read
% äöüßáéçèê∂åëæƒÏ€ then change the encoding setting in your editor, not the line below. If your editor
% does not support utf8 use another editor!
% ****************************************************************************************************
\PassOptionsToPackage{utf8}{inputenc}
  \usepackage{inputenc}

\PassOptionsToPackage{T1}{fontenc} % T2A for cyrillics
  \usepackage{fontenc}

% ****************************************************************************************************
% 1. Configure classicthesis for your needs here, e.g., remove "drafting" below
% in order to deactivate the time-stamp on the pages
% (see ClassicThesis.pdf for more information):
% ****************************************************************************************************
\PassOptionsToPackage{
  drafting=true,    % print version information on the bottom of the pages
  tocaligned=false, % the left column of the toc will be aligned (no indentation)
  dottedtoc=false,  % page numbers in ToC flushed right
  eulerchapternumbers=true, % use AMS Euler for chapter font (otherwise Palatino)
  linedheaders=false,       % chaper headers will have line above and beneath
  floatperchapter=true,     % numbering per chapter for all floats (i.e., Figure 1.1)
  eulermath=false,  % use awesome Euler fonts for mathematical formulae (only with pdfLaTeX)
  beramono=true,    % toggle a nice monospaced font (w/ bold)
  palatino=true,    % deactivate standard font for loading another one, see the last section at the end of this file for suggestions
  style=classicthesis % classicthesis, arsclassica
}{classicthesis}


% ****************************************************************************************************
% 2. Personal data and user ad-hoc commands (insert your own data here)
% ****************************************************************************************************
\newcommand{\myTitle}{Matemáticas I\xspace}
\newcommand{\myCourse}{Ingenierías de la rama industrial}
\newcommand{\mySubtitle}{para Ingeniería Industrial\xspace}
%\newcommand{\myDegree}{Doktor-Ingenieur (Dr.-Ing.)\xspace}
\newcommand{\myName}{Diego García Zamora y Nuria V. Castillo González\xspace}
\newcommand{\myDepartment}{Departamento de Matemáticas\xspace}
\newcommand{\myUni}{Universidad de Jaén\xspace}
\newcommand{\myLocation}{Jaén\xspace}
\newcommand{\myTime}{Septiembre de 2026\xspace}
\newcommand{\myYear}{2026--2027\xspace}
\newcommand{\myPosition}{Profesor Ayudante Doctor\xspace}



%\newcommand{\myVersion}{\classicthesis}

% ********************************************************************
% Setup, finetuning, and useful commands
% ********************************************************************
\providecommand{\mLyX}{L\kern-.1667em\lower.25em\hbox{Y}\kern-.125emX\@}
\newcommand{\ie}{i.\,e.}
\newcommand{\Ie}{I.\,e.}
\newcommand{\eg}{e.\,g.}
\newcommand{\Eg}{E.\,g.}
% ****************************************************************************************************


% ****************************************************************************************************
% 3. Loading some handy packages
% ****************************************************************************************************

%********************************************************************
\PassOptionsToPackage{american}{babel} % change this to your language(s), main language last
% Spanish languages need extra options in order to work with this template
\PassOptionsToPackage{spanish,es-lcroman}{babel}
\usepackage{babel}
\addto\captionsamerican{\renewcommand{\contentsname}{Índice General}}
\addto\captionsamerican{\renewcommand{\listfigurename}{Lista de figuras}}
\addto\captionsamerican{\renewcommand{\figurename}{Fig.}}
\addto\captionsamerican{\renewcommand{\listtablename}{Lista de tablas}}
\addto\captionsamerican{\renewcommand{\tablename}{Tabla}}
\addto\captionsamerican{\renewcommand{\partname}{Parte}}
%\addto\captionsamerican{\renewcommand{\bibname}{Referencias}}
%\renewcommand{\bibname}{Referencias}
%\addto\captionsspanish{  \renewcommand{\contentsname}{Whatever}}

\usepackage{csquotes}
\PassOptionsToPackage{%
  %backend=biber,bibencoding=utf8, %instead of bibtex
  backend=bibtex8,bibencoding=ascii,%
  language=spanish,%
  %style=numeric-comp,%
  style=authoryear-comp, % Author 1999, 2010
  %bibstyle=authoryear,dashed=false, % dashed: substitute rep. author with ---
  sorting=nyt, % name, year, title
  maxbibnames=10, % default: 3, et al.
  %backref=true,%
  natbib=true % natbib compatibility mode (\citep and \citet still work)
}{biblatex}
    \usepackage{biblatex}

\PassOptionsToPackage{fleqn}{amsmath}       % math environments and more by the AMS
  \usepackage{amsmath}

% ********************************************************************
% General useful packages
% ********************************************************************
\usepackage{graphicx} %
\usepackage{scrhack} % fix warnings when using KOMA with listings package
\usepackage{xspace} % to get the spacing after macros right
\PassOptionsToPackage{printonlyused,smaller}{acronym}
  \usepackage{acronym} % nice macros for handling all acronyms in the thesis
  %\renewcommand{\bflabel}[1]{{#1}\hfill} % fix the list of acronyms --> no longer working
  %\renewcommand*{\acsfont}[1]{\textsc{#1}}
  %\renewcommand*{\aclabelfont}[1]{\acsfont{#1}}
  %\def\bflabel#1{{#1\hfill}}
  \def\bflabel#1{{\acsfont{#1}\hfill}}
  \def\aclabelfont#1{\acsfont{#1}}
% ****************************************************************************************************
%\usepackage{pgfplots} % External TikZ/PGF support (thanks to Andreas Nautsch)
%\usetikzlibrary{external}
%\tikzexternalize[mode=list and make, prefix=ext-tikz/]
% ****************************************************************************************************


% ****************************************************************************************************
% 4. Setup floats: tables, (sub)figures, and captions
% ****************************************************************************************************
\usepackage{tabularx} % better tables
  \setlength{\extrarowheight}{3pt} % increase table row height
\newcommand{\tableheadline}[1]{\multicolumn{1}{l}{\spacedlowsmallcaps{#1}}}
\newcommand{\myfloatalign}{\centering} % to be used with each float for alignment
\usepackage{subfig}
% ****************************************************************************************************


% ****************************************************************************************************
% 5. Setup code listings
% ****************************************************************************************************
\usepackage{listings}
%\lstset{emph={trueIndex,root},emphstyle=\color{BlueViolet}}%\underbar} % for special keywords
\lstset{language=[LaTeX]Tex,%C++,
  morekeywords={PassOptionsToPackage,selectlanguage},
  keywordstyle=\color{RoyalBlue},%\bfseries,
  basicstyle=\small\ttfamily,
  %identifierstyle=\color{NavyBlue},
  commentstyle=\color{Green}\ttfamily,
  stringstyle=\rmfamily,
  numbers=none,%left,%
  numberstyle=\scriptsize,%\tiny
  stepnumber=5,
  numbersep=8pt,
  showstringspaces=false,
  breaklines=true,
  %frameround=ftff,
  %frame=single,
  belowcaptionskip=.75\baselineskip
  %frame=L
}
% ****************************************************************************************************




% ****************************************************************************************************
% 6. Last calls before the bar closes
% ****************************************************************************************************
% ********************************************************************
% Her Majesty herself
% ********************************************************************
\usepackage{classicthesis}


% ********************************************************************
% Fine-tune hyperreferences (hyperref should be called last)
% ********************************************************************
\hypersetup{%
  %draft, % hyperref's draft mode, for printing see below
  colorlinks=true, linktocpage=true, pdfstartpage=3, pdfstartview=FitV,%
  % uncomment the following line if you want to have black links (e.g., for printing)
  %colorlinks=false, linktocpage=false, pdfstartpage=3, pdfstartview=FitV, pdfborder={0 0 0},%
  breaklinks=true, pageanchor=true,%
  pdfpagemode=UseNone, %
  % pdfpagemode=UseOutlines,%
  plainpages=false, bookmarksnumbered, bookmarksopen=true, bookmarksopenlevel=1,%
  hypertexnames=true, pdfhighlight=/O,%nesting=true,%frenchlinks,%
  urlcolor=CTurl, linkcolor=CTlink, citecolor=CTcitation, %pagecolor=RoyalBlue,%
  %urlcolor=Black, linkcolor=Black, citecolor=Black, %pagecolor=Black,%
  pdftitle={Matemáticas I para Ingeniería Industrial - Apuntes y Ejercicios},%
  pdfauthor={\myName, \myUni, \myDepartment	},%\textcopyright
  pdfsubject={},%
  pdfkeywords={},%
  pdfcreator={pdfLaTeX},%
  pdfproducer={LaTeX with hyperref and classicthesis}%
}


% ********************************************************************
% Setup autoreferences (hyperref and babel)
% ********************************************************************
% There are some issues regarding autorefnames
% http://www.tex.ac.uk/cgi-bin/texfaq2html?label=latexwords
% you have to redefine the macros for the
% language you use, e.g., american, ngerman
% (as chosen when loading babel/AtBeginDocument)
% ********************************************************************
\makeatletter
\@ifpackageloaded{babel}%
  {%
    \addto\extrasamerican{%
      \renewcommand*{\figureautorefname}{Figure}%
      \renewcommand*{\tableautorefname}{Table}%
      \renewcommand*{\partautorefname}{Part}%
      \renewcommand*{\chapterautorefname}{Chapter}%
      \renewcommand*{\sectionautorefname}{Section}%
      \renewcommand*{\subsectionautorefname}{Section}%
      \renewcommand*{\subsubsectionautorefname}{Section}%
    }%
    \addto\extrasngerman{%
      \renewcommand*{\paragraphautorefname}{Absatz}%
      \renewcommand*{\subparagraphautorefname}{Unterabsatz}%
      \renewcommand*{\footnoteautorefname}{Fu\"snote}%
      \renewcommand*{\FancyVerbLineautorefname}{Zeile}%
      \renewcommand*{\theoremautorefname}{Theorem}%
      \renewcommand*{\appendixautorefname}{Anhang}%
      \renewcommand*{\equationautorefname}{Gleichung}%
      \renewcommand*{\itemautorefname}{Punkt}%
    }%
      % Fix to getting autorefs for subfigures right (thanks to Belinda Vogt for changing the definition)
      \providecommand{\subfigureautorefname}{\figureautorefname}%
    }{\relax}
\makeatother


% ********************************************************************
% Development Stuff
% ********************************************************************
\listfiles
%\PassOptionsToPackage{l2tabu,orthodox,abort}{nag}
%  \usepackage{nag}
%\PassOptionsToPackage{warning, all}{onlyamsmath}
%  \usepackage{onlyamsmath}


% ****************************************************************************************************
% 7. Further adjustments (experimental)
% ****************************************************************************************************
% ********************************************************************
% Changing the text area
% ********************************************************************
%\areaset[current]{312pt}{761pt} % 686 (factor 2.2) + 33 head + 42 head \the\footskip
%\setlength{\marginparwidth}{7em}%
%\setlength{\marginparsep}{2em}%

% ********************************************************************
% Using different fonts
% ********************************************************************
%\usepackage[oldstylenums]{kpfonts} % oldstyle notextcomp
% \usepackage[osf]{libertine}
%\usepackage[light,condensed,math]{iwona}
%\renewcommand{\sfdefault}{iwona}
%\usepackage{lmodern} % <-- no osf support :-(
%\usepackage{cfr-lm} %
%\usepackage[urw-garamond]{mathdesign} <-- no osf support :-(
%\usepackage[default,osfigures]{opensans} % scale=0.95
%\usepackage[sfdefault]{FiraSans}
% \usepackage[opticals,mathlf]{MinionPro} % onlytext
% ********************************************************************
%\usepackage[largesc,osf]{newpxtext}
%\linespread{1.05} % a bit more for Palatino
% Used to fix these:
% https://bitbucket.org/amiede/classicthesis/issues/139/italics-in-pallatino-capitals-chapter
% https://bitbucket.org/amiede/classicthesis/issues/45/problema-testatine-su-classicthesis-style
% ********************************************************************
% ****************************************************************************************************

% ****************************************************************************************************
% If you like the classicthesis, then I would appreciate a postcard.
% My address can be found in the file ClassicThesis.pdf. A collection
% of the postcards I received so far is available online at
% http://postcards.miede.de
% ****************************************************************************************************
\usepackage{float}
\usepackage{multirow}
% ****************************************************************************************************
% 0. Set the encoding of your files. UTF-8 is the only sensible encoding nowadays. If you can't read
% äöüßáéçèê∂åëæƒÏ€ then change the encoding setting in your editor, not the line below. If your editor
% does not support utf8 use another editor!
% ****************************************************************************************************
\PassOptionsToPackage{utf8}{inputenc}
  \usepackage{inputenc}

\PassOptionsToPackage{T1}{fontenc} % T2A for cyrillics
  \usepackage{fontenc}

% ****************************************************************************************************
% 1. Configure classicthesis for your needs here, e.g., remove "drafting" below
% in order to deactivate the time-stamp on the pages
% (see ClassicThesis.pdf for more information):
% ****************************************************************************************************
\PassOptionsToPackage{
  drafting=true,    % print version information on the bottom of the pages
  tocaligned=false, % the left column of the toc will be aligned (no indentation)
  dottedtoc=false,  % page numbers in ToC flushed right
  eulerchapternumbers=true, % use AMS Euler for chapter font (otherwise Palatino)
  linedheaders=false,       % chaper headers will have line above and beneath
  floatperchapter=true,     % numbering per chapter for all floats (i.e., Figure 1.1)
  eulermath=false,  % use awesome Euler fonts for mathematical formulae (only with pdfLaTeX)
  beramono=true,    % toggle a nice monospaced font (w/ bold)
  palatino=true,    % deactivate standard font for loading another one, see the last section at the end of this file for suggestions
  style=classicthesis % classicthesis, arsclassica
}{classicthesis}


% ****************************************************************************************************
% 2. Personal data and user ad-hoc commands (insert your own data here)
% ****************************************************************************************************
\newcommand{\myTitle}{Matemáticas I\xspace}
\newcommand{\myCourse}{Ingenierías de la rama industrial}
\newcommand{\mySubtitle}{para Ingeniería Industrial\xspace}
%\newcommand{\myDegree}{Doktor-Ingenieur (Dr.-Ing.)\xspace}
\newcommand{\myName}{Diego García Zamora y Nuria V. Castillo González\xspace}
\newcommand{\myDepartment}{Departamento de Matemáticas\xspace}
\newcommand{\myUni}{Universidad de Jaén\xspace}
\newcommand{\myLocation}{Jaén\xspace}
\newcommand{\myTime}{Septiembre de 2026\xspace}
\newcommand{\myYear}{2026--2027\xspace}
\newcommand{\myPosition}{Profesor Ayudante Doctor\xspace}



%\newcommand{\myVersion}{\classicthesis}

% ********************************************************************
% Setup, finetuning, and useful commands
% ********************************************************************
\providecommand{\mLyX}{L\kern-.1667em\lower.25em\hbox{Y}\kern-.125emX\@}
\newcommand{\ie}{i.\,e.}
\newcommand{\Ie}{I.\,e.}
\newcommand{\eg}{e.\,g.}
\newcommand{\Eg}{E.\,g.}
% ****************************************************************************************************


% ****************************************************************************************************
% 3. Loading some handy packages
% ****************************************************************************************************

%********************************************************************
\PassOptionsToPackage{american}{babel} % change this to your language(s), main language last
% Spanish languages need extra options in order to work with this template
\PassOptionsToPackage{spanish,es-lcroman}{babel}
\usepackage{babel}
\addto\captionsamerican{\renewcommand{\contentsname}{Índice General}}
\addto\captionsamerican{\renewcommand{\listfigurename}{Lista de figuras}}
\addto\captionsamerican{\renewcommand{\figurename}{Fig.}}
\addto\captionsamerican{\renewcommand{\listtablename}{Lista de tablas}}
\addto\captionsamerican{\renewcommand{\tablename}{Tabla}}
\addto\captionsamerican{\renewcommand{\partname}{Parte}}
%\addto\captionsamerican{\renewcommand{\bibname}{Referencias}}
%\renewcommand{\bibname}{Referencias}
%\addto\captionsspanish{  \renewcommand{\contentsname}{Whatever}}

\usepackage{csquotes}
\PassOptionsToPackage{%
  %backend=biber,bibencoding=utf8, %instead of bibtex
  backend=bibtex8,bibencoding=ascii,%
  language=spanish,%
  %style=numeric-comp,%
  style=authoryear-comp, % Author 1999, 2010
  %bibstyle=authoryear,dashed=false, % dashed: substitute rep. author with ---
  sorting=nyt, % name, year, title
  maxbibnames=10, % default: 3, et al.
  %backref=true,%
  natbib=true % natbib compatibility mode (\citep and \citet still work)
}{biblatex}
    \usepackage{biblatex}

\PassOptionsToPackage{fleqn}{amsmath}       % math environments and more by the AMS
  \usepackage{amsmath}

% ********************************************************************
% General useful packages
% ********************************************************************
\usepackage{graphicx} %
\usepackage{scrhack} % fix warnings when using KOMA with listings package
\usepackage{xspace} % to get the spacing after macros right
\PassOptionsToPackage{printonlyused,smaller}{acronym}
  \usepackage{acronym} % nice macros for handling all acronyms in the thesis
  %\renewcommand{\bflabel}[1]{{#1}\hfill} % fix the list of acronyms --> no longer working
  %\renewcommand*{\acsfont}[1]{\textsc{#1}}
  %\renewcommand*{\aclabelfont}[1]{\acsfont{#1}}
  %\def\bflabel#1{{#1\hfill}}
  \def\bflabel#1{{\acsfont{#1}\hfill}}
  \def\aclabelfont#1{\acsfont{#1}}
% ****************************************************************************************************
%\usepackage{pgfplots} % External TikZ/PGF support (thanks to Andreas Nautsch)
%\usetikzlibrary{external}
%\tikzexternalize[mode=list and make, prefix=ext-tikz/]
% ****************************************************************************************************


% ****************************************************************************************************
% 4. Setup floats: tables, (sub)figures, and captions
% ****************************************************************************************************
\usepackage{tabularx} % better tables
  \setlength{\extrarowheight}{3pt} % increase table row height
\newcommand{\tableheadline}[1]{\multicolumn{1}{l}{\spacedlowsmallcaps{#1}}}
\newcommand{\myfloatalign}{\centering} % to be used with each float for alignment
\usepackage{subfig}
% ****************************************************************************************************


% ****************************************************************************************************
% 5. Setup code listings
% ****************************************************************************************************
\usepackage{listings}
%\lstset{emph={trueIndex,root},emphstyle=\color{BlueViolet}}%\underbar} % for special keywords
\lstset{language=[LaTeX]Tex,%C++,
  morekeywords={PassOptionsToPackage,selectlanguage},
  keywordstyle=\color{RoyalBlue},%\bfseries,
  basicstyle=\small\ttfamily,
  %identifierstyle=\color{NavyBlue},
  commentstyle=\color{Green}\ttfamily,
  stringstyle=\rmfamily,
  numbers=none,%left,%
  numberstyle=\scriptsize,%\tiny
  stepnumber=5,
  numbersep=8pt,
  showstringspaces=false,
  breaklines=true,
  %frameround=ftff,
  %frame=single,
  belowcaptionskip=.75\baselineskip
  %frame=L
}
% ****************************************************************************************************




% ****************************************************************************************************
% 6. Last calls before the bar closes
% ****************************************************************************************************
% ********************************************************************
% Her Majesty herself
% ********************************************************************
\usepackage{classicthesis}


% ********************************************************************
% Fine-tune hyperreferences (hyperref should be called last)
% ********************************************************************
\hypersetup{%
  %draft, % hyperref's draft mode, for printing see below
  colorlinks=true, linktocpage=true, pdfstartpage=3, pdfstartview=FitV,%
  % uncomment the following line if you want to have black links (e.g., for printing)
  %colorlinks=false, linktocpage=false, pdfstartpage=3, pdfstartview=FitV, pdfborder={0 0 0},%
  breaklinks=true, pageanchor=true,%
  pdfpagemode=UseNone, %
  % pdfpagemode=UseOutlines,%
  plainpages=false, bookmarksnumbered, bookmarksopen=true, bookmarksopenlevel=1,%
  hypertexnames=true, pdfhighlight=/O,%nesting=true,%frenchlinks,%
  urlcolor=CTurl, linkcolor=CTlink, citecolor=CTcitation, %pagecolor=RoyalBlue,%
  %urlcolor=Black, linkcolor=Black, citecolor=Black, %pagecolor=Black,%
  pdftitle={Matemáticas I para Ingeniería Industrial - Apuntes y Ejercicios},%
  pdfauthor={\myName, \myUni, \myDepartment	},%\textcopyright
  pdfsubject={},%
  pdfkeywords={},%
  pdfcreator={pdfLaTeX},%
  pdfproducer={LaTeX with hyperref and classicthesis}%
}


% ********************************************************************
% Setup autoreferences (hyperref and babel)
% ********************************************************************
% There are some issues regarding autorefnames
% http://www.tex.ac.uk/cgi-bin/texfaq2html?label=latexwords
% you have to redefine the macros for the
% language you use, e.g., american, ngerman
% (as chosen when loading babel/AtBeginDocument)
% ********************************************************************
\makeatletter
\@ifpackageloaded{babel}%
  {%
    \addto\extrasamerican{%
      \renewcommand*{\figureautorefname}{Figure}%
      \renewcommand*{\tableautorefname}{Table}%
      \renewcommand*{\partautorefname}{Part}%
      \renewcommand*{\chapterautorefname}{Chapter}%
      \renewcommand*{\sectionautorefname}{Section}%
      \renewcommand*{\subsectionautorefname}{Section}%
      \renewcommand*{\subsubsectionautorefname}{Section}%
    }%
    \addto\extrasngerman{%
      \renewcommand*{\paragraphautorefname}{Absatz}%
      \renewcommand*{\subparagraphautorefname}{Unterabsatz}%
      \renewcommand*{\footnoteautorefname}{Fu\"snote}%
      \renewcommand*{\FancyVerbLineautorefname}{Zeile}%
      \renewcommand*{\theoremautorefname}{Theorem}%
      \renewcommand*{\appendixautorefname}{Anhang}%
      \renewcommand*{\equationautorefname}{Gleichung}%
      \renewcommand*{\itemautorefname}{Punkt}%
    }%
      % Fix to getting autorefs for subfigures right (thanks to Belinda Vogt for changing the definition)
      \providecommand{\subfigureautorefname}{\figureautorefname}%
    }{\relax}
\makeatother


% ********************************************************************
% Development Stuff
% ********************************************************************
\listfiles
%\PassOptionsToPackage{l2tabu,orthodox,abort}{nag}
%  \usepackage{nag}
%\PassOptionsToPackage{warning, all}{onlyamsmath}
%  \usepackage{onlyamsmath}


% ****************************************************************************************************
% 7. Further adjustments (experimental)
% ****************************************************************************************************
% ********************************************************************
% Changing the text area
% ********************************************************************
%\areaset[current]{312pt}{761pt} % 686 (factor 2.2) + 33 head + 42 head \the\footskip
%\setlength{\marginparwidth}{7em}%
%\setlength{\marginparsep}{2em}%

% ********************************************************************
% Using different fonts
% ********************************************************************
%\usepackage[oldstylenums]{kpfonts} % oldstyle notextcomp
% \usepackage[osf]{libertine}
%\usepackage[light,condensed,math]{iwona}
%\renewcommand{\sfdefault}{iwona}
%\usepackage{lmodern} % <-- no osf support :-(
%\usepackage{cfr-lm} %
%\usepackage[urw-garamond]{mathdesign} <-- no osf support :-(
%\usepackage[default,osfigures]{opensans} % scale=0.95
%\usepackage[sfdefault]{FiraSans}
% \usepackage[opticals,mathlf]{MinionPro} % onlytext
% ********************************************************************
%\usepackage[largesc,osf]{newpxtext}
%\linespread{1.05} % a bit more for Palatino
% Used to fix these:
% https://bitbucket.org/amiede/classicthesis/issues/139/italics-in-pallatino-capitals-chapter
% https://bitbucket.org/amiede/classicthesis/issues/45/problema-testatine-su-classicthesis-style
% ********************************************************************
% ****************************************************************************************************

\addbibresource{Bibliography.bib}

% Paquetes matemáticos y gráficos de uso general
\usepackage{amsmath,amsthm,amsfonts,amssymb}
\usepackage{stmaryrd}
\usepackage{graphicx}
\usepackage{tikz}
\usepackage{pgfplots}
\pgfplotsset{compat=1.18}
\usetikzlibrary{arrows.meta,calc,math}
\usepackage[export]{adjustbox}
\usepackage[shortlabels]{enumitem}
\usepackage{listings}
\usepackage[most]{tcolorbox}

% Código visual para los elementos destacados. La paleta parte de los tonos
% sobrios de classicthesis y mantiene fondos muy claros para no sobrecargar
% visualmente las páginas.
\definecolor{BoxResult}{HTML}{596D36}      % verde oliva: resultados
\definecolor{BoxDefinition}{HTML}{5D7180}  % azul pizarra: definiciones
\definecolor{BoxExercise}{HTML}{8B624B}    % tierra: ejercicios resueltos
\definecolor{BoxExample}{HTML}{9A7B43}     % ocre: ejemplos

\tcbset{
  mathboxbase/.style={
    enhanced,
    breakable,
    sharp corners=south,
    arc=1.2mm,
    boxrule=0.35pt,
    leftrule=2.2pt,
    left=2.2mm,right=2.2mm,top=1.6mm,bottom=1.6mm,
    before skip=8pt plus 2pt minus 2pt,
    after skip=8pt plus 2pt minus 2pt,
    fonttitle=\bfseries,
  },
  resultbox/.style={
    mathboxbase,
    colframe=BoxResult,
    colback=BoxResult!4!white,
  },
  definitionbox/.style={
    mathboxbase,
    colframe=BoxDefinition,
    colback=BoxDefinition!4!white,
  },
  solvedexercisebox/.style={
    mathboxbase,
    colframe=BoxExercise,
    colback=BoxExercise!4!white,
    colbacklower=BoxExercise!2!white,
    coltitle=BoxExercise,
    colbacktitle=BoxExercise!7!white,
    boxed title style={boxrule=0pt,sharp corners,arc=0pt},
    attach boxed title to top left={xshift=1.5mm,yshift*=-1mm},
    segmentation style={draw=BoxExercise!45,line width=.35pt},
  },
  examplebox/.style={
    mathboxbase,
    colframe=BoxExample,
    colback=BoxExample!4!white,
    colbacklower=BoxExample!2!white,
    coltitle=BoxExample,
    colbacktitle=BoxExample!7!white,
    boxed title style={boxrule=0pt,sharp corners,arc=0pt},
    attach boxed title to top left={xshift=1.5mm,yshift*=-1mm},
    segmentation style={draw=BoxExample!45,line width=.35pt},
  },
}

% Los entornos de resultados y definiciones conservan su numeración y su
% sintaxis; tcolorbox aporta únicamente el tratamiento visual.
\tcolorboxenvironment{theorem}{resultbox}
\tcolorboxenvironment{proposition}{resultbox}
\tcolorboxenvironment{lemma}{resultbox}
\tcolorboxenvironment{corollary}{resultbox}
\tcolorboxenvironment{definition}{definitionbox}

\NewDocumentEnvironment{solvedexercise}{o}
  {%
    \refstepcounter{exercise}%
    \begin{tcolorbox}[
      solvedexercisebox,
      title={Ejercicio resuelto \theexercise\IfValueT{#1}{: #1}}
    ]%
  }
  {\end{tcolorbox}}

% Símbolos habituales
\newcommand{\R}{\mathbb{R}}
\newcommand{\N}{\mathbb{N}}
\newcommand{\C}{\mathbb{C}}
\newcommand{\Q}{\mathbb{Q}}
\newcommand{\Z}{\mathbb{Z}}
\newcommand{\tq}{\text{ : }}
\newcommand{\pt}{\;\forall\;}

% Marca para ejercicios recomendados para Wolfram Mathematica
\newcommand{\M}{\textbf{\textsf{[M]}}\xspace}

% Entornos, manteniendo la filosofía de la plantilla
\newtheorem{theorem}{Teorema}[chapter]
\newtheorem{proposition}{Proposición}[chapter]
\newtheorem{lemma}{Lema}[chapter]
\newtheorem{corollary}{Corolario}[chapter]
\newtheorem{definition}{Definición}[chapter]
\newcounter{example}[chapter]
\renewcommand{\theexample}{\thechapter.\arabic{example}}
\NewDocumentEnvironment{example}{o}
  {%
    \refstepcounter{example}%
    \begin{tcolorbox}[
      examplebox,
      title={Ejemplo \theexample\IfValueT{#1}{: #1}}
    ]%
  }
  {\end{tcolorbox}}
\newtheorem{exercise}{Ejercicio}[chapter]
\renewcommand*{\proofname}{Demostración}
\newenvironment{resolution}
  {\begin{itemize}\item[]\textbf{Resolución:}}
  {\end{itemize}}

% Código de Wolfram Mathematica. Se usará solo cuando aporte valor didáctico.
\lstdefinelanguage{Wolfram}{
  sensitive=true,
  morecomment=[s]{(*}{*)},
  morestring=[b]"
}
\lstnewenvironment{mathematicacode}
  {\lstset{language=Wolfram,basicstyle=\small\ttfamily,breaklines=true,showstringspaces=false}}
  {}

\begin{document}
\frenchspacing
\raggedbottom
\selectlanguage{spanish}
\pagenumbering{roman}
\pagestyle{plain}

\cfoot{\hrulefill\\ {D. García Zamora y N. Castillo González\hfill Curso \myYear}}

% Frontmatter
\include{FrontBackmatter/Titlepage}
\include{FrontBackmatter/Titleback}
\include{FrontBackmatter/Introduccion}
\cleardoublepage
\include{FrontBackmatter/Contents}

% Mainmatter
\cleardoublepage
\pagestyle{scrheadings}
\pagenumbering{arabic}

\chapter{Números reales y complejos}
\label{ch:numeros}

Las magnitudes que aparecen en ingeniería se expresan habitualmente mediante números reales: longitudes, temperaturas, caudales, concentraciones, potencias o tiempos son algunos ejemplos. Sin embargo, existen situaciones en las que los números reales no son suficientes. En particular, el estudio de fenómenos oscilatorios y de corriente alterna conduce de forma natural a los números complejos. En este capítulo repasaremos las propiedades de los números reales que utilizaremos a lo largo del curso y presentaremos las distintas formas de representar y operar con números complejos.

\section{Números reales}
\label{sec:reales}

\subsection{Conjuntos numéricos, orden e intervalos}

Recordemos los principales conjuntos numéricos:
\[
\N=\{1,2,3,\ldots\},\qquad
\Z=\{0,\pm1,\pm2,\ldots\},
\]
\[
\Q=\left\{\frac{p}{q}:p\in\Z,\;q\in\N\right\},
\qquad
\R=\text{conjunto de los números reales}.
\]
Se verifican las inclusiones
\[
\N\subset\Z\subset\Q\subset\R.
\]
Los números de \(\R\setminus\Q\) reciben el nombre de \emph{irracionales}. Por ejemplo,
\(\sqrt{2}\), \(\pi\) y \(e\) son irracionales.

A lo largo de los apuntes utilizaremos, entre otros, los símbolos
\[
\begin{array}{ll}
\forall\;	\text{(para todo)}, & \exists\;	\text{(existe)},\\[2mm]
\in\;	\text{(pertenece)}, & \subseteq\;	\text{(está contenido en)}.
\end{array}
\]

El conjunto \(\R\) está ordenado: dados \(x,y\in\R\), podemos comparar ambos números mediante las relaciones \(<\), \(\leq\), \(>\) y \(\geq\). Las reglas elementales para trabajar con desigualdades se utilizarán continuamente. En particular,
\begin{align*}
 x\leq y &\Longrightarrow x+z\leq y+z,\\
 x\leq y,\ z>0 &\Longrightarrow xz\leq yz,\\
 x\leq y,\ z<0 &\Longrightarrow xz\geq yz.
\end{align*}
La última propiedad es especialmente importante: al multiplicar o dividir una desigualdad por un número negativo, el sentido de la desigualdad cambia.

\begin{example}
Resolver la desigualdad
\[
3-2x<7.
\]
\tcblower
\textbf{Resolución.}\par\smallskip
Restando \(3\) en ambos miembros obtenemos \(-2x<4\). Al dividir entre \(-2\), la desigualdad cambia de sentido:
\[
x>-2.
\]
\end{example}

Los intervalos permiten describir de forma compacta subconjuntos de la recta real. Dados \(a,b\in\R\), con \(a<b\), escribiremos
\begin{align*}
[a,b]&=\{x\in\R:a\leq x\leq b\},\\
]a,b]&=\{x\in\R:a<x\leq b\},\\
[a,b[&=\{x\in\R:a\leq x<b\},\\
]a,b[&=\{x\in\R:a<x<b\}.
\end{align*}
También utilizaremos intervalos no acotados, como
\[
[a,+\infty[=\{x\in\R:x\geq a\},
\qquad
]-\infty,b[=\{x\in\R:x<b\}.
\]
El símbolo \(\infty\) no representa un número real y, por tanto, nunca se incluye como extremo de un intervalo.

\begin{example}
Una pieza cilíndrica debe fabricarse con diámetro nominal \(20\,\mathrm{mm}\) y una tolerancia de \(0.03\,\mathrm{mm}\). Si \(d\) es el diámetro real de la pieza, la condición de aceptación puede escribirse como
\[
19.97\leq d\leq20.03,
\]
es decir,
\[
d\in[19.97,20.03].
\]
La notación mediante intervalos traduce directamente una especificación de fabricación en una condición matemática.
\end{example}

Dados dos subconjuntos \(A,B\subseteq\R\), utilizaremos además
\[
A\cup B=\{x:x\in A\text{ o }x\in B\},
\qquad
A\cap B=\{x:x\in A\text{ y }x\in B\}.
\]
Estas operaciones aparecerán de forma natural al estudiar dominios de funciones y conjuntos de soluciones de desigualdades.

\begin{solvedexercise}
Determinar el conjunto
\[
A=[-2,4]\cap\bigl(]-\infty,1[\cup[3,+\infty[\bigr).
\]
\tcblower
\textbf{Resolución.}\par\smallskip
Distribuyendo la intersección,
\[
A=\bigl([-2,4]\cap]-\infty,1[\bigr)\cup
\bigl([-2,4]\cap[3,+\infty[\bigr),
\]
y, por tanto,
\[
A=[-2,1[\cup[3,4].
\]
\end{solvedexercise}

\subsection{Valor absoluto y distancia}

\begin{definition}[Valor absoluto]
Para \(x\in\R\), se define el valor absoluto de \(x\) por
\[
|x|=
\begin{cases}
 x, & x\geq0,\\
 -x, & x<0.
\end{cases}
\]
\end{definition}

El valor absoluto de un número real es su distancia al origen. En particular,
\[
|x|\geq0,
\qquad
|x|=0\Longleftrightarrow x=0,
\qquad
|-x|=|x|.
\]
Además, para cualesquiera \(x,y\in\R\),
\[
|xy|=|x|\,|y|,
\qquad
|x+y|\leq |x|+|y|.
\]
La segunda desigualdad recibe el nombre de \emph{desigualdad triangular}.

El valor absoluto permite expresar distancias en la recta real.

\begin{definition}[Distancia en \(\R\)]
Dados \(x,y\in\R\), la distancia entre \(x\) e \(y\) es
\[
d(x,y)=|x-y|.
\]
\end{definition}

De esta interpretación se obtienen equivalencias muy útiles. Si \(a>0\), entonces
\[
|x|<a\Longleftrightarrow -a<x<a,
\qquad
|x|\leq a\Longleftrightarrow -a\leq x\leq a.
\]
Más generalmente,
\[
|x-c|<a\Longleftrightarrow c-a<x<c+a,
\]
es decir, \(|x-c|<a\) describe los puntos cuya distancia a \(c\) es menor que \(a\).

\begin{example}
La temperatura de operación de un reactor debe mantenerse a menos de \(2.5\,^{\circ}\mathrm{C}\) de su valor nominal \(T_0=350\,^{\circ}\mathrm{C}\). La condición puede escribirse como
\[
|T-350|<2.5.
\]
Equivalentemente,
\[
347.5<T<352.5.
\]
La formulación mediante valor absoluto expresa directamente una tolerancia alrededor de un valor de referencia.
\end{example}

\begin{solvedexercise}
Resolver
\[
|2x-3|\leq5.
\]
\tcblower
\textbf{Resolución.}\par\smallskip
Usando la equivalencia anterior,
\[
-5\leq2x-3\leq5.
\]
Sumando \(3\),
\[
-2\leq2x\leq8,
\]
y dividiendo entre \(2\),
\[
-1\leq x\leq4.
\]
Por tanto, el conjunto solución es \([-1,4]\).
\end{solvedexercise}

También será útil la desigualdad triangular en la forma
\[
\bigl||x|-|y|\bigr|\leq|x-y|.
\]
Esta expresión indica que la diferencia entre las distancias de dos puntos al origen nunca puede superar la distancia entre ambos puntos.

\subsection{Potencias, raíces y notación científica}

Las reglas usuales de las potencias serán empleadas de forma constante. Para \(a\neq0\) y \(m,n\in\Z\),
\[
a^m a^n=a^{m+n},
\qquad
\frac{a^m}{a^n}=a^{m-n},
\qquad
(a^m)^n=a^{mn}.
\]
Si \(n\in\N\), la raíz \(n\)-ésima real de un número \(a\) es un número \(x\) que satisface \(x^n=a\). Cuando \(n\) es par, una raíz real sólo existe para \(a\geq0\). Además, el símbolo \(\sqrt[n]{a}\) representa la raíz real no negativa cuando \(n\) es par.

Por ejemplo,
\[
\sqrt{9}=3,
\qquad
x^2=9\Longleftrightarrow x=\pm3.
\]
No debe confundirse el valor de la raíz cuadrada principal con el conjunto de soluciones de una ecuación cuadrática.

En aplicaciones de ingeniería es frecuente trabajar con magnitudes muy grandes o muy pequeñas. La notación científica escribe un número no nulo en la forma
\[
a\,10^n,\qquad 1\leq |a|<10,\quad n\in\Z.
\]
Así,
\[
0.000047=4.7\,10^{-5},
\qquad
2\,350\,000=2.35\,10^6.
\]

\begin{example}
Una corriente de \(35\,\mu\mathrm{A}\) equivale a
\[
35\,10^{-6}\,\mathrm{A}=3.5\,10^{-5}\,\mathrm{A}.
\]
Una potencia de \(2.4\,\mathrm{MW}\) equivale a
\[
2.4\,10^6\,\mathrm{W}.
\]
La notación científica permite operar con ambas magnitudes sin perder de vista su orden de magnitud.
\end{example}

\section{Números complejos}
\label{sec:complejos}

La ecuación
\[
x^2+1=0
\]
no tiene solución en \(\R\), pues \(x^2\geq0\) para todo \(x\in\R\). Para resolver este tipo de ecuaciones ampliamos el sistema numérico introduciendo la unidad imaginaria.

\begin{definition}[Unidad imaginaria]
La unidad imaginaria, denotada por \(i\), se caracteriza por
\[
i^2=-1.
\]
\end{definition}

\begin{definition}[Número complejo]
Un número complejo es una expresión de la forma
\[
z=a+bi,
\]
donde \(a,b\in\R\). El conjunto de los números complejos se denota por
\[
\C=\{a+bi:a,b\in\R\}.
\]
El número \(a\) es la \emph{parte real} de \(z\), denotada por \(\operatorname{Re}(z)\), y \(b\) es la \emph{parte imaginaria}, denotada por \(\operatorname{Im}(z)\).
\end{definition}

Dos números complejos son iguales si y sólo si coinciden sus partes real e imaginaria:
\[
a+bi=c+di\Longleftrightarrow a=c\text{ y }b=d.
\]
Los números reales se identifican con los complejos cuya parte imaginaria es cero, de modo que
\[
\R\subset\C.
\]

\subsection{Forma binómica, conjugado y operaciones}

La expresión \(z=a+bi\) recibe el nombre de \emph{forma binómica}. La suma y el producto se realizan aplicando las reglas algebraicas usuales y utilizando \(i^2=-1\):
\begin{align*}
(a+bi)+(c+di)&=(a+c)+(b+d)i,\\
(a+bi)(c+di)&=(ac-bd)+(ad+bc)i.
\end{align*}

\begin{example}
Sean \(z=3-2i\) y \(w=1+4i\). Entonces
\[
z+w=4+2i,
\]
y
\[
zw=(3-2i)(1+4i)=3+12i-2i-8i^2=11+10i.
\]
\end{example}

Las potencias de \(i\) se repiten con período \(4\):
\[
i^0=1,\qquad i^1=i,\qquad i^2=-1,\qquad i^3=-i,\qquad i^4=1.
\]
Por tanto, para calcular \(i^n\) basta conocer el resto de dividir \(n\) entre \(4\).

\begin{example}
Como \(241=4\cdot60+1\),
\[
i^{241}=i.
\]
\end{example}

\begin{definition}[Conjugado]
Si \(z=a+bi\), el conjugado de \(z\) es
\[
\overline{z}=a-bi.
\]
\end{definition}

El conjugado verifica
\[
\overline{z+w}=\overline z+\overline w,
\qquad
\overline{zw}=\overline z\,\overline w,
\qquad
z\overline z=a^2+b^2.
\]
Esta última identidad permite dividir números complejos.

\begin{proposition}
Sea \(z=a+bi\neq0\). Entonces
\[
\frac1z=\frac{\overline z}{z\overline z}
=\frac{a-bi}{a^2+b^2}.
\]
\end{proposition}

\begin{example}
Calcular
\[
\frac{1+2i}{3-4i}.
\]
Multiplicamos numerador y denominador por el conjugado del denominador:
\[
\frac{1+2i}{3-4i}
=\frac{(1+2i)(3+4i)}{(3-4i)(3+4i)}
=\frac{-5+10i}{25}
=-\frac15+\frac25i.
\]
\end{example}

\begin{solvedexercise}
Encontrar \(x,y\in\R\) tales que
\[
3x+2yi-xi+5y=7+5i.
\]
\tcblower
\textbf{Resolución.}\par\smallskip
Agrupando parte real e imaginaria,
\[
(3x+5y)+(2y-x)i=7+5i.
\]
La igualdad de números complejos equivale al sistema
\[
\begin{cases}
3x+5y=7,\\
-x+2y=5.
\end{cases}
\]
De la segunda ecuación, \(x=2y-5\). Sustituyendo en la primera,
\[
3(2y-5)+5y=7\Longrightarrow11y=22\Longrightarrow y=2,
\]
y, por tanto, \(x=-1\).
\end{solvedexercise}

\subsection{Representación geométrica, módulo y argumento}

A cada número complejo \(z=a+bi\) le asociamos el punto \((a,b)\) del plano. El eje horizontal se denomina \emph{eje real} y el vertical \emph{eje imaginario}. Esta identificación permite interpretar un número complejo como un punto o como un vector del plano.

\begin{center}
\begin{tikzpicture}[scale=0.9]
  \draw[-latex] (-3.2,0)--(4.2,0) node[right] {$\operatorname{Re}$};
  \draw[-latex] (0,-3.0)--(0,3.4) node[above] {$\operatorname{Im}$};
  \draw[-latex,thick] (0,0)--(3,2) node[midway,above left] {$z$};
  \draw[dashed] (3,0)--(3,2)--(0,2);
  \fill (3,2) circle (1.8pt) node[above right] {$a+bi$};
  \node[below] at (3,0) {$a$};
  \node[left] at (0,2) {$b$};
  \draw (0.7,0) arc[start angle=0,end angle=33.69,radius=0.7];
  \node at (0.95,0.25) {$\theta$};
\end{tikzpicture}
\end{center}

\begin{definition}[Módulo]
El módulo de \(z=a+bi\) es
\[
|z|=\sqrt{a^2+b^2}.
\]
\end{definition}

Geométricamente, \(|z|\) es la distancia de \(z\) al origen. Más generalmente, la distancia entre \(z,w\in\C\) es
\[
|z-w|.
\]
El módulo conserva muchas de las propiedades del valor absoluto real:
\[
|z|\geq0,
\qquad
|z|=0\Longleftrightarrow z=0,
\qquad
|zw|=|z|\,|w|,
\]
y se verifica la desigualdad triangular
\[
|z+w|\leq|z|+|w|.
\]
Además,
\[
|z|^2=z\overline z.
\]

\begin{example}
El conjunto
\[
\{z\in\C:|z-(2+3i)|\leq1\}
\]
está formado por los puntos cuya distancia al punto \(2+3i\) es menor o igual que \(1\). Por tanto, representa el disco cerrado de centro \((2,3)\) y radio \(1\).
\end{example}

\begin{definition}[Argumento]
Sea \(z\in\C\), \(z\neq0\). Un argumento de \(z\) es cualquier ángulo \(\theta\in\R\) tal que
\[
z=|z|\bigl(\cos\theta+i\sin\theta\bigr).
\]
El argumento principal, denotado por \(\arg(z)\), es el único argumento perteneciente a \(]-\pi,\pi]\).
\end{definition}

Si \(\theta\) es un argumento de \(z\), entonces todos sus argumentos son
\[
\theta+2k\pi,\qquad k\in\Z.
\]
Para calcular correctamente el argumento no basta, en general, con usar \(\arctan(b/a)\): es necesario tener en cuenta el cuadrante en el que se encuentra el punto \((a,b)\).

\begin{example}
Consideremos
\[
z=-2+2\sqrt3\,i.
\]
Su módulo es
\[
|z|=\sqrt{(-2)^2+(2\sqrt3)^2}=4.
\]
Como \(z\) está en el segundo cuadrante y
\[
\tan\theta=\frac{2\sqrt3}{-2}=-\sqrt3,
\]
el argumento principal es
\[
\arg(z)=\frac{2\pi}{3}.
\]
Por tanto,
\[
z=4\left(\cos\frac{2\pi}{3}+i\sin\frac{2\pi}{3}\right).
\]
\end{example}

\subsection{Forma polar y exponencial. Producto, cociente y potencias}

Si \(z\neq0\), podemos escribir
\[
z=\rho(\cos\theta+i\sin\theta),
\qquad \rho=|z|>0,
\]
que es la \emph{forma trigonométrica} o \emph{polar} de \(z\). También utilizaremos la notación
\[
z=\rho e^{i\theta},
\]
basada en la fórmula de Euler
\[
e^{i\theta}=\cos\theta+i\sin\theta.
\]
En este capítulo la expresión exponencial se empleará como una forma cómoda de representar números complejos; el estudio de la función exponencial se realizará en el capítulo siguiente para variable real.

La forma polar simplifica notablemente productos y cocientes. Si
\[
z=\rho e^{i\theta},
\qquad
w=\sigma e^{i\varphi},
\]
entonces
\[
zw=\rho\sigma e^{i(\theta+\varphi)},
\qquad
\frac zw=\frac{\rho}{\sigma}e^{i(\theta-\varphi)},\quad w\neq0.
\]
Así, al multiplicar números complejos se multiplican los módulos y se suman los argumentos; al dividir, se dividen los módulos y se restan los argumentos.

Esta propiedad tiene una interpretación geométrica inmediata: multiplicar por un número complejo \(\rho e^{i\theta}\) combina una dilatación de factor \(\rho\) y una rotación de ángulo \(\theta\).

\begin{example}
Multiplicar por \(i=e^{i\pi/2}\) rota un vector \(90^\circ\) en sentido antihorario sin modificar su módulo. Por ejemplo,
\[
i(2+i)=-1+2i.
\]
El punto \((2,1)\) pasa al punto \((-1,2)\), que es precisamente su rotación de \(90^\circ\) alrededor del origen.
\end{example}

\begin{proposition}[Fórmula de De Moivre]
Sea
\[
z=\rho(\cos\theta+i\sin\theta)
\]
y sea \(n\in\Z\). Entonces
\[
z^n=\rho^n\bigl(\cos(n\theta)+i\sin(n\theta)\bigr),
\]
siempre que la expresión esté definida.
\end{proposition}

\begin{example}
Calcular \((1+i)^8\).

Como
\[
1+i=\sqrt2\,e^{i\pi/4},
\]
por De Moivre,
\[
(1+i)^8=(\sqrt2)^8e^{i2\pi}=16.
\]
El cálculo en forma polar evita desarrollar una potencia de orden elevado en forma binómica.
\end{example}

La representación compleja es especialmente útil para magnitudes sinusoidales. En Ingeniería Eléctrica suele emplearse la letra \(j\) en lugar de \(i\), para evitar confusión con la intensidad de corriente. Matemáticamente ambas notaciones representan la misma unidad imaginaria, \(j^2=-1\).

\begin{example}[Impedancia en corriente alterna]
Una impedancia
\[
Z=4+3i\;\Omega
\]
tiene módulo
\[
|Z|=5\;\Omega
\]
y argumento
\[
\arg(Z)=\arctan\frac34\approx0.644\;\mathrm{rad}\approx36.87^\circ.
\]
Por tanto,
\[
Z=5e^{i0.644}\;\Omega.
\]
La forma binómica separa las componentes resistiva y reactiva; la forma polar muestra directamente el módulo de la impedancia y el desfase que introduce.
\end{example}

\begin{solvedexercise}
Una corriente sinusoidal se representa mediante el fasor
\[
I=2e^{-i\pi/6}\;\mathrm{A}
\]
y atraviesa una impedancia
\[
Z=5e^{i\pi/4}\;\Omega.
\]
Calcular el fasor de tensión \(V=ZI\) en forma polar y binómica.
\tcblower
\textbf{Resolución.}\par\smallskip
En forma polar,
\[
V=ZI=10e^{i(\pi/4-\pi/6)}=10e^{i\pi/12}\;\mathrm{V}.
\]
Pasando a forma binómica,
\[
V=10\left(\cos\frac{\pi}{12}+i\sin\frac{\pi}{12}\right).
\]
Como
\[
\cos\frac{\pi}{12}=\frac{\sqrt6+\sqrt2}{4},
\qquad
\sin\frac{\pi}{12}=\frac{\sqrt6-\sqrt2}{4},
\]
obtenemos
\[
V=\frac52(\sqrt6+\sqrt2)
+i\frac52(\sqrt6-\sqrt2)\;\mathrm{V}.
\]
Numéricamente,
\[
V\approx9.659+2.588i\;\mathrm{V}.
\]
\end{solvedexercise}

\subsection{Raíces de números complejos}

Sea \(w=\rho e^{i\theta}\neq0\). Resolver
\[
z^n=w
\]
consiste en buscar todos los números complejos \(z=re^{i\varphi}\) tales que
\[
r^ne^{in\varphi}=\rho e^{i\theta}.
\]
Debe cumplirse
\[
r=\sqrt[n]{\rho}
\]
y
\[
n\varphi=\theta+2k\pi,
\qquad k\in\Z.
\]
Por tanto, existen exactamente \(n\) raíces distintas.

\begin{proposition}[Raíces \(n\)-ésimas]
Sea \(w=\rho e^{i\theta}\neq0\), con \(\rho>0\), y \(n\in\N\). Las soluciones de
\[
z^n=w
\]
son
\[
z_k=\sqrt[n]{\rho}\,e^{i(\theta+2k\pi)/n},
\qquad k=0,1,\ldots,n-1.
\]
\end{proposition}

Todas las raíces tienen el mismo módulo y sus argumentos difieren en \(2\pi/n\). Geométricamente son los vértices de un polígono regular centrado en el origen.

\begin{example}
Resolver
\[
z^3=-8i.
\]
Escribimos
\[
-8i=8e^{-i\pi/2}.
\]
Las raíces cúbicas tienen módulo \(2\) y argumentos
\[
\frac{-\pi/2+2k\pi}{3},\qquad k=0,1,2.
\]
Por tanto,
\[
z_0=2e^{-i\pi/6},\qquad
z_1=2e^{i\pi/2},\qquad
z_2=2e^{i7\pi/6}.
\]
En forma binómica,
\[
z_0=\sqrt3-i,
\qquad
z_1=2i,
\qquad
z_2=-\sqrt3-i.
\]
\end{example}

\begin{center}
\begin{tikzpicture}[scale=1.0]
  \draw[-latex] (-2.7,0)--(2.7,0) node[right] {$\operatorname{Re}$};
  \draw[-latex] (0,-2.7)--(0,2.7) node[above] {$\operatorname{Im}$};
  \draw[dashed] (0,0) circle (2);
  \foreach \x/\y/\lab in {1.732/-1/$z_0$,0/2/$z_1$,-1.732/-1/$z_2$}{
    \fill (\x,\y) circle (2pt) node[above right] {\lab};
    \draw[thin] (0,0)--(\x,\y);
  }
\end{tikzpicture}
\end{center}

\begin{example}[Raíces conjugadas]
La ecuación
\[
x^2-6x+10=0
\]
tiene discriminante
\[
\Delta=36-40=-4.
\]
Por tanto,
\[
x=\frac{6\pm\sqrt{-4}}2=3\pm i.
\]
Como los coeficientes del polinomio son reales, las raíces complejas aparecen conjugadas. En efecto,
\[
(x-(3+i))(x-(3-i))=x^2-6x+10.
\]
\end{example}

\paragraph{Comprobación con Mathematica.}
En \textit{Wolfram Mathematica} la unidad imaginaria se escribe con la letra mayúscula \texttt{I}. Por ejemplo, las raíces de la ecuación anterior pueden obtenerse mediante
\begin{mathematicacode}
Solve[z^3 == -8 I, z]
\end{mathematicacode}
Para consultar módulo y argumento principal se utilizan, respectivamente,
\begin{mathematicacode}
Abs[-2 + 2 Sqrt[3] I]
Arg[-2 + 2 Sqrt[3] I]
\end{mathematicacode}
No es necesario memorizar ahora la sintaxis: los comandos utilizados durante el curso se recopilan en el glosario final.

\section*{Ejercicios propuestos}
\label{sec:ejercicios-t1}

Los ejercicios siguientes combinan cálculo elemental, razonamiento y problemas de aplicación. Los marcados con \M son especialmente adecuados para ser resueltos o comprobados también con \textit{Wolfram Mathematica}.

\begin{enumerate}[label=\textbf{\arabic*.},leftmargin=*]

\item Clasificar los siguientes números indicando a cuáles de los conjuntos \(\N\), \(\Z\), \(\Q\) y \(\R\) pertenecen:
\[
-7,\qquad \frac{13}{5},\qquad \sqrt{49},\qquad \sqrt3,\qquad 0,\qquad \pi-3.
\]

\item Escribir como unión de intervalos los conjuntos siguientes:
\begin{enumerate}[label=\alph*)]
\item \(\{x\in\R:-2\leq x<4\}\).
\item \(\{x\in\R:x<1\text{ o }x\geq5\}\).
\item \([-3,6]\setminus]0,2[\).
\item \(]-\infty,4]\cap[-1,+\infty[\).
\end{enumerate}

\item Resolver las siguientes desigualdades:
\begin{enumerate}[label=\alph*)]
\item \(5-3x\geq2x+10\).
\item \((x-2)(x+3)\leq0\).
\item \(\dfrac{x-1}{x+2}>0\).
\end{enumerate}

\item Resolver:
\begin{enumerate}[label=\alph*)]
\item \(|x-3|<5\).
\item \(|2x+1|\geq3\).
\item \(|x-2|+|x+1|\leq5\).
\end{enumerate}

\item Una resistencia de \(220\,\Omega\) se fabrica con tolerancia del \(5\%\). Expresar mediante un intervalo el conjunto de valores admisibles de su resistencia real \(R\). Escribir también la misma condición mediante una desigualdad con valor absoluto.

\item El diámetro de un eje debe satisfacer \(d=40.00\pm0.02\,\mathrm{mm}\). Una medición da \(d=40.018\,\mathrm{mm}\). Determinar si la pieza cumple la especificación y expresar la condición de aceptación mediante valor absoluto.

\item Demostrar, para \(x,y\in\R\), que
\[
\bigl||x|-|y|\bigr|\leq|x-y|.
\]
Indicar cuándo puede alcanzarse la igualdad.

\item Escribir en notación científica y en unidades del SI:
\begin{enumerate}[label=\alph*)]
\item \(36\,\mu\mathrm{m}\).
\item \(4.7\,\mathrm{GJ}\).
\item \(0.85\,\mathrm{mA}\).
\item \(125\,\mathrm{nm}\).
\end{enumerate}

\item Dados
\[
z_1=2i,\qquad z_2=1+3i,\qquad z_3=\frac1{2-i},\qquad z_4=5-2i,
\]
calcular:
\begin{enumerate}[label=\alph*)]
\item \(z_2+z_4\),
\item \(z_1-z_3\),
\item \(z_2z_3\),
\item \(z_4/z_2\),
\item \(z_2\overline{z_2}\),
\item \(|z_4|\,|z_3|\).
\end{enumerate}

\item Demostrar las propiedades
\[
\overline{z+w}=\overline z+\overline w,
\qquad
\overline{zw}=\overline z\,\overline w,
\qquad
z\overline z=|z|^2.
\]

\item Escribir en forma polar y exponencial:
\[
1-i,\qquad -4,\qquad 8i,\qquad -1+i,\qquad \sqrt3-i,
\qquad -2+2\sqrt3\,i.
\]

\item Escribir en forma binómica:
\[
2e^{i\pi/3},\qquad e^{i\pi},\qquad
3e^{-i\pi/6},\qquad
7e^{i3\pi/2},\qquad
10e^{i\pi/4}.
\]

\item Calcular, dando el resultado final en forma binómica:
\begin{enumerate}[label=\alph*)]
\item \((2-2i)^7\).
\item \((1+i)^8\).
\item \(i^{241}\).
\item \(\left(\dfrac{1-i}{1+i}\right)^{12}\).
\end{enumerate}

\item Comprobar las identidades:
\begin{enumerate}[label=\alph*)]
\item \((1-i)^4=-4\).
\item \(\dfrac{5i}{2+i}=1+2i\).
\item \((-1+i)^7=-8(1+i)\).
\end{enumerate}

\item Describir geométricamente los siguientes conjuntos:
\begin{enumerate}[label=\alph*)]
\item \(\{z\in\C:|z|=1\}\).
\item \(\{z\in\C:|z|<2\}\).
\item \(\{z\in\C:|z-(2+3i)|\leq1\}\).
\item \(\{z\in\C:|z-2i|>3\}\).
\item \(\{z\in\C:z+\overline z=4\}\).
\end{enumerate}

\item Hallar \(a\in\R\) para que el número \(z=a+3i\) tenga módulo \(5\). Interpretar geométricamente las soluciones.

\item Encontrar \(x,y\in\R\) tales que
\[
(2x-y)+(x+3y)i=7-i.
\]

\item Dar la expresión del polinomio mónico de segundo grado con coeficientes reales cuyas raíces son
\[
8+3i\qquad\text{y}\qquad8-3i.
\]
Generalizar el resultado para las raíces \(a\pm bi\).

\item Resolver y expresar las soluciones en forma polar:
\begin{enumerate}[label=\alph*)]
\item \(z^3+8i=0\).
\item \(z^4=16\).
\item \(z^5+32i=0\).
\item \((z^2-7)(z^3+i)=0\).
\end{enumerate}

\item Los afijos de seis números complejos forman un hexágono regular centrado en el origen. Uno de sus vértices es \(z_1=1+i\). Hallar los otros cinco vértices y expresarlos en forma binómica.

\item Sea
\[
w=9e^{i3\pi/2}.
\]
\begin{enumerate}[label=\alph*)]
\item Resolver \(z^3-w=0\) y expresar las soluciones en forma trigonométrica.
\item Calcular la parte real y la parte imaginaria de \(w^{2022}\) sin desarrollar potencias en forma binómica.
\end{enumerate}

\item Sean \(z,w\in\C\setminus\{0\}\). Razonar las siguientes cuestiones:
\begin{enumerate}[label=\alph*)]
\item Si \(w=\overline z\), ¿qué relación existe entre sus módulos y entre sus argumentos?
\item Si \(w=-z\), ¿qué relación existe entre sus módulos y entre sus argumentos?
\item Si \(|z|=|w|=3\), ¿dónde están situados sus afijos?
\item Si \(w^{10}=z\), ¿qué relación existe entre \(|w|\) y \(|z|\)?
\end{enumerate}

\item Una impedancia eléctrica vale
\[
Z=6-8i\;\Omega.
\]
\begin{enumerate}[label=\alph*)]
\item Obtener su módulo y argumento principal.
\item Expresarla en forma polar.
\item Si \(I=3e^{i\pi/6}\,\mathrm{A}\), calcular \(V=ZI\) en forma polar y binómica.
\end{enumerate}

\item En un circuito en serie se tienen
\[
Z_1=4+3i\;\Omega,
\qquad
Z_2=2-5i\;\Omega.
\]
Calcular la impedancia total \(Z=Z_1+Z_2\), su módulo y su argumento. Interpretar el resultado en términos de las componentes real e imaginaria.

\item Una magnitud armónica se representa mediante la amplitud compleja
\[
A=12e^{i\pi/3}.
\]
La amplitud se modifica mediante una transformación que multiplica por
\[
H=0.8e^{-i\pi/6}.
\]
Determinar la nueva amplitud compleja \(HA\), su módulo y su desfase respecto de \(A\). Esta forma de cálculo aparece, por ejemplo, al estudiar la respuesta de sistemas lineales a señales sinusoidales.

\item \M Utilizar Mathematica para comprobar los resultados del ejercicio 13. Emplear únicamente operaciones algebraicas y, cuando sea necesario, \texttt{Re}, \texttt{Im}, \texttt{Abs}, \texttt{Arg} y \texttt{N}.

\item \M Resolver con Mathematica
\[
z^6=64i
\]
mediante \texttt{Solve}. Para cada solución, calcular numéricamente su módulo y su argumento principal y comprobar que los argumentos están uniformemente espaciados.

\item \M Resolver con Mathematica las ecuaciones del ejercicio 19 y comparar las soluciones obtenidas con las calculadas a mano. La comprobación con ordenador debe realizarse después de haber planteado matemáticamente la forma polar de las raíces.

\end{enumerate}

\cleardoublepage
\chapter{Funciones y continuidad}
\label{ch:funciones}

En ingeniería es habitual describir una magnitud a partir de otra: la longitud de una barra cambia con la temperatura, la concentración de un reactivo evoluciona con el tiempo, la tensión de una red varía periódicamente y la potencia producida por una máquina depende de sus condiciones de funcionamiento. El lenguaje matemático que permite expresar estas dependencias es el de las funciones. En este capítulo estudiaremos las funciones reales de variable real, sus principales familias, el concepto de límite y la continuidad. Estos conceptos constituyen, además, la base sobre la que se construyen el cálculo diferencial y el cálculo integral.

\section{Funciones reales de variable real}
\label{sec:funciones-reales}

\subsection{Concepto de función, dominio, recorrido y gráfica}

Una variable es una cantidad que puede tomar distintos valores. Cuando los valores de una variable están determinados por los de otra, aparece de forma natural el concepto de función.

\begin{example}
El área de una sección circular depende de su radio. Si el radio es $r>0$, entonces
\[
A(r)=\pi r^2.
\]
En este modelo, $r$ es la variable independiente y $A$ la variable dependiente.
\end{example}

\begin{definition}[Función real de variable real]
Sean $D\subseteq\R$ y $f$ una regla que asigna a cada $x\in D$ un único número real $f(x)$. Diremos que
\[
f:D\longrightarrow\R
\]
es una \emph{función real de variable real}. El conjunto $D$ se denomina \emph{dominio} de $f$ y se denota $\operatorname{Dom}(f)$.
\end{definition}

El conjunto de valores que realmente toma la función se denomina \emph{recorrido} o \emph{imagen}:
\[
\operatorname{Im}(f)=\{f(x):x\in\operatorname{Dom}(f)\}.
\]
La notación
\[
y=f(x)
\]
indica que $x$ es la variable independiente y que $y$ depende de $x$.

Una función puede describirse mediante una fórmula, una tabla, una gráfica o una regla verbal. En aplicaciones, estas representaciones se complementan. Una tabla de medidas experimentales puede sugerir una gráfica y, a partir de ella, un modelo algebraico aproximado.

\begin{example}[Curva de calibración]
Un sensor de temperatura proporciona, dentro de su intervalo de trabajo, una tensión aproximadamente lineal
\[
V(T)=0.010T+0.50,
\]
donde $T$ se mide en grados Celsius y $V$ en voltios. La fórmula permite calcular la salida del sensor para cualquier temperatura admisible. Por ejemplo,
\[
V(80)=1.30\ \mathrm{V}.
\]
Si el fabricante especifica que el modelo es válido para $-20\leq T\leq150$, el dominio físico del modelo es $[-20,150]$, aunque la expresión algebraica $0.010T+0.50$ esté definida para todo $T\in\R$.
\end{example}

Esta distinción entre \emph{dominio matemático} y \emph{dominio de validez de un modelo} será importante a lo largo del curso.

\begin{definition}[Gráfica de una función]
La gráfica de una función $f:D\to\R$ es el conjunto
\[
G_f=\{(x,f(x)):x\in D\}\subseteq\R^2.
\]
\end{definition}

Una curva del plano representa una función de $x$ si toda recta vertical corta a la curva, como máximo, en un punto. Esta observación se conoce como \emph{criterio de la recta vertical}.

En muchas ocasiones la fórmula de una función determina su dominio de forma implícita. Cuando no se indique otra cosa, entenderemos por dominio el mayor subconjunto de $\R$ para el que la expresión tiene sentido.

Las restricciones más frecuentes son:
\begin{itemize}
\item un denominador no puede anularse;
\item el radicando de una raíz de índice par debe ser no negativo;
\item el argumento de un logaritmo debe ser estrictamente positivo.
\end{itemize}

\begin{example}
Determinar el dominio de
\[
f(x)=\frac{\sqrt{x+1}}{x-2}.
\]
Para que la raíz esté definida debe cumplirse $x+1\geq0$, es decir, $x\geq-1$. Además, el denominador exige $x\neq2$. Por tanto,
\[
\operatorname{Dom}(f)=[-1,2[\,\cup\,]2,+\infty[.
\]
\end{example}

\begin{solvedexercise}
Determinar el dominio de
\[
g(x)=\ln\!\left(\frac{x-1}{x+2}\right).
\]
\tcblower
\textbf{Resolución.}\par\smallskip
El argumento del logaritmo debe ser positivo:
\[
\frac{x-1}{x+2}>0.
\]
Los puntos que pueden cambiar el signo son $x=-2$ y $x=1$. El cociente es positivo en
\[
]-\infty,-2[\,\cup\,]1,+\infty[.
\]
Por tanto,
\[
\operatorname{Dom}(g)=]-\infty,-2[\,\cup\,]1,+\infty[.
\]
\end{solvedexercise}

Además del dominio y del recorrido, algunas propiedades globales proporcionan información cualitativa muy útil.

\begin{definition}[Paridad]
Sea $f:D\to\R$ y supongamos que $x\in D\Rightarrow -x\in D$.
\begin{itemize}
\item $f$ es \emph{par} si $f(-x)=f(x)$ para todo $x\in D$.
\item $f$ es \emph{impar} si $f(-x)=-f(x)$ para todo $x\in D$.
\end{itemize}
\end{definition}

La gráfica de una función par es simétrica respecto del eje vertical, mientras que la de una función impar es simétrica respecto del origen.

\begin{definition}[Periodicidad]
Una función $f$ es \emph{periódica} si existe $T>0$ tal que
\[
f(x+T)=f(x)
\]
siempre que $x$ y $x+T$ pertenezcan al dominio. Un valor $T$ con esta propiedad se denomina período. Cuando existe un período positivo mínimo, éste se denomina \emph{período fundamental}.
\end{definition}

La periodicidad aparece de forma natural en fenómenos oscilatorios. Por ejemplo, una tensión sinusoidal ideal puede escribirse como
\[
v(t)=V_0\sin(\omega t+\varphi),
\]
y tiene período
\[
T=\frac{2\pi}{\omega}.
\]

\begin{definition}[Monotonía]
Sea $I$ un intervalo. Una función $f$ es \emph{creciente} en $I$ si
\[
x_1<x_2\Longrightarrow f(x_1)\leq f(x_2),
\]
y \emph{estrictamente creciente} si la desigualdad es estricta. Análogamente, $f$ es \emph{decreciente} en $I$ si
\[
x_1<x_2\Longrightarrow f(x_1)\geq f(x_2).
\]
\end{definition}

\begin{definition}[Acotación]
Sea $A\subseteq\operatorname{Dom}(f)$.
\begin{itemize}
\item $f$ está acotada superiormente en $A$ si existe $M\in\R$ tal que $f(x)\leq M$ para todo $x\in A$.
\item $f$ está acotada inferiormente en $A$ si existe $m\in\R$ tal que $f(x)\geq m$ para todo $x\in A$.
\item $f$ está acotada en $A$ si está acotada superior e inferiormente.
\end{itemize}
\end{definition}

Por ejemplo, $\sin x$ y $\cos x$ están acotadas entre $-1$ y $1$, una propiedad que tiene interpretación inmediata cuando representan magnitudes sinusoidales normalizadas.

\subsection{Funciones elementales}

Las funciones que aparecerán con mayor frecuencia pueden construirse a partir de unas pocas familias elementales.

\paragraph{Funciones polinómicas y racionales.}
Una función polinómica tiene la forma
\[
p(x)=a_nx^n+a_{n-1}x^{n-1}+\cdots+a_1x+a_0,
\qquad a_n\neq0,
\]
y está definida para todo $x\in\R$. Una función racional es un cociente
\[
r(x)=\frac{p(x)}{q(x)},
\]
donde $p$ y $q$ son polinomios y
\[
\operatorname{Dom}(r)=\{x\in\R:q(x)\neq0\}.
\]

Las funciones lineales y cuadráticas son casos particulares de funciones polinómicas. Una relación lineal
\[
y=mx+b
\]
representa una dependencia con tasa de variación constante. Una función cuadrática
\[
y=ax^2+bx+c,
\qquad a\neq0,
\]
tiene como gráfica una parábola.

\begin{example}[Deformación térmica]
En un intervalo moderado de temperaturas, la longitud de una barra puede modelarse por
\[
L(T)=L_0\bigl(1+\alpha(T-T_0)\bigr).
\]
Como función de $T$, se trata de una función lineal. El parámetro $\alpha$ mide cuánto cambia proporcionalmente la longitud por unidad de temperatura. El modelo no pretende describir temperaturas arbitrarias: su dominio físico es el intervalo en el que la aproximación lineal sea válida.
\end{example}

\paragraph{Funciones potenciales y radicales.}
Una función potencial tiene la forma
\[
f(x)=x^r,
\]
con exponente real $r$, en el dominio en el que la potencia esté definida como número real. Son especialmente frecuentes las potencias enteras y las raíces.

Por ejemplo,
\[
f(x)=\sqrt{x}
\]
tiene dominio $[0,+\infty[$ y recorrido $[0,+\infty[$, mientras que
\[
g(x)=\frac{1}{x}
\]
tiene dominio $\R\setminus\{0\}$.

\paragraph{Funciones exponenciales y logarítmicas.}
Sea $a>0$, $a\neq1$. La función exponencial de base $a$ es
\[
f(x)=a^x.
\]
Su dominio es $\R$, su recorrido es $]0,+\infty[$ y nunca se anula. Si $a>1$ es estrictamente creciente y si $0<a<1$ es estrictamente decreciente.

La base exponencial más importante es
\[
e=2.718281828\ldots,
\]
y escribiremos
\[
\exp(x)=e^x.
\]

La función logarítmica de base $a$ se define como la inversa de la exponencial:
\[
\log_a(x)=y\Longleftrightarrow a^y=x.
\]
Su dominio es $]0,+\infty[$ y su recorrido es $\R$. Para la base $e$ utilizaremos el logaritmo natural
\[
\ln x=\log_e x.
\]
Se verifican, para $x,y>0$,
\[
\ln(xy)=\ln x+\ln y,
\qquad
\ln\left(\frac{x}{y}\right)=\ln x-\ln y,
\qquad
\ln(x^r)=r\ln x.
\]

\begin{example}[Decaimiento de una concentración]
En un modelo de cinética de primer orden, la concentración de un reactivo puede escribirse como
\[
C(t)=C_0e^{-kt},\qquad t\geq0,
\]
donde $C_0>0$ es la concentración inicial y $k>0$ una constante característica del proceso. La función es positiva y decreciente, y satisface
\[
\frac{C(t)}{C_0}=e^{-kt}.
\]
Si se quiere conocer el instante en que la concentración se reduce a la mitad, se resuelve
\[
e^{-kt}=\frac12,
\]
de donde
\[
t=\frac{\ln2}{k}.
\]
El logaritmo aparece, por tanto, de forma natural al despejar una variable situada en un exponente.
\end{example}

\paragraph{Funciones trigonométricas.}
Las funciones seno y coseno tienen dominio $\R$, recorrido $[-1,1]$ y período fundamental $2\pi$:
\[
\sin(x+2\pi)=\sin x,
\qquad
\cos(x+2\pi)=\cos x.
\]
Además,
\[
\sin^2x+\cos^2x=1.
\]
La tangente se define por
\[
\tan x=\frac{\sin x}{\cos x}
\]
y su dominio es
\[
\R\setminus\left\{\frac{\pi}{2}+k\pi:k\in\Z\right\}.
\]

Los argumentos de las funciones trigonométricas se medirán en radianes. Recordemos que
\[
\pi\ \text{radianes}=180^\circ.
\]

\begin{example}[Movimiento armónico]
La posición de una masa sometida a una oscilación ideal puede modelarse por
\[
x(t)=0.04\cos(10t)\quad \text{m}.
\]
La amplitud es $0.04\,\mathrm{m}$ y el período de la oscilación es
\[
T=\frac{2\pi}{10}=\frac{\pi}{5}\ \mathrm{s}.
\]
La propia forma de la función permite leer propiedades físicas sin realizar ningún cálculo diferencial.
\end{example}

Las funciones trigonométricas inversas se obtienen restringiendo adecuadamente los dominios. Utilizaremos
\[
\arcsin:[-1,1]\to\left[-\frac\pi2,\frac\pi2\right],
\]
\[
\arccos:[-1,1]\to[0,\pi],
\qquad
\arctan:\R\to\left]-\frac\pi2,\frac\pi2\right[.
\]

En la Figura~\ref{fig:elementales-t2} se muestran algunas funciones elementales que servirán de referencia durante el curso.

\begin{figure}[H]
\centering
\begin{tikzpicture}
\begin{axis}[
  width=.82\textwidth,height=6.5cm,
  axis lines=middle,
  xmin=-3.2,xmax=3.2,ymin=-2.2,ymax=4.2,
  xlabel={$x$},ylabel={$y$},
  samples=160,
  legend style={at={(0.02,0.98)},anchor=north west,draw=none,fill=none},
  domain=-3:3
]
\addplot[thick,blue] {x};
\addlegendentry{$x$}
\addplot[thick,red] {x^2};
\addlegendentry{$x^2$}
\addplot[thick,green!50!black] {exp(x)};
\addlegendentry{$e^x$}
\addplot[thick,purple,domain=0.06:3] {ln(x)};
\addlegendentry{$\ln x$}
\end{axis}
\end{tikzpicture}
\caption{Algunas funciones elementales.}
\label{fig:elementales-t2}
\end{figure}

\paragraph{Transformaciones de gráficas.}
A partir de la gráfica de una función $f$ pueden obtenerse otras mediante transformaciones sencillas. Si $c>0$:
\[
f(x)+c
\]
desplaza la gráfica $c$ unidades hacia arriba,
\[
f(x-c)
\]
la desplaza $c$ unidades hacia la derecha,
\[
-f(x)
\]
la refleja respecto del eje horizontal y
\[
f(-x)
\]
la refleja respecto del eje vertical.

Estas reglas son útiles para reconocer gráficas sin construir largas tablas de valores.

\subsection{Composición y función inversa}

Si $f$ y $g$ son funciones reales, podemos definir, en los puntos donde tenga sentido,
\[
(f+g)(x)=f(x)+g(x),\qquad
(fg)(x)=f(x)g(x),
\]
\[
\left(\frac fg\right)(x)=\frac{f(x)}{g(x)},\qquad g(x)\neq0.
\]
El dominio de la nueva función se obtiene intersectando los dominios de las funciones que intervienen y añadiendo las restricciones que correspondan.

La composición formaliza dependencias encadenadas.

\begin{definition}[Composición]
Sean $g:D_g\to\R$ y $f:D_f\to\R$. La función compuesta $f\circ g$ se define por
\[
(f\circ g)(x)=f(g(x))
\]
para aquellos $x\in D_g$ tales que $g(x)\in D_f$.
\end{definition}

\begin{example}[Cadena de conversión]
Un transmisor proporciona una tensión $V$ en función de una presión $p$ mediante
\[
V(p)=0.02p+1,
\]
y un sistema de adquisición convierte esa tensión en una lectura digital
\[
N(V)=1000V.
\]
La lectura en función de la presión es la composición
\[
(N\circ V)(p)=1000(0.02p+1)=20p+1000.
\]
El modelo completo se obtiene componiendo los modelos de cada etapa.
\end{example}

En general,
\[
f\circ g\neq g\circ f.
\]

\begin{definition}[Función inyectiva]
Una función $f:D\to\R$ es \emph{inyectiva} si
\[
f(x_1)=f(x_2)\Longrightarrow x_1=x_2.
\]
Equivalentemente, valores distintos de la variable independiente producen valores distintos de la función.
\end{definition}

Si $f$ es inyectiva, puede definirse una función inversa sobre su imagen.

\begin{definition}[Función inversa]
Sea $f:D\to\operatorname{Im}(f)$ una función inyectiva. Su función inversa
\[
f^{-1}:\operatorname{Im}(f)\to D
\]
se caracteriza por
\[
f^{-1}(f(x))=x,
\qquad
f(f^{-1}(y))=y.
\]
\end{definition}

No debe confundirse $f^{-1}(x)$ con $1/f(x)$. Las gráficas de $f$ y $f^{-1}$ son simétricas respecto de la recta $y=x$.

\begin{proposition}
Toda función estrictamente monótona en un intervalo es inyectiva en ese intervalo y, por tanto, admite función inversa sobre su imagen.
\end{proposition}

Por ejemplo, $e^x$ y $\ln x$ son funciones inversas entre sí. La función $x^2$ no es inyectiva en $\R$, pero sí lo es al restringirla a $[0,+\infty[$; en ese intervalo su inversa es $\sqrt{x}$.

\paragraph{Comprobación con Mathematica.}
Para representar una función en un intervalo utilizaremos fundamentalmente \texttt{Plot}. Por ejemplo,
\begin{mathematicacode}
Plot[{Exp[x], Log[x]}, {x, 0.1, 2}]
\end{mathematicacode}
La representación gráfica es una herramienta de exploración y comprobación; no sustituye al análisis del dominio ni a la justificación de las propiedades de la función.

\section{Límites}
\label{sec:limites}

El concepto de límite permite describir el comportamiento de una función cuando la variable independiente se aproxima a un punto, aunque la función no esté definida en ese punto. Esta idea hará posible definir en el capítulo siguiente la velocidad instantánea, la pendiente de una curva y, en general, la derivada.

\subsection{Límite en un punto y límites laterales}

De forma intuitiva, decir que $f(x)$ se aproxima a $L$ cuando $x$ se aproxima a $a$ significa que podemos hacer $f(x)$ tan próxima a $L$ como queramos tomando $x$ suficientemente próxima a $a$, sin exigir $x=a$.

\begin{definition}[Límite de una función]
Sea $f:D\subseteq\R\to\R$ y sea $a$ un punto al que pueden aproximarse puntos de $D$ distintos de $a$. Diremos que
\[
\lim_{x\to a}f(x)=L
\]
si para todo $\varepsilon>0$ existe $\delta>0$ tal que
\[
x\in D,\quad 0<|x-a|<\delta
\quad\Longrightarrow\quad
|f(x)-L|<\varepsilon.
\]
\end{definition}

La definición no exige que $f(a)$ exista. El límite describe lo que sucede \emph{alrededor} de $a$, no necesariamente en $a$.

\begin{example}
Sea
\[
f(x)=\frac{x^2-1}{x-1},\qquad x\neq1.
\]
Como
\[
f(x)=x+1\qquad (x\neq1),
\]
cuando $x$ se aproxima a $1$, $f(x)$ se aproxima a $2$. Por tanto,
\[
\lim_{x\to1}\frac{x^2-1}{x-1}=2,
\]
aunque la expresión original no esté definida en $x=1$.
\end{example}

En la recta real podemos aproximarnos a un punto desde valores menores o mayores.

\begin{definition}[Límites laterales]
Diremos que
\[
\lim_{x\to a^-}f(x)=L
\]
si $f(x)$ se aproxima a $L$ cuando $x$ se aproxima a $a$ con $x<a$. Análogamente,
\[
\lim_{x\to a^+}f(x)=L
\]
se refiere a la aproximación con $x>a$.
\end{definition}

\begin{theorem}
Existe $\displaystyle\lim_{x\to a}f(x)=L$ si y sólo si existen los dos límites laterales y
\[
\lim_{x\to a^-}f(x)=\lim_{x\to a^+}f(x)=L.
\]
\end{theorem}

\begin{example}
Consideremos
\[
f(x)=
\begin{cases}
2x+1, & x<1,\\
x^2+2, & x\geq1.
\end{cases}
\]
Entonces
\[
\lim_{x\to1^-}f(x)=3,
\qquad
\lim_{x\to1^+}f(x)=3.
\]
Por tanto,
\[
\lim_{x\to1}f(x)=3.
\]
En este caso, además, $f(1)=3$; más adelante veremos que la función es continua en $x=1$.
\end{example}

Los límites se comportan bien respecto de las operaciones algebraicas.

\begin{proposition}[Álgebra de límites]
Si
\[
\lim_{x\to a}f(x)=A,
\qquad
\lim_{x\to a}g(x)=B,
\]
entonces, siempre que las expresiones estén definidas,
\begin{align*}
\lim_{x\to a}(f(x)\pm g(x))&=A\pm B,\\
\lim_{x\to a}f(x)g(x)&=AB,\\
\lim_{x\to a}\frac{f(x)}{g(x)}&=\frac AB,\qquad B\neq0.
\end{align*}
Además, para las potencias y composiciones de funciones continuas pueden aplicarse las reglas habituales de sustitución.
\end{proposition}

En particular, si una función elemental es continua en $a$, su límite en $a$ se obtiene sustituyendo $x=a$.

\subsection{Límites infinitos, límites en el infinito y asíntotas}

En ocasiones los valores de una función crecen sin cota cuando $x$ se aproxima a un punto. Utilizaremos entonces la notación de límite infinito.

\begin{definition}[Límite infinito]
Escribiremos
\[
\lim_{x\to a}f(x)=+\infty
\]
si para todo $M>0$ existe $\delta>0$ tal que
\[
x\in D,\quad 0<|x-a|<\delta\Longrightarrow f(x)>M.
\]
La definición de límite $-\infty$ es análoga.
\end{definition}

El símbolo $+\infty$ no representa un número real. La notación anterior expresa un comportamiento de la función.

Si al menos uno de los límites laterales en $a$ es infinito, la recta
\[
x=a
\]
se denomina \emph{asíntota vertical} de la gráfica.

\begin{example}
Para
\[
f(x)=\frac{1}{x-2}
\]
se tiene
\[
\lim_{x\to2^-}f(x)=-\infty,
\qquad
\lim_{x\to2^+}f(x)=+\infty.
\]
Por tanto, $x=2$ es una asíntota vertical.
\end{example}

También podemos estudiar el comportamiento para valores de $x$ arbitrariamente grandes en valor absoluto.

\begin{definition}[Límite en el infinito]
Supongamos que $D$ no está acotado superiormente. Diremos que
\[
\lim_{x\to+\infty}f(x)=L
\]
si para todo $\varepsilon>0$ existe $M>0$ tal que
\[
x\in D,\quad x>M\Longrightarrow |f(x)-L|<\varepsilon.
\]
Análogamente se define el límite cuando $x\to-\infty$ si $D$ no está acotado inferiormente.
\end{definition}

Si
\[
\lim_{x\to+\infty}f(x)=L
\quad\text{o}\quad
\lim_{x\to-\infty}f(x)=L,
\]
la recta $y=L$ es una \emph{asíntota horizontal} por el lado correspondiente.

\begin{example}[Respuesta estacionaria]
Consideremos
\[
T(t)=T_a+(T_0-T_a)e^{-kt},\qquad k>0,
\]
como modelo idealizado de la temperatura de un cuerpo que intercambia calor con un ambiente a temperatura $T_a$. Puesto que
\[
\lim_{t\to+\infty}e^{-kt}=0,
\]
se obtiene
\[
\lim_{t\to+\infty}T(t)=T_a.
\]
La recta $T=T_a$ es una asíntota horizontal del modelo y expresa matemáticamente que la temperatura se aproxima al equilibrio con el ambiente.
\end{example}

Una asíntota puede ser también oblicua.

\begin{definition}[Asíntota oblicua]
La recta $y=mx+n$, con $m\neq0$, es una asíntota oblicua de $f$ cuando $x\to+\infty$ si
\[
\lim_{x\to+\infty}\bigl(f(x)-(mx+n)\bigr)=0.
\]
Si existen los límites correspondientes, sus coeficientes pueden obtenerse mediante
\[
m=\lim_{x\to+\infty}\frac{f(x)}x,
\qquad
n=\lim_{x\to+\infty}\bigl(f(x)-mx\bigr).
\]
Se procede análogamente cuando $x\to-\infty$.
\end{definition}

\begin{example}
Sea
\[
f(x)=\frac{x^2-2x+4}{x-2}.
\]
Dividiendo polinomios,
\[
f(x)=x+\frac{4}{x-2}.
\]
Como
\[
\lim_{x\to\pm\infty}\frac{4}{x-2}=0,
\]
la recta
\[
y=x
\]
es una asíntota oblicua tanto cuando $x\to+\infty$ como cuando $x\to-\infty$. Además, $x=2$ es una asíntota vertical.
\end{example}

\subsection{Cálculo de límites e indeterminaciones}

Cuando la sustitución directa produce un valor real y todas las operaciones son válidas, el límite queda determinado inmediatamente. Sin embargo, ciertas expresiones no permiten decidir el resultado por sustitución. Las más habituales son
\[
\frac00,\qquad \frac{\infty}{\infty},\qquad
\infty-\infty,\qquad 0\cdot\infty,
\qquad 1^\infty,
\]
y también pueden aparecer $0^0$ o $\infty^0$. Se denominan \emph{formas indeterminadas}: dos expresiones distintas con la misma forma pueden tener límites diferentes.

No existe una única técnica para resolverlas. La transformación adecuada depende de la estructura de la expresión.

\paragraph{Factorización y simplificación.}
Cuando aparece $0/0$ en cocientes algebraicos, puede ser posible factorizar y cancelar un factor que sólo se anula en el punto al que tiende $x$.

\begin{example}
\[
\lim_{x\to1}\frac{x^3-x^2+x-1}{x-1}
=\lim_{x\to1}\frac{(x-1)(x^2+1)}{x-1}
=\lim_{x\to1}(x^2+1)=2.
\]
La cancelación es válida en el cálculo del límite porque se trabaja con $x\neq1$ suficientemente próximo a $1$.
\end{example}

\paragraph{Racionalización.}
Cuando intervienen raíces, multiplicar por la expresión conjugada puede eliminar la indeterminación.

\begin{example}
\[
\lim_{x\to0}\frac{\sqrt{1+x}-1}{x}
=\lim_{x\to0}
\frac{1}{\sqrt{1+x}+1}
=\frac12.
\]
\end{example}

\paragraph{Comparación de órdenes de crecimiento.}
Para $x\to+\infty$ resulta útil recordar que, para $a>1$ y $m,n\in\N$,
\[
a^x\gg x^n\gg (\ln x)^m,
\]
es decir, las exponenciales crecientes dominan a las potencias y éstas dominan a los logaritmos. Entre polinomios, el término de mayor grado determina el comportamiento en el infinito.

Por tanto, si
\[
r(x)=\frac{a_nx^n+\cdots+a_0}{b_mx^m+\cdots+b_0},
\qquad a_n b_m\neq0,
\]
entonces, cuando $x\to\pm\infty$, basta comparar $n$ y $m$ para determinar el comportamiento dominante. Si $n=m$, el límite es $a_n/b_m$; si $n<m$, el límite es $0$; y si $n>m$, el cociente no tiene límite finito en general.

\begin{example}
\[
\lim_{x\to+\infty}\frac{5x^2+3x-1}{-2x^2+7}
=-\frac52.
\]
Dividiendo numerador y denominador entre $x^2$ se obtiene directamente el cociente de los coeficientes principales.
\end{example}

\paragraph{La indeterminación $1^\infty$.}
El número $e$ puede caracterizarse mediante
\[
e=\lim_{n\to\infty}\left(1+\frac1n\right)^n.
\]
Más generalmente, si
\[
f(x)\to1,
\qquad
g(x)\to\pm\infty,
\]
y existe el límite indicado, entonces
\[
\lim_{x\to a}f(x)^{g(x)}
=
\exp\left(\lim_{x\to a}g(x)(f(x)-1)\right).
\]

\begin{example}
Calcular
\[
\lim_{x\to+\infty}\left(1+\frac{3}{x}\right)^x.
\]
Tomando
\[
f(x)=1+\frac3x,
\qquad g(x)=x,
\]
se tiene
\[
\lim_{x\to+\infty}g(x)(f(x)-1)
=\lim_{x\to+\infty}x\frac3x=3,
\]
y, por tanto,
\[
\lim_{x\to+\infty}\left(1+\frac3x\right)^x=e^3.
\]
\end{example}

Una herramienta especialmente útil en este tema es la equivalencia de infinitésimos.

\begin{definition}[Infinitésimo]
Una función $f$ es un infinitésimo cuando $x\to a$ si
\[
\lim_{x\to a}f(x)=0.
\]
\end{definition}

\begin{definition}[Infinitésimos equivalentes]
Sean $f$ y $g$ infinitésimos cuando $x\to a$. Diremos que son equivalentes, y escribiremos
\[
f(x)\sim g(x)\qquad(x\to a),
\]
si
\[
\lim_{x\to a}\frac{f(x)}{g(x)}=1.
\]
\end{definition}

Cuando $u\to0$, las equivalencias fundamentales que utilizaremos son
\[
\sin u\sim u,
\qquad
\tan u\sim u,
\qquad
\ln(1+u)\sim u,
\]
\[
e^u-1\sim u,
\qquad
1-\cos u\sim\frac{u^2}{2}.
\]
Más generalmente, para $a>0$,
\[
a^u-1\sim u\ln a.
\]

Estas equivalencias pueden sustituirse dentro de productos y cocientes en el cálculo de límites, pero no de forma indiscriminada dentro de sumas y diferencias.

\begin{solvedexercise}
Calcular
\[
\lim_{x\to0}\frac{\sin(3x)\,\ln(1+2x)}{x^2}.
\]
\tcblower
\textbf{Resolución.}\par\smallskip
Cuando $x\to0$,
\[
\sin(3x)\sim3x,
\qquad
\ln(1+2x)\sim2x.
\]
Por tanto,
\[
\frac{\sin(3x)\ln(1+2x)}{x^2}
\sim
\frac{(3x)(2x)}{x^2}=6,
\]
y
\[
\boxed{\displaystyle
\lim_{x\to0}\frac{\sin(3x)\ln(1+2x)}{x^2}=6.}
\]
\end{solvedexercise}

\begin{example}[Pequeñas oscilaciones]
En mecánica y electricidad aparecen aproximaciones de pequeño ángulo. Si $\theta$ se mide en radianes y $|\theta|$ es pequeño,
\[
\sin\theta\approx\theta.
\]
La justificación matemática local es precisamente
\[
\frac{\sin\theta}{\theta}\longrightarrow1
\qquad(\theta\to0).
\]
La aproximación no afirma que las dos funciones sean iguales; afirma que su error relativo se hace pequeño cerca de cero.
\end{example}

\paragraph{Comprobación con Mathematica.}
El comando básico para calcular límites será \texttt{Limit}. Por ejemplo,
\begin{mathematicacode}
Limit[Sin[3 x] Log[1 + 2 x]/x^2, x -> 0]
\end{mathematicacode}
y para límites en el infinito se utiliza \texttt{Infinity}:
\begin{mathematicacode}
Limit[(5 x^2 + 3 x - 1)/(-2 x^2 + 7), x -> Infinity]
\end{mathematicacode}
El cálculo con ordenador debe servir para verificar el resultado y explorar el comportamiento de la función, no para sustituir la identificación de la indeterminación y la elección de una técnica adecuada.

\section{Continuidad}
\label{sec:continuidad}

La continuidad formaliza una idea geométrica sencilla: una pequeña variación de la variable independiente produce una pequeña variación de la función. En modelos físicos, la continuidad suele ser una hipótesis razonable mientras no se atraviesen cambios de régimen, conmutaciones, impactos o transiciones de fase que introduzcan saltos u otras singularidades.

\subsection{Continuidad}

\begin{definition}[Continuidad]
Sea $f:D\subseteq\R\to\R$ y $a\in D$. Diremos que $f$ es \emph{continua en $a$} si
\[
\lim_{x\to a}f(x)=f(a).
\]
Equivalentemente, para todo $\varepsilon>0$ existe $\delta>0$ tal que
\[
x\in D,\quad |x-a|<\delta
\Longrightarrow
|f(x)-f(a)|<\varepsilon.
\]
\end{definition}

Para comprobar la continuidad en un punto interior del dominio pueden separarse tres condiciones:
\begin{enumerate}
\item $f(a)$ está definida;
\item existe $\displaystyle\lim_{x\to a}f(x)$;
\item $\displaystyle\lim_{x\to a}f(x)=f(a)$.
\end{enumerate}

\begin{definition}[Continuidad en un intervalo]
Una función es continua en un intervalo abierto si es continua en cada uno de sus puntos. En un intervalo cerrado $[a,b]$ se exige continuidad en $]a,b[$, continuidad por la derecha en $a$ y continuidad por la izquierda en $b$.
\end{definition}

Las funciones polinómicas son continuas en $\R$. Las racionales son continuas donde el denominador no se anula. Las funciones exponenciales, trigonométricas y las funciones potenciales son continuas en sus dominios; los logaritmos son continuos en $]0,+\infty[$.

\begin{proposition}[Operaciones con funciones continuas]
Si $f$ y $g$ son continuas en $a$, entonces también lo son
\[
f+g,\qquad f-g,\qquad fg,
\]
y
\[
\frac fg
\]
si $g(a)\neq0$. Además, si $g$ es continua en $a$ y $f$ es continua en $g(a)$, entonces $f\circ g$ es continua en $a$.
\end{proposition}

\begin{solvedexercise}
Determinar el valor de $k$ para que
\[
f(x)=
\begin{cases}
x^2, & x\leq1,\\
3x+k, & x>1
\end{cases}
\]
sea continua en $\R$.
\tcblower
\textbf{Resolución.}\par\smallskip
Cada expresión es continua en su intervalo. El único punto que debe estudiarse es $x=1$. Tenemos
\[
f(1)=1,
\qquad
\lim_{x\to1^-}f(x)=1,
\]
y
\[
\lim_{x\to1^+}f(x)=3+k.
\]
La continuidad exige
\[
3+k=1,
\]
de donde
\[
\boxed{k=-2.}
\]
\end{solvedexercise}

Las discontinuidades pueden clasificarse atendiendo al comportamiento de los límites.

\begin{itemize}
\item Una discontinuidad es \emph{evitable} si existe un límite finito en $a$ pero $f(a)$ no está definida o no coincide con ese límite. Puede eliminarse redefiniendo $f(a)$.
\item Hay una discontinuidad de \emph{salto} si existen límites laterales finitos pero son distintos.
\item Hay una discontinuidad \emph{infinita} si al menos uno de los límites laterales es infinito.
\end{itemize}

\begin{example}
La función
\[
f(x)=\frac{x^2-4}{x-2},\qquad x\neq2,
\]
presenta una discontinuidad evitable en $x=2$, pues
\[
\lim_{x\to2}f(x)=\lim_{x\to2}(x+2)=4.
\]
Definiendo $f(2)=4$ se obtiene una función continua en todo $\R$.
\end{example}

\subsection{Teoremas fundamentales de las funciones continuas}

La continuidad permite obtener conclusiones globales a partir de información local o de los valores en los extremos de un intervalo.

\begin{theorem}[Teorema del valor intermedio]
Sea $f:[a,b]\to\R$ continua. Si $k$ está comprendido entre $f(a)$ y $f(b)$, entonces existe $c\in[a,b]$ tal que
\[
f(c)=k.
\]
\end{theorem}

Una consecuencia fundamental es el Teorema de Bolzano.

\begin{theorem}[Teorema de Bolzano]
Sea $f:[a,b]\to\R$ continua y supongamos que
\[
f(a)f(b)<0.
\]
Entonces existe al menos un $c\in]a,b[$ tal que
\[
f(c)=0.
\]
\end{theorem}

Bolzano garantiza la \emph{existencia} de una raíz, pero no dice que sea única ni proporciona su valor exacto. Este hecho será la base de algunos métodos numéricos de resolución de ecuaciones.

\begin{example}[Punto de operación]
Consideremos la ecuación
\[
x^3+2x-5=0.
\]
Sea
\[
f(x)=x^3+2x-5.
\]
Como $f$ es un polinomio, es continua. Además,
\[
f(1)=-2<0,
\qquad
f(2)=7>0.
\]
Por el Teorema de Bolzano existe al menos una solución en $]1,2[$. Antes de calcularla numéricamente ya sabemos que el punto de operación descrito por la ecuación existe en ese intervalo.
\end{example}

\begin{theorem}[Teorema de Weierstrass]
Sea $f:[a,b]\to\R$ continua. Entonces $f$ está acotada en $[a,b]$ y alcanza un máximo absoluto y un mínimo absoluto en dicho intervalo.
\end{theorem}

Este resultado es importante en ingeniería porque muchas variables de diseño están restringidas a intervalos cerrados. Si una magnitud objetivo depende continuamente de una variable de diseño, sabemos que dentro del intervalo admisible existen valores donde se alcanzan sus extremos. En el capítulo de cálculo diferencial aprenderemos a localizarlos.

\begin{theorem}[Conservación del signo]
Sea $f$ continua en $a$ y supongamos que $f(a)\neq0$. Entonces existe $\delta>0$ tal que $f(x)$ tiene el mismo signo que $f(a)$ para todo $x$ del dominio con
\[
|x-a|<\delta.
\]
\end{theorem}

Esta propiedad permite estudiar desigualdades mediante los ceros y los puntos de discontinuidad de una función. Entre dos puntos consecutivos de este tipo, una función continua que no se anula conserva su signo.

\begin{example}
Resolver
\[
\frac{x-1}{x-5}\geq0.
\]
Los puntos relevantes son $x=1$, donde el numerador se anula, y $x=5$, donde la función no está definida. Estos puntos dividen la recta en
\[
]-\infty,1[,\qquad ]1,5[,\qquad ]5,+\infty[.
\]
La función es continua y no se anula en cada uno de esos intervalos, por lo que conserva el signo. Probando un punto de cada intervalo se obtiene
\[
\frac{x-1}{x-5}>0\quad\text{en}\quad]-\infty,1[\cup]5,+\infty[.
\]
Como en $x=1$ la función vale $0$, la solución es
\[
]-\infty,1]\cup]5,+\infty[.
\]
\end{example}

\section*{Ejercicios propuestos}
\label{sec:ejercicios-t2}

Los ejercicios siguientes recogen el nivel de cálculo y razonamiento que se utilizará en el curso. Algunos proceden, con ajustes de notación, de las relaciones de ejercicios empleadas en cursos anteriores; otros completan contenidos que allí aparecían con menor frecuencia. Los marcados con \M son especialmente adecuados para ser resueltos o comprobados también con \textit{Wolfram Mathematica}.

\begin{enumerate}[label=\textbf{\arabic*.},leftmargin=*]

\item Determinar el dominio de las funciones:
\begin{enumerate}[label=\alph*)]
\item $f(x)=\dfrac{x+1}{x^2-4}$.
\item $g(x)=\sqrt{x^2-2}$.
\item $h(x)=\dfrac{x^2-4}{x+2}$.
\item $j(x)=\ln(x^2+x-6)$.
\item $k(x)=e^{1/(x-3)}$.
\item $m(x)=\dfrac{\sqrt{x+1}}{e^x-1}$.
\item $s(x)=\dfrac{x^3-4x}{x^2+1}$.
\item $r(x)=\sqrt{\dfrac{x-2}{x+1}}$.
\end{enumerate}

\item Para cada función, estudiar si es par, impar o ninguna de las dos:
\[
f(x)=x^4-3x^2+1,
\qquad
g(x)=x^3-4x,
\]
\[
h(x)=e^x,
\qquad
r(x)=\frac{x}{1+x^2}.
\]

\item Determinar el período fundamental, cuando exista, de
\[
f(x)=\sin(3x),
\qquad
g(x)=\cos\left(\frac{x}{2}\right),
\qquad
h(x)=2\sin(4x-1).
\]

\item Una señal eléctrica viene dada por
\[
v(t)=230\sqrt2\sin(100\pi t+\pi/6)\ \mathrm{V}.
\]
Determinar su amplitud, su período y su frecuencia. Calcular $v(0)$ y $v(0.01)$.

\item Una concentración evoluciona según
\[
C(t)=1.8e^{-0.12t}\ \mathrm{mol/L},\qquad t\geq0.
\]
\begin{enumerate}[label=\alph*)]
\item Determinar $C(5)$.
\item Calcular el instante en que $C(t)=0.9$.
\item Razonar, sin derivar, si la función es creciente o decreciente.
\item Indicar su recorrido físico para $t\geq0$.
\end{enumerate}

\item Sean
\[
f(x)=\sqrt{x+1},\qquad g(x)=x^2-4.
\]
Determinar $f\circ g$, $g\circ f$ y sus dominios. Comprobar que las composiciones no coinciden.

\item Determinar, cuando sea posible, la función inversa y especificar dominio y recorrido:
\begin{enumerate}[label=\alph*)]
\item $f(x)=3x-5$.
\item $f(x)=e^{2x}$.
\item $f(x)=x^2$, restringida a $[0,+\infty[$.
\item $f(x)=\dfrac{x-1}{x+2}$.
\end{enumerate}

\item A partir de la gráfica de $y=f(x)$, describir geométricamente las transformaciones que producen
\[
f(x)+3,
\qquad f(x-2),
\qquad -f(x),
\qquad f(-x),
\qquad 2f(x).
\]
Aplicarlas después a $f(x)=x^2$ y esbozar las gráficas resultantes.

\item Calcular, si existen, los siguientes límites:
\begin{enumerate}[label=\alph*)]
\item $\displaystyle\lim_{x\to\pi/2}\sin x$.
\item $\displaystyle\lim_{x\to0}\frac{\sqrt{1+x}-1}{x}$.
\item $\displaystyle\lim_{x\to2}\frac{x^2-4}{x-2}$.
\item $\displaystyle\lim_{x\to-2}\frac{x^3+8}{x+2}$.
\item $\displaystyle\lim_{x\to1}\frac{x^4-1}{x-1}$.
\item $\displaystyle\lim_{x\to3}\frac{\sqrt{x+1}-2}{x-3}$.
\end{enumerate}

\item Calcular los límites laterales y decidir si existe el límite bilateral:
\begin{enumerate}[label=\alph*)]
\item $\displaystyle\lim_{x\to2}\frac{1}{x-2}$.
\item $\displaystyle\lim_{x\to0}\frac{|x|}{x}$.
\item $\displaystyle\lim_{x\to1}f(x)$, donde
\[
f(x)=\begin{cases}
x^2+1,&x<1,\\3x-1,&x\geq1.
\end{cases}
\]
\item $\displaystyle\lim_{x\to0}g(x)$, donde
\[
g(x)=\begin{cases}
x^2-2,&x<0,\\1,&x=0,\\3x,&x>0.
\end{cases}
\]
\end{enumerate}

\item Calcular:
\begin{enumerate}[label=\alph*)]
\item $\displaystyle\lim_{x\to+\infty}\frac{5x+3x^2-4}{-3x^4+10x+6}$.
\item $\displaystyle\lim_{x\to+\infty}\frac{x^3-16x}{9x^3+x+19}$.
\item $\displaystyle\lim_{x\to+\infty}\frac{x^6+x^3}{9x^5+7x}$.
\item $\displaystyle\lim_{x\to-\infty}(9-5x^3+7x^2)$.
\item $\displaystyle\lim_{x\to+\infty}e^{-x}$.
\item $\displaystyle\lim_{x\to-\infty}7^x$.
\end{enumerate}

\item Determinar todas las asíntotas verticales, horizontales y oblicuas de:
\begin{enumerate}[label=\alph*)]
\item $f(x)=\dfrac{x-1}{x+1}$.
\item $g(x)=\dfrac{x^2-2x+4}{x-2}$.
\item $h(x)=\dfrac{2x^2+1}{x^2-9}$.
\item $r(x)=\dfrac{x^3}{x^2+1}$.
\end{enumerate}

\item Calcular:
\begin{enumerate}[label=\alph*)]
\item $\displaystyle\lim_{x\to+\infty}\left(1+\frac{2}{x}\right)^x$.
\item $\displaystyle\lim_{x\to0}(1+3x)^{1/x}$.
\item $\displaystyle\lim_{x\to1}\left(\frac{x^2+x+1}{x+2}\right)^{1/(x-1)}$.
\item $\displaystyle\lim_{x\to0^+}(1+x)^{1/x^2}$.
\end{enumerate}

\item Utilizando equivalencias de infinitésimos, calcular:
\begin{enumerate}[label=\alph*)]
\item $\displaystyle\lim_{x\to0}\frac{\tan x\,\sin x}{x^2}$.
\item $\displaystyle\lim_{x\to0}\frac{\sin x(1-\cos x)}{\ln^3(1+x)}$.
\item $\displaystyle\lim_{x\to0}\frac{\ln(1+5x)}{2x}$.
\item $\displaystyle\lim_{x\to0}\frac{\sin^2(x/3)}{x^2}$.
\item $\displaystyle\lim_{x\to0}\frac{e^{5x}-1}{10x}$.
\item $\displaystyle\lim_{x\to0}\frac{1-\cos(4x)}{\sin^2(2x)}$.
\end{enumerate}

\item Calcular mediante un cambio de variable adecuado:
\begin{enumerate}[label=\alph*)]
\item $\displaystyle\lim_{x\to0}\frac{\sqrt{1+x}-1}{\sqrt[3]{1+x}-1}$.
\item $\displaystyle\lim_{x\to1}\frac{\sqrt{x}-1}{x-1}$.
\item $\displaystyle\lim_{x\to64}\frac{\sqrt[3]{x}-4}{\sqrt[6]{x}-2}$.
\item $\displaystyle\lim_{x\to1}\frac{\sqrt[3]{x}-1}{\sqrt[4]{x}-1}$.
\end{enumerate}

\item Sean
\[
f(x)=4-x^2,
\qquad
g(x)=\sqrt{x+1}.
\]
Calcular, cuando tengan sentido,
\[
\lim_{x\to4}f(x),
\qquad
\lim_{x\to1/2}g(x),
\qquad
\lim_{x\to4}g(f(x)),
\qquad
\lim_{x\to4}f(g(x)).
\]
Explicar cualquier dificultad relacionada con los dominios.

\item Estudiar la continuidad en todo su dominio de:
\begin{enumerate}[label=\alph*)]
\item $f(x)=\dfrac{1}{x^2-4}$.
\item $g(x)=\dfrac{x^2-4x}{x+2}$.
\item $h(x)=\ln(x^2-1)$.
\item $r(x)=\sqrt{4-x^2}$.
\end{enumerate}

\item Estudiar la continuidad de
\[
f(x)=
\begin{cases}
x-1,&x<1,\\
1,&x=1,\\
\ln x,&1<x<e^3,\\
e^x,&x\geq e^3.
\end{cases}
\]
Clasificar las discontinuidades que aparezcan.

\item Estudiar la continuidad de
\[
f(x)=
\begin{cases}
x+1,&-2\leq x\leq0,\\
1,&0<x\leq1,\\
x^2-2,&1<x\leq3.
\end{cases}
\]

\item Determinar $m\in\R$ para que
\[
f(x)=
\begin{cases}
\dfrac{x+1}{x^2-x-2},&x\neq-1,\\
m,&x=-1
\end{cases}
\]
sea continua en $x=-1$, si ello es posible.

\item Determinar $a\in\R$ para que
\[
f(x)=
\begin{cases}
\dfrac{x^2-a^2}{x-a},&x\neq a,\\
a^2,&x=a
\end{cases}
\]
sea continua en $x=a$.

\item La función
\[
f(x)=\frac{1}{x-2}
\]
no se anula nunca, aunque $f(1)$ y $f(3)$ tienen signos distintos. Explicar por qué este hecho no contradice el Teorema de Bolzano aplicado al intervalo $[1,3]$.

\item Utilizar el Teorema de Bolzano para demostrar que cada ecuación tiene al menos una solución en el intervalo indicado:
\begin{enumerate}[label=\alph*)]
\item $x^3+2x-1=0$ en $[0,1]$.
\item $\cos^2x-2\sin(x/4)=0$ en $[0,2]$.
\item $e^{-x}=x$ en $[0,1]$.
\end{enumerate}

\item Una máquina presenta una potencia útil aproximada por
\[
P(n)=\frac{120n}{n+500}\ \mathrm{kW},\qquad n\geq0,
\]
donde $n$ representa una variable proporcional a la velocidad de giro dentro del régimen de validez del modelo.
\begin{enumerate}[label=\alph*)]
\item Determinar $P(500)$ y $P(1500)$.
\item Calcular $\displaystyle\lim_{n\to+\infty}P(n)$ e interpretarlo dentro del modelo.
\item ¿Qué asíntota horizontal presenta la función?
\item Explicar por qué el límite matemático no implica que sea físicamente razonable hacer $n\to\infty$ en una máquina real.
\end{enumerate}

\item En un proceso térmico se modela la diferencia de temperatura con el ambiente mediante
\[
\Delta T(t)=80e^{-0.08t},\qquad t\geq0.
\]
\begin{enumerate}[label=\alph*)]
\item Calcular $\displaystyle\lim_{t\to+\infty}\Delta T(t)$.
\item Determinar el tiempo necesario para que $\Delta T(t)<5$.
\item Interpretar el resultado del límite y la desigualdad.
\end{enumerate}

\item Una característica idealizada de un dispositivo está dada por
\[
I(V)=
\begin{cases}
0,&V<0.7,\\
0.02(V-0.7),&V\geq0.7,
\end{cases}
\]
donde $I$ se mide en amperios y $V$ en voltios. Estudiar la continuidad en $V=0.7$ e interpretar qué significa matemáticamente que el modelo sea continuo en el umbral.

\item Para $n,m\in\N$, calcular
\[
\lim_{x\to1}\frac{1-x^n}{1-x^m}.
\]
Justificar el resultado sin utilizar la regla de L'Hôpital.

\item Calcular
\[
\lim_{x\to0^+}(1+x)^{1/x^2}
\]
y decidir si el límite es finito o infinito. Justificar cada transformación realizada.

\item Estudiar la continuidad en su dominio de
\[
f(x)=
\begin{cases}
\dfrac{x^3}{x^2-4},&-2<x\leq0,\\[2mm]
\dfrac{\ln x}{\ln(\sin x)},&x>0.
\end{cases}
\]
Prestar especial atención a las restricciones de dominio antes de estudiar los puntos de unión.

\item \M Representar con Mathematica, en intervalos adecuados, las funciones
\[
x^2,\qquad e^x,\qquad \ln x,\qquad \sin x,\qquad \frac1{x-2}.
\]
Antes de ejecutar \texttt{Plot}, indicar dominio, simetrías, periodicidad y asíntotas que puedan deducirse analíticamente. Comparar después con las gráficas.

\item \M Comprobar con \texttt{Limit} los resultados de los ejercicios 9, 11 y 14. Para cada límite que produzca una forma indeterminada por sustitución directa, anotar primero la forma y el procedimiento utilizado a mano.

\item \M Para
\[
f(x)=\frac{x^2-2x+4}{x-2},
\]
utilizar \texttt{Plot} y \texttt{Limit} para comprobar las asíntotas obtenidas analíticamente en el ejercicio 12. La respuesta debe incluir las ecuaciones de las asíntotas y no sólo la salida del programa.

\item \M Considerar
\[
f(x)=x^3+2x-5.
\]
Representar la función en $[1,2]$ y comprobar gráficamente la existencia de una raíz. Utilizar después \texttt{Solve} para obtener las soluciones de la ecuación y \texttt{N} para aproximar numéricamente la raíz real. Explicar qué parte de la conclusión proporciona el Teorema de Bolzano y qué parte proporciona el cálculo computacional.

\end{enumerate}

\cleardoublepage
\chapter{Cálculo diferencial}
\label{ch:diferencial}

En el capítulo anterior estudiamos cómo describir una función y cómo analizar su comportamiento mediante límites y continuidad. El cálculo diferencial añade una pregunta esencial: \emph{¿cómo está cambiando una magnitud?} La derivada permite responderla localmente y, a partir de esa información, deducir propiedades globales de una función.

Esta idea admite varios registros que deben manejarse de forma coordinada. Analíticamente, la derivada se obtiene como un límite y se representa mediante una nueva función $f'$. Gráficamente, $f'(x)$ es la pendiente de la recta tangente a la gráfica de $f$ en el punto de abscisa $x$. Verbalmente, expresa la tasa instantánea de cambio de la magnitud representada por $f$. A lo largo del capítulo se trabajará de forma sistemática el paso entre estos tres registros. En ingeniería, esta traducción es tan importante como el propio cálculo simbólico: no basta con obtener una derivada; hay que saber qué información proporciona sobre el sistema que se está estudiando.

\section{Derivada y diferenciabilidad}
\label{sec:derivada-diferenciabilidad}

\subsection{Tasa de variación y pendiente de la recta tangente}

Sea $y=f(x)$ una magnitud que depende de otra magnitud $x$. Si $x$ cambia desde $x_0$ hasta $x_0+h$, el incremento de la variable independiente es
\[
\Delta x=h,
\]
y el correspondiente incremento de la función es
\[
\Delta y=f(x_0+h)-f(x_0).
\]

\begin{definition}[Tasa de variación media]
La \emph{tasa de variación media} de $f$ entre $x_0$ y $x_0+h$, con $h\neq0$, es
\[
\frac{\Delta y}{\Delta x}
=
\frac{f(x_0+h)-f(x_0)}{h}.
\]
\end{definition}

Desde el punto de vista geométrico, este cociente es la pendiente de la recta secante que pasa por
\[
P=(x_0,f(x_0)),
\qquad
Q=(x_0+h,f(x_0+h)).
\]
Si hacemos que $Q$ se aproxime a $P$, es decir, $h\to0$, la secante tiende, cuando existe el límite correspondiente, a la recta tangente.

\begin{figure}[htbp]
\centering
\begin{tikzpicture}
\begin{axis}[
    width=.76\textwidth,
    height=7cm,
    axis lines=middle,
    xmin=-.2,xmax=4.3,
    ymin=-.5,ymax=10.5,
    xlabel={$x$},ylabel={$y$},
    samples=150,
    xtick={1,3},
    xticklabels={$x_0$,$x_0+h$},
    ytick=\empty,
    clip=false
]
\addplot[thick,domain=0:4] {0.8*x^2+0.5};
\addplot[dashed,domain=.25:4.05] {3.2*x-1.9};
\addplot[domain=.25:3.75] {1.6*x-0.3};
\addplot[only marks,mark=*] coordinates {(1,1.3) (3,7.7)};
\node[anchor=north east] at (axis cs:1,1.3) {$P$};
\node[anchor=south west] at (axis cs:3,7.7) {$Q$};
\node[anchor=south west] at (axis cs:3.45,8.34) {secante};
\node[anchor=north west] at (axis cs:2.55,3.78) {tangente en $P$};
\end{axis}
\end{tikzpicture}
\caption{La pendiente de las secantes se aproxima a la pendiente de la tangente cuando $Q$ se acerca a $P$.}
\label{fig:secante-tangente}
\end{figure}

Esta misma construcción aparece al estudiar movimiento. Si $s(t)$ indica la posición de un móvil en el instante $t$, entonces
\[
\frac{s(t_0+h)-s(t_0)}{h}
\]
es la velocidad media durante el intervalo de tiempo $[t_0,t_0+h]$. Cuando el intervalo se hace arbitrariamente pequeño, obtenemos la velocidad instantánea.

\begin{definition}[Derivada en un punto]
Sea $f:D\subseteq\R\to\R$ y sea $x_0$ un punto interior de $D$. La \emph{derivada de $f$ en $x_0$} es, si existe y es finita, el límite
\[
f'(x_0)=\lim_{h\to0}\frac{f(x_0+h)-f(x_0)}{h}.
\]
Cuando este límite existe, se dice que $f$ es \emph{derivable} o \emph{diferenciable} en $x_0$.
\end{definition}

El cociente
\[
\frac{f(x_0+h)-f(x_0)}{h}
\]
se denomina \emph{cociente incremental}. La derivada es, por tanto, el límite de las tasas de variación medias cuando el incremento de la variable independiente tiende a cero.

\begin{example}[De la velocidad media a la instantánea]
La posición de un carro móvil a lo largo de una guía viene dada por
\[
s(t)=t^2+2t\qquad (s\text{ en metros},\ t\text{ en segundos}).
\]
La velocidad media entre $t=2$ y $t=2+h$ es
\[
\frac{s(2+h)-s(2)}{h}
=
\frac{(2+h)^2+2(2+h)-8}{h}
=
6+h.
\]
Al hacer $h\to0$,
\[
v(2)=s'(2)=6\ \mathrm{m/s}.
\]
La misma cantidad tiene tres lecturas: analíticamente es $s'(2)=6$; gráficamente es la pendiente de la tangente a la gráfica posición--tiempo en $t=2$; verbalmente significa que, en ese instante, la posición aumenta aproximadamente $6$ metros por cada segundo adicional.
\end{example}

\begin{solvedexercise}
Calcular, utilizando exclusivamente la definición, la derivada de
\[
f(x)=x^2
\]
en un punto arbitrario $x_0$.
\tcblower
\textbf{Resolución.}\par\smallskip
Partimos del cociente incremental:
\[
\frac{f(x_0+h)-f(x_0)}{h}
=
\frac{(x_0+h)^2-x_0^2}{h}
=
\frac{2x_0h+h^2}{h}
=2x_0+h.
\]
Por tanto,
\[
f'(x_0)=\lim_{h\to0}(2x_0+h)=2x_0.
\]
Como $x_0$ es arbitrario, podemos escribir
\[
f'(x)=2x.
\]
\end{solvedexercise}

\subsection{Interpretación geométrica, física y unidades}

Si $f$ es derivable en $x_0$, la pendiente de la recta tangente a la gráfica $y=f(x)$ en el punto $(x_0,f(x_0))$ es $f'(x_0)$. Por tanto, su ecuación es
\[
y-f(x_0)=f'(x_0)(x-x_0).
\]

\begin{definition}[Recta tangente]
Si $f$ es derivable en $x_0$, la recta tangente a la gráfica de $f$ en $(x_0,f(x_0))$ es
\[
\boxed{\ y=f(x_0)+f'(x_0)(x-x_0).\ }
\]
\end{definition}

\begin{example}
Para
\[
f(x)=x^2+1,
\]
tenemos $f'(x)=2x$. En $x_0=2$,
\[
f(2)=5,
\qquad
f'(2)=4.
\]
La recta tangente es
\[
y-5=4(x-2),
\qquad\text{es decir,}\qquad
 y=4x-3.
\]
El signo positivo de la pendiente indica, además, que la función está creciendo alrededor de $x=2$.
\end{example}

La unidad de la derivada se obtiene dividiendo la unidad de la magnitud dependiente entre la unidad de la variable independiente. Esta observación es esencial para interpretar resultados en problemas aplicados.

\begin{example}[Velocidad de calentamiento]
La temperatura de una pieza durante un proceso térmico se modela por
\[
T(t)=20+80\bigl(1-e^{-0.05t}\bigr),
\]
donde $T$ se mide en grados Celsius y $t$ en minutos. Su derivada es
\[
T'(t)=4e^{-0.05t}.
\]
Por tanto, $T'(10)$ se mide en $^{\circ}\mathrm{C}/\mathrm{min}$ y expresa la velocidad instantánea con la que cambia la temperatura a los diez minutos. Como $T'(t)>0$, la temperatura aumenta; como el valor de $T'(t)$ disminuye con el tiempo, el calentamiento se hace progresivamente más lento.
\end{example}

Cuando $f(x_0)\neq0$, también puede resultar útil comparar la tasa de cambio con el tamaño de la propia magnitud.

\begin{definition}[Tasa relativa de cambio]
La \emph{tasa relativa de cambio} de $f$ en $x_0$ es
\[
\frac{f'(x_0)}{f(x_0)}.
\]
Multiplicada por $100$, se interpreta como porcentaje instantáneo de cambio por unidad de la variable independiente.
\end{definition}

Por ejemplo, si la potencia generada por una instalación es $P(t)$, el cociente $P'(t)/P(t)$ indica la variación instantánea de la potencia en relación con el nivel de producción existente en ese instante.

\subsection{Función derivada, notación y diferenciabilidad}

Si $f$ es derivable en todos los puntos de un intervalo, podemos considerar una nueva función que a cada punto le asigna la pendiente de la tangente correspondiente.

\begin{definition}[Función derivada]
La \emph{función derivada} de $f$ es
\[
f':x\longmapsto f'(x).
\]
\end{definition}

Utilizaremos indistintamente las notaciones
\[
f'(x),
\qquad
y',
\qquad
\frac{dy}{dx},
\qquad
\frac{df}{dx}.
\]
Cuando la variable independiente es el tiempo también es habitual, especialmente en mecánica, la notación de Newton
\[
\dot y(t)=\frac{dy}{dt}.
\]

La derivabilidad es una propiedad más fuerte que la continuidad.

\begin{theorem}[Derivabilidad implica continuidad]
Si $f$ es derivable en $x_0$, entonces $f$ es continua en $x_0$.
\end{theorem}

\begin{proof}
Para $x=x_0+h$,
\[
f(x_0+h)-f(x_0)
=
\frac{f(x_0+h)-f(x_0)}{h}\,h.
\]
Si $f$ es derivable en $x_0$, el primer factor tiende a $f'(x_0)$ y el segundo a $0$. Por tanto,
\[
\lim_{h\to0}\bigl(f(x_0+h)-f(x_0)\bigr)=0,
\]
de donde
\[
\lim_{x\to x_0}f(x)=f(x_0).
\]
\end{proof}

El recíproco es falso: una función puede ser continua y no derivable. El ejemplo más sencillo es $f(x)=|x|$ en $x=0$. Sus derivadas laterales son
\[
f'_-(0)=-1,
\qquad
f'_+(0)=1,
\]
y, al no coincidir, $f'(0)$ no existe. Gráficamente aparece un punto anguloso.

También pueden aparecer tangentes verticales. En ese caso el cociente incremental puede tender a $+\infty$ o $-\infty$, pero, con la definición anterior, no diremos que existe una derivada real finita.

\begin{solvedexercise}
Estudiar la derivabilidad de
\[
f(x)=
\begin{cases}
ax+1, & x\leq1,\\
x^2+b, & x>1,
\end{cases}
\]
y determinar $a$ y $b$ para que sea derivable en todo $\R$.
\tcblower
\textbf{Resolución.}\par\smallskip
Sólo hay que estudiar $x=1$. Para que la función sea derivable debe ser primero continua:
\[
a+1=1+b,
\]
es decir,
\[
b=a.
\]
Las derivadas laterales son
\[
f'_-(1)=a,
\qquad
f'_+(1)=2.
\]
La derivabilidad exige $a=2$, y por continuidad $b=2$. Por tanto,
\[
\boxed{a=b=2}.
\]
\end{solvedexercise}

\paragraph{Comprobación con Mathematica.}
La función \texttt{D} calcula derivadas simbólicas. Por ejemplo,
\begin{mathematicacode}
D[x^2 Exp[-x], x]
\end{mathematicacode}
Para visualizar simultáneamente una función y su recta tangente en $x=2$ puede utilizarse, una vez calculada la derivada,
\begin{mathematicacode}
f[x_] := x^2 + 1
Plot[{f[x], f[2] + f'[2] (x - 2)}, {x, -1, 4}]
\end{mathematicacode}
El gráfico sirve para comprobar la interpretación geométrica; la pendiente y la ecuación de la tangente deben haberse obtenido y justificado analíticamente.

\section{Cálculo de derivadas}
\label{sec:calculo-derivadas}

La definición permite calcular cualquier derivada siempre que el límite correspondiente pueda resolverse. Sin embargo, repetir ese proceso para cada función sería poco eficiente. A partir de la definición se obtienen reglas que permiten derivar directamente la mayoría de las funciones elementales.

\subsection{Reglas básicas y derivadas elementales}

Sean $f$ y $g$ funciones derivables y $\lambda\in\R$. Se verifican las reglas
\[
(f\pm g)'=f'\pm g',
\]
\[
(\lambda f)'=\lambda f',
\]
\[
(fg)'=f'g+fg',
\]
y, siempre que $g(x)\neq0$,
\[
\left(\frac{f}{g}\right)'
=
\frac{f'g-fg'}{g^2}.
\]

La siguiente tabla recoge las derivadas que utilizaremos con mayor frecuencia:

\begin{center}
\renewcommand{\arraystretch}{1.5}
\begin{tabular}{@{}ll@{}}
\toprule
Función & Derivada \\
\midrule
$c$ & $0$ \\
$x^n$ & $nx^{n-1}$ \\
$e^x$ & $e^x$ \\
$a^x$ & $a^x\ln a$ \\
$\ln x$ & $1/x$ \\
$\log_a x$ & $1/(x\ln a)$ \\
$\sin x$ & $\cos x$ \\
$\cos x$ & $-\sin x$ \\
$\tan x$ & $1/\cos^2x$ \\
$\arcsin x$ & $1/\sqrt{1-x^2}$ \\
$\arccos x$ & $-1/\sqrt{1-x^2}$ \\
$\arctan x$ & $1/(1+x^2)$ \\
\bottomrule
\end{tabular}
\end{center}

Para potencias reales, la regla
\[
\frac{d}{dx}x^r=rx^{r-1}
\]
se aplica en los intervalos en los que la función y la expresión resultante estén definidas.

\begin{solvedexercise}
Derivar
\[
f(x)=\frac{x^2e^x}{1+x^2}.
\]
\tcblower
\textbf{Resolución.}\par\smallskip
Tomamos
\[
u(x)=x^2e^x,
\qquad
v(x)=1+x^2.
\]
Entonces
\[
u'(x)=2xe^x+x^2e^x=e^x(2x+x^2),
\qquad
v'(x)=2x.
\]
Por la regla del cociente,
\[
f'(x)=
\frac{e^x(2x+x^2)(1+x^2)-2x^3e^x}{(1+x^2)^2}.
\]
Sacando factor común $e^x$ y simplificando,
\[
\boxed{
f'(x)=\frac{e^x(x^4+x^2+2x)}{(1+x^2)^2}
}.
\]
La forma no simplificada también es correcta y, en algunos problemas, puede ser más útil para interpretar el origen de cada término.
\end{solvedexercise}

\subsection{Funciones compuestas y regla de la cadena}

Muchas magnitudes dependen de otras a través de varias etapas. Si $u=g(x)$ e $y=f(u)$, entonces
\[
y=(f\circ g)(x)=f(g(x)).
\]

\begin{theorem}[Regla de la cadena]
Si $g$ es derivable en $x$ y $f$ es derivable en $g(x)$, entonces $f\circ g$ es derivable en $x$ y
\[
(f\circ g)'(x)=f'(g(x))g'(x).
\]
En notación de Leibniz,
\[
\frac{dy}{dx}=\frac{dy}{du}\frac{du}{dx}.
\]
\end{theorem}

La regla de la cadena expresa matemáticamente la propagación de una tasa de cambio a través de dependencias sucesivas.

\begin{example}[Variación de una resistencia con la temperatura]
Supongamos que la resistencia eléctrica de un conductor depende de la temperatura mediante
\[
R(T)=R_0\bigl(1+\alpha(T-T_0)\bigr),
\]
y que la temperatura evoluciona con el tiempo según $T=T(t)$. Entonces
\[
\frac{dR}{dt}
=
\frac{dR}{dT}\frac{dT}{dt}
=
R_0\alpha\,T'(t).
\]
La fórmula separa dos efectos: la sensibilidad de la resistencia a la temperatura y la velocidad con la que cambia la temperatura.
\end{example}

\begin{example}
Para
\[
f(x)=\sin\!\bigl((1+x^2)^3\bigr),
\]
aplicamos sucesivamente la regla de la cadena:
\[
f'(x)
=
\cos\!\bigl((1+x^2)^3\bigr)\,3(1+x^2)^2\,2x.
\]
Por tanto,
\[
\boxed{
f'(x)=6x(1+x^2)^2\cos\!\bigl((1+x^2)^3\bigr)
}.
\]
\end{example}

\subsection{Derivación implícita y logarítmica}

No siempre una relación entre variables viene dada explícitamente como $y=f(x)$. Si $x$ e $y$ están ligados por una ecuación
\[
F(x,y)=0,
\]
podemos derivar ambos miembros respecto de $x$, teniendo en cuenta que $y=y(x)$.

\begin{example}[Derivación implícita]
Sea
\[
x^2+xy+y^2=7.
\]
Derivando respecto de $x$,
\[
2x+y+xy'+2yy'=0.
\]
Agrupando los términos que contienen $y'$,
\[
(x+2y)y'=-(2x+y),
\]
y, siempre que $x+2y\neq0$,
\[
\boxed{
y'=-\frac{2x+y}{x+2y}
}.
\]
No ha sido necesario despejar previamente $y$.
\end{example}

\begin{solvedexercise}
Mostrar que las tangentes en el origen a las curvas
\[
5y-2x+y^3-x^2y=0
\]
y
\[
2y+5x+x^4-x^3y^2=0
\]
son perpendiculares.
\tcblower
\textbf{Resolución.}\par\smallskip
Para la primera curva, derivando implícitamente,
\[
5y'-2+3y^2y'-2xy-x^2y'=0.
\]
En $(0,0)$,
\[
5y'(0)-2=0,
\qquad
y'(0)=\frac25.
\]
Para la segunda,
\[
2y'+5+4x^3-3x^2y^2-2x^3yy'=0.
\]
En el origen,
\[
2y'(0)+5=0,
\qquad
y'(0)=-\frac52.
\]
El producto de las pendientes es
\[
\frac25\left(-\frac52\right)=-1,
\]
por lo que las tangentes son perpendiculares.
\end{solvedexercise}

La derivación logarítmica resulta especialmente útil cuando la variable aparece simultáneamente en una base y en un exponente, o cuando una expresión está formada por productos y cocientes de muchas potencias.

Si $y=f(x)>0$, tomamos logaritmos:
\[
\ln y=\ln f(x),
\]
derivamos implícitamente y despejamos $y'$.

\begin{example}[Derivación logarítmica]
Sea
\[
y=x^{\sin x},
\qquad x>0.
\]
Tomando logaritmos,
\[
\ln y=\sin x\,\ln x.
\]
Derivando,
\[
\frac{y'}{y}
=
\cos x\,\ln x+\frac{\sin x}{x}.
\]
Por tanto,
\[
\boxed{
y'=x^{\sin x}\left(\cos x\,\ln x+\frac{\sin x}{x}\right)
}.
\]
\end{example}

\subsection{Derivadas sucesivas}

Si $f'$ es derivable, su derivada se denomina \emph{segunda derivada} de $f$ y se denota
\[
f''(x)=\frac{d^2f}{dx^2}.
\]
De forma sucesiva se definen
\[
f'''(x),
\qquad
f^{(4)}(x),
\qquad
\ldots,
\qquad
f^{(n)}(x)=\frac{d^nf}{dx^n}.
\]

Las derivadas de orden superior describen cambios de las propias tasas de cambio. Si $s(t)$ es la posición de un móvil,
\[
v(t)=s'(t)
\]
es su velocidad y
\[
a(t)=s''(t)
\]
es su aceleración. Por tanto, $s$, $s'$ y $s''$ no son tres objetos aislados: representan distintos niveles de variación de un mismo fenómeno.

\begin{solvedexercise}
La posición de un mecanismo viene dada por
\[
s(t)=2\sin(3t),
\]
con $s$ en centímetros y $t$ en segundos. Calcular velocidad y aceleración, y relacionar sus signos con el movimiento cuando $t=\pi/18$.
\tcblower
\textbf{Resolución.}\par\smallskip
Tenemos
\[
v(t)=s'(t)=6\cos(3t),
\qquad
a(t)=s''(t)=-18\sin(3t).
\]
Para $t=\pi/18$ se tiene $3t=\pi/6$, y por tanto
\[
v\!\left(\frac{\pi}{18}\right)=3\sqrt3>0,
\]
\[
a\!\left(\frac{\pi}{18}\right)=-9<0.
\]
La posición está aumentando en ese instante, pero la velocidad está disminuyendo. El sistema se mueve en el sentido positivo mientras experimenta una aceleración en sentido contrario al movimiento.
\end{solvedexercise}

\section{Teoremas fundamentales y aproximación}
\label{sec:teoremas-aproximacion}

Las reglas anteriores permiten calcular derivadas. Los resultados de esta sección explican qué consecuencias tienen esas derivadas sobre el comportamiento de una función y permiten utilizar información local para obtener conclusiones globales.

\subsection{Teoremas de Rolle y del valor medio}

Comenzamos con una condición necesaria para que una función derivable tenga un extremo interior.

\begin{theorem}[Teorema de Fermat]
Sea $f$ derivable en un punto interior $x_0$ de su dominio. Si $f$ tiene en $x_0$ un máximo o un mínimo relativo, entonces
\[
f'(x_0)=0.
\]
\end{theorem}

El recíproco no es cierto: $f'(x_0)=0$ no garantiza que exista un extremo. Por ejemplo, $f(x)=x^3$ satisface $f'(0)=0$ y, sin embargo, es estrictamente creciente alrededor de $0$.

\begin{theorem}[Teorema de Rolle]
Sea $f:[a,b]\to\R$ continua en $[a,b]$ y derivable en $]a,b[$. Si
\[
f(a)=f(b),
\]
entonces existe al menos un punto $c\in]a,b[$ tal que
\[
f'(c)=0.
\]
\end{theorem}

Geométricamente, si una curva derivable parte y termina a la misma altura, debe tener al menos una tangente horizontal entre ambos extremos.

\begin{theorem}[Teorema del valor medio]
Sea $f:[a,b]\to\R$ continua en $[a,b]$ y derivable en $]a,b[$. Entonces existe al menos un punto $c\in]a,b[$ tal que
\[
f'(c)=\frac{f(b)-f(a)}{b-a}.
\]
\end{theorem}

El teorema afirma que, bajo las hipótesis indicadas, alguna tasa instantánea de cambio coincide con la tasa media de cambio en todo el intervalo.

\begin{example}[Interpretación cinemática]
Un vehículo recorre $180$ km en dos horas y su posición es una función continua del tiempo y derivable durante el trayecto. Su velocidad media es
\[
\frac{180}{2}=90\ \mathrm{km/h}.
\]
El Teorema del Valor Medio garantiza que existe al menos un instante del recorrido en el que la velocidad instantánea es exactamente $90\ \mathrm{km/h}$.
\end{example}

\begin{solvedexercise}
Comprobar que la función
\[
f(x)=x^2-4x+3
\]
satisface las hipótesis del Teorema de Rolle en $[1,3]$ y hallar todos los puntos cuya existencia garantiza el teorema.
\tcblower
\textbf{Resolución.}\par\smallskip
La función es polinómica, luego es continua en $[1,3]$ y derivable en $]1,3[$. Además,
\[
f(1)=0=f(3).
\]
Por Rolle existe al menos $c\in]1,3[$ con $f'(c)=0$. Como
\[
f'(x)=2x-4,
\]
se obtiene
\[
2c-4=0,
\qquad
\boxed{c=2}.
\]
\end{solvedexercise}

\subsection{Regla de L'Hôpital}

La Regla de L'Hôpital permite resolver determinadas indeterminaciones de tipo $0/0$ o $\infty/\infty$ utilizando derivadas.

\begin{theorem}[Regla de L'Hôpital]
Sean $f$ y $g$ derivables en un entorno perforado de $a$, con $g'(x)\neq0$ en dicho entorno. Supongamos que
\[
\lim_{x\to a}f(x)=\lim_{x\to a}g(x)=0
\]
o que $|f(x)|$ y $|g(x)|$ tienden a infinito. Si existe
\[
\lim_{x\to a}\frac{f'(x)}{g'(x)}=L,
\]
con $L\in\R\cup\{\pm\infty\}$, entonces, bajo estas hipótesis,
\[
\lim_{x\to a}\frac{f(x)}{g(x)}=L.
\]
El resultado también se aplica a límites laterales y a límites cuando $x\to\pm\infty$.
\end{theorem}

La regla sólo debe utilizarse después de comprobar que se está ante una de las indeterminaciones permitidas. No es una regla general para derivar numerador y denominador en cualquier cociente.

\begin{example}
Calcular
\[
\lim_{x\to0}\frac{e^{2x}-1}{x}.
\]
La sustitución directa produce $0/0$. Aplicando L'Hôpital,
\[
\lim_{x\to0}\frac{e^{2x}-1}{x}
=
\lim_{x\to0}\frac{2e^{2x}}{1}
=2.
\]
\end{example}

Las indeterminaciones $0\cdot\infty$, $\infty-\infty$ y las indeterminaciones exponenciales deben transformarse previamente en un cociente de tipo $0/0$ o $\infty/\infty$, o bien tratarse tomando logaritmos cuando corresponda.

\begin{solvedexercise}
Calcular
\[
\lim_{x\to0^+}x\ln x.
\]
\tcblower
\textbf{Resolución.}\par\smallskip
La expresión es de tipo $0\cdot(-\infty)$. La escribimos como
\[
x\ln x=\frac{\ln x}{1/x}.
\]
Ahora tenemos una indeterminación $(-\infty)/(+\infty)$. Aplicando L'Hôpital,
\[
\lim_{x\to0^+}\frac{1/x}{-1/x^2}
=
\lim_{x\to0^+}(-x)=0.
\]
Por tanto,
\[
\boxed{\lim_{x\to0^+}x\ln x=0}.
\]
\end{solvedexercise}

\subsection{Aproximación lineal y polinomio de Taylor}

La derivabilidad implica que, a escala suficientemente pequeña, una función se comporta aproximadamente como una función lineal. Si $f$ es derivable en $x_0$, entonces
\[
f(x_0+h)=f(x_0)+f'(x_0)h+o(h),
\qquad h\to0.
\]
Por ello, para $x$ próximo a $x_0$,
\[
\boxed{
f(x)\approx f(x_0)+f'(x_0)(x-x_0)
}.
\]
El miembro derecho es precisamente la ordenada de la recta tangente.

\begin{example}[Estimación de una variación]
El radio de una pieza circular pasa de $10$ cm a $10.1$ cm. Para
\[
A(r)=\pi r^2,
\]
la variación exacta es
\[
\Delta A=A(10.1)-A(10)=2.01\pi\ \mathrm{cm}^2.
\]
La aproximación lineal proporciona
\[
\Delta A\approx A'(10)\,\Delta r
=20\pi\cdot0.1
=2\pi\ \mathrm{cm}^2.
\]
La diferencia es pequeña porque el incremento del radio es pequeño respecto de su valor inicial.
\end{example}

La aproximación puede mejorarse incorporando derivadas de orden superior.

\begin{definition}[Polinomio de Taylor]
Supongamos que $f$ tiene derivadas hasta orden $n$ en un entorno de $x_0$. El \emph{polinomio de Taylor de orden $n$ de $f$ centrado en $x_0$} es
\[
P_n(x)
=
\sum_{k=0}^n\frac{f^{(k)}(x_0)}{k!}(x-x_0)^k.
\]
Es decir,
\begin{align*}
P_n(x)={}&f(x_0)+f'(x_0)(x-x_0)
+\frac{f''(x_0)}{2!}(x-x_0)^2\\
&+\cdots
+\frac{f^{(n)}(x_0)}{n!}(x-x_0)^n.
\end{align*}
\end{definition}

Cuando $x_0=0$ se habla también de \emph{polinomio de Maclaurin}.

\begin{theorem}[Taylor con resto de Lagrange]
Si $f$ tiene $n+1$ derivadas continuas en un intervalo que contiene a $x_0$ y a $x$, existe un punto $\xi$ entre $x_0$ y $x$ tal que
\[
f(x)=P_n(x)+R_n(x),
\]
donde
\[
R_n(x)=\frac{f^{(n+1)}(\xi)}{(n+1)!}(x-x_0)^{n+1}.
\]
\end{theorem}

Este resto permite estimar el error cometido al sustituir una función por su polinomio de Taylor.

\begin{example}
Obtener el polinomio de Maclaurin de orden $3$ de $e^x$. Como todas las derivadas de $e^x$ coinciden con $e^x$,
\[
f^{(k)}(0)=1
\]
para todo $k$. Por tanto,
\[
\boxed{
P_3(x)=1+x+\frac{x^2}{2}+\frac{x^3}{6}
}.
\]
Cerca de $x=0$, este polinomio reproduce no sólo el valor y la pendiente de $e^x$, sino también sus dos siguientes niveles de variación.
\end{example}

\begin{solvedexercise}
Obtener el polinomio de Maclaurin de orden $5$ de $\sin x$ y utilizarlo para aproximar $\sin(0.2)$.
\tcblower
\textbf{Resolución.}\par\smallskip
Las derivadas de $\sin x$ se repiten cíclicamente. En $x=0$,
\begin{align*}
f(0)&=0, & f'(0)&=1, & f''(0)&=0,\\
f'''(0)&=-1, & f^{(4)}(0)&=0, & f^{(5)}(0)&=1.
\end{align*}
Por tanto,
\[
P_5(x)=x-\frac{x^3}{3!}+\frac{x^5}{5!}.
\]
Así,
\[
\sin(0.2)
\approx
0.2-\frac{0.2^3}{6}+\frac{0.2^5}{120}
\approx0.1986693.
\]
\end{solvedexercise}

\paragraph{Comprobación con Mathematica.}
El comando mínimo para obtener desarrollos de Taylor será \texttt{Series}. Por ejemplo,
\begin{mathematicacode}
Series[Sin[x], {x, 0, 5}]
\end{mathematicacode}
La salida incluye un término de orden que indica el resto. El objetivo de la práctica no es sustituir el cálculo de las derivadas, sino comprobar el polinomio obtenido y comparar gráficamente la función con su aproximación en distintos intervalos.

\section{La derivada y la forma de una función}
\label{sec:forma-derivada}

La derivada no debe entenderse únicamente como una expresión que se calcula. Su utilidad principal consiste en traducir información entre distintas representaciones de una función. En esta sección estudiaremos de forma explícita las relaciones
\[
f\quad\longleftrightarrow\quad f'\quad\longleftrightarrow\quad f'',
\]
y aprenderemos a pasar del registro analítico al gráfico y al verbal, y viceversa.

\subsection{Monotonía y signo de la primera derivada}

Recordemos que una función es creciente en un intervalo si valores mayores de $x$ producen valores mayores o iguales de $f(x)$, y decreciente si ocurre lo contrario.

\begin{theorem}[Criterio de monotonía]
Sea $f$ continua en un intervalo $I$ y derivable en su interior.
\begin{itemize}
\item Si $f'(x)>0$ para todo $x$ interior a $I$, entonces $f$ es estrictamente creciente en $I$.
\item Si $f'(x)<0$ para todo $x$ interior a $I$, entonces $f$ es estrictamente decreciente en $I$.
\item Si $f'(x)=0$ para todo $x$ interior a $I$, entonces $f$ es constante en $I$.
\end{itemize}
Además, si $f$ es creciente y derivable en un punto interior $x$, entonces $f'(x)\geq0$; si es decreciente, $f'(x)\leq0$.
\end{theorem}

La información puede leerse en tres registros:

\begin{center}
\renewcommand{\arraystretch}{1.45}
\begin{tabularx}{\textwidth}{@{}>{\bfseries}lXXX@{}}
\toprule
 & Registro analítico & Registro gráfico & Registro verbal \\
\midrule
$f'(x)>0$ & La derivada es positiva & La tangente tiene pendiente positiva & $f$ aumenta al aumentar $x$ \\
$f'(x)<0$ & La derivada es negativa & La tangente tiene pendiente negativa & $f$ disminuye al aumentar $x$ \\
$f'(x)=0$ & Punto estacionario & Tangente horizontal & La variación instantánea es nula \\
$|f'(x)|$ grande & Gran valor absoluto de la derivada & Gráfica localmente muy inclinada & La magnitud cambia rápidamente \\
$|f'(x)|$ pequeño & Derivada próxima a cero & Gráfica localmente poco inclinada & La magnitud cambia lentamente \\
\bottomrule
\end{tabularx}
\end{center}

La fila $f'(x)=0$ requiere atención: una tangente horizontal no implica necesariamente un máximo o un mínimo. Puede tratarse, por ejemplo, de un punto de inflexión horizontal.

\begin{example}[Del registro analítico al gráfico]
Sea
\[
f(x)=x^3-3x.
\]
Entonces
\[
f'(x)=3x^2-3=3(x-1)(x+1).
\]
El signo de $f'$ es
\[
\begin{array}{c|ccccc}
x & ]-\infty,-1[ & -1 & ]-1,1[ & 1 & ]1,+\infty[\\
\hline
f'(x) & + & 0 & - & 0 & +
\end{array}
\]
Por tanto, $f$ crece en $]-\infty,-1[$, decrece en $]-1,1[$ y vuelve a crecer en $]1,+\infty[$. Sin representar todavía la función, sabemos que la gráfica debe subir, alcanzar una tangente horizontal en $x=-1$, bajar hasta otra tangente horizontal en $x=1$ y volver a subir.
\end{example}

\begin{figure}[htbp]
\centering
\begin{minipage}{.48\textwidth}
\centering
\begin{tikzpicture}
\begin{axis}[
 width=\linewidth,height=6.2cm,
 axis lines=middle,xmin=-2.3,xmax=2.3,ymin=-4,ymax=4,
 xlabel={$x$},ylabel={$f(x)$},samples=150,
 xtick={-1,1},ytick=\empty
]
\addplot[thick,domain=-2.05:2.05] {x^3-3*x};
\addplot[only marks,mark=*] coordinates {(-1,2) (1,-2)};
\end{axis}
\end{tikzpicture}
\end{minipage}\hfill
\begin{minipage}{.48\textwidth}
\centering
\begin{tikzpicture}
\begin{axis}[
 width=\linewidth,height=6.2cm,
 axis lines=middle,xmin=-2.3,xmax=2.3,ymin=-4,ymax=9,
 xlabel={$x$},ylabel={$f'(x)$},samples=150,
 xtick={-1,1},ytick={0}
]
\addplot[thick,domain=-2.05:2.05] {3*x^2-3};
\end{axis}
\end{tikzpicture}
\end{minipage}
\caption{Una misma información en dos registros: la gráfica de $f$ crece exactamente donde la gráfica de $f'$ está por encima del eje horizontal.}
\label{fig:f-y-derivada}
\end{figure}

La Figura~\ref{fig:f-y-derivada} debe leerse en ambos sentidos. Si conocemos $f$, la inclinación de su gráfica permite anticipar el signo y el tamaño de $f'$. Si conocemos $f'$, su posición respecto del eje horizontal permite reconstruir los intervalos de crecimiento y decrecimiento de $f$.

\begin{solvedexercise}[De la derivada a la función]
Sea $f$ una función derivable cuya derivada es
\[
f'(x)=6x(x+1)(x-2).
\]
Sin calcular una expresión para $f$, estudiar los intervalos de crecimiento y decrecimiento y clasificar sus extremos relativos.
\tcblower
\textbf{Resolución.}\par\smallskip
Los puntos donde $f'$ se anula son
\[
x=-1,\qquad x=0,\qquad x=2.
\]
Estudiamos el signo del producto:
\[
\begin{array}{c|ccccccc}
x & ]-\infty,-1[ & -1 & ]-1,0[ & 0 & ]0,2[ & 2 & ]2,+\infty[\\
\hline
f'(x) & - & 0 & + & 0 & - & 0 & +
\end{array}
\]
Por tanto,
\begin{itemize}
\item $f$ decrece en $]-\infty,-1[$;
\item $f$ crece en $]-1,0[$;
\item $f$ decrece en $]0,2[$;
\item $f$ crece en $]2,+\infty[$.
\end{itemize}
El cambio de signo $-\to+$ en $x=-1$ indica un mínimo relativo; el cambio $+\to-$ en $x=0$ indica un máximo relativo; y el cambio $-\to+$ en $x=2$ indica otro mínimo relativo. Todo se ha deducido a partir de $f'$ sin necesidad de conocer $f$.
\end{solvedexercise}

\subsection{Extremos relativos y absolutos}

\begin{definition}[Extremos relativos]
Sea $x_0$ un punto del dominio de $f$.
\begin{itemize}
\item $f$ tiene un \emph{máximo relativo} en $x_0$ si existe un entorno de $x_0$ en el que $f(x)\leq f(x_0)$.
\item $f$ tiene un \emph{mínimo relativo} en $x_0$ si existe un entorno de $x_0$ en el que $f(x)\geq f(x_0)$.
\end{itemize}
\end{definition}

\begin{definition}[Punto crítico]
Un punto $x_0$ interior al dominio de $f$ es \emph{crítico} si
\[
f'(x_0)=0
\]
o si $f'(x_0)$ no existe.
\end{definition}

Los extremos relativos pueden aparecer en puntos donde la derivada no existe. Por ejemplo, $f(x)=|x|$ tiene un mínimo relativo y absoluto en $x=0$, aunque no es derivable en ese punto. Por ello, al buscar extremos no basta con resolver $f'(x)=0$; hay que tener en cuenta también los puntos no derivables del dominio.

\begin{theorem}[Criterio de la primera derivada]
Sea $x_0$ un punto crítico de una función continua.
\begin{itemize}
\item Si $f'$ cambia de signo de positivo a negativo al atravesar $x_0$, entonces $f$ tiene un máximo relativo en $x_0$.
\item Si $f'$ cambia de signo de negativo a positivo al atravesar $x_0$, entonces $f$ tiene un mínimo relativo en $x_0$.
\item Si $f'$ no cambia de signo, el punto crítico no es un extremo relativo.
\end{itemize}
\end{theorem}

\begin{definition}[Extremos absolutos]
La función $f$ alcanza un \emph{máximo absoluto} en $x_M$ si
\[
f(x)\leq f(x_M)
\]
para todo $x$ de su dominio. Análogamente, alcanza un \emph{mínimo absoluto} en $x_m$ si
\[
f(x)\geq f(x_m)
\]
para todo $x$ de su dominio.
\end{definition}

\begin{theorem}[Teorema de Weierstrass]
Si $f$ es continua en un intervalo cerrado y acotado $[a,b]$, entonces alcanza en dicho intervalo un máximo absoluto y un mínimo absoluto.
\end{theorem}

Para encontrar los extremos absolutos de una función continua en $[a,b]$ se evalúa la función en los puntos críticos interiores y en los extremos $a$ y $b$, y se comparan los valores obtenidos.

\begin{solvedexercise}
Estudiar los extremos absolutos de
\[
f(x)=\frac{x+1}{x^2+1}
\]
en el intervalo $[-1,1/2]$.
\tcblower
\textbf{Resolución.}\par\smallskip
La función es continua en todo $\R$, luego alcanza máximo y mínimo absolutos en el intervalo. Su derivada es
\[
f'(x)=\frac{x^2+1-2x(x+1)}{(x^2+1)^2}
=
\frac{1-2x-x^2}{(x^2+1)^2}.
\]
Los puntos críticos satisfacen
\[
1-2x-x^2=0,
\]
es decir,
\[
x=-1\pm\sqrt2.
\]
Sólo $-1+\sqrt2$ pertenece a $[-1,1/2]$. Evaluamos
\[
f(-1)=0,
\]
\[
f\left(\frac12\right)=\frac{6}{5},
\]
y
\[
f(-1+\sqrt2)=\frac{1+\sqrt2}{2}.
\]
Como
\[
\frac{1+\sqrt2}{2}>\frac65>0,
\]
el mínimo absoluto es $0$, alcanzado en $x=-1$, y el máximo absoluto es $(1+\sqrt2)/2$, alcanzado en $x=-1+\sqrt2$.
\end{solvedexercise}

\subsection{Segunda derivada, convexidad y puntos de inflexión}

La primera derivada informa sobre si $f$ crece o decrece. La segunda derivada informa sobre cómo cambia la pendiente de $f$.

\begin{definition}[Convexidad y concavidad]
Sea $f$ dos veces derivable en un intervalo $I$.
\begin{itemize}
\item Si $f''(x)\geq0$ en $I$, diremos que $f$ es \emph{convexa} en $I$.
\item Si $f''(x)\leq0$ en $I$, diremos que $f$ es \emph{cóncava} en $I$.
\end{itemize}
\end{definition}

En términos de $f'$, si $f''>0$, las pendientes de las tangentes aumentan; si $f''<0$, disminuyen. Esta relación permite leer la convexidad de $f$ directamente en la gráfica de $f'$.

\begin{center}
\renewcommand{\arraystretch}{1.45}
\begin{tabularx}{\textwidth}{@{}>{\bfseries}lXXX@{}}
\toprule
 & Registro analítico & Registro gráfico & Registro verbal \\
\midrule
$f''(x)>0$ & $f'$ aumenta & Las pendientes se hacen mayores & La tasa de cambio de $f$ aumenta \\
$f''(x)<0$ & $f'$ disminuye & Las pendientes se hacen menores & La tasa de cambio de $f$ disminuye \\
$f''(x)=0$ & Candidato a cambio de convexidad & Posible cambio de curvatura & La variación de la tasa puede cambiar de signo \\
\bottomrule
\end{tabularx}
\end{center}

\begin{definition}[Punto de inflexión]
Un punto $x_0$ es un \emph{punto de inflexión} de $f$ si la función cambia de convexidad al atravesarlo.
\end{definition}

Si $f''$ es continua y $x_0$ es un punto de inflexión, entonces necesariamente $f''(x_0)=0$. Sin embargo, $f''(x_0)=0$ por sí solo no garantiza un punto de inflexión. Por ejemplo, para $f(x)=x^4$ se tiene $f''(0)=0$, pero $f$ es convexa a ambos lados de $0$.

\begin{theorem}[Criterio de la segunda derivada para extremos]
Sea $f$ dos veces derivable y supongamos que $f'(x_0)=0$.
\begin{itemize}
\item Si $f''(x_0)>0$, entonces $f$ tiene un mínimo relativo estricto en $x_0$.
\item Si $f''(x_0)<0$, entonces $f$ tiene un máximo relativo estricto en $x_0$.
\item Si $f''(x_0)=0$, el criterio no decide.
\end{itemize}
\end{theorem}

\begin{example}[Lectura verbal de la segunda derivada]
Supongamos que la temperatura de un horno aumenta durante todo un proceso, pero los operarios observan que \emph{cada minuto aumenta menos que el anterior}. Si $T(t)$ representa la temperatura, la frase se traduce en
\[
T'(t)>0
\]
porque la temperatura aumenta, y
\[
T''(t)<0
\]
porque la tasa de aumento disminuye. Gráficamente, $T$ debe ser creciente y cóncava. A su vez, la gráfica de $T'$ debe estar por encima del eje horizontal y ser decreciente.

Este ejemplo muestra que una descripción verbal suficientemente precisa permite anticipar simultáneamente la forma de $T$ y la de $T'$ sin disponer todavía de una fórmula explícita.
\end{example}

\subsection{Lectura coordinada de $f$, $f'$ y $f''$}

Para analizar correctamente una función conviene evitar una confusión frecuente: el signo de $f'$ y el crecimiento de $f'$ son informaciones distintas. El primero describe la monotonía de $f$; el segundo describe su convexidad.

Por ejemplo:
\begin{itemize}
\item si $f'>0$ y $f''>0$, la función crece y lo hace cada vez más deprisa;
\item si $f'>0$ y $f''<0$, la función crece, pero cada vez más despacio;
\item si $f'<0$ y $f''>0$, la función decrece, pero su pendiente se hace menos negativa;
\item si $f'<0$ y $f''<0$, la función decrece y su pendiente se hace cada vez más negativa.
\end{itemize}

\begin{solvedexercise}[De un registro gráfico a uno analítico y verbal]
La siguiente gráfica representa la derivada $f'$ de una función $f$.

\begin{center}
\begin{tikzpicture}
\begin{axis}[
 width=.72\textwidth,height=6.5cm,
 axis lines=middle,
 xmin=-3.3,xmax=3.3,ymin=-4.5,ymax=5,
 xlabel={$x$},ylabel={$f'(x)$},
 xtick={-2,1},ytick={0},samples=150
]
\addplot[thick,domain=-3:3] {(x+2)*(x-1)};
\addplot[only marks,mark=*] coordinates {(-2,0) (1,0)};
\end{axis}
\end{tikzpicture}
\end{center}

Determinar, sin conocer una expresión de $f$:
\begin{enumerate}[label=\alph*)]
\item los intervalos de crecimiento y decrecimiento de $f$;
\item la naturaleza de los puntos críticos visibles;
\item los intervalos de convexidad y concavidad de $f$;
\item una descripción verbal coherente del comportamiento de $f$ cuando $x$ recorre el gráfico de izquierda a derecha.
\end{enumerate}
\tcblower
\textbf{Resolución.}\par\smallskip
La gráfica de $f'$ corta al eje en $x=-2$ y $x=1$.

Para $x<-2$, $f'(x)>0$, de modo que $f$ crece. Entre $-2$ y $1$, $f'(x)<0$, luego $f$ decrece. Para $x>1$, $f'(x)>0$, y $f$ vuelve a crecer. Por tanto, $x=-2$ es un máximo relativo y $x=1$ un mínimo relativo.

La gráfica de $f'$ es una parábola con vértice en
\[
x=-\frac12.
\]
Como $f'$ decrece para $x<-1/2$, se tiene $f''<0$ y $f$ es cóncava. Como $f'$ crece para $x>-1/2$, se tiene $f''>0$ y $f$ es convexa. Así, $x=-1/2$ es un punto de inflexión de $f$.

Verbalmente: la función empieza aumentando, pero su pendiente va disminuyendo; en $x=-2$ deja de aumentar y comienza a disminuir. Continúa perdiendo pendiente hasta $x=-1/2$. A partir de ahí la pendiente empieza a recuperarse, aunque todavía es negativa hasta $x=1$. En $x=1$ la función alcanza un mínimo relativo y después vuelve a crecer cada vez con mayor pendiente.
\end{solvedexercise}

Este tipo de razonamiento será especialmente importante en los ejercicios en los que se proporciona la gráfica de $f$ o de $f'$ sin una fórmula explícita.

\subsection{Estudio de la gráfica de una función}

Para realizar un estudio cualitativo completo de una función no existe una única receta rígida, pero suele ser útil considerar, en este orden, los siguientes elementos:
\begin{enumerate}
\item dominio, cortes con los ejes, simetrías y periodicidad;
\item límites en los extremos del dominio y asíntotas;
\item continuidad y puntos no derivables;
\item signo de $f'$: monotonía y extremos;
\item signo de $f''$: convexidad, concavidad y puntos de inflexión;
\item información adicional relevante del modelo.
\end{enumerate}

La representación final debe ser coherente con \emph{toda} la información obtenida. El objetivo no es dibujar una curva aproximada al azar, sino sintetizar en un único registro gráfico las conclusiones analíticas anteriores.

\begin{solvedexercise}
Realizar un estudio de forma y esbozar la gráfica de
\[
f(x)=\frac{2(x^2-9)}{x^2-4}.
\]
\tcblower
\textbf{Resolución.}\par\smallskip
El dominio es
\[
\R\setminus\{-2,2\}.
\]
La función es par. Además,
\[
\lim_{x\to\pm\infty}f(x)=2,
\]
por lo que $y=2$ es asíntota horizontal. En $x=\pm2$ existen asíntotas verticales.

La derivada es
\[
f'(x)=\frac{20x}{(x^2-4)^2}.
\]
Como el denominador es positivo en el dominio, el signo depende de $x$: la función decrece para $x<0$ en cada componente del dominio y crece para $x>0$. El único punto estacionario es $x=0$, donde
\[
f(0)=\frac92,
\]
y se obtiene un mínimo relativo en el intervalo $]-2,2[$.

La segunda derivada es
\[
f''(x)=-\frac{20(3x^2+4)}{(x^2-4)^3}.
\]
El numerador es siempre negativo. Por tanto, $f''>0$ en $]-2,2[$ y $f''<0$ en $]-\infty,-2[\cup]2,+\infty[$. No hay puntos de inflexión porque los cambios de convexidad se producen en $x=\pm2$, que no pertenecen al dominio.

Estos datos determinan el esbozo cualitativo: una rama central convexa con mínimo en $(0,9/2)$ y dos ramas exteriores cóncavas que se aproximan a $y=2$ cuando $|x|\to\infty$.
\end{solvedexercise}

\subsection{Problemas de optimización}

Optimizar consiste en determinar el valor de una variable que hace máxima o mínima una magnitud. La dificultad principal suele estar antes de derivar: hay que identificar correctamente la función objetivo, expresar todas las variables en términos de una sola y determinar el dominio físico del problema.

Un procedimiento habitual es:
\begin{enumerate}
\item identificar la magnitud que se quiere maximizar o minimizar;
\item formular las restricciones del problema;
\item expresar la función objetivo en una sola variable;
\item determinar el intervalo físicamente admisible;
\item localizar puntos críticos y comparar, cuando corresponda, con los extremos del intervalo;
\item interpretar la solución en el contexto original y comprobar sus unidades.
\end{enumerate}

\begin{solvedexercise}[Diseño de un depósito]
Se desea fabricar un depósito cilíndrico cerrado de volumen $1\ \mathrm{m}^3$. Determinar el radio y la altura que minimizan la superficie de material empleada.
\tcblower
\textbf{Resolución.}\par\smallskip
Si $r$ es el radio y $h$ la altura,
\[
\pi r^2h=1,
\]
de donde
\[
h=\frac{1}{\pi r^2},
\qquad r>0.
\]
La superficie total es
\[
S=2\pi r^2+2\pi rh.
\]
Sustituyendo la restricción,
\[
S(r)=2\pi r^2+\frac{2}{r}.
\]
Derivamos:
\[
S'(r)=4\pi r-\frac{2}{r^2}.
\]
La condición $S'(r)=0$ da
\[
4\pi r^3=2,
\qquad
r^3=\frac{1}{2\pi},
\]
por lo que
\[
\boxed{r=\left(\frac{1}{2\pi}\right)^{1/3}}.
\]
De la restricción,
\[
h=\frac{1}{\pi r^2}=2r.
\]
Además,
\[
S''(r)=4\pi+\frac{4}{r^3}>0,
\]
por lo que el punto crítico es un mínimo. El depósito que utiliza menos material tiene altura igual al diámetro.
\end{solvedexercise}

\begin{example}[Tendido de un cable]
Una planta eléctrica está situada en la ribera de un río de $1$ km de ancho. En la ribera opuesta, una fábrica se encuentra $2$ km aguas arriba del punto situado frente a la planta. Tender cable por tierra cuesta $3$ unidades monetarias por kilómetro y bajo el agua cuesta $5$.

Sea $x$ la distancia, medida sobre la ribera de la fábrica desde el punto enfrentado a la planta hasta el punto en el que emerge el cable. La longitud submarina es
\[
\sqrt{1+x^2},
\]
y la longitud terrestre restante es $2-x$, con $0\leq x\leq2$. El coste es
\[
C(x)=5\sqrt{1+x^2}+3(2-x).
\]
Entonces
\[
C'(x)=\frac{5x}{\sqrt{1+x^2}}-3.
\]
La condición $C'(x)=0$ conduce a
\[
\frac{5x}{\sqrt{1+x^2}}=3,
\]
y, al elevar al cuadrado,
\[
25x^2=9(1+x^2).
\]
Así,
\[
16x^2=9,
\qquad
x=\frac34,
\]
que pertenece al intervalo físico. Comparando con los extremos, se comprueba que el coste mínimo se obtiene haciendo emerger el cable $0.75$ km aguas arriba del punto enfrentado a la planta.
\end{example}

La solución matemática debe terminar siempre con una interpretación: un punto crítico fuera del dominio físico, una longitud negativa o una temperatura incompatible con el intervalo de operación no constituyen soluciones válidas del problema de ingeniería.

\paragraph{Comprobación con Mathematica.}
Para los estudios de forma no introduciremos comandos nuevos innecesarios. Se utilizarán principalmente \texttt{D}, \texttt{Solve} y \texttt{Plot}. Por ejemplo,
\begin{mathematicacode}
f[x_] := x^3 - 3 x
D[f[x], x]
Solve[D[f[x], x] == 0, x]
Plot[{f[x], D[f[x], x]}, {x, -2, 2}]
\end{mathematicacode}
La práctica debe comparar la información analítica con la representación. Una gráfica obtenida por ordenador no sustituye a la tabla de signos ni justifica por sí sola la existencia de un máximo, un mínimo o un punto de inflexión.

\section*{Ejercicios propuestos}
\label{sec:ejercicios-t3}

Los ejercicios siguientes combinan cálculo, interpretación y traducción entre registros. Algunos proceden, con ajustes de notación o contexto, de las relaciones utilizadas en cursos anteriores; otros completan aspectos que conviene trabajar con mayor intensidad, especialmente la relación entre una función y sus derivadas. Los marcados con \M son especialmente adecuados para ser resueltos o comprobados también con \textit{Wolfram Mathematica}.

\begin{enumerate}[label=\textbf{\arabic*.},leftmargin=*]

\item Calcular mediante la definición la derivada de las funciones siguientes en el punto indicado:
\begin{enumerate}[label=\alph*)]
\item $f(x)=x^2+1$ en $x_0=2$.
\item $g(x)=1/x$ en $x_0=1$.
\item $h(x)=\sqrt{x}$ en $x_0=4$.
\end{enumerate}

\item Para $f(x)=x+4/x$:
\begin{enumerate}[label=\alph*)]
\item calcular la ecuación de la secante por $(1,f(1))$ y $(8,f(8))$;
\item calcular la tangente en $(3,f(3))$;
\item interpretar geométricamente la diferencia entre ambas pendientes.
\end{enumerate}

\item Hallar la ecuación de la recta tangente a $y=1+x^2$ en $(0,1)$ y a $y=1/(2x^2)$ en el punto de abscisa $x=1$.

\item Determinar los puntos de la gráfica de
\[
f(x)=12x-x^3
\]
en los que la recta tangente es horizontal.

\item Encontrar los puntos de la gráfica de
\[
f(x)=x^3-5x+2
\]
en los que la recta tangente es perpendicular a la recta $y=\frac14x+3$.

\item Estudiar la derivabilidad de:
\begin{enumerate}[label=\alph*)]
\item $f(x)=|x^2-1|$;
\item $g(x)=|x^2-3x+2|$;
\item $h(x)=x^2-|x|+2$;
\item $p(x)=\begin{cases}\dfrac{x}{1+e^x},&x<0,\\0,&x=0,\\x^2+1,&x>0.\end{cases}$
\end{enumerate}

\item Determinar $k$ para que
\[
f(x)=
\begin{cases}
2x+5,&x\leq1,\\
x^2+k,&x>1
\end{cases}
\]
sea derivable en todo su dominio.

\item Determinar $a$ y $b$ para que
\[
f(x)=
\begin{cases}
ax^2+3x,&x\leq2,\\
x^2-bx-4,&x>2
\end{cases}
\]
sea derivable en $x=2$.

\item Derivar las funciones:
\begin{enumerate}[label=\alph*)]
\item $f(x)=\sin(3x)\cos(x^4)$;
\item $g(x)=e^{6-x^2}\cos(5x-1)$;
\item $h(x)=100\arctan(7x+8)$;
\item $p(x)=\dfrac{2x-4x^5}{2(x-5)}$;
\item $q(x)=\dfrac{\ln(4x+2)}{x}$;
\item $r(x)=\ln(x^2+1)$;
\item $s(x)=4^{\sin^2x}\ln\!\left(\dfrac{x^2-1}{x}\right)$.
\end{enumerate}

\item Aplicar derivación logarítmica para calcular la derivada de:
\begin{enumerate}[label=\alph*)]
\item $x^{\sin x}$, $x>0$;
\item $(3x)^{\cos x}$, $x>0$;
\item $\dfrac{(x^2+1)^3\sqrt{x-1}}{x^5}$, $x>1$.
\end{enumerate}

\item Mediante derivación implícita, hallar $dy/dx$ en:
\begin{enumerate}[label=\alph*)]
\item $x^2+y^2=25$;
\item $x^2+xy+y^2=7$;
\item $x^3+y^3=6xy$.
\end{enumerate}
Determinar en cada caso las pendientes de las tangentes en los puntos que se indiquen o que puedan deducirse fácilmente de la ecuación.

\item Una partícula se mueve según
\[
s(t)=t^3-6t^2+9t,
\qquad t\geq0.
\]
Calcular $v(t)$ y $a(t)$. Determinar cuándo la partícula está en reposo y describir, en esos instantes, el signo de la aceleración y el cambio de sentido del movimiento.

\item La potencia entregada por un sistema viene dada por
\[
P(t)=120\bigl(1-e^{-0.4t}\bigr),
\qquad t\geq0.
\]
Calcular $P'(t)$ y $P''(t)$ e interpretar verbalmente el signo de ambas derivadas. Describir, antes de representarla, la forma cualitativa de la gráfica de $P$.

\item Comprobar las hipótesis del Teorema de Rolle y hallar los puntos correspondientes para:
\begin{enumerate}[label=\alph*)]
\item $f(x)=x^2-4x+3$ en $[1,3]$;
\item $g(x)=\sin x$ en $[0,\pi]$;
\item $h(x)=x^3-x$ en $[-1,1]$.
\end{enumerate}

\item Aplicar el Teorema del Valor Medio a
\[
f(x)=\sqrt{x}
\]
en $[1,4]$. Interpretar geométricamente el resultado.

\item Un proceso productivo acumula $500$ unidades entre las horas $t=2$ y $t=7$. Bajo las hipótesis adecuadas de continuidad y derivabilidad, explicar qué puede garantizarse sobre la tasa instantánea de producción en algún instante del intervalo. Distinguir claramente qué información proporciona el Teorema del Valor Medio y qué información no proporciona.

\item Calcular mediante la Regla de L'Hôpital:
\begin{enumerate}[label=\alph*)]
\item $\displaystyle\lim_{x\to\infty}x\sin(1/x)$;
\item $\displaystyle\lim_{x\to0^+}x\ln x$;
\item $\displaystyle\lim_{x\to0}\frac{e^{2x}-1}{x}$;
\item $\displaystyle\lim_{x\to0}\frac{x-\sin x}{x^3}$;
\item $\displaystyle\lim_{x\to1}\frac{\sin(\pi/x)}{x-1}$;
\item $\displaystyle\lim_{x\to\infty}\frac{x^2}{e^x}$.
\end{enumerate}

\item Resolver las indeterminaciones exponenciales:
\begin{enumerate}[label=\alph*)]
\item $\displaystyle\lim_{x\to\infty}x^{1/x}$;
\item $\displaystyle\lim_{x\to0}(1+x)^{1/x}$;
\item $\displaystyle\lim_{x\to0}(1+x^2)^{1/x^2}$.
\end{enumerate}
Justificar la transformación previa antes de utilizar L'Hôpital.

\item Hallar los polinomios de Maclaurin de los órdenes indicados:
\begin{enumerate}[label=\alph*)]
\item $e^x$, orden $4$;
\item $\sin x$, orden $5$;
\item $\cos x$, orden $6$;
\item $\ln(1+x)$, orden $4$.
\end{enumerate}

\item Obtener el polinomio de Taylor de orden $3$ de
\[
f(x)=\sqrt{x}
\]
centrado en $x_0=4$ y utilizarlo para aproximar $\sqrt{4.1}$.

\item Utilizar la aproximación lineal para estimar:
\begin{enumerate}[label=\alph*)]
\item $\sqrt{100.4}$;
\item $1/1.02$;
\item el incremento del volumen de una esfera cuando su radio pasa de $20$ cm a $20.1$ cm.
\end{enumerate}
Comparar, cuando sea posible, con el valor exacto.

\item \M Para $f(x)=e^x$, utilizar \texttt{Series} para obtener los polinomios de Maclaurin de órdenes $1$, $2$, $3$ y $5$. Representar con \texttt{Plot} la función y los cuatro polinomios en $[-1,1]$ y después en $[-3,3]$. Explicar por qué la calidad de la aproximación depende del intervalo.

\item Estudiar la monotonía y los extremos relativos de
\[
f(x)=\frac13x^3-x^2-3x+4
\]
en todo $\R$. Repetir después el estudio de los extremos absolutos en $[0,1]$, $[-2,4]$ y $[0,4]$.

\item Hallar los extremos relativos y absolutos de
\[
f(x)=x^2-|x|+2.
\]
Prestar atención al punto en el que la derivada no existe.

\item Estudiar los extremos absolutos de
\[
f(x)=3+2xe^{-0.1x}
\]
en $[0,40]$.

\item Determinar $a$ y $b$ para que
\[
f(x)=ax^2+bx
\]
alcance un máximo relativo de valor $8$ en $x=3$.

\item Sea
\[
f'(x)=6x(x+1)(x-2).
\]
Sin obtener $f(x)$:
\begin{enumerate}[label=\alph*)]
\item estudiar la monotonía de $f$;
\item determinar sus extremos relativos;
\item estudiar la convexidad de $f$ utilizando $f''$;
\item localizar los puntos de inflexión que puedan deducirse;
\item realizar un esbozo cualitativo de una posible función $f$, sabiendo además que $f(0)=1$.
\end{enumerate}

\item Una función derivable $f$ satisface la siguiente descripción: es decreciente hasta $x=-1$, creciente entre $x=-1$ y $x=2$, y vuelve a decrecer para $x>2$. Además, la pendiente aumenta hasta $x=0$ y disminuye a partir de $x=0$.
\begin{enumerate}[label=\alph*)]
\item traducir la información a condiciones sobre los signos de $f'$ y $f''$;
\item indicar los extremos relativos y los posibles puntos de inflexión;
\item dibujar un esbozo cualitativo de $f$;
\item dibujar un esbozo compatible de $f'$.
\end{enumerate}

\item La temperatura de un reactor aumenta desde el inicio del proceso. Durante los primeros minutos, la velocidad de calentamiento también aumenta; después de cierto instante, la temperatura sigue aumentando, pero cada vez más lentamente.
\begin{enumerate}[label=\alph*)]
\item describir verbalmente el cambio de signo de $T''$;
\item representar cualitativamente una posible gráfica de $T$;
\item representar una posible gráfica de $T'$;
\item identificar qué característica del proceso correspondería a un punto de inflexión de $T$.
\end{enumerate}

\item La siguiente función representa la derivada de una función desconocida:
\[
g(x)=f'(x)=x^2-2x-3.
\]
Analizar $f$ en términos de monotonía, extremos y convexidad sin integrar $g$. Explicar qué datos adicionales serían necesarios para determinar una única función $f$.

\item \M Elegir una función polinómica de grado $4$ con al menos dos puntos críticos. Calcular $f'$ y $f''$ analíticamente y con \texttt{D}. Representar en gráficos separados $f$, $f'$ y $f''$. Elaborar una tabla que relacione ceros y signos de las derivadas con crecimiento, extremos y convexidad de $f$.

\item Dada una gráfica de una función $f$, un estudiante afirma: ``como la gráfica está por encima del eje $OX$, su derivada es positiva''. Explicar por qué la afirmación es falsa y proporcionar dos contraejemplos: uno gráfico y otro analítico.

\item Dada una gráfica de $f'$ que permanece por debajo del eje $OX$ pero es estrictamente creciente, describir con precisión qué puede afirmarse sobre $f$ y sobre $f''$. Ilustrar la respuesta con una función concreta.

\item Realizar un estudio de forma y un esbozo de las gráficas:
\begin{enumerate}[label=\alph*)]
\item $f(x)=|x^2-9|$;
\item $g(x)=\dfrac{2(x^2-9)}{x^2-4}$;
\item $h(x)=xe^{-x}$.
\end{enumerate}
El estudio debe incluir dominio, continuidad, asíntotas cuando existan, monotonía, extremos, convexidad y puntos de inflexión.

\item \M Para cada función del ejercicio anterior, utilizar \texttt{Plot} sólo después de completar el estudio analítico. Señalar cualquier rasgo de la gráfica que pudiera interpretarse incorrectamente si se observara únicamente una ventana de representación inadecuada.

\item Se desea construir una caja sin tapa, de base cuadrada y volumen $50\ \mathrm{cm}^3$. Determinar sus dimensiones para minimizar la cantidad de material empleado.

\item Un envase cilíndrico cerrado debe tener superficie total $80\ \mathrm{cm}^2$. Determinar radio y altura para que su volumen sea máximo.

\item Una pista consta de un rectángulo y dos semicírculos en sus extremos. Si su perímetro total es $400$ m, determinar las dimensiones que maximizan el área rectangular interior.

\item Un alambre de $10$ m se corta en dos trozos. Con uno se forma un cuadrado y con el otro un triángulo equilátero. Determinar cómo debe realizarse el corte para que la suma de las áreas sea mínima. Analizar también qué ocurre con el máximo dentro del intervalo físicamente admisible.

\item En el diseño de una conducción, el coste por metro bajo tierra es distinto del coste por metro en superficie. Formular y resolver un problema de optimización análogo al tendido de cable estudiado en el texto, utilizando una separación lateral de $3$ km y costes de $7$ y $4$ unidades monetarias por kilómetro, respectivamente. Definir con claridad la variable de decisión y el intervalo físicamente posible.

\item Una planta química utiliza un depósito cilíndrico cerrado de volumen $20\ \mathrm{m}^3$. El material de las bases cuesta un $40\%$ más por metro cuadrado que el material lateral. Formular la función de coste en términos del radio y determinar la relación óptima entre altura y radio.

\item La potencia útil de una máquina en función de su velocidad angular se aproxima, dentro de su rango de operación, por
\[
P(\omega)=a\omega-b\omega^3,
\qquad a,b>0,
\qquad \omega\geq0.
\]
Determinar la velocidad angular que maximiza la potencia y explicar cómo cambia la potencia antes y después de ese punto. Expresar la potencia máxima en función de $a$ y $b$.

\item En un proceso de calentamiento,
\[
T(t)=20+10\bigl(1-e^{-0.1t}\bigr),
\qquad t\geq0.
\]
\begin{enumerate}[label=\alph*)]
\item calcular $T'$ y $T''$;
\item interpretar físicamente ambas derivadas;
\item decidir si existe un instante interior en el que la tasa de calentamiento sea máxima;
\item explicar la respuesta utilizando simultáneamente el registro analítico, el gráfico y el verbal.
\end{enumerate}

\item Una línea de producción tiene producción acumulada
\[
Q(t)=7t^2-t^3,
\qquad 0\leq t\leq7.
\]
\begin{enumerate}[label=\alph*)]
\item calcular la tasa instantánea de producción $Q'(t)$;
\item determinar cuándo aumenta y cuándo disminuye la producción acumulada;
\item localizar el máximo de $Q$ en el intervalo;
\item estudiar cuándo la tasa de producción $Q'$ aumenta o disminuye;
\item interpretar la diferencia entre las preguntas anteriores.
\end{enumerate}

\item Un elemento elástico presenta, dentro de su rango de operación $0\leq x\leq L$, una fuerza
\[
F(x)=kx^2,
\qquad k>0.
\]
Calcular $F'$ y $F''$ e interpretar físicamente ambas expresiones. Explicar por qué no existe un máximo interior de $F$ y dónde se alcanza el máximo en el intervalo de operación.

\item \M Seleccionar cuatro ejercicios de derivación simbólica de esta relación y comprobar los resultados mediante \texttt{D}. En al menos uno de ellos utilizar \texttt{Simplify} para comparar dos expresiones algebraicamente distintas de la misma derivada. La entrega debe contener primero el desarrollo manual.

\item \M Resolver con \texttt{D}, \texttt{Solve} y \texttt{Plot} el estudio de
\[
f(x)=10+3xe^{-0.2x}.
\]
Antes de utilizar el ordenador, predecir verbalmente el comportamiento de la función para $x\geq0$ y justificar después qué partes de la predicción confirma cada comando.

\item \M Para una función $f$ elegida por el estudiante, construir en Mathematica una gráfica de $f$ y otra de $f'$ en el mismo intervalo. Señalar al menos dos puntos o intervalos en los que pueda verse directamente la correspondencia entre pendiente de $f$ y valor de $f'$, y otros dos en los que pueda relacionarse el crecimiento de $f'$ con la convexidad de $f$.

\end{enumerate}

\cleardoublepage
\chapter{Cálculo integral}
\label{ch:integral}

En el capítulo anterior estudiamos la derivada como herramienta para describir una tasa instantánea de cambio. El cálculo integral aborda, en cierto sentido, el problema inverso: a partir de una tasa de variación, reconstruir una magnitud acumulada. Si conocemos la velocidad de un móvil, podemos recuperar su desplazamiento; si conocemos el caudal que atraviesa una conducción, podemos calcular el volumen transportado; si conocemos la intensidad de corriente, podemos determinar la carga que ha circulado; y si conocemos una densidad lineal, podemos obtener la masa total de una pieza.

Esta relación entre \emph{tasa} y \emph{acumulación} constituye una de las ideas centrales del cálculo. La integral aparece también al medir áreas, volúmenes y otras magnitudes geométricas. El Teorema Fundamental del Cálculo establece que ambos puntos de vista están íntimamente ligados con la derivación.

\section{Primitivas e integral indefinida}
\label{sec:primitivas}

\subsection{Primitivas y reglas básicas}

Comenzamos con el problema inverso de la derivación. Conocida una función $f$, buscamos una función cuya derivada sea precisamente $f$.

\begin{definition}[Primitiva]
Sea $f:I\to\R$ una función definida en un intervalo $I$. Una función $F:I\to\R$ se llama \emph{primitiva} o \emph{antiderivada} de $f$ en $I$ si
\[
F'(x)=f(x),\qquad \forall x\in I.
\]
\end{definition}

Por ejemplo,
\[
F(x)=x^3
\]
es una primitiva de $f(x)=3x^2$. También lo son $x^3+4$, $x^3-\pi$ o, en general, $x^3+C$ para cualquier constante real $C$.

\begin{theorem}
Si $F$ es una primitiva de $f$ en un intervalo $I$, entonces todas las primitivas de $f$ en $I$ son de la forma
\[
F(x)+C,\qquad C\in\R.
\]
\end{theorem}

\begin{proof}
Si $G$ es otra primitiva de $f$, entonces
\[
(G-F)'=G'-F'=f-f=0.
\]
Por tanto, $G-F$ es constante en $I$, es decir, $G(x)=F(x)+C$.
\end{proof}

\begin{definition}[Integral indefinida]
El conjunto de todas las primitivas de una función $f$ se denomina \emph{integral indefinida} de $f$ y se escribe
\[
\int f(x)\,dx=F(x)+C,
\]
donde $F'(x)=f(x)$ y $C\in\R$ es la \emph{constante de integración}.
\end{definition}

La integral indefinida no es, por tanto, una única función, sino una familia de funciones que difieren en una constante.

La derivación y la integración indefinida son operaciones inversas en el sentido de que
\[
\frac{d}{dx}\left(\int f(x)\,dx\right)=f(x)
\]
y
\[
\int F'(x)\,dx=F(x)+C.
\]

De las reglas de derivación se obtienen inmediatamente reglas de integración. Si $\alpha,\beta\in\R$ y $f,g$ admiten primitivas, entonces
\[
\int \bigl(\alpha f(x)+\beta g(x)\bigr)\,dx
=
\alpha\int f(x)\,dx+\beta\int g(x)\,dx.
\]

Entre las primitivas elementales que utilizaremos con mayor frecuencia están
\[
\int x^n\,dx=\frac{x^{n+1}}{n+1}+C,\qquad n\neq-1,
\]
\[
\int \frac1x\,dx=\ln|x|+C,
\]
\[
\int e^x\,dx=e^x+C,
\qquad
\int a^x\,dx=\frac{a^x}{\ln a}+C\quad (a>0,\ a\neq1),
\]
\[
\int \cos x\,dx=\sin x+C,
\qquad
\int \sin x\,dx=-\cos x+C,
\]
\[
\int \frac{1}{1+x^2}\,dx=\arctan x+C,
\qquad
\int \frac{1}{\sqrt{1-x^2}}\,dx=\arcsin x+C.
\]

Estas fórmulas deben entenderse en intervalos donde las expresiones estén definidas. Por ejemplo, la derivada de $\ln|x|$ es $1/x$ tanto para $x>0$ como para $x<0$.

\begin{example}
Calcular
\[
\int \left(4x^3-\frac{3}{x^2}+2e^x\right)dx.
\]
Aplicando linealidad,
\[
\int4x^3\,dx=x^4,
\qquad
\int -3x^{-2}\,dx=\frac{3}{x},
\qquad
\int2e^x\,dx=2e^x.
\]
Por tanto,
\[
\boxed{\int \left(4x^3-\frac{3}{x^2}+2e^x\right)dx
=x^4+\frac3x+2e^x+C.}
\]
Una comprobación inmediata consiste en derivar el resultado.
\end{example}

\begin{solvedexercise}
La velocidad de un carro que se desplaza sobre una guía es
\[
v(t)=6t^2-4t+3\qquad (\mathrm{m/s}).
\]
Sabiendo que $s(0)=5$ m, determinar su posición $s(t)$.
\tcblower
\textbf{Resolución.}\par\smallskip
Como $s'(t)=v(t)$,
\[
s(t)=\int(6t^2-4t+3)\,dt
=2t^3-2t^2+3t+C.
\]
La condición inicial permite determinar la constante:
\[
5=s(0)=C.
\]
Por tanto,
\[
\boxed{s(t)=2t^3-2t^2+3t+5.}
\]
La constante de integración no es un detalle algebraico: contiene la información sobre la posición inicial, que la velocidad por sí sola no determina.
\end{solvedexercise}

\subsection{De una tasa a una cantidad acumulada}

En muchos modelos de ingeniería no se conoce directamente una cantidad $Q(t)$, sino su tasa de cambio $Q'(t)$. La integración permite reconstruir $Q$ salvo una constante, que debe determinarse mediante una condición inicial.

Si
\[
Q'(t)=r(t),
\]
entonces
\[
Q(t)=\int r(t)\,dt=R(t)+C,
\]
donde $R$ es una primitiva de $r$. La interpretación depende del contexto:

\begin{itemize}
\item si $Q$ es posición, $r$ es velocidad;
\item si $Q$ es carga eléctrica, $r$ es intensidad de corriente;
\item si $Q$ es volumen almacenado, $r$ es caudal neto;
\item si $Q$ es masa producida, $r$ es tasa de producción;
\item si $Q$ es energía transferida, $r$ puede ser una potencia.
\end{itemize}

\begin{example}[Producción acumulada]
Una línea de producción trabaja, durante un cierto régimen, a una tasa
\[
r(t)=40+2t-\frac15t^2
\]
toneladas por hora, donde $t$ se mide en horas desde el inicio del turno. Si $Q(t)$ es la producción acumulada y $Q(0)=0$, entonces
\[
Q'(t)=r(t)
\]
y, por tanto,
\[
Q(t)=40t+t^2-\frac1{15}t^3+C.
\]
Como $Q(0)=0$, resulta $C=0$ y
\[
\boxed{Q(t)=40t+t^2-\frac1{15}t^3.}
\]
Obsérvese la diferencia entre ambas magnitudes: $r(t)$ se mide en toneladas por hora y $Q(t)$ en toneladas.
\end{example}

\section{Métodos de integración}
\label{sec:metodos-integracion}

Las fórmulas inmediatas permiten integrar funciones sencillas. Sin embargo, en la mayoría de los problemas es necesario transformar el integrando hasta reconocer una forma conocida. No existe un algoritmo elemental único que resuelva todas las integrales. Por ello resulta esencial identificar la estructura de la función y elegir el método adecuado.

\subsection{Cambio de variable e integración por partes}

\subsubsection{Cambio de variable}

La regla de la cadena para derivadas afirma que
\[
\frac{d}{dx}F(g(x))=F'(g(x))g'(x).
\]
Leída en sentido inverso conduce al método de sustitución.

\begin{theorem}[Cambio de variable]
Si $F'=f$ y $g$ es diferenciable, entonces
\[
\int f(g(x))g'(x)\,dx=F(g(x))+C.
\]
Equivalentemente, haciendo
\[
u=g(x),\qquad du=g'(x)\,dx,
\]
se transforma la integral en
\[
\int f(u)\,du.
\]
\end{theorem}

La dificultad principal no consiste en escribir $u=g(x)$, sino en reconocer una composición y comprobar que el resto del integrando contiene, salvo una constante, la derivada de la función interior.

\begin{example}
Calcular
\[
\int 2x\,e^{x^2+1}\,dx.
\]
Tomamos
\[
u=x^2+1,
\qquad du=2x\,dx.
\]
Entonces
\[
\int 2x\,e^{x^2+1}\,dx
=
\int e^u\,du
=e^u+C,
\]
y deshaciendo el cambio,
\[
\boxed{\int 2x\,e^{x^2+1}\,dx=e^{x^2+1}+C.}
\]
\end{example}

\begin{example}
Calcular
\[
\int \frac{x}{1+x^2}\,dx.
\]
Como la derivada de $1+x^2$ es $2x$, tomamos $u=1+x^2$. Entonces
\[
\int \frac{x}{1+x^2}\,dx
=
\frac12\int\frac{du}{u}
=
\frac12\ln|u|+C.
\]
Puesto que $1+x^2>0$,
\[
\boxed{\int \frac{x}{1+x^2}\,dx=\frac12\ln(1+x^2)+C.}
\]
\end{example}

En integrales definidas es posible cambiar también los límites de integración. Si $u=g(x)$ y $x$ varía de $a$ a $b$, entonces $u$ varía de $g(a)$ a $g(b)$:
\[
\int_a^b f(g(x))g'(x)\,dx
=
\int_{g(a)}^{g(b)}f(u)\,du.
\]

\subsubsection{Integración por partes}

La regla del producto
\[
(uv)'=u'v+uv'
\]
puede escribirse como
\[
u v'= (uv)'-u'v.
\]
Integrando ambos miembros obtenemos la fórmula de integración por partes.

\begin{theorem}[Integración por partes]
Si $u$ y $v$ son funciones diferenciables,
\[
\boxed{\int u\,dv=uv-\int v\,du.}
\]
\end{theorem}

Este método es útil cuando el integrando es un producto en el que uno de los factores se simplifica al derivarlo y el otro puede integrarse fácilmente. También permite integrar funciones como $\ln x$, $\arctan x$ o $\arcsin x$ escribiéndolas como un producto por $1$.

\begin{solvedexercise}
Calcular
\[
\int x e^{2x}\,dx.
\]
\tcblower
\textbf{Resolución.}\par\smallskip
Tomamos
\[
u=x,
\qquad dv=e^{2x}\,dx.
\]
Entonces
\[
du=dx,
\qquad v=\frac12e^{2x}.
\]
Por tanto,
\[
\int xe^{2x}\,dx
=
\frac{x}{2}e^{2x}-\frac12\int e^{2x}\,dx
=
\frac{x}{2}e^{2x}-\frac14e^{2x}+C.
\]
Así,
\[
\boxed{\int xe^{2x}\,dx=\frac{e^{2x}}4(2x-1)+C.}
\]
\end{solvedexercise}

En ocasiones hay que aplicar el método varias veces. En otras, la integral original reaparece después de integrar por partes y puede despejarse algebraicamente.

\begin{example}
Sea
\[
I=\int e^x\cos x\,dx.
\]
Integrando por partes dos veces se obtiene
\[
I=e^x\sin x-\int e^x\sin x\,dx
\]
a partir de la primera aplicación, y
\[
\int e^x\sin x\,dx=-e^x\cos x+I
\]
a partir de la segunda. Sustituyendo,
\[
I=e^x\sin x+e^x\cos x-I,
\]
de modo que
\[
\boxed{\int e^x\cos x\,dx=\frac{e^x}{2}(\sin x+\cos x)+C.}
\]
\end{example}

\subsection{Funciones racionales}

Una función racional tiene la forma
\[
R(x)=\frac{P(x)}{Q(x)},
\]
donde $P$ y $Q$ son polinomios y $Q\not\equiv0$. Su integración se basa en reducir la fracción a sumas de expresiones elementales.

Si $\deg P\geq\deg Q$, se realiza primero la división de polinomios:
\[
\frac{P(x)}{Q(x)}=S(x)+\frac{R(x)}{Q(x)},
\qquad \deg R<\deg Q.
\]
A continuación se factoriza el denominador real en factores lineales y cuadráticos irreducibles y se utiliza una descomposición en fracciones simples.

Por ejemplo, si
\[
Q(x)=(x-a)(x-b),\qquad a\neq b,
\]
se busca
\[
\frac{P(x)}{(x-a)(x-b)}
=
\frac{A}{x-a}+\frac{B}{x-b}.
\]
Si aparecen factores repetidos, deben incluirse todas sus potencias. Por ejemplo,
\[
\frac{P(x)}{(x-a)^2(x-b)}
=
\frac{A}{x-a}+\frac{B}{(x-a)^2}+\frac{C}{x-b}.
\]
Para un factor cuadrático irreducible $x^2+px+q$ aparece un numerador lineal:
\[
\frac{Ax+B}{x^2+px+q}.
\]

\begin{solvedexercise}
Calcular
\[
\int\frac{2x}{(x-2)(x+1)}\,dx.
\]
\tcblower
\textbf{Resolución.}\par\smallskip
Buscamos constantes $A$ y $B$ tales que
\[
\frac{2x}{(x-2)(x+1)}
=
\frac{A}{x-2}+\frac{B}{x+1}.
\]
Multiplicando por $(x-2)(x+1)$,
\[
2x=A(x+1)+B(x-2).
\]
Tomando $x=2$ se obtiene $4=3A$, luego $A=4/3$. Tomando $x=-1$ se obtiene $-2=-3B$, luego $B=2/3$. Por tanto,
\[
\int\frac{2x}{(x-2)(x+1)}\,dx
=
\frac43\ln|x-2|+\frac23\ln|x+1|+C.
\]
\end{solvedexercise}

Cuando el denominador contiene factores repetidos, el \emph{método de Hermite} permite separar la derivada de una función racional de una fracción con denominador sin factores repetidos. En términos generales, para una función racional propia $P/Q$ puede buscarse una descomposición de la forma
\[
\frac{P(x)}{Q(x)}
=
R'(x)+\frac{S(x)}{T(x)},
\]
donde $R$ es racional y $T$ no contiene factores repetidos. Así,
\[
\int\frac{P(x)}{Q(x)}\,dx
=
R(x)+\int\frac{S(x)}{T(x)}\,dx.
\]
El método evita desarrollar por completo fracciones simples asociadas a potencias elevadas de factores irreducibles.

\begin{example}[Reducción de Hermite]
Calcular
\[
\int\frac{dx}{(1+x^2)^2}.
\]
Observamos que
\[
\frac{1}{(1+x^2)^2}
=
\frac{d}{dx}\left(\frac{x}{2(1+x^2)}\right)
+
\frac{1}{2(1+x^2)}.
\]
Por tanto,
\[
\int\frac{dx}{(1+x^2)^2}
=
\frac{x}{2(1+x^2)}+\frac12\arctan x+C.
\]
La parte racional procede de reducir el factor repetido $(1+x^2)^2$ y la integral restante contiene ya un denominador sin repetición.
\end{example}

\subsection{Integrales trigonométricas e irracionales}

\subsubsection{Integrales trigonométricas}

Las identidades trigonométricas permiten transformar productos y potencias hasta reconocer primitivas elementales. Entre las identidades que utilizaremos están
\[
\sin^2x+\cos^2x=1,
\qquad
1+\tan^2x=\sec^2x,
\]
\[
\sin^2x=\frac{1-\cos2x}{2},
\qquad
\cos^2x=\frac{1+\cos2x}{2}.
\]

Para integrales de la forma
\[
\int\sin^m x\cos^n x\,dx
\]
son útiles las siguientes ideas:

\begin{itemize}
\item si la potencia del seno es impar, se reserva un factor $\sin x$ y se utiliza $\sin^2x=1-\cos^2x$;
\item si la potencia del coseno es impar, se reserva un factor $\cos x$ y se utiliza $\cos^2x=1-\sin^2x$;
\item si ambas potencias son pares, se aplican las fórmulas de ángulo doble.
\end{itemize}

\begin{example}
Calcular
\[
\int\sin^3x\cos^2x\,dx.
\]
Escribimos
\[
\sin^3x=\sin x(1-\cos^2x).
\]
Con $u=\cos x$, $du=-\sin x\,dx$,
\[
\int\sin^3x\cos^2x\,dx
=-\int(1-u^2)u^2\,du
=-\frac{u^3}{3}+\frac{u^5}{5}+C.
\]
Por tanto,
\[
\boxed{\int\sin^3x\cos^2x\,dx
=-\frac{\cos^3x}{3}+\frac{\cos^5x}{5}+C.}
\]
\end{example}

Para potencias de tangente y secante se emplea la identidad $1+\tan^2x=\sec^2x$ y las derivadas
\[
(\tan x)'=\sec^2x,
\qquad
(\sec x)'=\sec x\tan x.
\]
Un caso clásico que requiere integración por partes es
\[
\int\sec^3x\,dx
=
\frac12\sec x\tan x+\frac12\ln|\sec x+\tan x|+C.
\]

\subsubsection{Integrales irracionales}

Las integrales que contienen radicales pueden, en ciertos casos, transformarse mediante una sustitución que elimine la raíz. Si aparece una raíz de una expresión lineal, por ejemplo $\sqrt{ax+b}$, suele ser útil tomar
\[
t=\sqrt{ax+b}.
\]

\begin{example}
Calcular
\[
\int\frac{\sqrt{2x+1}}{x+1}\,dx.
\]
Tomamos
\[
t=\sqrt{2x+1},
\qquad t^2=2x+1,
\qquad x=\frac{t^2-1}{2},
\qquad dx=t\,dt.
\]
Además,
\[
x+1=\frac{t^2+1}{2}.
\]
Por tanto,
\[
\int\frac{\sqrt{2x+1}}{x+1}\,dx
=
\int\frac{2t^2}{t^2+1}\,dt
=
2\int\left(1-\frac1{1+t^2}\right)dt,
\]
de donde
\[
\boxed{\int\frac{\sqrt{2x+1}}{x+1}\,dx
=2\sqrt{2x+1}-2\arctan\!\sqrt{2x+1}+C.}
\]
\end{example}

Para radicales cuadráticos son especialmente útiles las sustituciones trigonométricas
\[
\sqrt{a^2-x^2}:\quad x=a\sin t,
\]
\[
\sqrt{a^2+x^2}:\quad x=a\tan t,
\]
\[
\sqrt{x^2-a^2}:\quad x=a\sec t.
\]
La elección se justifica por las identidades $1-\sin^2t=\cos^2t$, $1+\tan^2t=\sec^2t$ y $\sec^2t-1=\tan^2t$.

\begin{solvedexercise}
Calcular
\[
\int\sqrt{16-x^2}\,dx.
\]
\tcblower
\textbf{Resolución.}\par\smallskip
Tomamos $x=4\sin t$. Entonces
\[
dx=4\cos t\,dt,
\qquad
\sqrt{16-x^2}=4\cos t
\]
en el intervalo correspondiente. Así,
\[
\int\sqrt{16-x^2}\,dx
=16\int\cos^2t\,dt
=8\int(1+\cos2t)\,dt.
\]
Por tanto,
\[
8t+4\sin2t+C.
\]
Como $t=\arcsin(x/4)$ y
\[
\sin2t=2\sin t\cos t
=2\frac{x}{4}\frac{\sqrt{16-x^2}}{4},
\]
obtenemos
\[
\boxed{\int\sqrt{16-x^2}\,dx
=\frac{x}{2}\sqrt{16-x^2}+8\arcsin\left(\frac{x}{4}\right)+C.}
\]
\end{solvedexercise}

La elección del método de integración forma parte del problema. Antes de comenzar los cálculos conviene preguntarse: ¿hay una composición clara?, ¿es un producto adecuado para partes?, ¿es una función racional?, ¿puede utilizarse una identidad trigonométrica?, ¿un cambio de variable elimina el radical? Estas preguntas evitan muchos desarrollos innecesarios.

\section{Integral definida y Teorema Fundamental del Cálculo}
\label{sec:integral-definida}

\subsection{Sumas de Riemann e integral definida}

La idea de integral definida surge de un procedimiento de aproximación. Supongamos que $f:[a,b]\to\R$ es continua y, para comenzar, que $f(x)\geq0$ en $[a,b]$. Queremos medir la región comprendida entre la gráfica de $f$, el eje horizontal y las rectas $x=a$ y $x=b$.

Dividimos el intervalo $[a,b]$ mediante una partición
\[
P=\{a=x_0<x_1<\cdots<x_n=b\}.
\]
Cada subintervalo $[x_{i-1},x_i]$ tiene longitud
\[
\Delta x_i=x_i-x_{i-1}.
\]
En cada uno de ellos elegimos un punto $c_i\in[x_{i-1},x_i]$. El rectángulo de base $\Delta x_i$ y altura $f(c_i)$ tiene área
\[
f(c_i)\Delta x_i.
\]
La suma de las áreas de todos estos rectángulos es
\[
S(P)=\sum_{i=1}^{n} f(c_i)\Delta x_i,
\]
y recibe el nombre de \emph{suma de Riemann} asociada a la partición $P$ y a los puntos $c_i$.

\begin{figure}[htbp]
\centering
\begin{tikzpicture}[x=1.25cm,y=1.0cm]
  % ejes sin puntas de flecha para evitar conflictos con babel-spanish
  \draw (-0.15,0)--(5.35,0);
  \draw (0,-0.15)--(0,4.8);
  \node[right] at (5.35,0) {$x$};
  \node[above] at (0,4.8) {$y$};

  % rectángulos de Riemann (extremos derechos)
  \fill[gray!22] (0.5,0) rectangle (1.25,{0.35*1.25*1.25+0.7});
  \fill[gray!22] (1.25,0) rectangle (2.0,{0.35*2.0*2.0+0.7});
  \fill[gray!22] (2.0,0) rectangle (2.75,{0.35*2.75*2.75+0.7});
  \fill[gray!22] (2.75,0) rectangle (3.5,{0.35*3.5*3.5+0.7});
  \fill[gray!22] (3.5,0) rectangle (4.25,{0.35*4.25*4.25+0.7});

  \draw[gray] (0.5,0) rectangle (1.25,{0.35*1.25*1.25+0.7});
  \draw[gray] (1.25,0) rectangle (2.0,{0.35*2.0*2.0+0.7});
  \draw[gray] (2.0,0) rectangle (2.75,{0.35*2.75*2.75+0.7});
  \draw[gray] (2.75,0) rectangle (3.5,{0.35*3.5*3.5+0.7});
  \draw[gray] (3.5,0) rectangle (4.25,{0.35*4.25*4.25+0.7});

  % curva
  \draw[thick] plot[domain=.5:4.25,samples=120] (\x,{0.35*\x*\x+0.7});

  \node[below] at (.5,0) {$a$};
  \node[below] at (4.25,0) {$b$};
  \node[above left] at (3.9,{0.35*3.9*3.9+0.7}) {$y=f(x)$};
\end{tikzpicture}
\caption{Aproximación mediante una suma de Riemann. Al refinar la partición, los rectángulos se adaptan cada vez mejor a la región bajo la curva.}
\label{fig:riemann}
\end{figure}

La aproximación mejora cuando todos los subintervalos son pequeños. Para medir su tamaño utilizamos la \emph{norma} o \emph{malla} de la partición,
\[
\|P\|=\max_{1\leq i\leq n}\Delta x_i.
\]
Decir que $\|P\|\to0$ significa que el mayor de los subintervalos de la partición tiende a tener longitud cero. En ese proceso, el número de rectángulos aumenta y sus anchuras disminuyen.

Lo importante es que, si $f$ es continua, las sumas de Riemann se aproximan a un mismo número independientemente de los puntos $c_i$ elegidos dentro de cada subintervalo.

\begin{definition}[Integral definida]
Sea $f:[a,b]\to\R$. Decimos que $f$ es \emph{integrable} en $[a,b]$ si existe un número real $I$ tal que
\[
\sum_{i=1}^{n}f(c_i)\Delta x_i\longrightarrow I
\qquad\text{cuando}\qquad
\|P\|\longrightarrow0,
\]
para cualquier elección de los puntos $c_i\in[x_{i-1},x_i]$. En ese caso escribimos
\[
\boxed{I=\int_a^b f(x)\,dx}
\]
y llamamos a $I$ \emph{integral definida} de $f$ entre $a$ y $b$.
\end{definition}

Toda función continua en un intervalo cerrado es integrable en él. Cuando $f\geq0$, la definición anterior formaliza la intuición geométrica: la integral es el límite de las áreas de los rectángulos utilizados para aproximar la región bajo la gráfica. Cuando $f$ toma valores negativos, los términos correspondientes de la suma de Riemann son negativos; por ello, la integral representa \emph{área con signo}, no necesariamente área geométrica.

\begin{example}
Consideremos $f(x)=x$ en $[0,1]$ y dividamos el intervalo en $n$ partes iguales. Si tomamos como $c_i$ el extremo derecho de cada subintervalo, entonces
\[
\Delta x=\frac1n,
\qquad
c_i=\frac{i}{n}.
\]
La suma de Riemann es
\[
S_n=\sum_{i=1}^{n}\frac{i}{n}\frac1n
=\frac{1}{n^2}\sum_{i=1}^{n}i
=\frac{1}{n^2}\frac{n(n+1)}{2}
=\frac{n+1}{2n}.
\]
Por tanto,
\[
\int_0^1 x\,dx
=\lim_{n\to\infty}S_n
=\frac12.
\]
Geométricamente, este valor coincide con el área del triángulo de base $1$ y altura $1$.
\end{example}

\subsection{Por qué la integral representa el área bajo una curva}

La definición mediante sumas de Riemann permite aproximar el área mediante rectángulos. Sin embargo, existe otra forma de entender la relación entre integral y área que conecta directamente este capítulo con el cálculo diferencial.

Supongamos, en primer lugar, que $f:[a,b]\to\R$ es continua y positiva. Para cada $t\in[a,b]$, definimos
\[
A(t)=\text{área comprendida bajo la gráfica de $f$ entre $a$ y $t$}.
\]
Así, $A(a)=0$ y el área total que queremos calcular es $A(b)$.


Aumentemos ahora $t$ una pequeña cantidad positiva $\Delta t$. Entonces
\[
A(t+\Delta t)-A(t)
\]
es el área de la estrecha región situada bajo la curva entre $t$ y $t+\Delta t$.

Si, para fijar ideas, $f$ es creciente en ese pequeño intervalo, dicha región queda comprendida entre dos rectángulos de base $\Delta t$: uno de altura $f(t)$ y otro de altura $f(t+\Delta t)$. Por tanto,
\[
f(t)\,\Delta t
\leq
A(t+\Delta t)-A(t)
\leq
f(t+\Delta t)\,\Delta t.
\]
Al dividir entre $\Delta t>0$, obtenemos
\[
f(t)
\leq
\frac{A(t+\Delta t)-A(t)}{\Delta t}
\leq
f(t+\Delta t).
\]

\begin{figure}[htbp]
\centering
\begin{tikzpicture}[x=1.5cm,y=.9cm]
  \draw[-latex] (-.2,0)--(4.5,0) node[right] {$x$};
  \draw[-latex] (0,-.2)--(0,6.2) node[above] {$y$};
  \fill[gray!25] (2.1,0) -- plot[domain=2.1:3.0,samples=50] (\x,{0.28*\x*\x+1.1}) -- (3.0,0) -- cycle;
  \draw[thick] plot[domain=.3:4,samples=100] (\x,{0.28*\x*\x+1.1});
  \draw[dashed] (2.1,0)--(2.1,{0.28*2.1*2.1+1.1});
  \draw[dashed] (3.0,0)--(3.0,{0.28*3.0*3.0+1.1});
  \draw[densely dotted] (2.1,{0.28*2.1*2.1+1.1}) -- (3.0,{0.28*2.1*2.1+1.1});
  \draw[densely dotted] (2.1,{0.28*3.0*3.0+1.1}) -- (3.0,{0.28*3.0*3.0+1.1});
  \node[below] at (2.1,0) {$t$};
  \node[below] at (3.0,0) {$t+\Delta t$};
  \node at (2.57,1.05) {$\Delta A$};
  \node[left] at (2.08,{0.28*2.1*2.1+1.1}) {$f(t)$};
  \node[left] at (2.08,{0.28*3.0*3.0+1.1}) {$f(t+\Delta t)$};
\end{tikzpicture}
\caption{El incremento de área $\Delta A=A(t+\Delta t)-A(t)$ queda acotado mediante rectángulos.}
\label{fig:incremento-area}
\end{figure}

Como $f$ es continua,
\[
\lim_{\Delta t\to0}f(t+\Delta t)=f(t).
\]
Por el teorema del sándwich,
\[
\lim_{\Delta t\to0}
\frac{A(t+\Delta t)-A(t)}{\Delta t}
=f(t).
\]
Pero el límite de la izquierda es, precisamente, la derivada de $A$. Por tanto,
\[
\boxed{A'(t)=f(t).}
\]

Hemos obtenido una conclusión fundamental: \emph{la función que mide el área acumulada bajo una curva tiene como derivada la propia función que determina la altura de la curva}. En consecuencia, $A$ es una primitiva de $f$.

Si $F$ es cualquier primitiva de $f$, entonces
\[
A(t)=F(t)+C.
\]
Como $A(a)=0$, se tiene
\[
0=F(a)+C,
\qquad\text{de donde}\qquad
C=-F(a).
\]
Así,
\[
A(t)=F(t)-F(a),
\]
y, tomando $t=b$,
\[
\boxed{A(b)=F(b)-F(a).}
\]

Este razonamiento explica por qué, cuando $f\geq0$,
\[
\boxed{\text{área bajo la gráfica de $f$ en $[a,b]$}=\int_a^b f(x)\,dx.}
\]
No se trata simplemente de una fórmula para calcular áreas: la relación surge porque la tasa instantánea de crecimiento del área acumulada es la altura $f(t)$ de la curva.

La hipótesis de monotonía utilizada en el argumento anterior solo simplifica la representación geométrica. Si $f$ es continua pero no monótona, en cada pequeño intervalo $[t,t+\Delta t]$ podemos considerar sus valores mínimo y máximo, $m_{\Delta t}$ y $M_{\Delta t}$. Entonces
\[
m_{\Delta t}\,\Delta t
\leq
A(t+\Delta t)-A(t)
\leq
M_{\Delta t}\,\Delta t.
\]
Por continuidad, tanto $m_{\Delta t}$ como $M_{\Delta t}$ tienden a $f(t)$ cuando $\Delta t\to0$, y se obtiene de nuevo $A'(t)=f(t)$.

Esta idea anticipa el Teorema Fundamental del Cálculo, que formaliza la relación entre derivación, acumulación e integración.

Una integral definida es un número. No debe confundirse con la integral indefinida, que representa una familia de primitivas.

Las propiedades fundamentales son:
\[
\int_a^a f(x)\,dx=0,
\qquad
\int_b^a f(x)\,dx=-\int_a^b f(x)\,dx,
\]
\[
\int_a^b f(x)\,dx
=
\int_a^c f(x)\,dx+\int_c^b f(x)\,dx,
\qquad c\in[a,b],
\]
\[
\int_a^b(\alpha f+\beta g)\,dx
=
\alpha\int_a^b f\,dx+\beta\int_a^b g\,dx.
\]
Además, si $f(x)\leq g(x)$ para todo $x\in[a,b]$, entonces
\[
\int_a^b f(x)\,dx\leq\int_a^b g(x)\,dx.
\]
En particular, si $f\geq0$,
\[
\int_a^b f(x)\,dx\geq0.
\]

La integral es una \emph{acumulación con signo}. Las partes de la gráfica situadas sobre el eje $OX$ contribuyen positivamente y las situadas bajo el eje contribuyen negativamente.

\begin{example}
Sea
\[
f(x)=x-1
\]
en $[0,3]$. Entonces
\[
\int_0^3(x-1)\,dx
=
\left[\frac{x^2}{2}-x\right]_0^3
=\frac32.
\]
Este valor no es el área geométrica total entre la gráfica y el eje $OX$, porque $f<0$ en $[0,1]$. El área geométrica es
\[
\int_0^1(1-x)\,dx+\int_1^3(x-1)\,dx
=\frac12+2=\frac52.
\]
\end{example}

\subsection{Teorema Fundamental y funciones de acumulación}

La conexión entre derivación e integración queda expresada por el Teorema Fundamental del Cálculo.

\begin{theorem}[Teorema Fundamental del Cálculo]
Sea $f$ continua en $[a,b]$ y definamos
\[
F(x)=\int_a^x f(t)\,dt.
\]
Entonces $F$ es derivable en $(a,b)$ y
\[
\boxed{F'(x)=f(x).}
\]
\end{theorem}

La función $F$ mide la acumulación de $f$ desde el punto fijo $a$ hasta el punto variable $x$. El teorema afirma que la tasa instantánea de crecimiento de esa acumulación es precisamente el valor que se está acumulando.

Esta relación permite pasar de un registro a otro:

\begin{center}
\begin{tabular}{p{.27\textwidth}p{.30\textwidth}p{.32\textwidth}}
\toprule
Información sobre $f$ & Información sobre $F(x)=\int_a^x f(t)\,dt$ & Lectura verbal \\
\midrule
$f(x)>0$ & $F'(x)>0$ & la acumulación aumenta \\
$f(x)<0$ & $F'(x)<0$ & la acumulación disminuye \\
$f(x)=0$ & $F'(x)=0$ & $F$ tiene un punto estacionario \\
$f$ creciente & $F''=f'>0$ & la pendiente de $F$ aumenta \\
$f$ decreciente & $F''=f'<0$ & la pendiente de $F$ disminuye \\
\bottomrule
\end{tabular}
\end{center}

Por tanto, una gráfica de $f$ contiene información directa sobre la forma de cualquier función de acumulación asociada, incluso sin calcular explícitamente la integral.

\begin{solvedexercise}
Sea
\[
f(t)=t^2-1
\]
y
\[
F(x)=\int_0^x(t^2-1)\,dt.
\]
Sin calcular inicialmente una fórmula para $F$, determinar sus intervalos de crecimiento y decrecimiento y sus extremos relativos. Después obtener $F$ analíticamente y comprobar las conclusiones.
\tcblower
\textbf{Resolución.}\par\smallskip
Por el Teorema Fundamental,
\[
F'(x)=f(x)=x^2-1.
\]
Así,
\[
F'(x)>0\quad\text{si}\quad x<-1\ \text{o}\ x>1,
\]
y
\[
F'(x)<0\quad\text{si}\quad -1<x<1.
\]
Por tanto, $F$ crece en $(-\infty,-1)$ y $(1,\infty)$ y decrece en $(-1,1)$. En $x=-1$ pasa de crecer a decrecer, luego hay un máximo relativo; en $x=1$ pasa de decrecer a crecer, luego hay un mínimo relativo.

Ahora calculamos
\[
F(x)=\int_0^x(t^2-1)\,dt
=
\left[\frac{t^3}{3}-t\right]_0^x
=
\frac{x^3}{3}-x,
\]
que confirma el análisis anterior. Obsérvese que para estudiar la monotonía de $F$ no era necesario encontrar esta expresión: bastaba estudiar el signo de $f$.
\end{solvedexercise}

La segunda parte del Teorema Fundamental proporciona un procedimiento eficaz para calcular integrales definidas.

\begin{theorem}[Regla de Barrow]
Si $f$ es continua en $[a,b]$ y $G$ es cualquier primitiva de $f$, entonces
\[
\boxed{\int_a^b f(x)\,dx=G(b)-G(a).}
\]
\end{theorem}

Utilizaremos frecuentemente la notación
\[
\left[G(x)\right]_a^b=G(b)-G(a).
\]

\begin{example}
Calcular
\[
\int_1^2(x^2+x+1)\,dx.
\]
Una primitiva es
\[
G(x)=\frac{x^3}{3}+\frac{x^2}{2}+x.
\]
Por tanto,
\[
\int_1^2(x^2+x+1)\,dx
=
\left[\frac{x^3}{3}+\frac{x^2}{2}+x\right]_1^2
=
\frac{29}{6}.
\]
\end{example}

Si el límite superior depende de otra variable, el Teorema Fundamental se combina con la regla de la cadena. Por ejemplo,
\[
F(x)=\int_a^{g(x)}f(t)\,dt
\]
implica
\[
\boxed{F'(x)=f(g(x))g'(x).}
\]
Análogamente,
\[
\frac{d}{dx}\left(\int_{u(x)}^{v(x)}f(t)\,dt\right)
=
f(v(x))v'(x)-f(u(x))u'(x).
\]

\begin{example}
Sea
\[
F(x)=\int_1^{x^2}\ln t\,dt,
\qquad x>0.
\]
Sin calcular la integral,
\[
F'(x)=\ln(x^2)\,2x=4x\ln x.
\]
\end{example}

\subsection{Interpretación de la integral definida}

La misma integral definida representa distintas magnitudes dependiendo de las unidades y del significado de $f$.

Si $f(x)\geq0$, el área entre la gráfica, el eje $OX$ y las rectas $x=a$, $x=b$ es
\[
A=\int_a^b f(x)\,dx.
\]
Si $f$ cambia de signo, el área geométrica total se obtiene integrando $|f|$ o dividiendo el intervalo por los ceros de $f$.

El área comprendida entre dos curvas $y=f(x)$ e $y=g(x)$, con $f\geq g$, es
\[
A=\int_a^b\bigl(f(x)-g(x)\bigr)\,dx.
\]
Si las curvas intercambian su posición, debe localizarse previamente sus puntos de corte y dividir el intervalo.

\begin{solvedexercise}
Calcular el área del recinto limitado por
\[
y=x^2
\qquad\text{y}\qquad
y=2x.
\]
\tcblower
\textbf{Resolución.}\par\smallskip
Los puntos de corte satisfacen
\[
x^2=2x,
\]
de donde $x=0$ y $x=2$. En $(0,2)$ se cumple $2x\geq x^2$. Por tanto,
\[
A=\int_0^2(2x-x^2)\,dx
=
\left[x^2-\frac{x^3}{3}\right]_0^2
=
\boxed{\frac43}.
\]
\end{solvedexercise}

Cuando una magnitud $r(t)$ representa una tasa, la integral
\[
\int_{t_1}^{t_2}r(t)\,dt
\]
representa el cambio neto de la cantidad asociada entre $t_1$ y $t_2$.

\begin{example}[Carga eléctrica]
Si la intensidad que atraviesa un conductor es
\[
I(t)=2+0.5t\qquad (\mathrm{A}),
\]
para $0\leq t\leq10$ s, la carga que atraviesa una sección del conductor durante esos diez segundos es
\[
Q=\int_0^{10}I(t)\,dt
=
\int_0^{10}(2+0.5t)\,dt
=45\ \mathrm{C}.
\]
Las unidades confirman la interpretación:
\[
\mathrm{A}\cdot\mathrm{s}=\mathrm{C}.
\]
\end{example}

Si $v(t)$ es una velocidad, entonces
\[
\int_{t_1}^{t_2}v(t)\,dt
\]
es el \emph{desplazamiento}, mientras que
\[
\int_{t_1}^{t_2}|v(t)|\,dt
\]
es la \emph{distancia total recorrida}. La diferencia entre ambas expresiones refleja el carácter con signo de la integral.

\begin{example}[Trabajo de una fuerza variable]
Una fuerza actúa en la dirección del movimiento y depende del desplazamiento según
\[
F(x)=100+20x\qquad (\mathrm{N}).
\]
El trabajo realizado al mover el punto de aplicación desde $x=0$ hasta $x=3$ m es
\[
W=\int_0^3F(x)\,dx
=
\int_0^3(100+20x)\,dx
=390\ \mathrm{J}.
\]
De nuevo, las unidades son coherentes: $\mathrm{N}\cdot\mathrm{m}=\mathrm{J}$.
\end{example}

El valor medio de una función integrable en $[a,b]$ se define como
\[
\boxed{f_{\mathrm{med}}=\frac1{b-a}\int_a^b f(x)\,dx.}
\]
Si $f$ es continua, existe al menos un punto $c\in[a,b]$ tal que
\[
f(c)=f_{\mathrm{med}}.
\]
Geométricamente, el rectángulo de base $b-a$ y altura $f_{\mathrm{med}}$ tiene la misma área con signo que la región representada por la integral.

También pueden expresarse mediante integrales algunas magnitudes geométricas que aparecen en modelos de diseño. Si $f\geq0$ es suficientemente regular en $[a,b]$, el volumen obtenido al girar la región bajo $y=f(x)$ alrededor del eje $OX$ es
\[
V=\pi\int_a^b f(x)^2\,dx.
\]
La longitud del arco de $y=f(x)$ entre $a$ y $b$ es
\[
L=\int_a^b\sqrt{1+\bigl(f'(x)\bigr)^2}\,dx,
\]
y el área de la superficie generada al girar la curva alrededor del eje $OX$ es
\[
S=2\pi\int_a^b f(x)\sqrt{1+\bigl(f'(x)\bigr)^2}\,dx.
\]
Estas fórmulas muestran cómo una magnitud global se construye sumando contribuciones locales cada vez más pequeñas.

\begin{example}[Volumen de un depósito idealizado]
El perfil de un recipiente de revolución viene dado por
\[
r(x)=\sqrt{1-\frac{x^2}{4}},
\qquad -2\leq x\leq2,
\]
donde $r$ y $x$ se miden en metros. Al girar el perfil alrededor del eje $x$, el volumen es
\[
V=\pi\int_{-2}^2r(x)^2\,dx
=
\pi\int_{-2}^2\left(1-\frac{x^2}{4}\right)dx
=
\frac{8\pi}{3}\ \mathrm{m}^3.
\]
\end{example}

Para comprobar integrales simbólicas con Mathematica utilizaremos principalmente \texttt{Integrate}. Cuando la integral definida no admita una expresión elemental sencilla o interese únicamente un valor numérico, puede emplearse \texttt{NIntegrate}. El objetivo de estas herramientas es comprobar o explorar cálculos, no sustituir la elección razonada del método de integración.

\begin{mathematicacode}
Integrate[x Exp[-x], x]
Integrate[x Exp[-x], {x, 0, 4}]
NIntegrate[Exp[-x^2], {x, 0, 1}]
\end{mathematicacode}

\section{Integrales impropias}
\label{sec:impropias}

La integral definida se ha introducido en intervalos cerrados y para funciones acotadas. En numerosos problemas aparecen intervalos infinitos o funciones que presentan una discontinuidad infinita. Las \emph{integrales impropias} extienden el concepto de integral definida a estas situaciones mediante límites.

\subsection{Intervalos infinitos y singularidades}

\begin{definition}[Integral impropia en un intervalo infinito]
Si $f$ es integrable en todo intervalo $[a,b]$ con $b>a$, se define
\[
\int_a^{\infty}f(x)\,dx
=
\lim_{b\to\infty}\int_a^b f(x)\,dx,
\]
si el límite existe y es finito. En ese caso la integral es \emph{convergente}; en caso contrario es \emph{divergente}.
\end{definition}

Análogamente,
\[
\int_{-\infty}^b f(x)\,dx
=
\lim_{a\to-\infty}\int_a^b f(x)\,dx.
\]
Para una integral sobre toda la recta se elige un punto $c$ y se exige la convergencia de ambas partes:
\[
\int_{-\infty}^{\infty}f(x)\,dx
=
\int_{-\infty}^{c}f(x)\,dx+
\int_c^{\infty}f(x)\,dx.
\]
No es correcto sustituir esta definición por un único límite simétrico, porque podría existir cancelación entre dos divergencias.

\begin{example}
Estudiar
\[
\int_1^{\infty}\frac{1}{x^2}\,dx.
\]
Por definición,
\[
\int_1^{\infty}\frac{1}{x^2}\,dx
=
\lim_{b\to\infty}\left[-\frac1x\right]_1^b
=
\lim_{b\to\infty}\left(1-\frac1b\right)=1.
\]
Por tanto, la integral es convergente.
\end{example}

En cambio,
\[
\int_1^{\infty}\frac1x\,dx
=
\lim_{b\to\infty}\ln b
=\infty,
\]
por lo que diverge. Aunque ambas funciones tienden a cero, lo hacen con distinta rapidez. Que el integrando tienda a cero es necesario para la convergencia en muchas situaciones, pero no es suficiente.

Un resultado básico es
\[
\boxed{\int_1^{\infty}\frac{1}{x^p}\,dx
\text{ converge si y sólo si }p>1.}
\]

También puede aparecer una singularidad en un extremo del intervalo.

\begin{definition}[Integral impropia por una singularidad]
Si $f$ no está acotada cuando $x\to a^+$, se define
\[
\int_a^b f(x)\,dx
=
\lim_{c\to a^+}\int_c^b f(x)\,dx,
\]
si el límite existe y es finito.
\end{definition}

Si la singularidad está en un punto interior $c\in(a,b)$, es obligatorio separar la integral:
\[
\int_a^b f(x)\,dx
=
\int_a^c f(x)\,dx+
\int_c^b f(x)\,dx,
\]
y ambas integrales impropias deben converger.

\begin{solvedexercise}
Estudiar la convergencia de
\[
\int_0^1\frac{1}{\sqrt{x}}\,dx.
\]
\tcblower
\textbf{Resolución.}\par\smallskip
La función no está acotada cuando $x\to0^+$. Por definición,
\[
\int_0^1x^{-1/2}\,dx
=
\lim_{a\to0^+}\int_a^1x^{-1/2}\,dx
=
\lim_{a\to0^+}\left[2\sqrt{x}\right]_a^1.
\]
Así,
\[
\int_0^1\frac{1}{\sqrt{x}}\,dx
=
\lim_{a\to0^+}(2-2\sqrt a)=2.
\]
La integral es convergente aunque el integrando tienda a $+\infty$ en el extremo izquierdo.
\end{solvedexercise}

En general,
\[
\boxed{\int_0^1\frac{1}{x^p}\,dx
\text{ converge si y sólo si }p<1.}
\]

Para funciones no negativas, los criterios de comparación permiten decidir la convergencia sin calcular exactamente la integral. Si
\[
0\leq f(x)\leq g(x)
\]
para $x$ suficientemente grande y $\int_a^\infty g$ converge, entonces $\int_a^\infty f$ también converge. Inversamente, si
\[
0\leq g(x)\leq f(x)
\]
y $\int_a^\infty g$ diverge, entonces $\int_a^\infty f$ diverge.

\begin{example}
Para $x\geq1$,
\[
0<\frac{1}{x^2+1}<\frac{1}{x^2}.
\]
Como
\[
\int_1^{\infty}\frac{1}{x^2}\,dx
\]
converge, también converge
\[
\int_1^{\infty}\frac{1}{x^2+1}\,dx.
\]
En este caso puede calcularse además exactamente:
\[
\int_1^{\infty}\frac{1}{x^2+1}\,dx
=
\left[\arctan x\right]_1^{\infty}
=
\frac{\pi}{4}.
\]
\end{example}

Las integrales impropias aparecen, por ejemplo, al modelar procesos que se prolongan indefinidamente. Si una potencia disipada viene dada por
\[
P(t)=P_0e^{-kt},\qquad k>0,
\]
la energía total disipada desde $t=0$ hasta tiempos arbitrariamente grandes es
\[
E=\int_0^{\infty}P_0e^{-kt}\,dt
=\frac{P_0}{k},
\]
que es finita pese a que el proceso idealizado se extiende hasta $t=\infty$.

\section*{Ejercicios propuestos}
\addcontentsline{toc}{section}{Ejercicios propuestos}

\begin{enumerate}

\item Calcular las siguientes integrales inmediatas:
\begin{enumerate}[label=\alph*)]
\item $\displaystyle\int e^{2x}\,dx$;
\item $\displaystyle\int\frac{10\pi}{x}\,dx$;
\item $\displaystyle\int x^{3/2}\,dx$;
\item $\displaystyle\int\left(4-5e^{-3x}-\cos2x+\sin5x\right)dx$;
\item $\displaystyle\int\frac{2e^{0.1x}}{(1+e^{0.1x})^2}\,dx$.
\end{enumerate}

\item Calcular, aplicando un cambio de variable adecuado cuando sea necesario:
\begin{enumerate}[label=\alph*)]
\item $\displaystyle\int(2x+1)e^{x^2+x}\,dx$;
\item $\displaystyle\int x\sqrt{1-x^2}\,dx$;
\item $\displaystyle\int\frac{x}{1-x^4}\,dx$;
\item $\displaystyle\int\frac{dx}{x\sqrt{1-(\ln x)^2}}$;
\item $\displaystyle\int\frac{e^{3x}}{\sqrt{4-e^{6x}}}\,dx$;
\item $\displaystyle\int \sin^5x\cos x\,dx$.
\end{enumerate}

\item Calcular mediante integración por partes:
\begin{enumerate}[label=\alph*)]
\item $\displaystyle\int\ln x\,dx$;
\item $\displaystyle\int xe^{2x}\,dx$;
\item $\displaystyle\int x^2e^{3x}\,dx$;
\item $\displaystyle\int x\sin x\,dx$;
\item $\displaystyle\int x^2\ln x\,dx$;
\item $\displaystyle\int\arctan x\,dx$;
\item $\displaystyle\int x^2\arctan x\,dx$;
\item $\displaystyle\int e^x\cos x\,dx$.
\end{enumerate}

\item Calcular las integrales racionales:
\begin{enumerate}[label=\alph*)]
\item $\displaystyle\int\frac{2x}{(x-2)(x+1)}\,dx$;
\item $\displaystyle\int\frac{2x}{(x-2)(x+1)^2}\,dx$;
\item $\displaystyle\int\frac{5x+10}{x(x+10)}\,dx$;
\item $\displaystyle\int\frac{3x+5}{x^3-x^2-x+1}\,dx$;
\item $\displaystyle\int\frac{x-1}{x^3+x}\,dx$;
\item $\displaystyle\int\frac{2x^3-4x^2-x-3}{x^2-2x-3}\,dx$;
\item $\displaystyle\int\frac{dx}{(1+x^2)^2}$.
\end{enumerate}
En los apartados con factores repetidos, comparar la descomposición completa en fracciones simples con una reducción de Hermite cuando proceda.

\item Calcular:
\begin{enumerate}[label=\alph*)]
\item $\displaystyle\int\cos^2x\,dx$;
\item $\displaystyle\int\sin^5x\,dx$;
\item $\displaystyle\int\cos^3x\,dx$;
\item $\displaystyle\int\sin^4x\,dx$;
\item $\displaystyle\int\sec x\,dx$;
\item $\displaystyle\int\sec^3x\,dx$;
\item $\displaystyle\int\sin^3x\cos^4x\,dx$.
\end{enumerate}

\item Calcular las integrales irracionales:
\begin{enumerate}[label=\alph*)]
\item $\displaystyle\int\sqrt{16-x^2}\,dx$;
\item $\displaystyle\int\sqrt{x^2+4}\,dx$;
\item $\displaystyle\int\sqrt{x^2-25}\,dx$;
\item $\displaystyle\int\frac{dx}{\sqrt{4-x^2}}$;
\item $\displaystyle\int\frac{dx}{\sqrt{x^2-9}}$;
\item $\displaystyle\int\frac{\sqrt{2x+1}}{x+1}\,dx$.
\end{enumerate}

\item Calcular las integrales definidas:
\begin{enumerate}[label=\alph*)]
\item $\displaystyle\int_1^2(x^2+x+1)\,dx$;
\item $\displaystyle\int_0^2e^{3x}\,dx$;
\item $\displaystyle\int_1^2xe^{x^2}\,dx$;
\item $\displaystyle\int_0^1x^{3/2}\,dx$;
\item $\displaystyle\int_{-1}^{1}\frac{x}{1+x^2}\,dx$;
\item $\displaystyle\int_0^1(e^x+e^{-x})\,dx$.
\end{enumerate}
Antes de integrar, aprovechar simetrías cuando sea posible.

\item Calcular el área limitada por la curva
\[
y=x-1,
\]
el eje $OX$ y las rectas $x=0$ y $x=5$. Distinguir entre integral con signo y área geométrica.

\item Calcular el área limitada por las curvas
\[
y=x+2
\qquad\text{y}\qquad
y=4-x^2.
\]
Representar primero la región y determinar analíticamente los puntos de corte.

\item Calcular el área comprendida entre
\[
y=1,\qquad y=x,\qquad y=1-x^2.
\]
Justificar cómo se divide el intervalo de integración.

\item Sea
\[
F(x)=\int_1^x\bigl(\sin^2t+\ln t\bigr)\,dt,
\qquad x>0.
\]
Calcular $F'(x)$ sin obtener previamente una primitiva del integrando.

\item Calcular la derivada de las funciones:
\begin{enumerate}[label=\alph*)]
\item $\displaystyle F(x)=\int_0^{x^2}(1+\cos^2t)\,dt$;
\item $\displaystyle G(x)=\int_{2x}^{x^2}\ln(1+t^2)\,dt$;
\item $\displaystyle H(x)=\int_{\sin x}^{e^x}\frac{dt}{1+t^2}$.
\end{enumerate}

\item Sea
\[
F(x)=\int_0^x(t^3-4t)\,dt.
\]
Sin calcular primero $F$:
\begin{enumerate}[label=\alph*)]
\item determinar los intervalos de crecimiento y decrecimiento de $F$;
\item localizar sus extremos relativos;
\item estudiar la convexidad de $F$;
\item localizar sus puntos de inflexión;
\item realizar un esbozo cualitativo de $F$ sabiendo que $F(0)=0$.
\end{enumerate}
Comprobar después las conclusiones obteniendo una expresión analítica de $F$.

\item Una función continua $f$ satisface: $f(x)>0$ para $x<1$, $f(1)=0$, $f(x)<0$ para $1<x<4$, $f(4)=0$ y $f(x)>0$ para $x>4$. Además, $f$ es creciente hasta $x=2$ y decreciente a partir de $x=2$. Para
\[
F(x)=\int_0^x f(t)\,dt,
\]
traducir esta información a crecimiento, decrecimiento, extremos y convexidad de $F$. Realizar un esbozo compatible de $F$ sin buscar una fórmula para $f$.

\item La velocidad de una partícula es
\[
v(t)=t^2-t-6,
\qquad 0\leq t\leq4.
\]
\begin{enumerate}[label=\alph*)]
\item calcular el desplazamiento durante los cuatro primeros segundos;
\item calcular la distancia total recorrida;
\item explicar por qué los dos resultados son distintos.
\end{enumerate}

\item La intensidad que circula por un dispositivo es
\[
I(t)=5e^{-0.2t}\ \mathrm{A}.
\]
Calcular la carga que atraviesa el dispositivo entre $t=0$ y $t=10$ s. Determinar también la carga total que circularía en el modelo si el proceso se prolongara indefinidamente.

\item En una conducción, el caudal viene dado por
\[
q(t)=20+4\sin\left(\frac{\pi t}{12}\right)
\qquad (\mathrm{m}^3/\mathrm{h}),
\]
durante $0\leq t\leq24$ h. Determinar el volumen total transportado y el caudal medio durante ese día.

\item Una planta química alimenta un reactor con una corriente cuya tasa neta de entrada de masa es
\[
r(t)=3te^{-0.5t}\qquad (\mathrm{kg/min}).
\]
Si inicialmente hay $20$ kg de sustancia en el reactor, obtener la masa $M(t)$ y calcular $M(5)$. Interpretar la constante de integración y las unidades de cada magnitud.

\item La potencia eléctrica suministrada a una resistencia durante un transitorio es
\[
P(t)=120e^{-0.4t}\qquad (\mathrm{W}).
\]
\begin{enumerate}[label=\alph*)]
\item calcular la energía suministrada durante los primeros $10$ s;
\item calcular la energía total que predice el modelo para $t\geq0$;
\item explicar por qué el segundo cálculo conduce de forma natural a una integral impropia.
\end{enumerate}

\item Una fuerza elástica no lineal viene dada por
\[
F(x)=800x+120x^3\qquad (\mathrm{N}),
\]
donde $x$ se mide en metros. Calcular el trabajo necesario para deformar el elemento desde $x=0$ hasta $x=0.1$ m.

\item Un eje de longitud $L=2$ m presenta una densidad lineal variable
\[
\rho(x)=4+x^2\qquad (\mathrm{kg/m}).
\]
Calcular su masa total y su densidad lineal media.

\item El perfil interior de un depósito de revolución está dado por
\[
r(x)=1+\frac{x}{4},
\qquad 0\leq x\leq4,
\]
en metros. Calcular el volumen del depósito idealizado al girar el perfil alrededor de su eje longitudinal.

\item Calcular la longitud del arco de
\[
y=\frac23x^{3/2},
\qquad 0\leq x\leq3.
\]

\item Estudiar la convergencia y, cuando sea posible, calcular el valor de:
\begin{enumerate}[label=\alph*)]
\item $\displaystyle\int_1^{\infty}\frac{dx}{x^2}$;
\item $\displaystyle\int_1^{\infty}\frac{dx}{x}$;
\item $\displaystyle\int_0^1\frac{dx}{\sqrt{x}}$;
\item $\displaystyle\int_0^1\frac{dx}{x^{3/2}}$;
\item $\displaystyle\int_0^{\infty}e^{-2x}\,dx$;
\item $\displaystyle\int_{-\infty}^{\infty}e^{-|x|}\,dx$.
\end{enumerate}

\item Determinar para qué valores de $p\in\R$ convergen:
\begin{enumerate}[label=\alph*)]
\item $\displaystyle\int_1^{\infty}\frac{dx}{x^p}$;
\item $\displaystyle\int_0^1\frac{dx}{x^p}$.
\end{enumerate}
Explicar por qué las condiciones obtenidas son diferentes.

\item Estudiar la convergencia de
\[
\int_1^{\infty}\frac{1}{x^2+x}\,dx
\]
mediante comparación y comprobar después el resultado calculando la integral.

\item Estudiar la integral
\[
\int_{-1}^{2}\frac{dx}{x-1}.
\]
Explicar por qué no es correcto aplicar directamente la regla de Barrow a una primitiva de $1/(x-1)$ en todo el intervalo.

\item Una función $f$ es positiva, continua y decreciente para $x\geq1$, y se sabe que
\[
\frac{1}{2x^2}\leq f(x)\leq\frac{3}{x^2}.
\]
Decidir si $\int_1^\infty f(x)\,dx$ converge y justificar la respuesta sin calcular una primitiva de $f$.

\item \M Seleccionar cuatro integrales de los ejercicios 2--6, resolverlas manualmente y comprobar el resultado mediante \texttt{Integrate}. En cada caso verificar la primitiva derivándola con \texttt{D}. Si Mathematica devuelve una expresión distinta de la obtenida a mano, utilizar \texttt{Simplify} para estudiar si son equivalentes en el dominio considerado.

\item \M Representar con \texttt{Plot} las funciones
\[
f(x)=x-1,
\qquad
g(x)=4-x^2,
\]
en intervalos adecuados y utilizar la gráfica para decidir cómo deben plantearse los ejercicios 8 y 9. Calcular después las áreas con \texttt{Integrate}. La solución entregada debe contener primero el planteamiento analítico de cada integral.

\item \M Para
\[
F(x)=\int_0^x(t^3-4t)\,dt,
\]
obtener una expresión con \texttt{Integrate}, derivarla con \texttt{D} y representar en gráficos separados el integrando y $F$. Relacionar los signos del integrando con la monotonía de $F$ y los cambios de monotonía del integrando con la convexidad de $F$.

\item \M Comparar numéricamente
\[
\int_0^1e^{-x^2}\,dx
\]
mediante \texttt{NIntegrate} con aproximaciones por sumas de rectángulos usando $n=10$, $50$ y $200$ subintervalos. No es necesario introducir nuevos comandos de Mathematica: las sumas pueden calcularse con una tabla preparada por el profesor o de forma manual/semiautomática. Interpretar qué ocurre al aumentar $n$.

\item \M Utilizar \texttt{NIntegrate} para aproximar
\[
\int_0^{\infty}e^{-x^2}\,dx.
\]
Comparar con la aproximación de $\int_0^b e^{-x^2}\,dx$ para $b=1,2,3,4$ y explicar verbalmente el significado de la convergencia observada.

\item (Examen) Calcular
\[
\int\frac{2x^3-4x^2-x-3}{x^2-2x-3}\,dx.
\]

\item (Examen) Calcular
\[
\int x^2\arctan x\,dx.
\]

\item (Examen) Calcular el área del recinto situado en el primer cuadrante y limitado por
\[
y=3x,
\qquad y=\sqrt{x+1},
\qquad y=1-x^2.
\]
Determinar cuidadosamente qué curvas delimitan cada parte del recinto antes de integrar.

\item (Examen) Calcular
\[
\int\frac{x}{\sqrt{4-x^2}}\,dx.
\]

\item (Examen) Sea
\[
f(x)=10+3xe^{-0.2x}.
\]
Calcular el área bajo la curva, el eje $OX$ y las rectas $x=0$ y $x=10$. Justificar previamente que $f(x)$ es positiva en ese intervalo.

\item Una empresa gestiona una reserva inicial de $100$ millones de unidades de material. La tasa de extracción es
\[
R(t)=3te^{-0.6t}
\]
millones de unidades por año, donde $t$ se mide en años desde el inicio del plan.
\begin{enumerate}[label=\alph*)]
\item interpretar $R(2)$ y sus unidades;
\item obtener la cantidad extraída durante los primeros dos años;
\item obtener una expresión para la reserva restante $S(t)$;
\item calcular la cantidad total que extraería el modelo si se mantuviera indefinidamente y discutir si agotaría la reserva.
\end{enumerate}

\end{enumerate}

\cleardoublepage
\chapter{Matrices}
\label{ch:matrices}

Los sistemas de ecuaciones lineales aparecen de forma natural cuando varias magnitudes están relacionadas simultáneamente. Un balance de materia puede imponer varias ecuaciones sobre caudales desconocidos; las leyes de Kirchhoff conducen a sistemas para corrientes y tensiones; las ecuaciones de equilibrio relacionan fuerzas y reacciones; y muchos modelos de producción, transporte o energía se expresan de manera compacta mediante tablas de coeficientes.

Las matrices proporcionan el lenguaje adecuado para organizar esos datos y operar con ellos. En este capítulo estudiaremos las operaciones matriciales, los determinantes y el rango, y utilizaremos estas herramientas para analizar y resolver sistemas de ecuaciones lineales. El objetivo no es únicamente ejecutar algoritmos: interesa también reconocer qué información aporta cada objeto matemático y cuándo un problema tiene una solución única, infinitas soluciones o ninguna.

\section{Matrices y operaciones}
\label{sec:matrices-operaciones}

\subsection{Concepto de matriz y tipos básicos}

\begin{definition}[Matriz]
Una \emph{matriz} de tamaño $m\times n$ es una disposición rectangular de $mn$ números reales organizada en $m$ filas y $n$ columnas. Se escribe
\[
A=(a_{ij})_{m\times n}
=
\begin{pmatrix}
 a_{11}&a_{12}&\cdots&a_{1n}\\
 a_{21}&a_{22}&\cdots&a_{2n}\\
 \vdots&\vdots&\ddots&\vdots\\
 a_{m1}&a_{m2}&\cdots&a_{mn}
\end{pmatrix},
\]
donde $a_{ij}$ es el elemento situado en la fila $i$ y la columna $j$.
\end{definition}

El conjunto de matrices reales de tamaño $m\times n$ se denota por $\mathcal M_{m\times n}(\R)$. Cuando $m=n$ se habla de una \emph{matriz cuadrada de orden $n$} y escribiremos simplemente $\mathcal M_n(\R)$.

Una matriz puede utilizarse para almacenar simultáneamente varias magnitudes. Por ejemplo, una planta fabrica dos productos $P_1$ y $P_2$ a partir de tres materias primas $M_1$, $M_2$ y $M_3$. Si las cantidades necesarias por unidad de producto son
\[
A=
\begin{pmatrix}
3&4\\
2&1\\
4&3
\end{pmatrix},
\]
las columnas describen los productos y las filas las materias primas. El elemento $a_{32}=3$ significa que cada unidad de $P_2$ consume tres unidades de $M_3$. La matriz no es sólo una tabla: más adelante veremos que sus operaciones reproducen operaciones reales sobre el proceso.

Entre los tipos de matrices que utilizaremos con frecuencia se encuentran los siguientes:

\begin{itemize}
\item una \emph{matriz fila} tiene una sola fila y una \emph{matriz columna} una sola columna;
\item la \emph{matriz nula}, denotada por $0$, tiene todos sus elementos iguales a cero;
\item una matriz cuadrada es \emph{diagonal} si todos los elementos fuera de la diagonal principal son nulos;
\item la \emph{matriz identidad} de orden $n$ es
\[
I_n=\begin{pmatrix}
1&0&\cdots&0\\
0&1&\cdots&0\\
\vdots&\vdots&\ddots&\vdots\\
0&0&\cdots&1
\end{pmatrix};
\]
\item una matriz es \emph{triangular superior} si $a_{ij}=0$ cuando $i>j$, y \emph{triangular inferior} si $a_{ij}=0$ cuando $i<j$;
\item una matriz cuadrada es \emph{simétrica} si $A^t=A$ y \emph{antisimétrica} si $A^t=-A$.
\end{itemize}

La suma de los elementos de la diagonal principal de una matriz cuadrada se denomina \emph{traza}:
\[
\operatorname{tr}(A)=a_{11}+a_{22}+\cdots+a_{nn}.
\]

\begin{example}
Para
\[
A=
\begin{pmatrix}
2&-1&0\\
-1&3&4\\
0&4&5
\end{pmatrix},
\]
tenemos $A^t=A$, por lo que $A$ es simétrica, y
\[
\operatorname{tr}(A)=2+3+5=10.
\]
Las matrices simétricas aparecerán de nuevo en el estudio de formas cuadráticas y diagonalización.
\end{example}

\subsection{Suma, producto por escalares y producto de matrices}

Dos matrices sólo pueden sumarse si tienen el mismo tamaño. Si $A=(a_{ij})$ y $B=(b_{ij})$ pertenecen a $\mathcal M_{m\times n}(\R)$, definimos
\[
A+B=(a_{ij}+b_{ij}).
\]
Si $\lambda\in\R$,
\[
\lambda A=(\lambda a_{ij}).
\]
Estas operaciones satisfacen las propiedades habituales de suma y producto por escalares: asociatividad, conmutatividad de la suma, existencia de matriz nula y opuesta, y distributividad.

\begin{example}[Variación de una producción]
Una fábrica registra la producción mensual de dos modelos con tres acabados mediante
\[
P=
\begin{pmatrix}
500&250&70\\
600&150&100
\end{pmatrix}.
\]
Si se planifica un aumento uniforme del $20\%$, la nueva producción prevista queda representada por
\[
1.2P=
\begin{pmatrix}
600&300&84\\
720&180&120
\end{pmatrix}.
\]
Si otra planta tiene una matriz de producción $Q$ con el mismo tamaño, la producción conjunta será $P+Q$.
\end{example}

El producto de matrices es diferente. Su definición está diseñada para encadenar relaciones lineales.

\begin{definition}[Producto de matrices]
Sean $A\in\mathcal M_{m\times n}(\R)$ y $B\in\mathcal M_{n\times p}(\R)$. El producto $AB$ es la matriz de tamaño $m\times p$ cuyo elemento $(i,j)$ es
\[
(AB)_{ij}=\sum_{k=1}^{n}a_{ik}b_{kj}.
\]
Por tanto, el elemento $(i,j)$ se obtiene multiplicando la fila $i$ de $A$ por la columna $j$ de $B$ y sumando los productos.
\end{definition}

La condición importante es que el número de columnas de la primera matriz coincida con el número de filas de la segunda:
\[
(m\times n)(n\times p)\longrightarrow(m\times p).
\]

\begin{example}[Consumo de materias primas]
Retomemos la matriz
\[
A=
\begin{pmatrix}
3&4\\
2&1\\
4&3
\end{pmatrix},
\]
que contiene el consumo de tres materias primas por unidad de dos productos. Si la producción semanal es
\[
x=\begin{pmatrix}20\\30\end{pmatrix},
\]
el consumo total de materias primas es
\[
Ax=
\begin{pmatrix}
3&4\\
2&1\\
4&3
\end{pmatrix}
\begin{pmatrix}20\\30\end{pmatrix}
=
\begin{pmatrix}
180\\70\\170
\end{pmatrix}.
\]
La primera componente, por ejemplo, significa que se consumen $180$ unidades de $M_1$ por semana.

Si además el vector fila de costes unitarios es
\[
c=\begin{pmatrix}6&10&12\end{pmatrix},
\]
entonces
\[
cA
=
\begin{pmatrix}86&70\end{pmatrix}
\]
proporciona el coste de materias primas por unidad de $P_1$ y $P_2$. Finalmente,
\[
cAx
\]
representa el coste total semanal. La asociación de productos matriciales reproduce así la secuencia física
\[
\text{producción}\longrightarrow\text{consumos}\longrightarrow\text{coste}.
\]
\end{example}

El producto matricial es asociativo y distributivo respecto de la suma:
\[
A(BC)=(AB)C,
\]
\[
A(B+C)=AB+AC,
\qquad
(A+B)C=AC+BC.
\]
Sin embargo, en general no es conmutativo:
\[
AB\neq BA.
\]
Incluso puede ocurrir que uno de los dos productos exista y el otro no.

\begin{example}
Sean
\[
A=\begin{pmatrix}1&2\\0&1\end{pmatrix},
\qquad
B=\begin{pmatrix}0&1\\1&0\end{pmatrix}.
\]
Entonces
\[
AB=\begin{pmatrix}2&1\\1&0\end{pmatrix},
\qquad
BA=\begin{pmatrix}0&1\\1&2\end{pmatrix},
\]
y por tanto $AB\neq BA$.
\end{example}

Para una matriz cuadrada $A$ se definen las potencias
\[
A^2=AA,
\qquad
A^3=AAA,
\qquad\ldots,
\qquad
A^0=I.
\]
Esta notación aparecerá de forma natural en modelos que evolucionan por etapas. Si un vector de estado $x_k$ cambia según
\[
x_{k+1}=Ax_k,
\]
entonces
\[
x_2=A^2x_0,
\qquad
x_k=A^kx_0.
\]
Más adelante, la diagonalización permitirá calcular potencias elevadas de matrices de manera eficiente.

\subsection{Matriz transpuesta, inversa y ecuaciones matriciales}

\begin{definition}[Matriz transpuesta]
Si $A=(a_{ij})\in\mathcal M_{m\times n}(\R)$, su \emph{matriz transpuesta} $A^t\in\mathcal M_{n\times m}(\R)$ se obtiene intercambiando filas por columnas:
\[
(A^t)_{ij}=a_{ji}.
\]
\end{definition}

Se cumplen
\[
(A^t)^t=A,
\qquad
(A+B)^t=A^t+B^t,
\qquad
(\lambda A)^t=\lambda A^t,
\]
y, de manera especialmente importante,
\[
(AB)^t=B^tA^t.
\]
Obsérvese que al transponer un producto se invierte el orden de los factores.

\begin{definition}[Matriz inversa]
Sea $A\in\mathcal M_n(\R)$. Se dice que $A$ es \emph{invertible} o \emph{regular} si existe una matriz $A^{-1}\in\mathcal M_n(\R)$ tal que
\[
AA^{-1}=A^{-1}A=I_n.
\]
Si no existe, $A$ se denomina \emph{singular}.
\end{definition}

La inversa, cuando existe, es única. Además,
\[
(A^{-1})^{-1}=A,
\qquad
(AB)^{-1}=B^{-1}A^{-1}
\]
si $A$ y $B$ son invertibles.

\begin{example}
Sea
\[
A=\begin{pmatrix}2&1\\1&1\end{pmatrix}.
\]
La matriz
\[
B=\begin{pmatrix}1&-1\\-1&2\end{pmatrix}
\]
satisface
\[
AB=BA=I_2.
\]
Por tanto,
\[
A^{-1}=B.
\]
En la siguiente sección obtendremos procedimientos sistemáticos para decidir si una matriz es invertible y calcular su inversa.
\end{example}

Las ecuaciones matriciales deben manipularse respetando el orden de los productos. Por ejemplo,
\[
AX=B
\]
con $A$ invertible implica
\[
A^{-1}AX=A^{-1}B
\quad\Longrightarrow\quad
X=A^{-1}B.
\]
En cambio,
\[
XA=B
\]
conduce a
\[
X=BA^{-1}.
\]
No podemos intercambiar estos resultados porque, en general, el producto matricial no es conmutativo.

\begin{solvedexercise}
Suponiendo que todas las matrices inversas que aparecen existen, despejar $X$ en
\[
A(X+C)=BX.
\]
\tcblower
\textbf{Resolución.}\par\smallskip
Desarrollamos el primer miembro:
\[
AX+AC=BX.
\]
Reunimos los términos que contienen $X$:
\[
AX-BX=-AC.
\]
Como $X$ aparece a la derecha de ambos factores,
\[
(A-B)X=-AC.
\]
Si $A-B$ es invertible,
\[
\boxed{X=-(A-B)^{-1}AC.}
\]
No sería correcto escribir $X=-AC(A-B)^{-1}$ sin una justificación adicional.
\end{solvedexercise}

\section{Determinantes y rango}
\label{sec:determinantes-rango}

\subsection{Determinantes}

A cada matriz cuadrada se le puede asociar un número, su determinante. Este número contiene información decisiva sobre la matriz: permite detectar si es invertible, interviene en el cálculo de la inversa y ayuda a estudiar el rango y los sistemas lineales.

\begin{definition}[Determinante de orden $2$]
Para
\[
A=\begin{pmatrix}a&b\\c&d\end{pmatrix},
\]
se define
\[
\det(A)=|A|=ad-bc.
\]
\end{definition}

Para matrices de orden superior utilizaremos menores y adjuntos.

\begin{definition}[Menor y adjunto]
Sea $A=(a_{ij})\in\mathcal M_n(\R)$. El \emph{menor complementario} $M_{ij}$ del elemento $a_{ij}$ es el determinante de la matriz de orden $n-1$ que se obtiene eliminando la fila $i$ y la columna $j$ de $A$. El \emph{adjunto} o \emph{cofactor} del elemento $a_{ij}$ es
\[
C_{ij}=(-1)^{i+j}M_{ij}.
\]
\end{definition}

El determinante puede desarrollarse por cualquier fila o columna. Por la fila $i$,
\[
\det(A)=\sum_{j=1}^n a_{ij}C_{ij},
\]
y por la columna $j$,
\[
\det(A)=\sum_{i=1}^n a_{ij}C_{ij}.
\]

Para una matriz de orden $3$,
\[
A=
\begin{pmatrix}
a&b&c\\
d&e&f\\
g&h&i
\end{pmatrix},
\]
el desarrollo por la primera fila da
\[
\det(A)
=a(ei-fh)-b(di-fg)+c(dh-eg).
\]
La regla de Sarrus es una forma mnemotécnica equivalente para matrices $3\times3$, pero no se generaliza a órdenes superiores.

\begin{example}
Calcular
\[
\det\begin{pmatrix}
1&2&-1\\
0&3&4\\
2&1&1
\end{pmatrix}.
\]
Desarrollando por la primera fila,
\begin{align*}
\det(A)
&=1\begin{vmatrix}3&4\\1&1\end{vmatrix}
-2\begin{vmatrix}0&4\\2&1\end{vmatrix}
-\begin{vmatrix}0&3\\2&1\end{vmatrix}\\
&=(3-4)-2(0-8)-(0-6)\\
&=-1+16+6=21.
\end{align*}
\end{example}

Las propiedades siguientes son fundamentales para simplificar cálculos.

\begin{theorem}[Propiedades básicas del determinante]
Sean $A,B\in\mathcal M_n(\R)$ y $\lambda\in\R$.
\begin{enumerate}[label=\alph*)]
\item $\det(I_n)=1$.
\item $\det(A^t)=\det(A)$.
\item $\det(AB)=\det(A)\det(B)$.
\item $\det(\lambda A)=\lambda^n\det(A)$.
\item Si $A$ es triangular, $\det(A)$ es el producto de los elementos de su diagonal.
\item Intercambiar dos filas cambia el signo del determinante.
\item Multiplicar una fila por $\lambda$ multiplica el determinante por $\lambda$.
\item Sumar a una fila un múltiplo de otra no cambia el determinante.
\item Si dos filas son iguales o proporcionales, el determinante es cero.
\end{enumerate}
\end{theorem}

Las mismas propiedades son válidas trabajando por columnas.

\begin{solvedexercise}
Sea $A\in\mathcal M_n(\R)$ una matriz idempotente, es decir, $A^2=A$. Determinar los valores posibles de $\det(A)$.
\tcblower
\textbf{Resolución.}\par\smallskip
Tomando determinantes,
\[
\det(A^2)=\det(A).
\]
Como
\[
\det(A^2)=\det(A)^2,
\]
obtenemos
\[
\det(A)^2=\det(A),
\]
y por tanto
\[
\det(A)\bigl(\det(A)-1\bigr)=0.
\]
Así,
\[
\boxed{\det(A)=0\quad\text{o}\quad\det(A)=1.}
\]
Este tipo de razonamiento utiliza las propiedades estructurales del determinante sin necesidad de conocer los elementos de $A$.
\end{solvedexercise}

\subsection{Determinante, inversa y regularidad}

El determinante caracteriza completamente la invertibilidad de una matriz cuadrada.

\begin{theorem}[Criterio de invertibilidad]
Sea $A\in\mathcal M_n(\R)$. Son equivalentes:
\begin{enumerate}[label=\roman*)]
\item $A$ es invertible;
\item $\det(A)\neq0$;
\item $\operatorname{rg}(A)=n$.
\end{enumerate}
\end{theorem}

La equivalencia con el rango se justificará después de introducirlo formalmente. Por ahora, el determinante proporciona ya un criterio operativo muy simple.

\begin{definition}[Matriz adjunta]
La \emph{matriz adjunta} de $A$ se define como la transpuesta de la matriz de sus cofactores:
\[
\operatorname{Adj}(A)=
\begin{pmatrix}
C_{11}&C_{21}&\cdots&C_{n1}\\
C_{12}&C_{22}&\cdots&C_{n2}\\
\vdots&\vdots&\ddots&\vdots\\
C_{1n}&C_{2n}&\cdots&C_{nn}
\end{pmatrix}.
\]
\end{definition}

Si $\det(A)\neq0$, entonces
\[
\boxed{A^{-1}=\frac{1}{\det(A)}\operatorname{Adj}(A).}
\]

Para matrices $2\times2$ obtenemos la conocida expresión
\[
\begin{pmatrix}a&b\\c&d\end{pmatrix}^{-1}
=
\frac{1}{ad-bc}
\begin{pmatrix}d&-b\\-c&a\end{pmatrix},
\qquad ad-bc\neq0.
\]

\begin{example}
Para
\[
A=\begin{pmatrix}2&1\\3&2\end{pmatrix},
\]
se tiene
\[
\det(A)=4-3=1,
\]
y por tanto
\[
A^{-1}=\begin{pmatrix}2&-1\\-3&2\end{pmatrix}.
\]
La comprobación $AA^{-1}=I_2$ permite detectar inmediatamente posibles errores de cálculo.
\end{example}

La fórmula mediante adjuntos es útil para matrices pequeñas y para razonamientos teóricos. Para matrices de mayor tamaño, la eliminación de Gauss--Jordan suele ser más eficiente: se aplican operaciones elementales por filas a la matriz ampliada
\[
(A\mid I_n)
\]
hasta obtener, si es posible,
\[
(I_n\mid A^{-1}).
\]
Si no puede reducirse el bloque izquierdo a la identidad, la matriz no es invertible.

\subsection{Rango y eliminación de Gauss}

El rango mide cuánta información lineal independiente contiene una matriz. Antes de estudiar formalmente la independencia lineal en el capítulo siguiente, podemos definirlo mediante determinantes de submatrices.

\begin{definition}[Menor de una matriz]
Un \emph{menor de orden $r$} de una matriz $A\in\mathcal M_{m\times n}(\R)$ es el determinante de cualquier submatriz cuadrada de orden $r$ obtenida seleccionando $r$ filas y $r$ columnas de $A$.
\end{definition}

\begin{definition}[Rango]
El \emph{rango} de una matriz $A$, denotado por $\operatorname{rg}(A)$, es el mayor orden de un menor no nulo de $A$. Para la matriz nula se define $\operatorname{rg}(0)=0$.
\end{definition}

En particular,
\[
0\leq\operatorname{rg}(A)\leq\min\{m,n\}.
\]
Si $A$ es cuadrada de orden $n$,
\[
\operatorname{rg}(A)=n
\quad\Longleftrightarrow\quad
\det(A)\neq0.
\]

El cálculo mediante menores es apropiado para matrices pequeñas y para estudiar parámetros, pero la herramienta general es la eliminación de Gauss.

\begin{definition}[Operaciones elementales por filas]
Las siguientes transformaciones se denominan \emph{operaciones elementales por filas}:
\begin{enumerate}[label=\roman*)]
\item intercambiar dos filas;
\item multiplicar una fila por un escalar no nulo;
\item sumar a una fila un múltiplo de otra fila.
\end{enumerate}
\end{definition}

Estas operaciones no modifican el rango de la matriz.

\begin{definition}[Matriz escalonada]
Una matriz está en \emph{forma escalonada} si:
\begin{itemize}
\item las filas nulas, si existen, están situadas al final;
\item el primer elemento no nulo de cada fila no nula, llamado \emph{pivote}, aparece a la derecha del pivote de la fila anterior;
\item debajo de cada pivote todos los elementos son cero.
\end{itemize}
\end{definition}

El método de eliminación de Gauss consiste en transformar una matriz, mediante operaciones elementales, en una matriz escalonada. El rango es entonces el número de filas no nulas, equivalentes al número de pivotes.

\begin{example}
Calcular el rango de
\[
A=
\begin{pmatrix}
1&2&-1&0\\
2&4&1&3\\
-1&-2&2&1
\end{pmatrix}.
\]
Aplicamos
\[
F_2\leftarrow F_2-2F_1,
\qquad
F_3\leftarrow F_3+F_1,
\]
y obtenemos
\[
\begin{pmatrix}
1&2&-1&0\\
0&0&3&3\\
0&0&1&1
\end{pmatrix}.
\]
Después,
\[
F_2\leftarrow F_2-3F_3,
\]
con lo que
\[
A\sim
\begin{pmatrix}
1&2&-1&0\\
0&0&0&0\\
0&0&1&1
\end{pmatrix}.
\]
Hay dos filas no nulas y, por tanto,
\[
\boxed{\operatorname{rg}(A)=2.}
\]
\end{example}

Cuando aparecen parámetros hay que evitar operaciones que dividan por una expresión cuyo valor podría ser cero sin estudiar antes ese caso.

\begin{example}[Rango dependiente de un parámetro]
Sea
\[
A_\lambda=
\begin{pmatrix}
1&-1&1\\
2&1&3\\
5&1&\lambda
\end{pmatrix}.
\]
Calculamos primero el determinante:
\[
\det(A_\lambda)
=
\begin{vmatrix}
1&-1&1\\
2&1&3\\
5&1&\lambda
\end{vmatrix}
=3\lambda-18=3(\lambda-6).
\]
Si $\lambda\neq6$, el determinante es no nulo y
\[
\operatorname{rg}(A_\lambda)=3.
\]
Para $\lambda=6$,
\[
A_6=
\begin{pmatrix}
1&-1&1\\
2&1&3\\
5&1&6
\end{pmatrix}.
\]
El menor
\[
\begin{vmatrix}1&-1\\2&1\end{vmatrix}=3\neq0,
\]
de modo que el rango no puede ser inferior a $2$. Como el determinante de orden $3$ es nulo,
\[
\boxed{\operatorname{rg}(A_6)=2.}
\]
En resumen,
\[
\operatorname{rg}(A_\lambda)=
\begin{cases}
3,&\lambda\neq6,\\
2,&\lambda=6.
\end{cases}
\]
\end{example}

El rango también tiene una interpretación útil en modelos de medida. Si tres sensores proporcionan tres combinaciones lineales de unas mismas magnitudes desconocidas, una matriz de coeficientes de rango $3$ indica que las tres ecuaciones contienen información suficiente e independiente para determinar, en principio, las tres magnitudes. Si el rango desciende a $2$, alguna medición es redundante respecto de las demás. Esta idea quedará formalizada en el capítulo de espacios vectoriales.

\section{Sistemas de ecuaciones lineales}
\label{sec:sistemas-lineales}

\subsection{Representación matricial y clasificación}

\begin{definition}[Sistema lineal]
Un \emph{sistema lineal} de $m$ ecuaciones con $n$ incógnitas $x_1,\ldots,x_n$ es un conjunto de ecuaciones de la forma
\[
\begin{cases}
a_{11}x_1+a_{12}x_2+\cdots+a_{1n}x_n=b_1,\\
a_{21}x_1+a_{22}x_2+\cdots+a_{2n}x_n=b_2,\\
\hspace{2cm}\vdots\\
a_{m1}x_1+a_{m2}x_2+\cdots+a_{mn}x_n=b_m.
\end{cases}
\]
\end{definition}

Definiendo
\[
A=
\begin{pmatrix}
a_{11}&\cdots&a_{1n}\\
\vdots&\ddots&\vdots\\
a_{m1}&\cdots&a_{mn}
\end{pmatrix},
\qquad
x=\begin{pmatrix}x_1\\\vdots\\x_n\end{pmatrix},
\qquad
b=\begin{pmatrix}b_1\\\vdots\\b_m\end{pmatrix},
\]
el sistema puede escribirse de forma compacta como
\[
\boxed{Ax=b.}
\]
La matriz
\[
(A\mid b)
\]
se denomina \emph{matriz ampliada} del sistema.

Una solución es un vector $x$ que satisface todas las ecuaciones simultáneamente. Según el número de soluciones, el sistema se clasifica como:

\begin{itemize}
\item \emph{compatible determinado} si tiene una única solución;
\item \emph{compatible indeterminado} si tiene infinitas soluciones;
\item \emph{incompatible} si no tiene solución.
\end{itemize}

El rango permite decidir cuál de estos casos se produce sin necesidad de resolver previamente el sistema.

\begin{theorem}[Teorema de Rouché--Frobenius]
Sea un sistema lineal $Ax=b$ con $n$ incógnitas. Entonces:
\begin{enumerate}[label=\roman*)]
\item el sistema es compatible si y sólo si
\[
\operatorname{rg}(A)=\operatorname{rg}(A\mid b);
\]
\item si es compatible, es determinado si y sólo si
\[
\operatorname{rg}(A)=\operatorname{rg}(A\mid b)=n;
\]
\item si es compatible y
\[
\operatorname{rg}(A)=\operatorname{rg}(A\mid b)=r<n,
\]
entonces es indeterminado y su solución general depende de $n-r$ parámetros libres.
\end{enumerate}
\end{theorem}

\begin{example}
Consideremos
\[
\begin{cases}
x+y+z=3,\\
2x-y+z=0,\\
3x+2z=3.
\end{cases}
\]
Su matriz de coeficientes y su ampliada son
\[
A=\begin{pmatrix}
1&1&1\\
2&-1&1\\
3&0&2
\end{pmatrix},
\qquad
(A\mid b)=
\left(
\begin{array}{ccc|c}
1&1&1&3\\
2&-1&1&0\\
3&0&2&3
\end{array}
\right).
\]
Como
\[
\det(A)=-3\neq0,
\]
tenemos $\operatorname{rg}(A)=3$. En consecuencia, también $\operatorname{rg}(A\mid b)=3$ y el sistema es compatible determinado.
\end{example}

\subsection{Eliminación de Gauss, Cramer y matriz inversa}

El método de Gauss resuelve un sistema aplicando a su matriz ampliada operaciones elementales por filas. Estas operaciones corresponden a transformaciones de las ecuaciones que no modifican el conjunto de soluciones.

\begin{example}[Resolución por Gauss]
Resolver
\[
\begin{cases}
2x+3y+2z+6t=10,\\
y+2z+t=2,\\
3x-3z+6t=9.
\end{cases}
\]
La matriz ampliada es
\[
\left(
\begin{array}{cccc|c}
2&3&2&6&10\\
0&1&2&1&2\\
3&0&-3&6&9
\end{array}
\right).
\]
Para evitar fracciones, realizamos primero
\[
F_3\leftarrow2F_3-3F_1,
\]
y obtenemos
\[
\left(
\begin{array}{cccc|c}
2&3&2&6&10\\
0&1&2&1&2\\
0&-9&-12&-6&-12
\end{array}
\right).
\]
Después,
\[
F_3\leftarrow F_3+9F_2,
\]
con lo que
\[
\left(
\begin{array}{cccc|c}
2&3&2&6&10\\
0&1&2&1&2\\
0&0&6&3&6
\end{array}
\right).
\]
El sistema tiene cuatro incógnitas y tres pivotes, de modo que queda un parámetro libre. Tomando $t=s$, la tercera ecuación da
\[
6z+3s=6
\quad\Longrightarrow\quad
z=1-\frac{s}{2}.
\]
La segunda,
\[
y+2z+s=2,
\]
conduce a $y=0$. Finalmente,
\[
2x+2z+6s=10,
\]
y por tanto
\[
x=4-\frac52s.
\]
La solución general es
\[
\boxed{(x,y,z,t)=\left(4-\frac52s,\,0,\,1-\frac12s,\,s\right),\qquad s\in\R.}
\]
El sistema es compatible indeterminado.
\end{example}

Cuando el sistema tiene el mismo número de ecuaciones que de incógnitas y la matriz de coeficientes es invertible, existen otros dos métodos especialmente sencillos.

Si
\[
Ax=b,
\qquad \det(A)\neq0,
\]
entonces
\[
\boxed{x=A^{-1}b.}
\]

La regla de Cramer expresa directamente cada incógnita mediante determinantes.

\begin{theorem}[Regla de Cramer]
Sea $Ax=b$ un sistema de $n$ ecuaciones con $n$ incógnitas y $\det(A)\neq0$. Para cada $i=1,\ldots,n$, sea $A_i$ la matriz que se obtiene sustituyendo la columna $i$ de $A$ por el vector $b$. Entonces la solución única viene dada por
\[
\boxed{x_i=\frac{\det(A_i)}{\det(A)},\qquad i=1,\ldots,n.}
\]
\end{theorem}

\begin{example}[Corrientes en dos mallas]
Tras aplicar las leyes de Kirchhoff a un circuito lineal se obtiene el sistema
\[
\begin{cases}
7I_1-2I_2=12,\\
-2I_1+5I_2=6.
\end{cases}
\]
La matriz de coeficientes es
\[
A=\begin{pmatrix}7&-2\\-2&5\end{pmatrix},
\qquad
\det(A)=35-4=31\neq0.
\]
El sistema tiene solución única. Por Cramer,
\[
I_1=
\frac{\begin{vmatrix}12&-2\\6&5\end{vmatrix}}{31}
=\frac{72}{31},
\qquad
I_2=
\frac{\begin{vmatrix}7&12\\-2&6\end{vmatrix}}{31}
=\frac{66}{31}.
\]
Por tanto,
\[
\boxed{I_1=\frac{72}{31}\ \mathrm A,\qquad I_2=\frac{66}{31}\ \mathrm A.}
\]
La resolución matemática comienza una vez formuladas las ecuaciones físicas; el álgebra lineal permite después decidir la unicidad y calcular las corrientes.
\end{example}

Cramer es útil para sistemas pequeños y para obtener fórmulas simbólicas, pero el método de Gauss es mucho más general: funciona también con matrices no cuadradas y con sistemas compatibles indeterminados o incompatibles.

\subsection{Sistemas con parámetros y sistemas homogéneos}

En muchos problemas los coeficientes dependen de un parámetro. En esos casos no basta con resolver para un valor concreto: hay que determinar para qué valores cambia el rango y, por tanto, la naturaleza del sistema.

\begin{solvedexercise}
Discutir según $\lambda\in\R$ el sistema
\[
\begin{cases}
\lambda x+y+z=1,\\
x+\lambda y+z=1,\\
x+y+\lambda z=1.
\end{cases}
\]
\tcblower
\textbf{Resolución.}\par\smallskip
La matriz de coeficientes es
\[
A_\lambda=
\begin{pmatrix}
\lambda&1&1\\
1&\lambda&1\\
1&1&\lambda
\end{pmatrix}.
\]
Su determinante es
\[
\det(A_\lambda)
=(\lambda-1)^2(\lambda+2).
\]
Por tanto, si
\[
\lambda\neq1,-2,
\]
el determinante es no nulo y el sistema tiene una única solución. Debido a la simetría de las ecuaciones, $x=y=z$, y al sustituir se obtiene
\[
\boxed{x=y=z=\frac{1}{\lambda+2}}.
\]

Para $\lambda=1$, las tres ecuaciones coinciden:
\[
x+y+z=1.
\]
Así,
\[
\operatorname{rg}(A_1)=\operatorname{rg}(A_1\mid b)=1<3,
\]
y el sistema es compatible indeterminado, con dos parámetros libres.

Para $\lambda=-2$, sumando las tres ecuaciones del sistema se obtiene
\[
0=3,
\]
lo cual es imposible. Por tanto,
\[
\operatorname{rg}(A_{-2})<\operatorname{rg}(A_{-2}\mid b),
\]
y el sistema es incompatible.
\end{solvedexercise}

Un caso especialmente importante es el de los sistemas homogéneos.

\begin{definition}[Sistema homogéneo]
Un sistema lineal es \emph{homogéneo} si todos sus términos independientes son cero:
\[
Ax=0.
\]
\end{definition}

Todo sistema homogéneo tiene al menos la solución
\[
x=0,
\]
llamada \emph{solución trivial}. Por tanto, nunca es incompatible.

\begin{theorem}
Sea $Ax=0$ un sistema homogéneo con $n$ incógnitas y $r=\operatorname{rg}(A)$. Entonces:
\begin{itemize}
\item si $r=n$, la única solución es la trivial;
\item si $r<n$, existen infinitas soluciones y la solución general depende de $n-r$ parámetros libres.
\end{itemize}
\end{theorem}

\begin{example}[Equilibrio con una familia de soluciones]
Consideremos
\[
\begin{cases}
2x+3y+z=0,\\
x+2y+5z=0.
\end{cases}
\]
La matriz de coeficientes tiene rango $2$ y hay $3$ incógnitas. Por tanto, el sistema homogéneo tiene infinitas soluciones con un parámetro libre. De la primera ecuación y la segunda puede obtenerse
\[
y=-9z,
\qquad
x=13z.
\]
Tomando $z=t$,
\[
\boxed{(x,y,z)=t(13,-9,1),\qquad t\in\R.}
\]
En el siguiente capítulo interpretaremos este conjunto de soluciones como un subespacio vectorial generado por el vector $(13,-9,1)$.
\end{example}

Los sistemas lineales aparecen también al formular balances. En una instalación química, por ejemplo, pueden desconocerse varios caudales y conocerse las cantidades totales de diferentes componentes. Cada balance de componente aporta una ecuación lineal. En mecánica, las ecuaciones de equilibrio de fuerzas y momentos generan sistemas para las reacciones desconocidas. En ingeniería eléctrica, las ecuaciones de nodos o mallas tienen la misma estructura algebraica. El procedimiento matemático es común: formular correctamente $A$, $x$ y $b$, analizar los rangos y sólo entonces resolver.

\begin{example}[Balance de mezclas]
Se desean mezclar dos corrientes $A$ y $B$. La corriente $A$ aporta $4$ unidades de un componente $C_1$ y $3$ de $C_2$ por unidad de caudal, mientras que $B$ aporta $2$ de $C_1$ y $5$ de $C_2$. Se necesitan en total $40$ unidades de $C_1$ y $60$ de $C_2$. Si $x$ e $y$ son los caudales de $A$ y $B$, respectivamente,
\[
\begin{cases}
4x+2y=40,\\
3x+5y=60.
\end{cases}
\]
En forma matricial,
\[
\begin{pmatrix}4&2\\3&5\end{pmatrix}
\begin{pmatrix}x\\y\end{pmatrix}
=
\begin{pmatrix}40\\60\end{pmatrix}.
\]
Como
\[
\det\begin{pmatrix}4&2\\3&5\end{pmatrix}=14\neq0,
\]
la mezcla es matemáticamente única. Resolviendo,
\[
\boxed{x=\frac{40}{7},\qquad y=\frac{60}{7}.}
\]
En un problema real habría que comprobar además que las soluciones sean físicamente admisibles, por ejemplo que los caudales no sean negativos.
\end{example}

\section*{Ejercicios propuestos}
\addcontentsline{toc}{section}{Ejercicios propuestos}

\begin{enumerate}[label=\textbf{\arabic*.},leftmargin=*]

\item Dadas
\[
A=\begin{pmatrix}1&3&2\\-1&9&2\\0&3&-1\end{pmatrix},
\qquad
B=\begin{pmatrix}5&1&-2\\0&6&0\\-1&-3&1\end{pmatrix},
\]
\[
F=\begin{pmatrix}1&5\\0&2\\4&-1\end{pmatrix},
\qquad
G=\begin{pmatrix}1&0&8\\2&-1&1\end{pmatrix},
\]
calcular, cuando sea posible,
\begin{enumerate}[label=\alph*)]
\item $2A-7B$;
\item $(FG)^t-A$;
\item $AB-BA$;
\item $ABA$.
\end{enumerate}
Justificar los productos que no estén definidos por incompatibilidad de dimensiones.

\item Sean
\[
A=\begin{pmatrix}2&3\\-1&1\end{pmatrix},
\qquad
B=\begin{pmatrix}1&9\\-3&k\end{pmatrix}.
\]
Determinar los valores de $k$ para los que $AB=BA$.

\item Encontrar dos matrices $A$ y $B$, ambas no nulas, tales que $AB=0$. Explicar por qué este hecho no contradice la propiedad de los números reales según la cual $ab=0$ implica $a=0$ o $b=0$.

\item Sean $A,B,C$ matrices de dimensiones compatibles. Decidir cuáles de las siguientes igualdades son siempre ciertas y justificar la respuesta:
\begin{enumerate}[label=\alph*)]
\item $A(B+C)=AB+AC$;
\item $(A+B)^2=A^2+2AB+B^2$;
\item $(AB)^t=A^tB^t$;
\item $(AB)^{-1}=A^{-1}B^{-1}$, cuando las inversas existen;
\item $A^2-B^2=(A-B)(A+B)$.
\end{enumerate}
Cuando una igualdad sea falsa, proporcionar un contraejemplo.

\item Una empresa produce tres componentes eléctricos $E_1,E_2,E_3$ en dos líneas de fabricación. La matriz
\[
P=\begin{pmatrix}
120&80&60\\
100&90&75
\end{pmatrix}
\]
recoge la producción diaria de cada línea. La primera fila corresponde a la línea 1 y la segunda a la línea 2.
\begin{enumerate}[label=\alph*)]
\item expresar matricialmente la producción total de cada componente;
\item si la producción de la línea 1 aumenta un $15\%$ y la de la línea 2 un $10\%$, construir la nueva matriz de producción;
\item si cada unidad de $E_1,E_2,E_3$ requiere respectivamente $2$, $3$ y $5$ minutos de inspección, expresar mediante un producto matricial el tiempo total de inspección de cada línea.
\end{enumerate}

\item Una planta química fabrica dos productos a partir de tres materias primas. El consumo por unidad de producto viene dado por
\[
A=\begin{pmatrix}3&4\\2&1\\4&3\end{pmatrix}.
\]
La producción semanal es $x=(20,30)^t$ y los costes unitarios de las materias primas son $c=(6,10,12)$.
\begin{enumerate}[label=\alph*)]
\item calcular $Ax$ e interpretar sus componentes;
\item calcular $cA$ e interpretar sus componentes;
\item calcular $cAx$ de dos formas distintas y explicar por qué los resultados coinciden.
\end{enumerate}

\item Una empresa de mantenimiento dispone de equipos en tres centros. El vector columna
\[
x_0=\begin{pmatrix}200\\300\\100\end{pmatrix}
\]
representa la distribución inicial. Cada semana, la distribución cambia según
\[
x_{k+1}=Px_k,
\qquad
P=\begin{pmatrix}
0.50&0.40&0.20\\
0.30&0.45&0.20\\
0.20&0.15&0.60
\end{pmatrix}.
\]
Calcular $x_1,x_2,x_3$ y escribir una expresión para $x_k$. Interpretar por qué las columnas de $P$ suman $1$.

\item Resolver las ecuaciones matriciales, suponiendo invertibles las matrices que sea necesario invertir:
\begin{enumerate}[label=\alph*)]
\item $AX+B^tX=2C$;
\item $AXA+AB=C$;
\item $A(X+C)=BX$;
\item $AX-B=3X$;
\item $XAB=C$.
\end{enumerate}
Prestar especial atención al orden de los factores.

\item Calcular los determinantes:
\begin{enumerate}[label=\alph*)]
\item $\begin{vmatrix}4&-2\\3&5\end{vmatrix}$;
\item $\begin{vmatrix}1&2&3\\0&-1&4\\2&1&0\end{vmatrix}$;
\item $\begin{vmatrix}2&1&0&0\\3&-1&4&0\\5&2&7&0\\1&3&2&6\end{vmatrix}$;
\item $\begin{vmatrix}1&2&0&0\\0&3&1&0\\0&0&-2&4\\0&0&0&5\end{vmatrix}$.
\end{enumerate}
Elegir en cada caso el procedimiento más eficiente.

\item Resolver
\[
\det\begin{pmatrix}
4&x&6\\
5&7&12\\
3&-1&x
\end{pmatrix}=0.
\]

\item Utilizando exclusivamente propiedades de determinantes, y sin desarrollar los determinantes elemento a elemento, resolver:
\begin{enumerate}[label=\alph*)]
\item si $A\in\mathcal M_n(\R)$ y $A^2=A$, determinar $\det(A)$;
\item si $\det(A)=5$ y $\det(B)=3$, calcular $\det(BA^t)$ y $\det(BAB)$;
\item si $A$ es invertible y $\det(A)=4$, calcular $\det(A^{-1})$;
\item expresar $\det(\lambda A)$ en función de $\lambda$, $n$ y $\det(A)$.
\end{enumerate}

\item Sea $A\in\mathcal M_n(\R)$ invertible. Demostrar que
\[
\det\bigl(\operatorname{Adj}(A)\bigr)=\det(A)^{n-1}.
\]
Indicar con claridad dónde se utiliza que $A$ es invertible.

\item Calcular la inversa, cuando exista, de:
\begin{enumerate}[label=\alph*)]
\item $\begin{pmatrix}2&1\\-1&1\end{pmatrix}$;
\item $\begin{pmatrix}1&3&2\\-1&1&5\\1&2&3\end{pmatrix}$;
\item $\begin{pmatrix}5&3&1\\0&2&-7\\3&4&-9\end{pmatrix}$;
\item $\begin{pmatrix}0&1&0\\1&0&0\\0&0&1\end{pmatrix}$.
\end{enumerate}
Utilizar en algunos apartados la fórmula de la adjunta y en otros Gauss--Jordan.

\item Determinar para qué valores de $\alpha$ y $\beta$ son invertibles
\[
A_\alpha=
\begin{pmatrix}
1&0&0\\
0&\alpha-1&2\\
0&-1&\alpha+1
\end{pmatrix},
\qquad
B_\beta=
\begin{pmatrix}
1&2&1\\
-1&\beta-1&2\\
2&3&\beta+1
\end{pmatrix}.
\]

\item Calcular la inversa de
\[
R(\theta)=
\begin{pmatrix}
\cos\theta&\sin\theta\\
-\sin\theta&\cos\theta
\end{pmatrix}
\]
y relacionarla con $R(-\theta)$ y con $R(\theta)^t$.

\item Calcular el rango de las matrices
\[
A=\begin{pmatrix}
0&-1&2&-1\\
3&-2&1&2
\end{pmatrix},
\qquad
B=\begin{pmatrix}
1&2&3\\
-1&1&1\\
3&3&5\\
0&3&4
\end{pmatrix}.
\]
Realizar el cálculo por Gauss y comprobarlo mediante menores no nulos adecuados.

\item Determinar, según $\lambda\in\R$, el rango de
\[
A_\lambda=
\begin{pmatrix}
1&-1&1\\
2&1&3\\
5&1&\lambda
\end{pmatrix}.
\]

\item Determinar $a$ y $b$ para que
\[
A=\begin{pmatrix}1&2&4\\a&4&b\end{pmatrix}
\]
tenga rango $1$.

\item Determinar para qué valores de $\alpha$ el rango de
\[
H_\alpha=
\begin{pmatrix}
\alpha&-1&0\\
1&\alpha^2&1+\alpha^2\\
2&2&2+\alpha
\end{pmatrix}
\]
es igual a $2$.

\item Sea
\[
A=\begin{pmatrix}
1&a&a^2\\
1&b&b^2\\
1&c&c^2
\end{pmatrix}.
\]
Calcular $\det(A)$ y deducir cuándo $A$ es invertible. Interpretar el resultado en función de la igualdad o diferencia entre $a,b,c$.

\item Estudiar el rango de
\[
A_{a,b}=
\begin{pmatrix}
1&a&1\\
2&2a&2\\
1&b&3
\end{pmatrix}
\]
según los valores de $a$ y $b$. Organizar el análisis por casos sin dividir nunca por una expresión paramétrica antes de saber si es nula.

\item Expresar como sistema de ecuaciones y resolver
\[
x\begin{pmatrix}2\\1\end{pmatrix}
-y\begin{pmatrix}-3\\5\end{pmatrix}
=2\begin{pmatrix}8\\11\end{pmatrix}.
\]

\item Resolver por eliminación de Gauss
\[
\begin{cases}
2x+3y+2z+6t=10,\\
y+2z+t=2,\\
3x-3z+6t=9.
\end{cases}
\]
Escribir la solución en forma paramétrica.

\item Discutir y resolver, cuando sea posible,
\[
\begin{cases}
x+y+z+t=0,\\
x-y+z-t=12,\\
x+y-z+t=-6,\\
x-y-z+t=-6.
\end{cases}
\]

\item Discutir y resolver
\[
\begin{cases}
2x-y-3z=2,\\
-x+2y-3z=7,\\
x+y+4z=-3.
\end{cases}
\]
Resolver después el mismo sistema mediante Cramer y mediante $x=A^{-1}b$, y comparar los tres procedimientos.

\item Discutir y resolver
\[
\begin{cases}
4x+y+z+w=6,\\
3x+7y-z+w=1,\\
7x+3y-5z+8w=-3,\\
x+y+z+2w=3.
\end{cases}
\]

\item Discutir según $k\in\R$ y resolver cuando sea posible
\[
\begin{cases}
x+y-z=1,\\
2x+3y+kz=3,\\
x+ky+3z=2.
\end{cases}
\]

\item Discutir según $\lambda\in\R$
\[
\begin{cases}
\lambda x+y+z=1,\\
x+\lambda y+z=1,\\
x+y+\lambda z=1.
\end{cases}
\]
Resolver el sistema en cada caso compatible.

\item Determinar los valores de $a$ para los que
\[
\begin{cases}
x+ay+3z=2,\\
x+y-z=1,\\
2x+3y+az=-1
\end{cases}
\]
\begin{enumerate}[label=\alph*)]
\item tiene solución única;
\item tiene más de una solución;
\item no tiene solución.
\end{enumerate}

\item Resolver los sistemas homogéneos:
\begin{enumerate}[label=\alph*)]
\item $\begin{cases}3w+5x-4y+2z=0,\\7w-2x+9y+3z=0;\end{cases}$
\item $\begin{cases}2x+3y-z=0,\\y+5z=0,\\3x-z=0;\end{cases}$
\item $\begin{cases}2x+3y+z=0,\\x+2y+5z=0.\end{cases}$
\end{enumerate}
En cada caso, indicar el rango y el número de parámetros libres.

\item Una mezcla industrial se prepara con dos corrientes $A$ y $B$. Cada unidad de $A$ contiene $4$ unidades de componente $C_1$ y $3$ de $C_2$, mientras que cada unidad de $B$ contiene $2$ de $C_1$ y $5$ de $C_2$.
\begin{enumerate}[label=\alph*)]
\item determinar si es posible obtener $40$ unidades de $C_1$ y $60$ de $C_2$;
\item si la composición de $A$ cambia y pasa a contener $a$ unidades de $C_1$, discutir según $a$ cuándo la mezcla objetivo es única, imposible o admite infinitas soluciones;
\item comentar qué restricciones físicas adicionales deberían imponerse a las soluciones matemáticas.
\end{enumerate}

\item En un nudo de una estructura actúan tres fuerzas de módulos desconocidos $T_1,T_2,T_3$. Tras proyectar las ecuaciones de equilibrio sobre tres direcciones independientes se obtiene
\[
\begin{cases}
T_1+T_2=10,\\
T_2+T_3=8,\\
T_1-T_3=2.
\end{cases}
\]
\begin{enumerate}[label=\alph*)]
\item escribir el sistema en forma matricial;
\item determinar su rango y discutir si las tres ecuaciones aportan información independiente;
\item resolverlo;
\item interpretar qué significaría físicamente obtener un valor negativo para uno de los módulos si el modelo hubiera definido $T_i$ como tracciones positivas.
\end{enumerate}

\item En una red eléctrica lineal, las corrientes de tres mallas satisfacen
\[
\begin{cases}
6I_1-2I_2-I_3=10,\\
-2I_1+5I_2-I_3=4,\\
-I_1-I_2+4I_3=6.
\end{cases}
\]
Determinar las corrientes mediante eliminación de Gauss. Comprobar la unicidad de la solución antes de resolver.

\item Un sistema energético combina tres fuentes $E_1,E_2,E_3$. Por unidad de energía suministrada, sus emisiones relativas de dos contaminantes y su coste vienen dados por las columnas de
\[
A=
\begin{pmatrix}
1&2&4\\
3&1&2\\
5&4&3
\end{pmatrix}.
\]
Se desea obtener un vector total
\[
b=\begin{pmatrix}18\\14\\31\end{pmatrix}.
\]
\begin{enumerate}[label=\alph*)]
\item formular el problema como $Ax=b$;
\item estudiar si existe una combinación única de las tres fuentes;
\item resolver el sistema;
\item discutir la admisibilidad del resultado si las componentes de $x$ representan cantidades que deben ser no negativas.
\end{enumerate}

\item Una planta produce dos aleaciones utilizando tres metales. Las columnas de
\[
A=\begin{pmatrix}
0.50&0.20\\
0.30&0.50\\
0.20&0.30
\end{pmatrix}
\]
indican las fracciones de cada metal en una tonelada de las aleaciones $L_1$ y $L_2$. Se dispone de un vector de metal utilizado
\[
b=\begin{pmatrix}40\\55\\35\end{pmatrix}.
\]
Determinar si puede existir una producción $(x_1,x_2)$ que consuma exactamente esas tres cantidades. Utilizar Rouché--Frobenius y explicar qué significa el resultado en términos de compatibilidad de los tres balances.

\item \M Introducir dos matrices pequeñas en Mathematica y comprobar las operaciones de los ejercicios 1 y 2. Utilizar únicamente las operaciones algebraicas habituales y comprobar con \texttt{Simplify} cualquier igualdad simbólica que contenga parámetros. No utilizar el ordenador para decidir previamente si un producto está definido.

\item \M Seleccionar tres matrices de los ejercicios 9--16. Calcular manualmente sus determinantes e inversas cuando existan y comprobar después con \texttt{Det} e \texttt{Inverse}. Explicar qué respuesta proporciona \texttt{Inverse} cuando la matriz es singular.

\item \M Para las matrices de los ejercicios 17--21, calcular primero el rango analíticamente y después comprobarlo con \texttt{MatrixRank}. En los casos con parámetros, estudiar por separado los valores críticos antes de pedir al ordenador un resultado general.

\item \M Resolver mediante \texttt{RowReduce} la matriz ampliada de los sistemas de los ejercicios 23--29. Identificar los pivotes y comprobar que la clasificación obtenida coincide con el Teorema de Rouché--Frobenius.

\item \M Resolver uno de los sistemas compatibles determinados mediante \texttt{LinearSolve} y mediante \texttt{Solve}. Comparar la forma de la entrada y de la salida de ambos comandos. Repetir con un sistema compatible indeterminado y discutir qué representación de la solución resulta más informativa.

\item (Examen) Para cada $\lambda\in\R$ se considera
\[
A_\lambda=
\begin{pmatrix}
\lambda&1&-2\\
1&\lambda&2\\
-2&2&\lambda-3
\end{pmatrix}.
\]
Calcular $\det(A_\lambda)$ y determinar los valores de $\lambda$ para los que $A_\lambda$ es regular.

\item (Examen) Sea
\[
A=
\begin{pmatrix}
1&0&1&0\\
0&1&-2&0\\
1&-2&5&0\\
0&0&0&6
\end{pmatrix}.
\]
Decidir razonadamente si $A$ es regular. Si lo es, indicar al menos dos procedimientos distintos que permitirían calcular $A^{-1}$ y elegir el más conveniente para esta matriz.

\item (Examen) Determinar los valores de $\alpha$ para los que
\[
H_\alpha=
\begin{pmatrix}
\alpha&-1&0\\
1&\alpha^2&1+\alpha^2\\
2&2&2+\alpha
\end{pmatrix}
\]
tiene rango exactamente $2$. La respuesta debe justificar tanto que el rango no sea $3$ como que no descienda por debajo de $2$.

\end{enumerate}

\cleardoublepage
\chapter{Espacios vectoriales y aplicaciones lineales}
\label{ch:espacios}

En el capítulo anterior hemos utilizado vectores como columnas de números y hemos comprobado que conceptos como el rango o la resolución de un sistema dependen, en último término, de si unas filas o columnas contienen información independiente. El lenguaje de los \emph{espacios vectoriales} permite formular estas ideas de manera general y entender qué tienen en común objetos aparentemente distintos, como vectores de $\R^n$, polinomios, matrices o señales.

En ingeniería, elegir un conjunto adecuado de variables independientes es una decisión esencial. Las componentes de una fuerza dependen del sistema de referencia; un estado físico puede describirse mediante distintas coordenadas; varios sensores pueden proporcionar información redundante; y una transformación lineal puede convertir magnitudes de entrada en respuestas de un sistema. Las nociones de dependencia lineal, base, dimensión, ortogonalidad, núcleo e imagen formalizan precisamente estas cuestiones.

A lo largo del capítulo trabajaremos principalmente en espacios vectoriales reales de dimensión finita. El objetivo no es acumular definiciones abstractas, sino aprender a pasar con soltura entre tres descripciones complementarias: ecuaciones, conjuntos de generadores y coordenadas respecto de una base.

\section{Espacios vectoriales y subespacios}
\label{sec:espacios-subespacios}

\subsection{Espacios vectoriales}

La suma de vectores y el producto de un vector por un número aparecen de forma natural en $\R^n$. La estructura de espacio vectorial conserva precisamente las propiedades algebraicas esenciales de estas dos operaciones.

\begin{definition}[Espacio vectorial]
Un \emph{espacio vectorial real} es un conjunto no vacío $V$ cuyos elementos, llamados \emph{vectores}, pueden sumarse entre sí y multiplicarse por números reales de manera que, para cualesquiera $u,v,w\in V$ y $\alpha,\beta\in\R$, se verifican:
\begin{enumerate}
\item $u+v\in V$ y $\alpha u\in V$;
\item $u+v=v+u$;
\item $(u+v)+w=u+(v+w)$;
\item existe un vector $0\in V$ tal que $u+0=u$;
\item para cada $u\in V$ existe $-u\in V$ tal que $u+(-u)=0$;
\item $\alpha(u+v)=\alpha u+\alpha v$;
\item $(\alpha+\beta)u=\alpha u+\beta u$;
\item $(\alpha\beta)u=\alpha(\beta u)$ y $1u=u$.
\end{enumerate}
\end{definition}

El ejemplo fundamental es $\R^n$, con las operaciones componente a componente. Sin embargo, no es el único.

\begin{example}
Son espacios vectoriales reales:
\begin{itemize}
\item $\R^n$;
\item el conjunto $\mathcal M_{m\times n}(\R)$ de matrices reales de tamaño $m\times n$;
\item el conjunto $\mathcal P_n$ de polinomios reales de grado menor o igual que $n$;
\item el conjunto de funciones reales continuas en un intervalo, con las operaciones habituales.
\end{itemize}
En cada caso, el vector nulo es diferente: en $\R^n$ es $(0,\ldots,0)$, en $\mathcal M_{m\times n}(\R)$ es la matriz nula y en $\mathcal P_n$ es el polinomio idénticamente nulo.
\end{example}

El uso de la palabra \emph{vector} es, por tanto, más amplio que el de una flecha en el plano o el espacio. Lo que importa es que los objetos puedan combinarse linealmente respetando las propiedades anteriores.

\subsection{Subespacios vectoriales}

En muchas situaciones no necesitamos utilizar todo el espacio, sino únicamente los vectores que satisfacen determinadas restricciones. Por ejemplo, los desplazamientos admisibles de un sistema pueden estar sujetos a ciertas ecuaciones lineales, o varias magnitudes pueden estar ligadas por balances homogéneos.

\begin{definition}[Subespacio vectorial]
Sea $V$ un espacio vectorial. Un subconjunto no vacío $W\subseteq V$ es un \emph{subespacio vectorial} de $V$ si, con las mismas operaciones de $V$, es también un espacio vectorial.
\end{definition}

No es necesario comprobar una a una todas las propiedades de la definición.

\begin{theorem}[Criterio de subespacio]
Sea $V$ un espacio vectorial y sea $W\subseteq V$, $W\neq\varnothing$. Entonces $W$ es un subespacio vectorial si y sólo si
\[
u,v\in W,\quad \alpha,\beta\in\R
\quad\Longrightarrow\quad
\alpha u+\beta v\in W.
\]
\end{theorem}

Una consecuencia inmediata es que todo subespacio debe contener al vector nulo. Esta observación permite descartar muchos conjuntos sin realizar más cálculos.

\begin{example}
Consideremos
\[
W=\{(x,y,z)\in\R^3:x+y-z=0\}.
\]
El vector nulo pertenece a $W$. Si $u=(x_1,y_1,z_1)$ y $v=(x_2,y_2,z_2)$ pertenecen a $W$, entonces
\[
x_1+y_1-z_1=0,
\qquad
x_2+y_2-z_2=0.
\]
Para $\alpha,\beta\in\R$,
\[
\begin{aligned}
&(\alpha x_1+\beta x_2)+(\alpha y_1+\beta y_2)
 -(\alpha z_1+\beta z_2)\\
&\qquad=\alpha(x_1+y_1-z_1)+\beta(x_2+y_2-z_2)=0.
\end{aligned}
\]
Por tanto, $W$ es un subespacio de $\R^3$.

En cambio,
\[
H=\{(x,y,z)\in\R^3:x+y-z=1\}
\]
no es un subespacio, pues $(0,0,0)\notin H$. Geométricamente, $W$ es un plano que pasa por el origen, mientras que $H$ es un plano paralelo a él que no pasa por el origen.
\end{example}

\begin{example}[Matrices simétricas]
Sea
\[
S=\{A\in\mathcal M_2(\R):A^t=A\}.
\]
Si $A,B\in S$ y $\alpha,\beta\in\R$, entonces
\[
(\alpha A+\beta B)^t
=\alpha A^t+\beta B^t
=\alpha A+\beta B.
\]
Luego $S$ es un subespacio de $\mathcal M_2(\R)$. Sus elementos tienen la forma
\[
\begin{pmatrix}
a&b\\
b&c
\end{pmatrix}
=
a\begin{pmatrix}1&0\\0&0\end{pmatrix}
+b\begin{pmatrix}0&1\\1&0\end{pmatrix}
+c\begin{pmatrix}0&0\\0&1\end{pmatrix}.
\]
Esta representación anticipa los conceptos de combinación lineal, sistema generador y base.
\end{example}

\subsection{Combinaciones lineales y sistemas generadores}

\begin{definition}[Combinación lineal]
Dados $v_1,\ldots,v_k\in V$, una \emph{combinación lineal} de estos vectores es cualquier vector de la forma
\[
\alpha_1v_1+\alpha_2v_2+\cdots+\alpha_kv_k,
\qquad
\alpha_1,\ldots,\alpha_k\in\R.
\]
\end{definition}

El conjunto formado por todas sus combinaciones lineales se denomina subespacio generado por $v_1,\ldots,v_k$ y se escribe
\[
\langle v_1,\ldots,v_k\rangle
=
\left\{
\alpha_1v_1+\cdots+\alpha_kv_k:
\alpha_i\in\R
\right\}.
\]
Es el menor subespacio vectorial que contiene a todos los vectores $v_1,\ldots,v_k$.

\begin{example}
Sean
\[
v_1=(1,2,-1),
\qquad
v_2=(0,1,1).
\]
Entonces
\[
\langle v_1,v_2\rangle
=
\{\alpha(1,2,-1)+\beta(0,1,1):\alpha,\beta\in\R\},
\]
es decir,
\[
\langle v_1,v_2\rangle
=
\{(\alpha,2\alpha+\beta,-\alpha+\beta):\alpha,\beta\in\R\}.
\]
Esta expresión constituye una descripción paramétrica del subespacio.
\end{example}

Un mismo subespacio puede describirse de maneras distintas. En $\R^n$ es frecuente disponer de:
\begin{itemize}
\item una descripción \emph{paramétrica}, mediante generadores;
\item una descripción \emph{implícita}, mediante un sistema de ecuaciones lineales homogéneas.
\end{itemize}
Pasar de una forma a otra es una destreza básica del capítulo.

\begin{solvedexercise}
Sea
\[
W=\{(x,y,z,t)\in\R^4:x+y+z+t=0,\quad y-2z-t=0\}.
\]
Hallar una descripción paramétrica de $W$ y un sistema de generadores.
\tcblower
\textbf{Resolución.}\par\smallskip
Resolvemos las ecuaciones tomando $z=s$ y $t=r$ como parámetros libres. De
\[
y-2z-t=0
\]
obtenemos
\[
y=2s+r.
\]
La primera ecuación proporciona
\[
x=-(y+z+t)=-(2s+r+s+r)=-3s-2r.
\]
Por tanto,
\[
(x,y,z,t)
=s(-3,2,1,0)+r(-2,1,0,1).
\]
Así,
\[
\boxed{W=\langle(-3,2,1,0),(-2,1,0,1)\rangle.}
\]
La descripción implícita pone de manifiesto las restricciones que deben satisfacer los vectores; la paramétrica muestra directamente cómo construir todos los elementos del subespacio.
\end{solvedexercise}

La conversión en sentido contrario puede hacerse eliminando los parámetros. Si
\[
W=\langle v_1,\ldots,v_k\rangle\subseteq\R^n,
\]
un vector $x\in\R^n$ pertenece a $W$ si y sólo si el sistema
\[
\alpha_1v_1+\cdots+\alpha_kv_k=x
\]
es compatible. Esta condición puede traducirse en ecuaciones lineales sobre las componentes de $x$ mediante eliminación de Gauss.

\subsection{Dependencia e independencia lineal}

No todos los generadores de un subespacio aportan información nueva. Si uno de ellos puede obtenerse a partir de los demás, es redundante.

\begin{definition}[Dependencia e independencia lineal]
Los vectores $v_1,\ldots,v_k\in V$ son \emph{linealmente independientes} si la igualdad
\[
\alpha_1v_1+\cdots+\alpha_kv_k=0
\]
implica necesariamente
\[
\alpha_1=\cdots=\alpha_k=0.
\]
Si existe una combinación lineal nula con algún coeficiente no nulo, los vectores son \emph{linealmente dependientes}.
\end{definition}

Equivalentemente, un conjunto finito de vectores es linealmente dependiente si alguno de ellos puede escribirse como combinación lineal de los restantes.

Cuando los vectores pertenecen a $\R^n$, podemos colocarlos como columnas de una matriz
\[
A=\begin{pmatrix}|&&|\\v_1&\cdots&v_k\\|&&|\end{pmatrix}.
\]
Entonces
\[
\alpha_1v_1+\cdots+\alpha_kv_k=0
\]
equivale al sistema homogéneo
\[
A\alpha=0.
\]
Por tanto, los vectores son linealmente independientes si y sólo si este sistema tiene únicamente la solución trivial. En particular,
\[
\boxed{v_1,\ldots,v_k\text{ son LI}\iff \operatorname{rg}(A)=k.}
\]
Esta relación enlaza directamente el concepto abstracto de independencia con el rango estudiado en el capítulo anterior.

\begin{example}
Consideremos
\[
v_1=(1,0,-1),\qquad
v_2=(0,-1,2),\qquad
v_3=(1,-1,1).
\]
Como
\[
v_3=v_1+v_2,
\]
los tres vectores son linealmente dependientes y
\[
\langle v_1,v_2,v_3\rangle
=
\langle v_1,v_2\rangle.
\]
Además, $v_1$ y $v_2$ no son proporcionales, por lo que son linealmente independientes. Hemos eliminado así un generador redundante sin cambiar el subespacio generado.
\end{example}

En un modelo experimental esta idea tiene una lectura directa. Si cada columna de una matriz contiene la respuesta de un sistema ante una entrada elemental y una de esas columnas es combinación de las demás, esa respuesta no añade una nueva dirección de comportamiento. Del mismo modo, si varias ecuaciones de medida son linealmente dependientes, alguna medición es redundante.

\section{Bases, coordenadas y estructura euclídea}
\label{sec:bases-euclidea}

\subsection{Bases y dimensión}

Un sistema generador describe todo el espacio, pero puede contener vectores redundantes. Una base elimina precisamente esa redundancia.

\begin{definition}[Base]
Una familia ordenada
\[
\mathcal B=(v_1,\ldots,v_n)
\]
es una \emph{base} de un espacio vectorial $V$ si:
\begin{enumerate}
\item $v_1,\ldots,v_n$ generan $V$;
\item $v_1,\ldots,v_n$ son linealmente independientes.
\end{enumerate}
\end{definition}

Una vez fijada una base, cada vector admite una representación única.

\begin{theorem}[Unicidad de las coordenadas]
Si $\mathcal B=(v_1,\ldots,v_n)$ es una base de $V$, para todo $v\in V$ existen únicos escalares $x_1,\ldots,x_n$ tales que
\[
v=x_1v_1+\cdots+x_nv_n.
\]
\end{theorem}

Los números $x_1,\ldots,x_n$ se denominan \emph{coordenadas de $v$ respecto de la base $\mathcal B$} y escribimos
\[
[v]_{\mathcal B}
=
\begin{pmatrix}
x_1\\ \vdots\\ x_n
\end{pmatrix}.
\]

\begin{definition}[Dimensión]
Si un espacio vectorial $V$ admite una base formada por $n$ vectores, se dice que $V$ tiene \emph{dimensión} $n$ y se escribe
\[
\dim(V)=n.
\]
Todas las bases de un espacio vectorial de dimensión finita tienen el mismo número de elementos.
\end{definition}

La base canónica de $\R^n$ es
\[
\mathcal E=(e_1,\ldots,e_n),
\]
donde $e_i$ tiene un $1$ en la posición $i$ y ceros en las restantes. Por tanto,
\[
\dim(\R^n)=n.
\]
También
\[
\dim\mathcal P_n=n+1,
\qquad
\dim\mathcal M_{m\times n}(\R)=mn.
\]

\begin{example}
Retomemos
\[
W=\{(x,y,z,t)\in\R^4:x+y+z+t=0,\ y-2z-t=0\}.
\]
Hemos obtenido
\[
W=\langle(-3,2,1,0),(-2,1,0,1)\rangle.
\]
Los dos generadores son linealmente independientes, pues no son proporcionales. Por tanto,
\[
\mathcal B_W=\bigl((-3,2,1,0),(-2,1,0,1)\bigr)
\]
es una base de $W$ y
\[
\boxed{\dim(W)=2.}
\]
La dimensión puede interpretarse como el número de parámetros independientes necesarios para describir un vector del subespacio.
\end{example}

Para un conjunto de vectores de $\R^n$ colocados como columnas de una matriz $A$, el rango proporciona la dimensión del subespacio generado:
\[
\boxed{
\dim\langle v_1,\ldots,v_k\rangle
=
\operatorname{rg}(A).
}
\]
Esto permite obtener una base seleccionando columnas linealmente independientes de la matriz original. Es importante escoger las columnas de la matriz original asociadas a las posiciones pivote, no las columnas de la matriz escalonada si lo que se desea es una base formada por los vectores dados.

\begin{solvedexercise}
Sea
\[
V=\{(4,-5,7),(3,-3,4),(1,1,-2),(2,-1,1)\}\subseteq\R^3.
\]
Determinar la dimensión del subespacio generado por $V$ y obtener una base formada por vectores de $V$.
\tcblower
\textbf{Resolución.}\par\smallskip
Colocamos los vectores como columnas:
\[
A=
\begin{pmatrix}
4&3&1&2\\
-5&-3&1&-1\\
7&4&-2&1
\end{pmatrix}.
\]
Aplicando eliminación de Gauss se obtiene un rango igual a $2$. Por ejemplo, las dos primeras columnas son linealmente independientes porque no son proporcionales, y las restantes pueden expresarse como combinaciones de ellas. En efecto,
\[
(1,1,-2)=-2(4,-5,7)+3(3,-3,4),
\]
\[
(2,-1,1)=-(4,-5,7)+2(3,-3,4).
\]
Por tanto,
\[
\boxed{\dim\langle V\rangle=2}
\]
y una posible base es
\[
\boxed{\mathcal B=\bigl((4,-5,7),(3,-3,4)\bigr).}
\]
Las coordenadas de los otros dos vectores respecto de esta base son, respectivamente,
\[
[(1,1,-2)]_{\mathcal B}=\begin{pmatrix}-2\\3\end{pmatrix},
\qquad
[(2,-1,1)]_{\mathcal B}=\begin{pmatrix}-1\\2\end{pmatrix}.
\]
\end{solvedexercise}

\subsection{Cambio de base}

Las coordenadas de un vector dependen de la base utilizada, aunque el vector sea el mismo. Este hecho es especialmente relevante en ingeniería: una fuerza física no cambia al girar el sistema de ejes, pero sí cambian sus componentes; del mismo modo, un estado puede describirse mediante variables originales o mediante combinaciones más convenientes de ellas.

Sean
\[
\mathcal B=(v_1,\ldots,v_n),
\qquad
\mathcal B'=(v'_1,\ldots,v'_n)
\]
dos bases de $V$. Definimos la matriz de cambio de base de $\mathcal B'$ a $\mathcal B$ como
\[
P_{\mathcal B\leftarrow\mathcal B'}
=
\begin{pmatrix}
|&&|\\
[v'_1]_{\mathcal B}&\cdots&[v'_n]_{\mathcal B}\\
|&&|
\end{pmatrix}.
\]
Entonces
\[
\boxed{
[v]_{\mathcal B}
=P_{\mathcal B\leftarrow\mathcal B'}[v]_{\mathcal B'}.
}
\]
Como ambas familias son bases, esta matriz es invertible y
\[
P_{\mathcal B'\leftarrow\mathcal B}
=
P_{\mathcal B\leftarrow\mathcal B'}^{-1}.
\]

\begin{example}
En $\R^2$ consideremos
\[
\mathcal E=((1,0),(0,1)),
\qquad
\mathcal B=((1,1),(1,-1)).
\]
La matriz que transforma coordenadas en la base $\mathcal B$ en coordenadas canónicas es
\[
P_{\mathcal E\leftarrow\mathcal B}
=
\begin{pmatrix}
1&1\\
1&-1
\end{pmatrix}.
\]
Para el vector $v=(4,2)$ buscamos $[v]_{\mathcal B}=(a,b)^t$. Debe cumplirse
\[
\begin{pmatrix}4\\2\end{pmatrix}
=
\begin{pmatrix}1&1\\1&-1\end{pmatrix}
\begin{pmatrix}a\\b\end{pmatrix}.
\]
Resolviendo,
\[
a=3,\qquad b=1,
\]
por lo que
\[
\boxed{[v]_{\mathcal B}=\begin{pmatrix}3\\1\end{pmatrix}.}
\]
En efecto,
\[
(4,2)=3(1,1)+(1,-1).
\]
\end{example}

Cuando la base es ortonormal, el cálculo de las coordenadas se simplifica considerablemente, como veremos a continuación.

\subsection{Producto escalar, ortogonalidad y proyección}

Además de sumar vectores y multiplicarlos por escalares, en $\R^n$ disponemos de una operación que permite medir longitudes y ángulos.

\begin{definition}[Producto escalar]
Dados $u=(u_1,\ldots,u_n)$ y $v=(v_1,\ldots,v_n)$ en $\R^n$, su \emph{producto escalar} es
\[
\langle u,v\rangle
=u\cdot v
=u_1v_1+\cdots+u_nv_n.
\]
\end{definition}

La norma asociada es
\[
\|u\|=\sqrt{u\cdot u},
\]
y la distancia entre dos vectores es
\[
d(u,v)=\|u-v\|.
\]
Si $u,v\neq0$, el ángulo $\theta\in[0,\pi]$ que forman satisface
\[
\cos\theta
=\frac{u\cdot v}{\|u\|\,\|v\|}.
\]

\begin{definition}[Ortogonalidad]
Dos vectores $u,v\in\R^n$ son \emph{ortogonales} si
\[
u\cdot v=0.
\]
Una familia es \emph{ortogonal} si sus vectores son ortogonales dos a dos, y es \emph{ortonormal} si, además, todos tienen norma $1$.
\end{definition}

Toda familia ortogonal formada por vectores no nulos es linealmente independiente. Por esta razón, una base ortonormal proporciona una representación especialmente estable y cómoda.

Si $u\neq0$, la proyección ortogonal de $v$ sobre la dirección generada por $u$ es
\[
\operatorname{proj}_{u}(v)
=
\frac{v\cdot u}{u\cdot u}\,u.
\]
Si $e$ es unitario,
\[
\operatorname{proj}_{e}(v)=(v\cdot e)e.
\]

\begin{example}[Componente de una fuerza en una dirección]
Una fuerza está representada por
\[
F=(6,2,-3)\ \mathrm{N},
\]
y una barra tiene la dirección
\[
u=(2,1,2).
\]
El vector unitario en la dirección de la barra es
\[
e=\frac{u}{\|u\|}=\frac{1}{3}(2,1,2).
\]
La componente escalar de $F$ sobre la barra es
\[
F\cdot e
=
(6,2,-3)\cdot\frac13(2,1,2)
=\frac{12+2-6}{3}
=\frac83\ \mathrm{N}.
\]
La proyección vectorial es
\[
\operatorname{proj}_{u}(F)
=(F\cdot e)e
=\frac{8}{9}(2,1,2).
\]
La operación es puramente algebraica, pero expresa la parte de la fuerza que actúa en una dirección físicamente relevante.
\end{example}

Si $\mathcal B=(e_1,\ldots,e_n)$ es una base ortonormal de $\R^n$, las coordenadas de cualquier vector $v$ se obtienen directamente mediante productos escalares:
\[
\boxed{
[v]_{\mathcal B}
=
\begin{pmatrix}
v\cdot e_1\\
\vdots\\
v\cdot e_n
\end{pmatrix}.
}
\]
Más generalmente, si $W\subseteq\R^n$ tiene una base ortonormal $(e_1,\ldots,e_k)$, la proyección ortogonal de $v$ sobre $W$ es
\[
\boxed{
\operatorname{proj}_{W}(v)
=\sum_{i=1}^{k}(v\cdot e_i)e_i.
}
\]
El vector
\[
v-\operatorname{proj}_{W}(v)
\]
es ortogonal a todos los vectores de $W$. De este modo se obtiene la descomposición
\[
v=v_W+v_W^{\perp},
\qquad
v_W\in W,
\quad
v_W^{\perp}\in W^{\perp}.
\]

\subsection{Método de ortogonalización de Gram--Schmidt}

No siempre una base viene dada en forma ortogonal. El método de Gram--Schmidt permite sustituir una familia linealmente independiente por otra ortogonal que genera exactamente el mismo subespacio.

\begin{theorem}[Proceso de Gram--Schmidt]
Sean $v_1,\ldots,v_k$ vectores linealmente independientes de $\R^n$. Definimos
\[
u_1=v_1,
\]
y sucesivamente
\[
u_j
=v_j-
\sum_{i=1}^{j-1}
\frac{v_j\cdot u_i}{u_i\cdot u_i}\,u_i,
\qquad j=2,\ldots,k.
\]
Entonces $u_1,\ldots,u_k$ son ortogonales y
\[
\langle u_1,\ldots,u_k\rangle
=
\langle v_1,\ldots,v_k\rangle.
\]
Normalizando,
\[
e_i=\frac{u_i}{\|u_i\|},
\]
se obtiene una base ortonormal del mismo subespacio.
\end{theorem}

La fórmula tiene una interpretación sencilla: de cada nuevo vector eliminamos las componentes que apuntan en las direcciones ya construidas.

\begin{solvedexercise}
Obtener mediante Gram--Schmidt una base ortonormal del subespacio
\[
W=\langle(1,1,0),(1,0,1)\rangle\subseteq\R^3.
\]
\tcblower
\textbf{Resolución.}\par\smallskip
Tomamos
\[
v_1=(1,1,0),
\qquad
v_2=(1,0,1).
\]
El primer vector ortogonal es
\[
u_1=v_1=(1,1,0).
\]
Para el segundo,
\[
u_2
=v_2-\frac{v_2\cdot u_1}{u_1\cdot u_1}u_1.
\]
Como
\[
v_2\cdot u_1=1,
\qquad
u_1\cdot u_1=2,
\]
tenemos
\[
u_2
=(1,0,1)-\frac12(1,1,0)
=\left(\frac12,-\frac12,1\right).
\]
Multiplicar $u_2$ por $2$ no altera la dirección ni la ortogonalidad, por lo que podemos utilizar
\[
\widetilde u_2=(1,-1,2).
\]
En efecto,
\[
u_1\cdot\widetilde u_2=1-1=0.
\]
Normalizamos:
\[
e_1=\frac{1}{\sqrt2}(1,1,0),
\qquad
 e_2=\frac{1}{\sqrt6}(1,-1,2).
\]
Así,
\[
\boxed{
\mathcal B_o=
\left(
\frac{1}{\sqrt2}(1,1,0),
\frac{1}{\sqrt6}(1,-1,2)
\right)
}
\]
es una base ortonormal de $W$.
\end{solvedexercise}

En problemas numéricos, trabajar con direcciones ortonormales evita que unas coordenadas queden artificialmente amplificadas por la longitud o la proximidad entre los vectores de una base. Por ello la ortogonalización aparece de forma recurrente en cálculo numérico, procesamiento de señales, ajuste de modelos y descomposición de movimientos.

\section{Aplicaciones lineales}
\label{sec:aplicaciones-lineales}

\subsection{Concepto de aplicación lineal}

Una aplicación lineal es una transformación que respeta las combinaciones lineales. Esta propiedad hace que conocer la respuesta de la transformación sobre una base sea suficiente para conocer su respuesta sobre cualquier vector.

\begin{definition}[Aplicación lineal]
Sean $V$ y $W$ espacios vectoriales reales. Una aplicación
\[
T:V\longrightarrow W
\]
es \emph{lineal} si para cualesquiera $u,v\in V$ y $\alpha,\beta\in\R$ se verifica
\[
T(\alpha u+\beta v)
=\alpha T(u)+\beta T(v).
\]
\end{definition}

Equivalentemente, basta comprobar
\[
T(u+v)=T(u)+T(v),
\qquad
T(\alpha u)=\alpha T(u).
\]
Toda aplicación lineal satisface necesariamente
\[
T(0)=0.
\]
Por tanto, una transformación que envíe el vector nulo a un vector no nulo no puede ser lineal.

\begin{example}
La aplicación
\[
T:\R^3\longrightarrow\R^2,
\qquad
T(x,y,z)=(x+y,y+z)
\]
es lineal. En efecto,
\[
T(x,y,z)
=
\begin{pmatrix}
1&1&0\\
0&1&1
\end{pmatrix}
\begin{pmatrix}x\\y\\z\end{pmatrix}.
\]
Toda transformación definida mediante un producto matricial $T(x)=Ax$ es lineal.

En cambio,
\[
S(x,y)=(x^2,y)
\]
no es lineal, pues el primer componente no respeta la multiplicación por escalares.
\end{example}

Los modelos matriciales del capítulo anterior son, por tanto, aplicaciones lineales. Si $x$ contiene cantidades de productos y $Ax$ los recursos consumidos, la propiedad
\[
A(\alpha x+\beta y)=\alpha Ax+\beta Ay
\]
expresa que, dentro de un modelo lineal, combinar dos planes de producción combina en la misma proporción sus consumos.

\subsection{Núcleo e imagen}

Dos subespacios asociados a una aplicación lineal describen gran parte de su comportamiento.

\begin{definition}[Núcleo e imagen]
Sea $T:V\to W$ una aplicación lineal.
\begin{itemize}
\item El \emph{núcleo} de $T$ es
\[
\ker(T)=\{v\in V:T(v)=0\}.
\]
\item La \emph{imagen} de $T$ es
\[
\operatorname{Im}(T)=\{T(v):v\in V\}.
\]
\end{itemize}
\end{definition}

Tanto $\ker(T)$ como $\operatorname{Im}(T)$ son subespacios vectoriales, respectivamente de $V$ y $W$.

Si $T(x)=Ax$, entonces
\[
\ker(T)=\{x:Ax=0\},
\]
de modo que calcular el núcleo equivale a resolver un sistema homogéneo. Por otro lado, la imagen está generada por las columnas de $A$:
\[
\operatorname{Im}(T)
=
\langle c_1,\ldots,c_n\rangle,
\]
donde $c_1,\ldots,c_n$ son las columnas de $A$.

\begin{solvedexercise}
Para
\[
T:\R^3\longrightarrow\R^2,
\qquad
T(x,y,z)=(x+y,y+z),
\]
determinar una base del núcleo y una base de la imagen.
\tcblower
\textbf{Resolución.}\par\smallskip
La matriz de $T$ en las bases canónicas es
\[
A=
\begin{pmatrix}
1&1&0\\
0&1&1
\end{pmatrix}.
\]
Para calcular el núcleo resolvemos
\[
\begin{cases}
x+y=0,\\
y+z=0.
\end{cases}
\]
Tomando $y=t$,
\[
x=-t,
\qquad
z=-t,
\]
por lo que
\[
\ker(T)=\langle(-1,1,-1)\rangle.
\]
Así,
\[
\boxed{\dim\ker(T)=1.}
\]

La imagen está generada por las columnas
\[
(1,0),\quad(1,1),\quad(0,1).
\]
Las dos primeras son linealmente independientes, de modo que
\[
\operatorname{Im}(T)=\R^2,
\qquad
\boxed{\dim\operatorname{Im}(T)=2.}
\]
En este ejemplo hay una dirección no nula de entrada que la transformación convierte en cero, pero cualquier vector de salida de $\R^2$ puede obtenerse a partir de alguna entrada.
\end{solvedexercise}

Las dimensiones del núcleo y la imagen no son independientes.

\begin{theorem}[Teorema de la dimensión]
Sea $T:V\to W$ una aplicación lineal y supongamos que $V$ tiene dimensión finita. Entonces
\[
\boxed{
\dim(V)=\dim\ker(T)+\dim\operatorname{Im}(T).
}
\]
\end{theorem}

Si $T(x)=Ax$ y $A$ tiene $n$ columnas,
\[
\dim\operatorname{Im}(T)=\operatorname{rg}(A),
\]
por lo que
\[
\boxed{
\dim\ker(T)=n-\operatorname{rg}(A).
}
\]
Esta igualdad es la formulación vectorial del número de parámetros libres de un sistema homogéneo.

\begin{definition}[Aplicación inyectiva y suprayectiva]
Una aplicación lineal $T:V\to W$ es
\begin{itemize}
\item \emph{inyectiva} si $T(u)=T(v)$ implica $u=v$;
\item \emph{suprayectiva} si $\operatorname{Im}(T)=W$;
\item \emph{biyectiva} si es simultáneamente inyectiva y suprayectiva.
\end{itemize}
\end{definition}

Para una aplicación lineal,
\[
T\text{ es inyectiva}
\iff
\ker(T)=\{0\}.
\]
Si $V$ y $W$ tienen la misma dimensión finita, inyectividad, suprayectividad y biyectividad son equivalentes.

En un modelo de medida $y=Ax$, un núcleo no trivial significa que existen variaciones de las variables de entrada que no producen ningún cambio en las medidas. En ese caso, a partir de esas medidas no puede recuperarse de forma única el estado original. Si, por el contrario, $\ker(T)=\{0\}$, entradas distintas producen salidas distintas.

\subsection{Matriz de una aplicación lineal}

Una aplicación lineal queda completamente determinada por la imagen de una base.

\begin{theorem}
Sea $\mathcal B=(v_1,\ldots,v_n)$ una base de $V$. Si $T:V\to W$ es lineal, conocer
\[
T(v_1),\ldots,T(v_n)
\]
determina $T$ sobre todo $V$.
\end{theorem}

En efecto, si
\[
v=x_1v_1+\cdots+x_nv_n,
\]
entonces
\[
T(v)=x_1T(v_1)+\cdots+x_nT(v_n).
\]

Sean ahora
\[
\mathcal B=(v_1,\ldots,v_n)
\]
una base de $V$ y
\[
\mathcal C=(w_1,\ldots,w_m)
\]
una base de $W$.

\begin{definition}[Matriz de una aplicación lineal respecto de dos bases]
La matriz de $T:V\to W$ respecto de la base $\mathcal B$ en el dominio y la base $\mathcal C$ en el codominio es la matriz
\[
[T]_{\mathcal C\leftarrow\mathcal B}
\]
cuyas columnas son las coordenadas de las imágenes de los vectores de $\mathcal B$ respecto de $\mathcal C$:
\[
[T]_{\mathcal C\leftarrow\mathcal B}
=
\begin{pmatrix}
|&&|\\
[T(v_1)]_{\mathcal C}&\cdots&[T(v_n)]_{\mathcal C}\\
|&&|
\end{pmatrix}.
\]
\end{definition}

Con esta convención,
\[
\boxed{
[T(v)]_{\mathcal C}
=
[T]_{\mathcal C\leftarrow\mathcal B}
[v]_{\mathcal B}.
}
\]

\begin{example}
Consideremos
\[
T:\R^2\longrightarrow\R^2,
\qquad
T(x,y)=(2x-y,x+y),
\]
y las bases
\[
\mathcal B=((1,1),(1,-1)),
\qquad
\mathcal C=((1,0),(1,1)).
\]
Calculamos primero
\[
T(1,1)=(1,2),
\qquad
T(1,-1)=(3,0).
\]
Para expresar $(1,2)$ en la base $\mathcal C$ escribimos
\[
(1,2)=a(1,0)+b(1,1)=(a+b,b),
\]
de donde $b=2$ y $a=-1$. Así,
\[
[T(1,1)]_{\mathcal C}
=\begin{pmatrix}-1\\2\end{pmatrix}.
\]
Análogamente,
\[
(3,0)=3(1,0)+0(1,1),
\]
por lo que
\[
[T(1,-1)]_{\mathcal C}
=\begin{pmatrix}3\\0\end{pmatrix}.
\]
Por tanto,
\[
\boxed{
[T]_{\mathcal C\leftarrow\mathcal B}
=
\begin{pmatrix}
-1&3\\
2&0
\end{pmatrix}.
}
\]
\end{example}

Si $\mathcal E$ denota la base canónica, la matriz habitual de $T(x)=Ax$ es precisamente
\[
[T]_{\mathcal E\leftarrow\mathcal E}=A.
\]

\subsection{Relación entre las matrices en bases distintas}

La expresión matricial de una aplicación cambia cuando cambian las coordenadas, pero la aplicación lineal es la misma. Sean $\mathcal B,\mathcal B'$ bases del dominio y $\mathcal C,\mathcal C'$ bases del codominio. Si
\[
A=[T]_{\mathcal C\leftarrow\mathcal B},
\]
entonces
\[
\boxed{
[T]_{\mathcal C'\leftarrow\mathcal B'}
=
P_{\mathcal C\leftarrow\mathcal C'}^{-1}
A
P_{\mathcal B\leftarrow\mathcal B'}.
}
\]
La fórmula simplemente encadena tres operaciones:
\[
\begin{array}{c}
\text{coordenadas en }\mathcal B'\\[-1mm]
\downarrow\\[-1mm]
\text{coordenadas en }\mathcal B
\end{array}
\quad\xrightarrow{\ T\ }\quad
\begin{array}{c}
\text{coordenadas en }\mathcal C\\[-1mm]
\downarrow\\[-1mm]
\text{coordenadas en }\mathcal C'
\end{array}.
\]

\begin{example}
En el ejemplo anterior, la matriz de $T$ en bases canónicas es
\[
A=
\begin{pmatrix}
2&-1\\
1&1
\end{pmatrix}.
\]
Además,
\[
P_{\mathcal E\leftarrow\mathcal B}
=
\begin{pmatrix}
1&1\\
1&-1
\end{pmatrix},
\qquad
P_{\mathcal E\leftarrow\mathcal C}
=
\begin{pmatrix}
1&1\\
0&1
\end{pmatrix}.
\]
Por tanto,
\[
[T]_{\mathcal C\leftarrow\mathcal B}
=
P_{\mathcal E\leftarrow\mathcal C}^{-1}
A
P_{\mathcal E\leftarrow\mathcal B}
=
\begin{pmatrix}
-1&3\\
2&0
\end{pmatrix},
\]
como habíamos obtenido calculando directamente las imágenes de los vectores de la base.
\end{example}

Cuando $T:V\to V$ es un endomorfismo y se utiliza la misma base para entrada y salida, si $A=[T]_{\mathcal B\leftarrow\mathcal B}$ y $A'=[T]_{\mathcal B'\leftarrow\mathcal B'}$, entonces
\[
\boxed{
A'=P^{-1}AP,
\qquad
P=P_{\mathcal B\leftarrow\mathcal B'}.
}
\]
Dos matrices relacionadas de esta forma se denominan \emph{semejantes}. Esta relación será el punto de partida del capítulo siguiente: diagonalizar una matriz consistirá en encontrar una base en la que la matriz de un endomorfismo adopte una forma diagonal.

\begin{example}[Un cambio de variables en un modelo lineal]
Supongamos que el estado de un sistema mecánico se describe mediante dos desplazamientos $x_1,x_2$. En lugar de ellos puede ser conveniente utilizar
\[
q_1=x_1+x_2,
\qquad
q_2=x_1-x_2,
\]
que representan un movimiento conjunto y un movimiento relativo. Los vectores
\[
(1,1),\qquad(1,-1)
\]
forman una base de $\R^2$. Cambiar de $(x_1,x_2)$ a $(q_1,q_2)$ no cambia el estado físico: cambia su representación. Si la evolución lineal del sistema está descrita por una matriz, esa matriz también cambia de acuerdo con la fórmula anterior. En el capítulo de diagonalización veremos cómo una elección adecuada de base puede desacoplar ciertos modelos.
\end{example}

\section*{Ejercicios propuestos}
\addcontentsline{toc}{section}{Ejercicios propuestos}
\label{sec:ejercicios-t6}

Los ejercicios siguientes combinan cálculo, razonamiento y traducción entre las distintas formas de describir un espacio vectorial o una aplicación lineal. Los primeros recuperan el nivel de los ejercicios utilizados en cursos anteriores y la colección se amplía para cubrir cambio de base, producto escalar, Gram--Schmidt y aplicaciones lineales. Los marcados con \M son especialmente adecuados para ser resueltos o comprobados también con \textit{Wolfram Mathematica}.

\begin{enumerate}[label=\textbf{\arabic*.},leftmargin=*]

\item Determinar razonadamente cuáles de los siguientes conjuntos son subespacios vectoriales del espacio indicado:
\begin{enumerate}[label=\alph*)]
\item $\{(a+1,a+1):a\in\R\}\subseteq\R^2$;
\item $\{(a+b,a+b+2):a,b\in\R\}\subseteq\R^2$;
\item $\{(x,y)\in\R^2:x=y\}$;
\item el conjunto de matrices simétricas de $\mathcal M_2(\R)$;
\item $\{p\in\mathcal P_3:p(1)=0\}$;
\item $\{p\in\mathcal P_3:p(1)=1\}$.
\end{enumerate}

\item Sea
\[
S=\{(x,y,z)\in\R^3:x+y-z=0\}.
\]
Demostrar que $S$ es un subespacio, obtener una base, calcular su dimensión y escribir sus ecuaciones paramétricas.

\item Dado
\[
W=\{(x,y,z,t)\in\R^4:x+y+z+t=0,\ y-2z-t=0\},
\]
hallar una base y la dimensión de $W$. Escribir todos sus elementos en forma paramétrica.

\item Obtener una base del subespacio
\[
H=\{(x,y,z,t)\in\R^4:2y+4z=0,\ x+2z=0\}.
\]
Determinar además unas ecuaciones implícitas equivalentes en las que aparezca el menor número posible de ecuaciones independientes.

\item Sea
\[
V=\{(4,-5,7),(3,-3,4),(1,1,-2),(2,-1,1)\}\subseteq\R^3.
\]
\begin{enumerate}[label=\alph*)]
\item Obtener una base de $\langle V\rangle$ formada por vectores de $V$.
\item Calcular $\dim\langle V\rangle$.
\item Decidir si $(1,-2,3)$ pertenece a $\langle V\rangle$.
\item Si pertenece, obtener sus coordenadas respecto de la base elegida.
\end{enumerate}

\item En $\R^4$ se considera
\[
V=\{(1,0,-1,2),(-1,1,0,1),(2,2,-4,10)\}.
\]
\begin{enumerate}[label=\alph*)]
\item Determinar una base y la dimensión de $\langle V\rangle$.
\item Escribir unas ecuaciones paramétricas del subespacio.
\item Estudiar si $w=(2,-1,-1,1)$ pertenece a él.
\end{enumerate}

\item Dado
\[
L=\langle(2,5,0,-2),(0,1,3,1)\rangle\subseteq\R^4,
\]
obtener unas ecuaciones implícitas de $L$, su dimensión y una base.

\item Dados
\[
u_1=(2,3,-1),\qquad
u_2=(0,5,2),\qquad
u_3=(6,-1,-7),
\]
responder razonadamente:
\begin{enumerate}[label=\alph*)]
\item ¿son $u_1,u_2,u_3$ linealmente independientes?;
\item determinar una base y la dimensión de $L=\langle u_1,u_2,u_3\rangle$;
\item decidir si $w=(1,2,0)$ pertenece a $L$;
\item obtener unas ecuaciones implícitas de $H=\langle u_1,u_2\rangle$;
\item relacionar las respuestas anteriores con el rango de la matriz que tiene $u_1,u_2,u_3$ como columnas.
\end{enumerate}

\item Determinar para qué valores de $a\in\R$ los vectores
\[
(1,a,1),\qquad(1,1,a),\qquad(a,1,1)
\]
son linealmente independientes. Para los valores restantes, obtener una relación de dependencia no trivial.

\item Sean
\[
\begin{aligned}
v_1&=(1,1,0,1),&\qquad v_2&=(2,1,1,0),\\
v_3&=(3,2,1,1),& v_4&=(0,1,-1,2).
\end{aligned}
\]
Obtener un subconjunto maximal de vectores linealmente independientes y una base del subespacio generado. Expresar los vectores eliminados como combinación lineal de la base elegida.

\item Determinar si
\[
\mathcal B=((1,1,0),(0,1,1),(1,0,1))
\]
es una base de $\R^3$. En caso afirmativo, calcular las coordenadas de $v=(3,2,1)$ respecto de $\mathcal B$.

\item En $\R^2$ se consideran las bases
\[
\mathcal B=((1,1),(1,-1)),
\qquad
\mathcal C=((2,1),(1,1)).
\]
\begin{enumerate}[label=\alph*)]
\item Obtener $P_{\mathcal E\leftarrow\mathcal B}$ y $P_{\mathcal E\leftarrow\mathcal C}$.
\item Calcular $P_{\mathcal C\leftarrow\mathcal B}$.
\item Si $[v]_{\mathcal B}=(2,-1)^t$, hallar $[v]_{\mathcal C}$ y las coordenadas canónicas de $v$.
\item Comprobar el resultado reconstruyendo el vector a partir de ambas bases.
\end{enumerate}

\item Sea
\[
\mathcal B=((1,0,1),(1,1,0),(0,1,1))
\]
una base de $\R^3$. Hallar el vector cuyas coordenadas en $\mathcal B$ son $(2,-1,3)^t$ y obtener las coordenadas de $(1,2,3)$ respecto de $\mathcal B$.

\item Determinar una base de $\R^4$ que contenga a los vectores
\[
(1,1,0,0),\qquad(0,1,1,0).
\]
Justificar que los vectores añadidos completan realmente una base. ¿Es única la respuesta?

\item Dados
\[
u=(2,-1,3),
\qquad
v=(1,4,-2),
\]
calcular $u\cdot v$, $\|u\|$, $\|v\|$ y el ángulo que forman. Determinar también la proyección de $u$ sobre $v$.

\item Una fuerza
\[
F=(8,-2,5)\ \mathrm{N}
\]
actúa sobre un elemento cuya dirección viene dada por $d=(1,2,2)$. Calcular la componente de $F$ en la dirección de $d$ y descomponer
\[
F=F_{\parallel}+F_{\perp},
\]
donde $F_{\parallel}$ es paralela a $d$ y $F_{\perp}$ es ortogonal a $d$.

\item Aplicar Gram--Schmidt a
\[
v_1=(1,1,0),
\qquad
v_2=(1,0,1),
\qquad
v_3=(0,1,1),
\]
y obtener una base ortonormal de $\R^3$. Comprobar explícitamente los productos escalares entre los vectores obtenidos.

\item Aplicar Gram--Schmidt a la base
\[
\mathcal B=((1,1,1),(1,1,0),(1,0,0))
\]
de $\R^3$. Comparar las coordenadas de $v=(2,1,3)$ en la base original y en la base ortonormal obtenida.

\item Determinar cuáles de las siguientes aplicaciones son lineales. En los casos lineales, obtener su matriz en las bases canónicas:
\begin{enumerate}[label=\alph*)]
\item $T:\R^2\to\R^3$, $T(x,y)=(x-y,2x+y,3y)$;
\item $S:\R^2\to\R^2$, $S(x,y)=(x+1,y)$;
\item $F:\R^3\to\R$, $F(x,y,z)=2x-y+4z$;
\item $G:\R^2\to\R^2$, $G(x,y)=(xy,y)$.
\end{enumerate}

\item Para
\[
T:\R^3\to\R^2,
\qquad
T(x,y,z)=(x+y,y+z),
\]
obtener bases de $\ker(T)$ y $\operatorname{Im}(T)$, calcular sus dimensiones y verificar el teorema de la dimensión. Decidir si $T$ es inyectiva y si es suprayectiva.

\item Sea
\[
T:\R^4\to\R^3,
\qquad
T(x_1,x_2,x_3,x_4)
=(x_1+x_2,\ x_2+x_3,\ x_3+x_4).
\]
\begin{enumerate}[label=\alph*)]
\item Obtener la matriz de $T$ en las bases canónicas.
\item Calcular una base del núcleo.
\item Calcular una base de la imagen.
\item Verificar el teorema de la dimensión.
\item Decidir si la salida $(1,0,1)$ puede producirse y, en caso afirmativo, encontrar todas las entradas que la generan.
\end{enumerate}

\item Sea $T:\mathcal P_2\to\mathcal P_2$ la aplicación
\[
T(p(x))=p'(x)+p(0).
\]
\begin{enumerate}[label=\alph*)]
\item Comprobar que $T$ es lineal.
\item Obtener su matriz respecto de la base $(1,x,x^2)$.
\item Determinar $\ker(T)$ y $\operatorname{Im}(T)$.
\item Verificar el teorema de la dimensión.
\end{enumerate}

\item Una instalación dispone de tres actuadores cuyas intensidades se recogen en
\[
x=(x_1,x_2,x_3)^t.
\]
Dos sensores proporcionan las medidas
\[
y_1=x_1+x_2-2x_3,
\qquad
y_2=2x_1+2x_2-4x_3.
\]
Interpretar el sistema como una aplicación lineal $T:\R^3\to\R^2$ y:
\begin{enumerate}[label=\alph*)]
\item determinar una base de $\ker(T)$ y explicar qué representan esos vectores en términos de las medidas;
\item obtener una base de $\operatorname{Im}(T)$;
\item discutir si las dos señales de los sensores aportan información independiente;
\item indicar cuántos grados de libertad del estado pueden reconstruirse a partir de estas medidas.
\end{enumerate}

\item Una planta química produce tres corrientes de proceso $P_1,P_2,P_3$. El consumo de dos recursos viene descrito por
\[
T(x)=Ax,
\qquad
A=
\begin{pmatrix}
2&1&3\\
1&2&3
\end{pmatrix},
\]
donde $x\in\R^3$ representa las cantidades producidas.
\begin{enumerate}[label=\alph*)]
\item determinar el núcleo y explicar qué significa matemáticamente que dos planes de producción difieran en un vector del núcleo;
\item calcular la dimensión de la imagen;
\item estudiar si cualquier vector de consumo de $\R^2$ puede obtenerse en el modelo lineal;
\item comentar por qué las restricciones físicas $x_i\geq0$ no forman parte de la estructura de espacio vectorial utilizada en el análisis algebraico.
\end{enumerate}

\item Para
\[
T:\R^2\to\R^2,
\qquad
T(x,y)=(2x-y,x+y),
\]
y las bases
\[
\mathcal B=((1,1),(1,-1)),
\qquad
\mathcal C=((1,0),(1,1)),
\]
obtener $[T]_{\mathcal C\leftarrow\mathcal B}$ de dos formas distintas:
\begin{enumerate}[label=\alph*)]
\item calculando las imágenes de los vectores de $\mathcal B$ y expresándolas en $\mathcal C$;
\item utilizando las matrices de cambio de base y la matriz canónica de $T$.
\end{enumerate}

\item Sea $T:\R^3\to\R^3$ la aplicación cuya matriz en la base canónica es
\[
A=
\begin{pmatrix}
1&1&0\\
0&1&1\\
1&0&1
\end{pmatrix}.
\]
Consideremos la base
\[
\mathcal B=((1,1,0),(0,1,1),(1,0,1)).
\]
Obtener $[T]_{\mathcal B\leftarrow\mathcal B}$ y comprobar que las dos matrices representan la misma transformación utilizando un vector elegido por el estudiante.

\item Sean $T:\R^2\to\R^3$ y $S:\R^3\to\R^2$ las aplicaciones
\[
T(x,y)=(x+y,x-y,2y),
\]
\[
S(x,y,z)=(x+z,y-z).
\]
Obtener las matrices canónicas de $T$, $S$, $S\circ T$ y $T\circ S$. Comprobar que
\[
[S\circ T]=[S][T],
\qquad
[T\circ S]=[T][S],
\]
cuando los productos están definidos.

\item \M Para los ejercicios 2--10, formar las matrices cuyas columnas sean los generadores y utilizar \texttt{MatrixRank} para comprobar las dimensiones obtenidas manualmente. En los ejercicios de pertenencia, comparar el rango de la matriz de generadores con el de la matriz ampliada.

\item \M Para las aplicaciones de los ejercicios 20--24, utilizar \texttt{NullSpace} para comprobar una base del núcleo y \texttt{MatrixRank} para obtener la dimensión de la imagen. Verificar en cada caso la igualdad
\[
\dim\ker(T)+\dim\operatorname{Im}(T)=\dim(V).
\]

\item \M Aplicar \texttt{Orthogonalize} a las familias de los ejercicios 17 y 18. Comparar el resultado del programa con el proceso manual de Gram--Schmidt y comprobar mediante productos escalares que la familia obtenida es ortonormal.

\item \M En los ejercicios 12, 13, 25 y 26, introducir las matrices de cambio de base y comprobar los cambios de coordenadas mediante productos matriciales. El objetivo es utilizar el ordenador para verificar el cálculo, no para sustituir la identificación de qué matriz transforma unas coordenadas en otras.

\item (Examen) Sea
\[
W=\{(x,y,z,t)\in\R^4:x+y-z+t=0,\quad 2x+y+t=0\}.
\]
\begin{enumerate}[label=\alph*)]
\item obtener una base y la dimensión de $W$;
\item decidir si $v=(1,-1,2,2)$ pertenece a $W$;
\item obtener unas ecuaciones paramétricas de $W$;
\item construir una base de $\R^4$ que contenga una base de $W$.
\end{enumerate}

\item (Examen) Sea
\[
T:\R^3\to\R^3,
\qquad
T(x,y,z)=(x+y+z,\ 2x+y,\ x+z).
\]
\begin{enumerate}[label=\alph*)]
\item obtener la matriz de $T$ en la base canónica;
\item calcular bases de $\ker(T)$ y $\operatorname{Im}(T)$;
\item decidir si $T$ es biyectiva;
\item para
\[
\mathcal B=((1,1,0),(1,0,1),(0,1,1)),
\]
obtener $[T]_{\mathcal B\leftarrow\mathcal B}$.
\end{enumerate}

\item (Examen) En $\R^3$ se considera el subespacio
\[
H=\langle(1,1,0),(1,0,1)\rangle.
\]
\begin{enumerate}[label=\alph*)]
\item obtener una base ortonormal de $H$ mediante Gram--Schmidt;
\item calcular la proyección ortogonal de $v=(2,-1,3)$ sobre $H$;
\item descomponer $v=v_H+v_H^{\perp}$ con $v_H\in H$ y $v_H^{\perp}$ ortogonal a $H$;
\item comprobar la ortogonalidad de la descomposición.
\end{enumerate}

\end{enumerate}

\cleardoublepage
\chapter{Diagonalización}
\label{ch:diagonalizacion}

En el capítulo anterior hemos visto que una aplicación lineal puede representarse mediante matrices diferentes cuando cambiamos de base. Esta libertad permite plantear una cuestión natural: \emph{¿podemos elegir una base en la que la matriz de una transformación sea especialmente sencilla?}

La situación ideal se produce cuando la matriz resulta diagonal. En ese caso, la transformación deja de mezclar las coordenadas: cada dirección básica se transforma simplemente multiplicándose por un número. Esta idea está detrás de muchas técnicas de ingeniería. Los modos normales de vibración de un sistema mecánico, la evolución de estados acoplados, ciertos modelos de transferencia de energía o la descomposición de una respuesta en direcciones independientes pueden entenderse a partir de valores y vectores propios.

En este capítulo estudiaremos primero la diagonalización por semejanza. Después veremos otra transformación de matrices, la \emph{congruencia}, que aparece de manera natural al estudiar formas cuadráticas. Aunque ambas conducen en determinados casos a matrices diagonales, representan operaciones matemáticas distintas y no deben confundirse.

\section{Valores y vectores propios}
\label{sec:valores-vectores-propios}

\subsection{Direcciones invariantes}

Sea $T:\R^n\to\R^n$ una aplicación lineal y sea $A$ su matriz en una determinada base. En general, al aplicar $T$ a un vector cambia tanto su módulo como su dirección. Sin embargo, existen direcciones especiales que se conservan.

\begin{definition}[Vector propio y valor propio]
Sea $A\in\mathcal M_n(\R)$. Un vector no nulo $v\in\R^n$ es un \emph{vector propio} de $A$ si existe un número $\lambda\in\R$ tal que
\[
Av=\lambda v.
\]
El número $\lambda$ se denomina \emph{valor propio} asociado a $v$.
\end{definition}

La igualdad $Av=\lambda v$ significa que la acción de $A$ sobre $v$ no cambia su dirección: únicamente multiplica el vector por el factor $\lambda$. Si $\lambda>1$, el vector se alarga; si $0<\lambda<1$, se contrae; si $\lambda<0$, además cambia de sentido; y si $\lambda=0$, el vector se transforma en el vector nulo.

\begin{example}
Consideremos
\[
A=\begin{pmatrix}
2&1\\
0&3
\end{pmatrix}.
\]
Para $v=(1,0)^t$ se tiene
\[
Av=\begin{pmatrix}2\\0\end{pmatrix}=2v,
\]
por lo que $v$ es un vector propio asociado al valor propio $\lambda=2$.

En cambio, para $w=(0,1)^t$,
\[
Aw=\begin{pmatrix}1\\3\end{pmatrix},
\]
que no es múltiplo de $w$. Por tanto, $w$ no es vector propio.
\end{example}

En problemas de ingeniería, un vector propio suele representar una configuración que conserva su forma bajo la transformación. Por ejemplo, si una matriz describe la evolución lineal de un estado, una condición inicial alineada con un vector propio evoluciona sin cambiar de dirección en el espacio de estados: solamente cambia su amplitud.

\subsection{Cálculo de los valores propios}

La ecuación
\[
Av=\lambda v
\]
puede escribirse como
\[
(A-\lambda I)v=0.
\]
Como buscamos soluciones no nulas, este sistema homogéneo debe ser compatible indeterminado. Por el Teorema de Rouché--Frobenius, esto sucede exactamente cuando
\[
\operatorname{rg}(A-\lambda I)<n,
\]
o, equivalentemente,
\[
\det(A-\lambda I)=0.
\]

\begin{definition}[Polinomio característico]
Sea $A\in\mathcal M_n(\R)$. El polinomio
\[
p_A(\lambda)=\det(A-\lambda I)
\]
se denomina \emph{polinomio característico} de $A$.
\end{definition}

\begin{theorem}
Un número $\lambda\in\R$ es valor propio de $A$ si y sólo si
\[
p_A(\lambda)=0.
\]
\end{theorem}

Por tanto, los valores propios reales de una matriz son las raíces reales de su polinomio característico.

\begin{example}
Sea
\[
A=\begin{pmatrix}
4&1\\
2&3
\end{pmatrix}.
\]
Entonces
\[
\begin{aligned}
p_A(\lambda)
&=\det\begin{pmatrix}4-\lambda&1\\2&3-\lambda\end{pmatrix}\\
&=(4-\lambda)(3-\lambda)-2\\
&=\lambda^2-7\lambda+10\\
&=(\lambda-2)(\lambda-5).
\end{aligned}
\]
Los valores propios son, por tanto,
\[
\boxed{\lambda_1=2,\qquad \lambda_2=5.}
\]
\end{example}

Una matriz real puede tener valores propios complejos. En este capítulo nos centraremos en la diagonalización sobre $\R$, por lo que necesitaremos valores y vectores propios reales. Cuando todos los valores propios sean reales, seguiremos trabajando dentro del espacio vectorial real que hemos utilizado hasta ahora.

\subsection{Espacios propios}

Una vez conocido un valor propio $\lambda$, sus vectores propios se obtienen resolviendo
\[
(A-\lambda I)v=0.
\]
El conjunto de soluciones es un subespacio vectorial.

\begin{definition}[Espacio propio]
Sea $\lambda$ un valor propio de $A$. El \emph{espacio propio} asociado a $\lambda$ es
\[
E_{\lambda}=\ker(A-\lambda I).
\]
Los vectores propios asociados a $\lambda$ son precisamente los vectores no nulos de $E_\lambda$.
\end{definition}

\begin{solvedexercise}
Calcular los valores propios y una base de cada espacio propio de
\[
A=\begin{pmatrix}
4&-1&1\\
0&3&0\\
1&-1&4
\end{pmatrix}.
\]
\tcblower
\textbf{Resolución.}\par\smallskip
Calculamos primero el polinomio característico:
\[
p_A(\lambda)=\det(A-\lambda I).
\]
Al desarrollar y factorizar se obtiene
\[
p_A(\lambda)=(3-\lambda)^2(5-\lambda).
\]
Por tanto, los valores propios son $\lambda=3$ y $\lambda=5$.

Para $\lambda=3$ resolvemos
\[
(A-3I)v=0,
\]
es decir,
\[
\begin{pmatrix}
1&-1&1\\
0&0&0\\
1&-1&1
\end{pmatrix}
\begin{pmatrix}x\\y\\z\end{pmatrix}=0.
\]
La única ecuación independiente es
\[
x-y+z=0,
\]
de modo que
\[
(x,y,z)=s(1,1,0)+t(-1,0,1).
\]
Así,
\[
E_3=\langle(1,1,0),(-1,0,1)\rangle.
\]

Para $\lambda=5$,
\[
(A-5I)v=0
\]
conduce a
\[
E_5=\langle(1,0,1)\rangle.
\]
En consecuencia, una base de cada espacio propio es
\[
\boxed{\mathcal B_3=((1,1,0),(-1,0,1)),\qquad
\mathcal B_5=((1,0,1)).}
\]
\end{solvedexercise}

El ejemplo anterior muestra una diferencia importante entre \emph{valor propio} y \emph{vector propio}: a un mismo valor propio pueden corresponder infinitos vectores propios, todos ellos contenidos en su espacio propio.

\subsection{Multiplicidad algebraica y geométrica}

\begin{definition}[Multiplicidad algebraica]
La \emph{multiplicidad algebraica} de un valor propio $\lambda$ es su multiplicidad como raíz del polinomio característico. La denotaremos por $m_a(\lambda)$.
\end{definition}

\begin{definition}[Multiplicidad geométrica]
La \emph{multiplicidad geométrica} de un valor propio $\lambda$ es la dimensión de su espacio propio:
\[
m_g(\lambda)=\dim E_\lambda.
\]
\end{definition}

Estas multiplicidades satisfacen
\[
1\leq m_g(\lambda)\leq m_a(\lambda).
\]
En el ejercicio anterior, $\lambda=3$ tiene multiplicidad algebraica $2$ y espacio propio de dimensión $2$, mientras que $\lambda=5$ tiene multiplicidad algebraica y geométrica iguales a $1$.

\begin{proposition}
Los vectores propios asociados a valores propios distintos son linealmente independientes.
\end{proposition}

Este resultado explica por qué una matriz con $n$ valores propios reales distintos es siempre diagonalizable: automáticamente proporciona $n$ direcciones propias linealmente independientes.

\section{Diagonalización por semejanza}
\label{sec:diagonalizacion-semejanza}

\subsection{Cambio de base y matriz diagonal}

Recordemos que si una aplicación lineal $T:\R^n\to\R^n$ tiene matriz $A$ en la base canónica y elegimos otra base $\mathcal B$, la nueva matriz está relacionada con $A$ mediante
\[
[T]_{\mathcal B\leftarrow\mathcal B}=P^{-1}AP,
\]
donde las columnas de $P$ son los vectores de la base $\mathcal B$ expresados en coordenadas canónicas.

Si escogemos una base formada por vectores propios
\[
\mathcal B=(v_1,\ldots,v_n),
\]
con
\[
Av_i=\lambda_i v_i,
\]
entonces la matriz de la transformación en esa base es
\[
D=\begin{pmatrix}
\lambda_1&0&\cdots&0\\
0&\lambda_2&\cdots&0\\
\vdots&\vdots&\ddots&\vdots\\
0&0&\cdots&\lambda_n
\end{pmatrix}.
\]
En otras palabras,
\[
P^{-1}AP=D
\qquad\Longleftrightarrow\qquad
A=PDP^{-1}.
\]

\begin{definition}[Matrices semejantes]
Dos matrices cuadradas $A$ y $B$ son \emph{semejantes} si existe una matriz regular $P$ tal que
\[
B=P^{-1}AP.
\]
\end{definition}

Matrices semejantes representan la misma aplicación lineal respecto de bases diferentes. Por ello tienen, entre otras propiedades, el mismo polinomio característico, los mismos valores propios, el mismo determinante y la misma traza.

\begin{definition}[Matriz diagonalizable]
Una matriz $A\in\mathcal M_n(\R)$ es \emph{diagonalizable sobre $\R$} si es semejante a una matriz diagonal, es decir, si existe una matriz regular $P$ y una matriz diagonal $D$ tales que
\[
A=PDP^{-1}.
\]
\end{definition}

\begin{theorem}[Criterio fundamental de diagonalización]
Una matriz $A\in\mathcal M_n(\R)$ es diagonalizable si y sólo si existe una base de $\R^n$ formada por vectores propios de $A$.
\end{theorem}

Por tanto, diagonalizar una matriz consiste esencialmente en encontrar suficientes vectores propios linealmente independientes.

\begin{solvedexercise}
Diagonalizar, si es posible, la matriz
\[
A=\begin{pmatrix}
4&-1&1\\
0&3&0\\
1&-1&4
\end{pmatrix}.
\]
\tcblower
\textbf{Resolución.}\par\smallskip
En la sección anterior obtuvimos
\[
E_3=\langle(1,1,0),(-1,0,1)\rangle,
\qquad
E_5=\langle(1,0,1)\rangle.
\]
Disponemos, por tanto, de tres vectores propios linealmente independientes. Tomamos
\[
P=\begin{pmatrix}
1&-1&1\\
1&0&0\\
0&1&1
\end{pmatrix}.
\]
Como las dos primeras columnas corresponden a $\lambda=3$ y la tercera a $\lambda=5$, definimos
\[
D=\begin{pmatrix}
3&0&0\\
0&3&0\\
0&0&5
\end{pmatrix}.
\]
Entonces
\[
AP=PD,
\]
y, puesto que $P$ es regular,
\[
\boxed{A=PDP^{-1}.}
\]
La matriz $A$ es diagonalizable.
\end{solvedexercise}

El orden de los valores propios en $D$ debe corresponderse con el orden de los vectores propios colocados como columnas de $P$. Cambiar simultáneamente ambos órdenes produce otra diagonalización igualmente válida.

\subsection{Criterios prácticos de diagonalización}

El criterio fundamental puede expresarse en términos de las multiplicidades.

\begin{theorem}
Sea $A\in\mathcal M_n(\R)$ y supongamos que su polinomio característico se descompone completamente en factores lineales reales. Entonces $A$ es diagonalizable si y sólo si, para cada valor propio $\lambda$,
\[
m_g(\lambda)=m_a(\lambda).
\]
\end{theorem}

Equivalentemente,
\[
\sum_{\lambda}\dim E_\lambda=n.
\]

\begin{corollary}
Si $A\in\mathcal M_n(\R)$ tiene $n$ valores propios reales distintos, entonces $A$ es diagonalizable.
\end{corollary}

La recíproca es falsa: una matriz puede ser diagonalizable aunque tenga valores propios repetidos, siempre que disponga de suficientes vectores propios independientes.

\begin{example}
Las matrices
\[
A=\begin{pmatrix}
2&0\\
0&2
\end{pmatrix},
\qquad
B=\begin{pmatrix}
2&1\\
0&2
\end{pmatrix}
\]
tienen el mismo único valor propio $\lambda=2$ con multiplicidad algebraica $2$.

Para $A$,
\[
E_2=\R^2,
\]
por lo que $m_g(2)=2$ y $A$ es diagonalizable.

Para $B$,
\[
B-2I=\begin{pmatrix}0&1\\0&0\end{pmatrix},
\]
y
\[
E_2=\langle(1,0)\rangle.
\]
Así, $m_g(2)=1<2=m_a(2)$ y $B$ no es diagonalizable.
\end{example}

\begin{solvedexercise}
Determinar para qué valores de $\alpha\in\R$ es diagonalizable la matriz
\[
M=\begin{pmatrix}
2&-1&1\\
0&\alpha&0\\
2&0&3
\end{pmatrix}.
\]
\tcblower
\textbf{Resolución.}\par\smallskip
El polinomio característico es
\[
p_M(\lambda)=(\alpha-\lambda)(4-\lambda)(1-\lambda).
\]
Si $\alpha\neq1,4$, la matriz tiene tres valores propios reales distintos y, por tanto, es diagonalizable.

Debemos estudiar por separado los casos en los que aparece un valor propio repetido.

Si $\alpha=1$, el valor propio $1$ tiene multiplicidad algebraica $2$. Al resolver
\[
(M-I)v=0
\]
se obtiene un espacio propio de dimensión $1$. Luego
\[
m_g(1)=1<2=m_a(1),
\]
y la matriz no es diagonalizable.

Si $\alpha=4$, el valor propio $4$ tiene multiplicidad algebraica $2$ y su espacio propio tiene dimensión $1$. Tampoco es diagonalizable.

Por tanto,
\[
\boxed{M\text{ es diagonalizable si y sólo si }\alpha\in\R\setminus\{1,4\}.}
\]
\end{solvedexercise}

Este tipo de ejercicio es importante: cuando aparecen parámetros no basta con calcular formalmente los valores propios. Hay que localizar los valores del parámetro que cambian sus multiplicidades y estudiar separadamente esos casos.

\subsection{¿Para qué sirve diagonalizar?}

Si
\[
A=PDP^{-1},
\]
entonces
\[
A^2=PDP^{-1}PDP^{-1}=PD^2P^{-1},
\]
y, en general,
\[
\boxed{A^k=PD^kP^{-1}.}
\]
Como elevar una matriz diagonal a una potencia consiste simplemente en elevar sus elementos diagonales, la diagonalización simplifica de manera drástica el cálculo de potencias de matrices.

\begin{example}[Intercambio térmico entre dos cuerpos]
Supongamos que, tras cada intervalo de tiempo, las temperaturas de dos cuerpos se actualizan según
\[
\begin{pmatrix}
T_{1,k+1}\\
T_{2,k+1}
\end{pmatrix}
=
\begin{pmatrix}
\frac45&\frac15\\[1mm]
\frac15&\frac45
\end{pmatrix}
\begin{pmatrix}
T_{1,k}\\
T_{2,k}
\end{pmatrix}.
\]
La matriz
\[
A=\begin{pmatrix}
\frac45&\frac15\\[1mm]
\frac15&\frac45
\end{pmatrix}
\]
tiene los valores propios
\[
\lambda_1=1,
\qquad
\lambda_2=\frac35,
\]
con vectores propios, respectivamente,
\[
v_1=(1,1),
\qquad
v_2=(1,-1).
\]
La dirección $v_1$ representa estados en los que ambos cuerpos tienen la misma temperatura; esa componente permanece constante. La dirección $v_2$ mide la diferencia de temperaturas y se multiplica en cada paso por $3/5$.

Si escribimos el estado inicial como
\[
T_0=c_1v_1+c_2v_2,
\]
entonces
\[
T_k=A^kT_0=c_1v_1+c_2\left(\frac35\right)^kv_2.
\]
Cuando $k\to\infty$,
\[
\left(\frac35\right)^k\to0,
\]
de modo que desaparece progresivamente la diferencia de temperaturas. La diagonalización no sólo facilita el cálculo: separa el modelo en dos \emph{modos} con comportamientos diferentes.
\end{example}

El mismo principio aparece en modelos mecánicos de vibración, en sistemas lineales discretos, en redes acopladas y en modelos de transferencia entre compartimentos: elegir una base de vectores propios permite estudiar por separado componentes que, en las coordenadas originales, aparecen mezcladas.

\subsection{Matrices simétricas y diagonalización ortogonal}

Las matrices simétricas tienen propiedades especialmente favorables.

\begin{theorem}[Teorema espectral]
Toda matriz real simétrica $A=A^t$ posee valores propios reales y admite una base ortonormal de $\R^n$ formada por vectores propios.
\end{theorem}

En consecuencia, si $A$ es simétrica, existe una matriz ortogonal $P$ tal que
\[
P^tP=I,
\qquad
P^{-1}=P^t,
\]
y
\[
\boxed{P^tAP=D,}
\]
donde $D$ es diagonal.

A esta operación se la denomina \emph{diagonalización ortogonal}. Es especialmente importante en ingeniería porque muchas matrices asociadas a energía, inercia, rigidez o covarianza son simétricas.

\begin{example}[Modos de un sistema mecánico sencillo]
Consideremos la matriz simétrica
\[
K=\begin{pmatrix}
2&-1\\
-1&2
\end{pmatrix}.
\]
Sus valores propios son
\[
\lambda_1=1,
\qquad
\lambda_2=3,
\]
con vectores propios ortogonales
\[
v_1=(1,1),
\qquad
v_2=(1,-1).
\]
Tras normalizarlos obtenemos
\[
u_1=\frac1{\sqrt2}(1,1),
\qquad
u_2=\frac1{\sqrt2}(1,-1).
\]
La matriz
\[
P=\frac1{\sqrt2}
\begin{pmatrix}
1&1\\
1&-1
\end{pmatrix}
\]
es ortogonal y satisface
\[
P^tKP=\begin{pmatrix}1&0\\0&3\end{pmatrix}.
\]
Las dos direcciones propias representan dos patrones independientes de deformación del sistema: uno en el que ambas componentes evolucionan conjuntamente y otro en el que lo hacen en sentidos opuestos.
\end{example}

\section{Formas cuadráticas y diagonalización por congruencia}
\label{sec:formas-cuadraticas}

\subsection{Formas cuadráticas}

Las expresiones cuadráticas en varias variables aparecen de manera natural al describir energías, errores, distancias o funciones de coste aproximadas. En dos variables, una expresión general de segundo grado homogénea es
\[
Q(x,y)=ax^2+2bxy+cy^2.
\]
Puede escribirse matricialmente como
\[
Q(x,y)=
\begin{pmatrix}x&y\end{pmatrix}
\begin{pmatrix}a&b\\b&c\end{pmatrix}
\begin{pmatrix}x\\y\end{pmatrix}.
\]

\begin{definition}[Forma cuadrática]
Una \emph{forma cuadrática} en $\R^n$ es una aplicación $Q:\R^n\to\R$ de la forma
\[
Q(x)=x^tAx,
\]
donde $A\in\mathcal M_n(\R)$ puede tomarse simétrica.
\end{definition}

La matriz simétrica asociada a una forma cuadrática es única. Los términos cuadrados ocupan la diagonal y cada coeficiente de un término cruzado $x_ix_j$ se reparte por mitades en las posiciones $(i,j)$ y $(j,i)$.

\begin{example}
La forma cuadrática
\[
Q(x,y,z)=3x^2+2y^2+5z^2+4xy-6xz+2yz
\]
tiene matriz asociada
\[
A=\begin{pmatrix}
3&2&-3\\
2&2&1\\
-3&1&5
\end{pmatrix},
\]
pues
\[
Q(x,y,z)=
\begin{pmatrix}x&y&z\end{pmatrix}
A
\begin{pmatrix}x\\y\\z\end{pmatrix}.
\]
Obsérvese que, por ejemplo, las dos entradas $a_{12}=a_{21}=2$ producen conjuntamente el término $4xy$.
\end{example}

Si una expresión cuadrática se proporciona inicialmente mediante una matriz no simétrica $B$, podemos sustituirla por su parte simétrica
\[
A=\frac{B+B^t}{2},
\]
porque
\[
x^tBx=x^tAx.
\]
Por ello, en el estudio de formas cuadráticas podemos trabajar siempre con matrices simétricas.

\subsection{Cambio de variables y congruencia}

Supongamos que realizamos un cambio lineal de variables
\[
x=Py,
\]
donde $P$ es una matriz regular. Entonces
\[
Q(x)=x^tAx=(Py)^tA(Py)=y^tP^tAPy.
\]
Por tanto, la matriz de la forma cuadrática en las nuevas coordenadas es
\[
\boxed{B=P^tAP.}
\]

\begin{definition}[Matrices congruentes]
Dos matrices cuadradas $A$ y $B$ son \emph{congruentes} si existe una matriz regular $P$ tal que
\[
B=P^tAP.
\]
\end{definition}

Esta expresión se parece a la relación de semejanza, pero representa una transformación distinta:
\[
\boxed{\text{Semejanza: }B=P^{-1}AP,}
\]
\[
\boxed{\text{Congruencia: }B=P^tAP.}
\]
La semejanza corresponde al cambio de base de una aplicación lineal de $V$ en sí mismo. La congruencia corresponde al cambio de variables de una forma cuadrática.

\begin{example}
Consideremos
\[
Q(x,y)=2x^2+4xy+5y^2,
\qquad
A=\begin{pmatrix}2&2\\2&5\end{pmatrix}.
\]
Introducimos el cambio
\[
\begin{pmatrix}x\\y\end{pmatrix}
=P\begin{pmatrix}u\\v\end{pmatrix}.
\]
Si elegimos $P$ formado por vectores propios ortonormales de $A$, el Teorema espectral garantiza que
\[
P^tAP=D
\]
es diagonal. Así, en las nuevas variables desaparece el término cruzado $uv$ y la forma cuadrática se escribe como suma de cuadrados ponderados.
\end{example}

Cuando $P$ es ortogonal, se cumple $P^{-1}=P^t$. En ese caso particular,
\[
P^{-1}AP=P^tAP,
\]
de modo que una diagonalización ortogonal de una matriz simétrica es simultáneamente una diagonalización por semejanza y por congruencia. Esta coincidencia es una propiedad especial del caso ortogonal y no debe utilizarse para identificar ambos conceptos en general.

\subsection{Diagonalización de una forma cuadrática}

Diagonalizar una forma cuadrática significa encontrar un cambio de variables para eliminar los términos cruzados. Si $A$ es la matriz simétrica de $Q$, buscamos una matriz regular $P$ tal que
\[
P^tAP=D,
\]
con $D$ diagonal. Entonces, si $x=Py$,
\[
Q(x)=y^tDy
=d_1y_1^2+\cdots+d_ny_n^2.
\]

Por el Teorema espectral, siempre podemos elegir $P$ ortogonal. En ese caso los coeficientes diagonales son los valores propios de $A$.

\begin{solvedexercise}
Reducir a forma diagonal mediante un cambio ortogonal la forma cuadrática
\[
Q(x,y)=5x^2-4xy+2y^2.
\]
\tcblower
\textbf{Resolución.}\par\smallskip
La matriz simétrica asociada es
\[
A=\begin{pmatrix}
5&-2\\
-2&2
\end{pmatrix}.
\]
Su polinomio característico es
\[
\det(A-\lambda I)
=(5-\lambda)(2-\lambda)-4
=\lambda^2-7\lambda+6,
\]
por lo que
\[
\lambda_1=1,
\qquad
\lambda_2=6.
\]

Para $\lambda_1=1$ podemos tomar el vector propio $(1,2)$, y para $\lambda_2=6$, el vector $(2,-1)$. Son ortogonales. Normalizando,
\[
u_1=\frac1{\sqrt5}(1,2),
\qquad
u_2=\frac1{\sqrt5}(2,-1).
\]
Tomamos
\[
P=\frac1{\sqrt5}
\begin{pmatrix}
1&2\\
2&-1
\end{pmatrix}.
\]
Entonces $P$ es ortogonal y
\[
P^tAP=
\begin{pmatrix}
1&0\\
0&6
\end{pmatrix}.
\]
Si
\[
\begin{pmatrix}x\\y\end{pmatrix}
=P\begin{pmatrix}u\\v\end{pmatrix},
\]
la forma cuadrática se convierte en
\[
\boxed{Q=u^2+6v^2.}
\]
El cambio de coordenadas ha eliminado el término cruzado.
\end{solvedexercise}

Esta reducción resulta especialmente útil cuando las variables originales están acopladas. Por ejemplo, una energía elástica de la forma
\[
E(x)=\frac12x^tKx
\]
puede contener términos cruzados que expresan interacción entre desplazamientos. Si $K$ es simétrica, una base ortonormal de vectores propios permite escribir la energía como suma de contribuciones independientes.


\subsection{Diagonalización por congruencia mediante completación de cuadrados}

La diagonalización ortogonal es especialmente cómoda cuando conocemos los valores y vectores propios de la matriz simétrica. Sin embargo, para diagonalizar una forma cuadrática por congruencia no es imprescindible utilizar una matriz ortogonal. También podemos realizar cambios lineales de variables obtenidos directamente al \emph{completar cuadrados}. Este procedimiento se conoce habitualmente como método de Lagrange.

La idea consiste en transformar sucesivamente la expresión hasta eliminar los términos cruzados. Cada sustitución lineal invertible equivale a una congruencia sobre la matriz asociada.

\begin{example}
Consideremos
\[
Q(x,y)=2x^2+4xy+3y^2.
\]
Agrupamos los términos que contienen $x$:
\[
Q(x,y)=2(x^2+2xy)+3y^2.
\]
Completando el cuadrado,
\[
x^2+2xy=(x+y)^2-y^2,
\]
y, por tanto,
\[
Q(x,y)=2(x+y)^2+y^2.
\]
Definimos las nuevas variables
\[
u=x+y,
\qquad
v=y.
\]
Entonces
\[
\boxed{Q=2u^2+v^2.}
\]

El cambio inverso es
\[
x=u-v,
\qquad
y=v,
\]
es decir,
\[
\begin{pmatrix}x\\y\end{pmatrix}
=
P\begin{pmatrix}u\\v\end{pmatrix},
\qquad
P=\begin{pmatrix}1&-1\\0&1\end{pmatrix}.
\]
La matriz original es
\[
A=\begin{pmatrix}2&2\\2&3\end{pmatrix},
\]
y se verifica
\[
P^tAP=
\begin{pmatrix}2&0\\0&1\end{pmatrix}.
\]
La matriz $P$ no es ortogonal, pero sí regular. Por tanto, hemos realizado una diagonalización por congruencia sin utilizar valores propios.
\end{example}

En dimensiones mayores se procede de forma análoga: se elige un término cuadrático con coeficiente no nulo, se completan cuadrados para eliminar los términos cruzados que contienen esa variable y se repite el proceso sobre las variables restantes. Si inicialmente no aparece ningún término cuadrático pero sí términos cruzados, puede realizarse primero un cambio elemental de variables que produzca uno.

Este método resulta útil para distinguir claramente las dos ideas del capítulo: la diagonalización por semejanza busca una base de vectores propios de una aplicación lineal, mientras que la diagonalización por congruencia busca unas variables en las que una forma cuadrática quede expresada sin términos cruzados.

\subsection{Clasificación de formas cuadráticas}

El signo de una forma cuadrática describe el comportamiento de $Q(x)$ para vectores no nulos.

\begin{definition}[Clasificación por el signo]
Sea $Q(x)=x^tAx$ una forma cuadrática real. Diremos que $Q$ es:
\begin{itemize}
\item \emph{definida positiva} si $Q(x)>0$ para todo $x\neq0$;
\item \emph{definida negativa} si $Q(x)<0$ para todo $x\neq0$;
\item \emph{semidefinida positiva} si $Q(x)\geq0$ para todo $x$ y existe algún $x\neq0$ con $Q(x)=0$;
\item \emph{semidefinida negativa} si $Q(x)\leq0$ para todo $x$ y existe algún $x\neq0$ con $Q(x)=0$;
\item \emph{indefinida} si toma valores positivos y negativos.
\end{itemize}
\end{definition}

Si $A$ es simétrica y la diagonalizamos ortogonalmente,
\[
P^tAP=\operatorname{diag}(\lambda_1,\ldots,\lambda_n),
\]
entonces
\[
Q=\lambda_1y_1^2+\cdots+\lambda_ny_n^2.
\]
Por tanto, la clasificación depende directamente del signo de los valores propios.

\begin{theorem}
Sea $A=A^t$. Entonces:
\begin{itemize}
\item $Q$ es definida positiva si y sólo si todos los valores propios de $A$ son positivos;
\item $Q$ es definida negativa si y sólo si todos son negativos;
\item $Q$ es semidefinida positiva si todos son no negativos y alguno es cero;
\item $Q$ es semidefinida negativa si todos son no positivos y alguno es cero;
\item $Q$ es indefinida si existen valores propios de signos distintos.
\end{itemize}
\end{theorem}

\begin{example}[Energía y estabilidad]
Una energía aproximada alrededor de una configuración de equilibrio puede tener la forma
\[
E(x)=\frac12x^tKx,
\]
donde $K$ es una matriz simétrica de rigidez. Si $K$ es definida positiva, cualquier desplazamiento no nulo produce un incremento positivo de energía. En cambio, si la forma es indefinida, existen direcciones en las que la expresión disminuye. La clasificación algebraica de la forma cuadrática tiene, por tanto, una interpretación directa sobre la respuesta del sistema.
\end{example}

Para matrices simétricas existe también un criterio que evita calcular explícitamente los valores propios.

\begin{theorem}[Criterio de Sylvester]
Sea $A=A^t\in\mathcal M_n(\R)$ y sean
\[
\Delta_k=\det A_k,
\]
donde $A_k$ es la submatriz principal formada por las primeras $k$ filas y columnas de $A$. Entonces $A$ es definida positiva si y sólo si
\[
\Delta_1>0,\quad\Delta_2>0,\quad\ldots,\quad\Delta_n>0.
\]
\end{theorem}

Para determinar si una forma es definida negativa puede aplicarse el criterio a $-A$, o equivalentemente comprobar la alternancia de signos
\[
\Delta_1<0,\quad\Delta_2>0,\quad\Delta_3<0,\ldots
\]
El criterio de Sylvester es especialmente eficiente cuando únicamente interesa decidir la definitud y no necesitamos obtener las direcciones propias.

\begin{solvedexercise}
Clasificar la forma cuadrática
\[
Q(x,y,z)=2x^2+2y^2+3z^2+2xy-2xz.
\]
\tcblower
\textbf{Resolución.}\par\smallskip
La matriz asociada es
\[
A=\begin{pmatrix}
2&1&-1\\
1&2&0\\
-1&0&3
\end{pmatrix}.
\]
Calculamos los menores principales:
\[
\Delta_1=2>0,
\]
\[
\Delta_2=\det\begin{pmatrix}2&1\\1&2\end{pmatrix}=3>0,
\]
y
\[
\Delta_3=\det A=7>0.
\]
Por el criterio de Sylvester,
\[
\boxed{Q\text{ es definida positiva}.}
\]
No ha sido necesario calcular los valores propios para responder a la pregunta.
\end{solvedexercise}

\subsection{Ley de inercia de Sylvester}

Una forma cuadrática puede diagonalizarse por congruencia de muchas maneras diferentes. Los valores concretos de la diagonal pueden cambiar si el cambio de variables no es ortogonal, pero hay una información que permanece invariante.

\begin{theorem}[Ley de inercia de Sylvester]
Todas las matrices diagonales congruentes con una matriz simétrica real $A$ tienen el mismo número de elementos diagonales positivos, negativos y nulos.
\end{theorem}

Por tanto, aunque una transformación por congruencia pueda modificar los coeficientes de los cuadrados, no puede cambiar la naturaleza de la forma: definida positiva, definida negativa, semidefinida o indefinida.

Este resultado permite entender por qué la clasificación de una forma cuadrática es una propiedad geométrica intrínseca y no depende del sistema de coordenadas elegido.

\section*{Ejercicios propuestos}
\addcontentsline{toc}{section}{Ejercicios propuestos}

Los ejercicios de esta relación combinan cálculo directo, razonamiento sobre diagonalización, problemas con parámetros, interpretación de modelos y formas cuadráticas. Los marcados con \M están recomendados para comprobar o ampliar el trabajo manual con Wolfram Mathematica.

\begin{enumerate}

\item Calcular los valores propios y obtener una base de cada espacio propio de las matrices
\[
A=\begin{pmatrix}2&1\\1&2\end{pmatrix},
\qquad
B=\begin{pmatrix}1&2\\0&3\end{pmatrix},
\qquad
C=\begin{pmatrix}4&0&0\\1&2&0\\-1&3&1\end{pmatrix}.
\]

\item Determinar cuáles de los siguientes vectores son vectores propios de
\[
A=\begin{pmatrix}
3&0&2\\
1&3&1\\
0&1&1
\end{pmatrix}:
\qquad
u=(2,3,1),\quad v=(-2,1,1),\quad w=(-1,0,1).
\]
En cada caso afirmativo, determinar el valor propio asociado.

\item Estudiar si la matriz
\[
A=\begin{pmatrix}
4&-1&1\\
0&3&0\\
1&-1&4
\end{pmatrix}
\]
es diagonalizable. En caso afirmativo, obtener matrices $P$ y $D$ tales que $A=PDP^{-1}$.

\item Sabemos que la matriz
\[
M=\begin{pmatrix}
x&1&u\\
y&2&v\\
z&-1&w
\end{pmatrix}
\]
admite como vectores propios
\[
(1,1,0),\qquad(-1,0,2),\qquad(0,1,-1).
\]
Determinar todos los elementos de $M$ y sus valores propios.

\item De la matriz
\[
A=\begin{pmatrix}
1&2&\alpha\\
2&1&\beta\\
2&2&\gamma
\end{pmatrix}
\]
se sabe que $\lambda=1$ es un valor propio y que $v=(1,1,1)$ es un vector propio asociado.
\begin{enumerate}[label=\alph*)]
\item Hallar $\alpha$, $\beta$ y $\gamma$.
\item Calcular los restantes valores propios.
\item Estudiar si la matriz resultante es diagonalizable.
\end{enumerate}

\item Determinar para qué valores de $\alpha\in\R$ la matriz
\[
M=\begin{pmatrix}
2&-1&1\\
0&\alpha&0\\
2&0&3
\end{pmatrix}
\]
es diagonalizable.

\item Dada
\[
A=\begin{pmatrix}
b&-2&6\\
0&3&-8\\
0&0&5
\end{pmatrix},
\]
estudiar para qué valores de $b\in\R$ es diagonalizable.

\item Dada la matriz
\[
A=\begin{pmatrix}
3&0&2&0\\
1&3&1&0\\
0&1&1&0\\
0&0&0&4
\end{pmatrix},
\]
responder a las siguientes cuestiones:
\begin{enumerate}[label=\alph*)]
\item determinar cuáles de los vectores
\[
\begin{aligned}
u&=(2,3,1,0), & v&=(0,1,0,1), & w&=(-2,1,1,0),\\
r&=(0,1,0,-1), & h&=(-1,0,1,0). &&
\end{aligned}
\]
son vectores propios;
\item calcular una base y la dimensión de
\[
H=\{x\in\R^4:Ax=4x\};
\]
\item decidir si $A$ es diagonalizable;
\item dada
\[
D=\begin{pmatrix}
4&0&0&0\\
0&2&0&0\\
0&0&4&0\\
0&0&0&1
\end{pmatrix},
\]
obtener, si es posible, una matriz regular $P$ tal que $A=PDP^{-1}$.
\end{enumerate}

\item Sea
\[
A=\begin{pmatrix}
1&1&0\\
0&1&0\\
0&0&2
\end{pmatrix}.
\]
Calcular sus valores propios y sus multiplicidades algebraicas y geométricas. Explicar con precisión por qué es o no diagonalizable.

\item Construir una matriz $3\times3$ que tenga valores propios $2,2,5$ y sea diagonalizable. Construir otra con los mismos valores propios que no sea diagonalizable. Justificar ambas respuestas sin limitarse a mostrar ejemplos numéricos.

\item Sean $A$ y $B$ matrices semejantes. Justificar cuáles de las siguientes magnitudes deben coincidir y cuáles no necesariamente:
\begin{enumerate}[label=\alph*)]
\item determinante;
\item traza;
\item rango;
\item elementos de la diagonal;
\item valores propios;
\item polinomio característico.
\end{enumerate}

\item Diagonalizar, si es posible,
\[
A=\begin{pmatrix}
2&1\\
1&2
\end{pmatrix}
\]
y utilizar la diagonalización para obtener una expresión cerrada de $A^n$.

\item Sea
\[
A=\begin{pmatrix}
3&1\\
0&2
\end{pmatrix}.
\]
Calcular $A^{10}$ utilizando diagonalización. Comparar el número de operaciones necesarias con el que exigiría multiplicar la matriz diez veces.

\item Una evolución discreta viene dada por
\[
x_{k+1}=Ax_k,
\qquad
A=\begin{pmatrix}
0.7&0.2\\
0.3&0.8
\end{pmatrix}.
\]
\begin{enumerate}[label=\alph*)]
\item Calcular valores y vectores propios de $A$.
\item Diagonalizar $A$.
\item Obtener una expresión para $x_k$ a partir de $x_0$.
\item Interpretar qué componente del estado permanece y cuál se atenúa al aumentar $k$.
\end{enumerate}

\item En un modelo térmico simplificado, las temperaturas de dos zonas satisfacen
\[
T_{k+1}=
\begin{pmatrix}
0.8&0.2\\
0.2&0.8
\end{pmatrix}T_k.
\]
Si
\[
T_0=\begin{pmatrix}80\\20\end{pmatrix},
\]
utilizar la diagonalización para obtener $T_k$ y determinar el estado límite. Interpretar físicamente los dos vectores propios.

\item Una matriz de rigidez simplificada de un sistema mecánico es
\[
K=\begin{pmatrix}
2&-1\\
-1&2
\end{pmatrix}.
\]
Obtener una base ortonormal de vectores propios y diagonalizar ortogonalmente $K$. Interpretar las direcciones propias como modos independientes de deformación.

\item Una red lineal de dos variables eléctricas está descrita por la matriz simétrica
\[
G=\begin{pmatrix}
5&-2\\
-2&5
\end{pmatrix}.
\]
Calcular sus direcciones propias y explicar qué ventaja aporta expresar un vector de estado en esa base cuando se aplica repetidamente la transformación $G$.

\item Determinar la matriz simétrica asociada a cada forma cuadrática:
\begin{enumerate}[label=\alph*)]
\item $Q(x,y)=3x^2+4xy+2y^2$;
\item $Q(x,y,z)=x^2+2y^2+4z^2-6xy+2xz$;
\item $Q(x_1,x_2,x_3)=2x_1x_2-4x_1x_3+6x_2x_3$.
\end{enumerate}

\item Para
\[
A=\begin{pmatrix}
2&1\\
1&3
\end{pmatrix},
\qquad
P=\begin{pmatrix}
1&1\\
0&1
\end{pmatrix},
\]
calcular $B=P^tAP$. Escribir la forma cuadrática $Q(x)=x^tAx$ en las nuevas variables determinadas por $x=Py$ y comprobar el resultado mediante sustitución directa.

\item Diagonalizar ortogonalmente la forma cuadrática
\[
Q(x,y)=5x^2-4xy+2y^2
\]
y escribir explícitamente el cambio de coordenadas que elimina el término cruzado.

\item Reducir por completación de cuadrados, sin calcular valores propios, la forma cuadrática
\[
Q(x,y,z)=2x^2+4xy+2xz+5y^2+2yz+2z^2.
\]
Escribir explícitamente el cambio lineal de variables empleado y la matriz $P$ que verifica $P^tAP=D$.

\item Diagonalizar ortogonalmente
\[
Q(x,y)=4x^2+4xy+y^2.
\]
Clasificar la forma y representar verbalmente qué indica la existencia o no de direcciones no nulas en las que $Q$ vale cero.

\item Clasificar mediante los valores propios las formas cuadráticas asociadas a
\[
A_1=\begin{pmatrix}2&1\\1&2\end{pmatrix},
\qquad
A_2=\begin{pmatrix}1&2\\2&1\end{pmatrix},
\]
\[
A_3=\begin{pmatrix}1&-1\\-1&1\end{pmatrix},
\qquad
A_4=\begin{pmatrix}-2&0\\0&-5\end{pmatrix}.
\]

\item Utilizar el criterio de Sylvester para clasificar
\[
Q(x,y,z)=3x^2+2y^2+4z^2+2xy-2yz.
\]
Comprobar después la clasificación mediante los valores propios de la matriz asociada.

\item La energía potencial elástica de un sistema idealizado viene dada por
\[
E(x_1,x_2)=\frac12\left(4x_1^2-4x_1x_2+3x_2^2\right).
\]
\begin{enumerate}[label=\alph*)]
\item Obtener la matriz simétrica $K$ tal que $E(x)=\frac12x^tKx$.
\item Determinar si la energía es positiva para todo desplazamiento no nulo.
\item Hallar unas coordenadas ortogonales en las que la energía no contenga términos cruzados.
\item Interpretar las nuevas direcciones como deformaciones desacopladas.
\end{enumerate}

\item Una aproximación cuadrática a una función de pérdidas de un dispositivo es
\[
L(x,y)=6x^2+4xy+3y^2.
\]
Determinar sus ejes principales mediante diagonalización ortogonal y decidir si $L$ puede tomar valores negativos. Explicar qué significa esta conclusión para una magnitud interpretada como pérdida.

\item Para la forma
\[
Q(x,y,z)=x^2+y^2-z^2+2xy,
\]
obtener la matriz asociada, calcular sus valores propios y clasificarla. Encontrar además dos vectores no nulos $u$ y $v$ tales que $Q(u)>0$ y $Q(v)<0$.

\item Razonar si las siguientes afirmaciones son verdaderas o falsas:
\begin{enumerate}[label=\alph*)]
\item Toda matriz con valores propios reales es diagonalizable.
\item Toda matriz con valores propios distintos es diagonalizable.
\item Toda matriz simétrica real es diagonalizable.
\item Si una matriz es diagonalizable, todos sus valores propios son distintos.
\item Dos matrices semejantes tienen los mismos valores propios.
\item Dos matrices congruentes tienen necesariamente los mismos valores propios.
\item Una diagonalización ortogonal de una matriz simétrica es a la vez una semejanza y una congruencia.
\end{enumerate}

\item Explicar con palabras propias la diferencia conceptual entre las expresiones
\[
P^{-1}AP
\qquad\text{y}\qquad
P^tAP.
\]
Indicar qué objeto matemático se está representando en cada caso y por qué confundir ambos cambios puede conducir a errores.

\item \M Para los ejercicios 1--8, utilizar \texttt{Eigenvalues} y \texttt{Eigenvectors} para comprobar los valores y vectores propios obtenidos manualmente. En los casos con parámetros, estudiar primero a mano qué valores requieren una discusión separada y utilizar después el ordenador únicamente como comprobación.

\item \M Para los ejercicios 12--17, construir manualmente $P$ y $D$ y comprobar en Mathematica las igualdades
\[
A=PDP^{-1}
\qquad\text{o}\qquad
P^tAP=D,
\]
según corresponda. Utilizar únicamente productos matriciales, \texttt{Inverse} y \texttt{Transpose}, evitando introducir comandos adicionales que automaticen por completo la diagonalización.

\item \M En los ejercicios 18--26, comprobar la matriz simétrica asociada y sus valores propios. Comparar la clasificación obtenida a partir de los signos de los valores propios con la obtenida manualmente mediante el criterio de Sylvester cuando sea aplicable.

\item (Examen) Sea
\[
A=\begin{pmatrix}
2&1&0\\
0&2&0\\
0&0&3
\end{pmatrix}.
\]
\begin{enumerate}[label=\alph*)]
\item Calcular el polinomio característico y los valores propios.
\item Determinar las multiplicidades algebraicas y geométricas.
\item Decidir razonadamente si $A$ es diagonalizable.
\item Si no lo es, modificar únicamente una entrada no diagonal de $A$ para obtener una matriz diagonalizable con los mismos valores propios y justificar la respuesta.
\end{enumerate}

\item (Examen) Sea
\[
A=\begin{pmatrix}
2&1&1\\
1&2&1\\
1&1&2
\end{pmatrix}.
\]
\begin{enumerate}[label=\alph*)]
\item Obtener sus valores propios y una base ortonormal de vectores propios.
\item Construir una matriz ortogonal $P$ tal que $P^tAP=D$.
\item Calcular $A^{20}$ utilizando la diagonalización.
\item Clasificar la forma cuadrática $Q(x)=x^tAx$.
\end{enumerate}

\item (Examen) Considérese la forma cuadrática
\[
Q(x,y,z)=2x^2+3y^2+2z^2+2xy-2yz.
\]
\begin{enumerate}[label=\alph*)]
\item Obtener su matriz simétrica asociada.
\item Clasificar la forma cuadrática.
\item Obtener un cambio ortogonal de variables que la transforme en una forma diagonal.
\item Indicar qué información sobre la forma permanece invariante si, en lugar de un cambio ortogonal, se utiliza cualquier cambio regular de variables.
\end{enumerate}

\end{enumerate}

\cleardoublepage
\chapter{Geometría}
\label{ch:geometria}

La geometría analítica permite traducir problemas geométricos a ecuaciones y, recíprocamente, interpretar geométricamente sistemas de ecuaciones. En ingeniería esta doble lectura es especialmente útil: la posición de un eje, la orientación de una superficie, la trayectoria ideal de una herramienta, la dirección de una fuerza o la separación mínima entre dos elementos pueden describirse mediante puntos, vectores, rectas y planos.

En el capítulo anterior hemos trabajado con espacios vectoriales, producto escalar y ortogonalidad. Ahora utilizaremos esas herramientas en $\R^2$ y $\R^3$ para estudiar la \emph{geometría afín} ---incidencia, paralelismo e intersecciones--- y la \emph{geometría euclídea} ---distancias, ángulos y perpendicularidad---. La intención es pasar con soltura entre diferentes registros: una descripción verbal de una situación, su representación vectorial y las ecuaciones que permiten resolverla.

\section{Puntos y vectores en el plano y el espacio}
\label{sec:geometria-vectores}

\subsection{Coordenadas, vectores y distancia entre puntos}

Fijado un sistema de referencia cartesiano, un punto del plano se representa mediante dos coordenadas,
\[
P=(p_1,p_2)\in\R^2,
\]
y un punto del espacio mediante tres,
\[
P=(p_1,p_2,p_3)\in\R^3.
\]
Conviene distinguir conceptualmente \emph{puntos} y \emph{vectores}. Un punto representa una posición; un vector representa un desplazamiento, una dirección o, en muchas aplicaciones, una magnitud física con módulo y dirección.

Dados dos puntos $P$ y $Q$, el vector que lleva de $P$ a $Q$ es
\[
\overrightarrow{PQ}=Q-P.
\]
En $\R^3$, si
\[
P=(p_1,p_2,p_3),\qquad Q=(q_1,q_2,q_3),
\]
entonces
\[
\overrightarrow{PQ}
=(q_1-p_1,q_2-p_2,q_3-p_3).
\]

\begin{definition}[Norma de un vector]
Para $v=(v_1,\ldots,v_n)\in\R^n$, su \emph{norma euclídea} es
\[
\|v\|=\sqrt{v_1^2+\cdots+v_n^2}.
\]
\end{definition}

La norma mide la longitud del vector. Si $v\neq0$, el vector
\[
\widehat v=\frac{v}{\|v\|}
\]
tiene norma $1$ y conserva la dirección y el sentido de $v$. Lo denominamos \emph{vector unitario} asociado a $v$.

\begin{definition}[Distancia entre dos puntos]
La distancia entre dos puntos $P,Q\in\R^n$ es
\[
d(P,Q)=\|Q-P\|.
\]
\end{definition}

Así, en el espacio,
\[
d(P,Q)
=\sqrt{(q_1-p_1)^2+(q_2-p_2)^2+(q_3-p_3)^2}.
\]

\begin{example}
El extremo de un actuador se desplaza desde
\[
P=(1,-2,0.5)\ \mathrm{m}
\]
hasta
\[
Q=(4,2,2.5)\ \mathrm{m}.
\]
El desplazamiento es
\[
\overrightarrow{PQ}=(3,4,2)\ \mathrm{m},
\]
y la distancia recorrida en línea recta entre ambas posiciones es
\[
\|\overrightarrow{PQ}\|
=\sqrt{3^2+4^2+2^2}
=\sqrt{29}\ \mathrm{m}.
\]
La dirección unitaria del desplazamiento es
\[
\widehat{PQ}=\frac{1}{\sqrt{29}}(3,4,2).
\]
Si el movimiento se realizara exactamente siguiendo esa dirección, cualquier posición de la trayectoria podría escribirse como
\[
P+t\widehat{PQ},\qquad t\geq0.
\]
\end{example}

La expresión anterior anticipa la ecuación vectorial de una recta: una posición inicial más un múltiplo de una dirección.

\subsection{Producto escalar, ángulos y proyecciones}

Recordemos que, dados $u,v\in\R^n$, su producto escalar es
\[
u\cdot v=u_1v_1+\cdots+u_nv_n.
\]
Además,
\[
u\cdot v=\|u\|\,\|v\|\cos\theta,
\]
donde $\theta\in[0,\pi]$ es el ángulo que forman los vectores no nulos $u$ y $v$. Por tanto,
\[
\cos\theta=\frac{u\cdot v}{\|u\|\,\|v\|}.
\]

\begin{proposition}[Criterio de ortogonalidad]
Dos vectores $u$ y $v$ son ortogonales si y sólo si
\[
u\cdot v=0.
\]
\end{proposition}

Esta igualdad permite traducir perpendicularidad geométrica en una ecuación algebraica.

\begin{definition}[Proyección ortogonal sobre una dirección]
Sean $u,v\in\R^n$, con $v\neq0$. La proyección ortogonal de $u$ sobre la dirección de $v$ es
\[
\operatorname{proj}_v(u)
=\frac{u\cdot v}{v\cdot v}\,v.
\]
\end{definition}

El vector $u$ puede descomponerse como
\[
u=u_{\parallel}+u_{\perp},
\]
donde
\[
u_{\parallel}=\operatorname{proj}_v(u),
\qquad
u_{\perp}=u-\operatorname{proj}_v(u),
\]
y $u_{\perp}\cdot v=0$.

\begin{solvedexercise}
Una fuerza
\[
F=(120,40,-20)\ \mathrm{N}
\]
actúa sobre un mecanismo cuyo desplazamiento permitido tiene la dirección
\[
d=(2,1,2).
\]
Calcular la componente vectorial de $F$ que actúa en la dirección del movimiento y la componente perpendicular.
\tcblower
\textbf{Resolución.}\par\smallskip
La proyección de $F$ sobre $d$ es
\[
F_{\parallel}
=\frac{F\cdot d}{d\cdot d}\,d.
\]
Como
\[
F\cdot d=120\cdot2+40\cdot1-20\cdot2=240,
\]
y
\[
d\cdot d=2^2+1^2+2^2=9,
\]
obtenemos
\[
F_{\parallel}=\frac{240}{9}(2,1,2)
=\left(\frac{160}{3},\frac{80}{3},\frac{160}{3}\right)\ \mathrm{N}.
\]
Por tanto,
\[
F_{\perp}=F-F_{\parallel}
=\left(\frac{200}{3},\frac{40}{3},-\frac{220}{3}\right)\ \mathrm{N}.
\]
Puede comprobarse que
\[
F_{\perp}\cdot d=0.
\]
La descomposición separa la parte de la fuerza que contribuye al movimiento permitido de la parte que actúa transversalmente sobre el mecanismo.
\end{solvedexercise}

En $\R^3$ también es útil describir una dirección mediante los ángulos que forma con los ejes coordenados. Si $v=(v_1,v_2,v_3)\neq0$ y $\alpha,\beta,\gamma$ son los ángulos con los ejes $OX$, $OY$ y $OZ$, respectivamente, entonces
\[
\cos\alpha=\frac{v_1}{\|v\|},
\qquad
\cos\beta=\frac{v_2}{\|v\|},
\qquad
\cos\gamma=\frac{v_3}{\|v\|},
\]
y, en consecuencia,
\[
\cos^2\alpha+\cos^2\beta+\cos^2\gamma=1.
\]

\subsection{Producto vectorial y producto mixto}

En $\R^3$ disponemos de una operación adicional que produce un vector perpendicular a otros dos.

\begin{definition}[Producto vectorial]
Sean
\[
u=(u_1,u_2,u_3),\qquad v=(v_1,v_2,v_3).
\]
Su \emph{producto vectorial} es
\[
u\times v=
\begin{vmatrix}
\mathbf i&\mathbf j&\mathbf k\\
u_1&u_2&u_3\\
v_1&v_2&v_3
\end{vmatrix},
\]
es decir,
\[
u\times v
=
(u_2v_3-u_3v_2,\,
 u_3v_1-u_1v_3,\,
 u_1v_2-u_2v_1).
\]
\end{definition}

El producto vectorial satisface:
\begin{itemize}
\item $u\times v$ es ortogonal a $u$ y a $v$;
\item $v\times u=-(u\times v)$;
\item $u\times v=0$ si y sólo si $u$ y $v$ son linealmente dependientes;
\item
\[
\|u\times v\|=\|u\|\,\|v\|\sin\theta.
\]
\end{itemize}

La última propiedad tiene una interpretación geométrica inmediata:
\[
\text{Área del paralelogramo generado por }u,v
=\|u\times v\|,
\]
y, por tanto,
\[
\text{Área del triángulo}
=\frac12\|u\times v\|.
\]

\begin{example}
Tres puntos de una placa son
\[
A=(1,0,0),\qquad B=(3,1,0),\qquad C=(2,0,2).
\]
Los vectores
\[
\overrightarrow{AB}=(2,1,0),
\qquad
\overrightarrow{AC}=(1,0,2)
\]
determinan la placa. Un vector perpendicular a ella es
\[
\overrightarrow{AB}\times\overrightarrow{AC}
=
\begin{vmatrix}
\mathbf i&\mathbf j&\mathbf k\\
2&1&0\\
1&0&2
\end{vmatrix}
=(2,-4,-1).
\]
Este vector normal será precisamente el que utilicemos más adelante para obtener la ecuación cartesiana del plano que contiene la placa.
\end{example}

\begin{definition}[Producto mixto]
Dados $u,v,w\in\R^3$, se denomina \emph{producto mixto} al número
\[
[u,v,w]=u\cdot(v\times w).
\]
\end{definition}

Se verifica
\[
[u,v,w]
=
\det
\begin{pmatrix}
u_1&u_2&u_3\\
v_1&v_2&v_3\\
w_1&w_2&w_3
\end{pmatrix}.
\]
Su valor absoluto es el volumen del paralelepípedo generado por $u,v,w$:
\[
V=|u\cdot(v\times w)|.
\]
Además,
\[
u,v,w\text{ son coplanarios}
\quad\Longleftrightarrow\quad
u\cdot(v\times w)=0.
\]

\begin{solvedexercise}
Una fuerza $F=(0,150,60)$ N actúa en el punto cuyo vector de posición respecto de un eje de referencia es
\[
r=(0.4,0,0.2)\ \mathrm{m}.
\]
Calcular el momento de la fuerza respecto del origen,
\[
M=r\times F.
\]
\tcblower
\textbf{Resolución.}\par\smallskip
Calculamos
\[
M=
\begin{vmatrix}
\mathbf i&\mathbf j&\mathbf k\\
0.4&0&0.2\\
0&150&60
\end{vmatrix}
=(-30,-24,60)\ \mathrm{N\,m}.
\]
El vector momento es perpendicular al plano generado por $r$ y $F$. Su módulo es
\[
\|M\|=\sqrt{(-30)^2+(-24)^2+60^2}
=\sqrt{5076}\ \mathrm{N\,m}.
\]
\end{solvedexercise}

\section{Rectas y planos}
\label{sec:rectas-planos}

\subsection{Rectas en el plano y en el espacio}

Una recta queda determinada por un punto y una dirección.

\begin{definition}[Ecuación vectorial de una recta]
Sea $A\in\R^n$ y sea $v\in\R^n$, $v\neq0$. La recta que pasa por $A$ y tiene vector director $v$ es
\[
r:\quad X=A+tv,\qquad t\in\R.
\]
\end{definition}

En el plano, si
\[
A=(a_1,a_2),\qquad v=(v_1,v_2),
\]
las ecuaciones paramétricas son
\[
\begin{cases}
x=a_1+tv_1,\\
y=a_2+tv_2,
\end{cases}
\qquad t\in\R.
\]
En el espacio,
\[
\begin{cases}
x=a_1+tv_1,\\
y=a_2+tv_2,\\
z=a_3+tv_3.
\end{cases}
\]

Cuando las componentes del vector director son no nulas, puede eliminarse el parámetro y escribirse la forma continua o simétrica
\[
\frac{x-a_1}{v_1}
=
\frac{y-a_2}{v_2}
\]
en el plano, y
\[
\frac{x-a_1}{v_1}
=
\frac{y-a_2}{v_2}
=
\frac{z-a_3}{v_3}
\]
en el espacio. Si alguna componente del vector director es cero, la coordenada correspondiente permanece constante y no debe dividirse por cero.

En $\R^2$ una recta también puede expresarse mediante una ecuación cartesiana
\[
r:\quad Ax+By+C=0,
\qquad (A,B)\neq(0,0).
\]
El vector
\[
n=(A,B)
\]
es normal a la recta. Por tanto, un vector director puede elegirse como
\[
v=(-B,A).
\]
Si $B\neq0$, la ecuación puede escribirse
\[
y=-\frac ABx-\frac CB.
\]
Así, la pendiente es $m=-A/B$.

\begin{example}
Hallar una ecuación de la recta del plano que pasa por
\[
P=(1,-2),\qquad Q=(5,1).
\]
Un vector director es
\[
v=Q-P=(4,3).
\]
Por tanto,
\[
(x,y)=(1,-2)+t(4,3),
\]
o, en forma paramétrica,
\[
\begin{cases}
x=1+4t,\\
y=-2+3t.
\end{cases}
\]
Un vector normal es $n=(3,-4)$, ya que
\[
(3,-4)\cdot(4,3)=0.
\]
La ecuación cartesiana es entonces
\[
3(x-1)-4(y+2)=0,
\]
es decir,
\[
\boxed{3x-4y-11=0.}
\]
\end{example}

Una recta del espacio también puede describirse como intersección de dos planos independientes:
\[
r:\quad
\begin{cases}
A_1x+B_1y+C_1z+D_1=0,\\
A_2x+B_2y+C_2z+D_2=0.
\end{cases}
\]
Si $n_1=(A_1,B_1,C_1)$ y $n_2=(A_2,B_2,C_2)$ no son paralelos, un vector director de la recta es
\[
v=n_1\times n_2.
\]

\begin{solvedexercise}
El eje de una tubería recta pasa por los puntos
\[
A=(1,2,0),\qquad B=(4,-1,6).
\]
Obtener sus ecuaciones vectorial, paramétricas y, cuando sea posible, continuas.
\tcblower
\textbf{Resolución.}\par\smallskip
Un vector director es
\[
v=B-A=(3,-3,6)=3(1,-1,2).
\]
Podemos utilizar la dirección simplificada $(1,-1,2)$. La ecuación vectorial es
\[
\boxed{(x,y,z)=(1,2,0)+t(1,-1,2).}
\]
Las ecuaciones paramétricas son
\[
\boxed{
\begin{cases}
x=1+t,\\
y=2-t,\\
z=2t.
\end{cases}}
\]
Como ninguna componente del vector director es nula,
\[
\boxed{x-1=-(y-2)=\frac z2.}
\]
Las tres expresiones describen exactamente el mismo conjunto de puntos.
\end{solvedexercise}

\subsection{Planos en el espacio}

Un plano puede determinarse mediante un punto y dos direcciones linealmente independientes.

\begin{definition}[Ecuación vectorial de un plano]
Sean $A\in\R^3$ y $u,v\in\R^3$ linealmente independientes. El plano que pasa por $A$ y tiene direcciones $u$ y $v$ es
\[
\pi:\quad X=A+su+tv,
\qquad s,t\in\R.
\]
\end{definition}

Si
\[
A=(a_1,a_2,a_3),\quad
u=(u_1,u_2,u_3),\quad
v=(v_1,v_2,v_3),
\]
obtenemos las ecuaciones paramétricas
\[
\begin{cases}
x=a_1+su_1+tv_1,\\
y=a_2+su_2+tv_2,\\
z=a_3+su_3+tv_3.
\end{cases}
\]

Como $u\times v$ es perpendicular a ambas direcciones, puede utilizarse como vector normal del plano.

\begin{definition}[Ecuación cartesiana de un plano]
Un plano de $\R^3$ puede escribirse
\[
\pi:\quad Ax+By+Cz+D=0,
\qquad (A,B,C)\neq(0,0,0).
\]
El vector
\[
n=(A,B,C)
\]
es un vector normal al plano.
\end{definition}

Si conocemos un punto $A=(a_1,a_2,a_3)$ y un vector normal $n=(n_1,n_2,n_3)$, la ecuación puede obtenerse directamente mediante
\[
n\cdot(X-A)=0,
\]
es decir,
\[
n_1(x-a_1)+n_2(y-a_2)+n_3(z-a_3)=0.
\]

Tres puntos no alineados $A,B,C$ determinan un plano. En este caso podemos tomar
\[
u=\overrightarrow{AB},
\qquad
v=\overrightarrow{AC},
\]
y calcular
\[
n=u\times v.
\]

\begin{solvedexercise}
Una superficie plana de captación está determinada, en un sistema de coordenadas local, por los puntos
\[
A=(0,0,2),\qquad B=(4,0,3),\qquad C=(0,3,4).
\]
Obtener una ecuación cartesiana del plano y un vector normal unitario.
\tcblower
\textbf{Resolución.}\par\smallskip
Tomamos
\[
u=B-A=(4,0,1),
\qquad
v=C-A=(0,3,2).
\]
Un vector normal es
\[
n=u\times v
=
\begin{vmatrix}
\mathbf i&\mathbf j&\mathbf k\\
4&0&1\\
0&3&2
\end{vmatrix}
=(-3,-8,12).
\]
Por tanto,
\[
-3(x-0)-8(y-0)+12(z-2)=0,
\]
y una ecuación cartesiana es
\[
\boxed{-3x-8y+12z-24=0.}
\]
Como
\[
\|n\|=\sqrt{9+64+144}=\sqrt{217},
\]
un vector normal unitario es
\[
\boxed{\widehat n=\frac1{\sqrt{217}}(-3,-8,12).}
\]
La orientación de una superficie puede describirse mediante su vector normal; por eso este tipo de cálculo aparece al estudiar superficies de montaje, paneles, placas y caras de elementos estructurales.
\end{solvedexercise}

\subsection{Posiciones relativas e intersecciones}

Las posiciones relativas se estudian combinando información sobre direcciones con la compatibilidad de sistemas de ecuaciones. Este punto conecta directamente la geometría con los capítulos de matrices y sistemas lineales.

\paragraph{Dos rectas en el plano.}
Sean
\[
r:A_1x+B_1y+C_1=0,
\qquad
s:A_2x+B_2y+C_2=0.
\]
El sistema formado por ambas ecuaciones puede ser:
\begin{itemize}
\item compatible determinado: las rectas son secantes;
\item compatible indeterminado: las rectas son coincidentes;
\item incompatible: las rectas son paralelas distintas.
\end{itemize}

Equivalentemente, si $v_r$ y $v_s$ son vectores directores, entonces las rectas son paralelas cuando $v_r$ y $v_s$ son proporcionales.

\paragraph{Recta y plano en el espacio.}
Sean
\[
r:X=A+tv,
\qquad
\pi:n\cdot(X-P)=0.
\]
La cantidad $n\cdot v$ determina la situación:
\begin{itemize}
\item si $n\cdot v\neq0$, la recta corta al plano en un único punto;
\item si $n\cdot v=0$ y $A\in\pi$, la recta está contenida en el plano;
\item si $n\cdot v=0$ y $A\notin\pi$, la recta es paralela al plano y no lo corta.
\end{itemize}

\begin{example}
Sea
\[
r:(x,y,z)=(1,0,2)+t(2,1,-1)
\]
y
\[
\pi: x-2y+z-5=0.
\]
El vector normal del plano es
\[
n=(1,-2,1).
\]
Como
\[
n\cdot v=(1,-2,1)\cdot(2,1,-1)=2-2-1=-1\neq0,
\]
la recta es secante al plano. Sustituyendo las ecuaciones paramétricas
\[
x=1+2t,\quad y=t,\quad z=2-t
\]
en la ecuación del plano,
\[
(1+2t)-2t+(2-t)-5=0,
\]
de donde
\[
t=-2.
\]
El punto de intersección es
\[
\boxed{(-3,-2,4).}
\]
\end{example}

\paragraph{Dos planos.}
Sean
\[
\pi_1:n_1\cdot X+d_1=0,
\qquad
\pi_2:n_2\cdot X+d_2=0.
\]
Si $n_1$ y $n_2$ no son proporcionales, los planos se cortan en una recta cuyo vector director puede tomarse como
\[
v=n_1\times n_2.
\]
Si los vectores normales son proporcionales, los planos son paralelos o coincidentes; para distinguir ambos casos basta comprobar si un punto de uno de ellos pertenece al otro.

\paragraph{Dos rectas en el espacio.}
Sean
\[
r:A+tu,
\qquad
s:B+\lambda v.
\]
Si $u$ y $v$ son proporcionales, las rectas son paralelas o coincidentes. Si no son proporcionales, resolvemos
\[
A+tu=B+\lambda v.
\]
Si el sistema tiene solución, las rectas son secantes. Si no tiene solución, son \emph{rectas cruzadas}: no son paralelas, pero tampoco se cortan.

Una caracterización útil de la coplanaridad es
\[
(B-A)\cdot(u\times v)=0.
\]
Cuando $u$ y $v$ no son paralelos, las rectas son secantes si y sólo si son coplanarias; si el producto mixto anterior es no nulo, son cruzadas.

\begin{solvedexercise}
Estudiar la posición relativa de
\[
r:(x,y,z)=(1,0,1)+t(1,2,-1)
\]
y
\[
s:(x,y,z)=(0,1,2)+\lambda(2,1,1).
\]
\tcblower
\textbf{Resolución.}\par\smallskip
Los vectores directores
\[
u=(1,2,-1),\qquad v=(2,1,1)
\]
no son proporcionales. Buscamos una posible intersección resolviendo
\[
(1,0,1)+t(1,2,-1)
=(0,1,2)+\lambda(2,1,1).
\]
Esto equivale a
\[
\begin{cases}
1+t=2\lambda,\\
2t=1+\lambda,\\
1-t=2+\lambda.
\end{cases}
\]
De la segunda ecuación,
\[
\lambda=2t-1.
\]
Sustituyendo en la primera,
\[
1+t=4t-2
\quad\Longrightarrow\quad
 t=1,
\]
y entonces $\lambda=1$. Sin embargo, la tercera ecuación daría
\[
0=3,
\]
lo cual es imposible. Por tanto, las rectas no se cortan.

Como tampoco son paralelas,
\[
\boxed{r\text{ y }s\text{ son rectas cruzadas}.}
\]
\end{solvedexercise}

\section{Problemas métricos}
\label{sec:problemas-metricos}

\subsection{Distancias}

La distancia entre objetos geométricos se define como la menor distancia entre sus puntos. Las fórmulas que utilizaremos proceden de proyecciones ortogonales y productos vectoriales.

\paragraph{Distancia de un punto a una recta del plano.}
Sea
\[
r:Ax+By+C=0
\]
y sea $P=(p_1,p_2)$. Entonces
\[
\boxed{
 d(P,r)=
 \frac{|Ap_1+Bp_2+C|}{\sqrt{A^2+B^2}}.}
\]

Esta expresión puede interpretarse proyectando el vector que une $P$ con un punto cualquiera de la recta sobre el vector normal $n=(A,B)$.

\begin{example}
Calcular la distancia del punto $P=(3,-1)$ a la recta
\[
r:2x-y-4=0.
\]
Aplicando la fórmula,
\[
d(P,r)
=\frac{|2\cdot3-(-1)-4|}{\sqrt{2^2+(-1)^2}}
=\frac{3}{\sqrt5}.
\]
Por tanto,
\[
\boxed{d(P,r)=\frac{3}{\sqrt5}.}
\]
\end{example}

\paragraph{Distancia de un punto a una recta del espacio.}
Sea
\[
r:X=A+tv
\]
y sea $P\in\R^3$. El área del paralelogramo generado por $P-A$ y $v$ es
\[
\|(P-A)\times v\|.
\]
También es igual a base por altura,
\[
\|v\|\,d(P,r).
\]
Por tanto,
\[
\boxed{
 d(P,r)=\frac{\|(P-A)\times v\|}{\|v\|}.}
\]

\begin{solvedexercise}
Calcular la distancia del punto
\[
P=(1,-2,3)
\]
a la recta
\[
r:(x,y,z)=(3,-1,0)+t(-1,1,1).
\]
\tcblower
\textbf{Resolución.}\par\smallskip
Tomamos
\[
A=(3,-1,0),
\qquad
v=(-1,1,1).
\]
Entonces
\[
P-A=(-2,-1,3)
\]
y
\[
(P-A)\times v
=
\begin{vmatrix}
\mathbf i&\mathbf j&\mathbf k\\
-2&-1&3\\
-1&1&1
\end{vmatrix}
=(-4,-1,-3).
\]
Por tanto,
\[
d(P,r)
=\frac{\sqrt{16+1+9}}{\sqrt{1+1+1}}
=\boxed{\sqrt{\frac{26}{3}}}.
\]
\end{solvedexercise}

\paragraph{Distancia de un punto a un plano.}
Sea
\[
\pi:Ax+By+Cz+D=0
\]
y $P=(p_1,p_2,p_3)$. Entonces
\[
\boxed{
 d(P,\pi)
 =\frac{|Ap_1+Bp_2+Cp_3+D|}{\sqrt{A^2+B^2+C^2}}.}
\]

De nuevo, el numerador mide ---salvo el factor $\|n\|$--- la componente del desplazamiento en la dirección normal al plano.

\begin{example}
Una carcasa metálica ocupa el plano
\[
\pi:2x-y+2z-6=0.
\]
Un sensor está situado en
\[
P=(3,1,4).
\]
La separación perpendicular entre el sensor y la superficie es
\[
d(P,\pi)
=\frac{|2\cdot3-1+2\cdot4-6|}{\sqrt{2^2+(-1)^2+2^2}}
=\frac{7}{3}.
\]
La distancia mínima es, por tanto,
\[
\boxed{\frac73}
\]
en las unidades del sistema de coordenadas utilizado.
\end{example}

\paragraph{Distancias entre objetos paralelos.}
Si dos rectas del plano son paralelas,
\[
r:Ax+By+C_1=0,
\qquad
s:Ax+By+C_2=0,
\]
entonces
\[
\boxed{
 d(r,s)=\frac{|C_2-C_1|}{\sqrt{A^2+B^2}}.}
\]
Análogamente, para dos planos paralelos
\[
\pi_1:Ax+By+Cz+D_1=0,
\qquad
\pi_2:Ax+By+Cz+D_2=0,
\]
se tiene
\[
\boxed{
 d(\pi_1,\pi_2)
 =\frac{|D_2-D_1|}{\sqrt{A^2+B^2+C^2}}.}
\]

\paragraph{Distancia entre dos rectas cruzadas.}
Sean
\[
r:A+tu,
\qquad
s:B+\lambda v,
\]
con $u\times v\neq0$. Si las rectas son cruzadas, el vector $u\times v$ es perpendicular a ambas direcciones. La distancia es la proyección de $B-A$ sobre esa dirección común:
\[
\boxed{
 d(r,s)
 =\frac{|(B-A)\cdot(u\times v)|}{\|u\times v\|}.}
\]
Si las rectas son secantes, el numerador vale cero y la fórmula proporciona, correctamente, distancia nula.

\subsection{Ángulos entre rectas y planos}

La medida de ángulos se reduce al producto escalar entre vectores directores o vectores normales.

\paragraph{Ángulo entre dos rectas.}
Si $u$ y $v$ son vectores directores de dos rectas, el ángulo geométrico entre ellas se toma habitualmente como el menor de los dos posibles, por lo que
\[
\boxed{
\cos\theta
=\frac{|u\cdot v|}{\|u\|\,\|v\|},
\qquad
0\leq\theta\leq\frac\pi2.}
\]

\paragraph{Ángulo entre dos planos.}
Si $n_1$ y $n_2$ son vectores normales de los planos, el ángulo entre los planos es el ángulo agudo entre sus normales:
\[
\boxed{
\cos\theta
=\frac{|n_1\cdot n_2|}{\|n_1\|\,\|n_2\|}.}
\]

\paragraph{Ángulo entre una recta y un plano.}
Sea $v$ un vector director de la recta y $n$ un vector normal al plano. Si $\alpha$ es el ángulo entre la recta y el plano, entonces
\[
\alpha+\theta=\frac\pi2,
\]
donde $\theta$ es el ángulo entre $v$ y $n$. Por tanto,
\[
\boxed{
\sin\alpha
=\frac{|v\cdot n|}{\|v\|\,\|n\|}.}
\]

\begin{solvedexercise}
Una conducción rectilínea tiene dirección
\[
v=(2,-1,2)
\]
y atraviesa una pared cuyo plano es
\[
\pi:x+2y+2z-8=0.
\]
Calcular el ángulo que forma la conducción con la pared.
\tcblower
\textbf{Resolución.}\par\smallskip
Un vector normal al plano es
\[
n=(1,2,2).
\]
Entonces
\[
v\cdot n=2-2+4=4,
\]
\[
\|v\|=3,
\qquad
\|n\|=3.
\]
Por tanto,
\[
\sin\alpha=\frac49,
\]
y
\[
\boxed{\alpha=\arcsin\left(\frac49\right).}
\]
Esta es la inclinación de la conducción respecto de la superficie, no respecto de la normal a la superficie.
\end{solvedexercise}

Las condiciones de paralelismo y perpendicularidad se obtienen como casos particulares. Por ejemplo:
\begin{itemize}
\item dos rectas son perpendiculares si sus vectores directores tienen producto escalar cero;
\item dos planos son perpendiculares si sus normales tienen producto escalar cero;
\item una recta es perpendicular a un plano si su vector director es paralelo al vector normal del plano;
\item una recta es paralela a un plano si su vector director es ortogonal al normal del plano.
\end{itemize}

\subsection{Proyecciones y puntos de mínima distancia}

Las fórmulas de distancia proporcionan un número, pero en muchos problemas interesa además localizar el punto que realiza esa distancia mínima.

\paragraph{Proyección de un punto sobre una recta.}
Sea
\[
r:X=A+tv
\]
y sea $P$ un punto. El pie de la perpendicular $H$ es un punto de la recta,
\[
H=A+t_0v,
\]
y debe satisfacer
\[
(P-H)\cdot v=0.
\]
Sustituyendo,
\[
(P-A-t_0v)\cdot v=0,
\]
de donde
\[
\boxed{
 t_0=\frac{(P-A)\cdot v}{v\cdot v}.}
\]
Así,
\[
\boxed{
 H=A+\frac{(P-A)\cdot v}{v\cdot v}\,v.}
\]

\paragraph{Proyección de un punto sobre un plano.}
Sea
\[
\pi:n\cdot(X-A)=0
\]
y $P$ un punto. El punto proyectado $H$ debe encontrarse sobre la recta perpendicular al plano que pasa por $P$:
\[
H=P-tn.
\]
Imponiendo $H\in\pi$,
\[
n\cdot(P-tn-A)=0,
\]
por lo que
\[
t=\frac{n\cdot(P-A)}{n\cdot n}.
\]
En consecuencia,
\[
\boxed{
 H=P-\frac{n\cdot(P-A)}{n\cdot n}\,n.}
\]

\begin{solvedexercise}
Una herramienta se encuentra en
\[
P=(4,1,5)
\]
y debe aproximarse perpendicularmente a la superficie
\[
\pi:2x-y+2z-7=0.
\]
Determinar el punto $H$ de la superficie al que llega siguiendo la trayectoria más corta y la longitud de dicha trayectoria.
\tcblower
\textbf{Resolución.}\par\smallskip
Tomamos
\[
n=(2,-1,2).
\]
Podemos utilizar cualquier punto del plano; por ejemplo, $A=(0,-7,0)$ satisface
\[
2\cdot0-(-7)+2\cdot0-7=0.
\]
Entonces
\[
P-A=(4,8,5),
\]
y
\[
n\cdot(P-A)=8-8+10=10,
\qquad
n\cdot n=9.
\]
Por tanto,
\[
H=P-\frac{10}{9}(2,-1,2)
=\left(\frac{16}{9},\frac{19}{9},\frac{25}{9}\right).
\]
La longitud de la trayectoria es
\[
\|P-H\|
=\frac{10}{9}\|n\|
=\frac{10}{3}.
\]
Así,
\[
\boxed{
H=\left(\frac{16}{9},\frac{19}{9},\frac{25}{9}\right),
\qquad
d(P,\pi)=\frac{10}{3}.}
\]
Como comprobación, al sustituir $H$ en la ecuación del plano se obtiene cero.
\end{solvedexercise}

Los problemas métricos suelen requerir combinar varias de las ideas del capítulo. Antes de utilizar una fórmula conviene identificar con precisión qué objetos intervienen y qué información los determina: puntos, vectores directores y normales. Después se traduce el problema a productos escalares, productos vectoriales o sistemas lineales.

\begin{example}[Distancia entre dos ejes no coplanarios]
Dos ejes de una instalación se modelizan mediante
\[
r:(x,y,z)=(0,0,0)+t(1,0,1),
\]
\[
s:(x,y,z)=(1,2,0)+\lambda(0,1,1).
\]
Como
\[
u=(1,0,1),\qquad v=(0,1,1),
\]
no son paralelos. Además,
\[
u\times v=(-1,-1,1).
\]
Con $A=(0,0,0)$ y $B=(1,2,0)$,
\[
(B-A)\cdot(u\times v)
=(1,2,0)\cdot(-1,-1,1)=-3\neq0,
\]
por lo que las rectas son cruzadas. Su distancia es
\[
d(r,s)
=\frac{3}{\sqrt3}
=\boxed{\sqrt3}.
\]
El mismo producto mixto que detecta que las rectas no son coplanarias proporciona, tras normalizar, la separación mínima entre ellas.
\end{example}

\section*{Ejercicios propuestos}
\addcontentsline{toc}{section}{Ejercicios propuestos}

Los ejercicios siguientes combinan cálculo vectorial, obtención de ecuaciones, estudio de posiciones relativas y problemas métricos. Algunos proceden del tipo de ejercicios utilizados en cursos anteriores y otros amplían la colección para trabajar de forma integrada los contenidos del capítulo. Los marcados con \M son especialmente adecuados para ser resueltos o comprobados también con \textit{Wolfram Mathematica}.

\begin{enumerate}

\item Dados los puntos
\[
A=(1,-4),\qquad B=(5,7),
\]
calcular $\overrightarrow{AB}$, la distancia $d(A,B)$ y un vector unitario con la dirección de $A$ hacia $B$.

\item Dados
\[
P=(1,2,3),\qquad Q=(5,-3,1),
\]
calcular la distancia entre ambos puntos. Determinar además el punto medio del segmento $PQ$.

\item Para
\[
u=(2,-1,3),\qquad v=(1,4,-2),
\]
calcular:
\begin{enumerate}[label=\alph*)]
\item $u\cdot v$;
\item las normas de $u$ y $v$;
\item el ángulo que forman;
\item la proyección de $u$ sobre $v$;
\item una descomposición $u=u_{\parallel}+u_{\perp}$ respecto de la dirección de $v$.
\end{enumerate}

\item Determinar el valor de $a\in\R$ para que los vectores
\[
u=(1,a,2),\qquad v=(2,-1,1)
\]
sean ortogonales. Para ese valor, calcular el área del paralelogramo que generan.

\item Dados
\[
u=(1,2,-1),\qquad v=(2,0,3),
\]
calcular $u\times v$ y comprobar que el resultado es ortogonal a ambos vectores.

\item Calcular el área del triángulo de vértices
\[
A=(1,0,2),\qquad B=(3,-1,1),\qquad C=(0,2,4).
\]

\item Estudiar si los vectores
\[
u=(1,2,0),\qquad v=(0,1,1),\qquad w=(2,5,1)
\]
son coplanarios. En caso contrario, calcular el volumen del paralelepípedo que generan.

\item Un brazo mecánico aplica una fuerza
\[
F=(80,30,-40)\ \mathrm{N}
\]
en un punto cuyo vector de posición respecto del eje de giro es
\[
r=(0.25,0.10,0)\ \mathrm{m}.
\]
Calcular el momento $M=r\times F$ y su módulo.

\item Una corriente de fluido tiene velocidad
\[
v=(3,4,0)\ \mathrm{m/s}
\]
y una sección de control posee vector normal unitario
\[
n=\frac{1}{3}(2,-1,2).
\]
Calcular la componente de la velocidad normal a la sección y la componente tangencial.

\item Hallar las ecuaciones vectorial y paramétricas de la recta del plano que pasa por
\[
P=(2,-1),\qquad Q=(-1,5),
\]
y obtener después una ecuación cartesiana.

\item Hallar la ecuación de la recta que pasa por $P=(3,1)$ y es perpendicular a
\[
r:2x-3y+4=0.
\]

\item Estudiar la posición relativa de las rectas
\[
r:2x-y+1=0,
\qquad
s:4x-2y-5=0.
\]
Calcular la distancia entre ellas cuando proceda.

\item Hallar las ecuaciones vectorial, paramétricas y continuas de la recta del espacio que pasa por
\[
P=(1,0,-1),\qquad Q=(2,1,-3).
\]

\item Hallar la recta que pasa por $A=(1,-3,0)$ y es paralela al vector $u\times v$, siendo
\[
u=(-1,-1,2),\qquad v=(6,0,0).
\]

\item La trayectoria ideal del extremo de un robot pasa por el punto
\[
A=(0.5,1,0)
\]
y tiene dirección
\[
v=(2,-1,3).
\]
Determinar el punto alcanzado para $t=2$ y decidir si
\[
P=(4.5,-1,6)
\]
pertenece a la trayectoria.

\item Obtener una ecuación cartesiana del plano que pasa por
\[
P=(0,1,-1),\qquad Q=(2,3,-5),\qquad R=(1,4,3).
\]

\item Hallar la ecuación del plano cuyo vector normal es
\[
n=(-1,2,-3)
\]
y que pasa por
\[
A=(2,0,4).
\]

\item Hallar las ecuaciones paramétricas y cartesianas del plano determinado por
\[
A=(1,-3,2),\qquad u=(2,1,0),\qquad v=(-1,0,3).
\]

\item Determinar el valor de $a$ para que el punto
\[
P=(1,a,2)
\]
pertenezca al plano
\[
\pi:2x-y+3z-5=0.
\]
Para ese valor, obtener una recta perpendicular a $\pi$ que pase por $P$.

\item Una placa de intercambio térmico está contenida en el plano que pasa por los puntos
\[
A=(0,0,0),\qquad B=(2,1,0),\qquad C=(0,2,1).
\]
Obtener un vector normal unitario a la placa y su ecuación cartesiana.

\item Estudiar la posición relativa de la recta
\[
r:(x,y,z)=(1,2,-1)+t(1,-1,2)
\]
y el plano
\[
\pi:2x+y-z-4=0.
\]
Si se cortan, calcular el punto de intersección.

\item Estudiar la posición relativa de
\[
r:(x,y,z)=(1,0,0)+t(1,2,1)
\]
y
\[
s:(x,y,z)=(0,1,1)+\lambda(2,4,2).
\]

\item Estudiar la posición relativa de
\[
r:(x,y,z)=(1,0,1)+t(1,1,0)
\]
y
\[
s:(x,y,z)=(0,1,2)+\lambda(0,1,1).
\]
Si son secantes, calcular el punto de corte; si son cruzadas, justificarlo mediante un producto mixto.

\item Estudiar la posición relativa de los planos
\[
\pi_1:x+y-z+1=0,
\qquad
\pi_2:2x-y+z-3=0.
\]
Si se cortan, obtener una ecuación paramétrica de la recta de intersección.

\item Estudiar, según el valor del parámetro $a$, la posición relativa de los planos
\[
\pi_1:x+2y-z=1,
\qquad
\pi_2:2x+4y-2z=a.
\]

\item Calcular la distancia del punto
\[
P=(2,-3)
\]
a la recta
\[
r:3x+4y-7=0.
\]
Obtener además el pie de la perpendicular desde $P$ a $r$.

\item Calcular la distancia del punto
\[
P=(1,-2,3)
\]
a la recta
\[
r:(x,y,z)=(3,-1,0)+t(-1,1,1).
\]
Obtener el punto de la recta más próximo a $P$.

\item Calcular la distancia del punto
\[
P=(1,2,3)
\]
al plano
\[
\pi:2x+3y-z=-7.
\]
Obtener también la proyección ortogonal de $P$ sobre $\pi$.

\item Calcular la distancia entre los planos paralelos
\[
\pi_1:2x-y+2z-3=0,
\qquad
\pi_2:4x-2y+4z+5=0.
\]
Antes de aplicar una fórmula, escribir ambos planos con el mismo vector normal.

\item Calcular la distancia entre las rectas
\[
r:(x,y,z)=(0,0,0)+t(1,1,0),
\]
\[
s:(x,y,z)=(1,0,2)+\lambda(1,-1,0).
\]
Estudiar primero su posición relativa.

\item Hallar el ángulo agudo entre las rectas cuyas direcciones son
\[
u=(1,2,2),\qquad v=(2,-1,2).
\]

\item Hallar el ángulo entre los planos
\[
\pi_1:x+2y-2z+1=0,
\qquad
\pi_2:2x-y+2z-3=0.
\]

\item Hallar el ángulo que forma la recta
\[
r:(x,y,z)=(1,0,2)+t(2,1,-2)
\]
con el plano
\[
\pi:x-y+2z=0.
\]

\item Una conducción debe atravesar perpendicularmente una pared
\[
\pi:2x-y+2z-9=0.
\]
Si el punto de partida es
\[
P=(1,2,0),
\]
obtener la recta de perforación, el punto donde alcanza la pared y la distancia recorrida hasta ella.

\item Una cámara de inspección está situada en
\[
P=(5,2,3)
\]
y debe orientarse hacia el punto más próximo del eje de una tubería
\[
r:(x,y,z)=(1,0,1)+t(2,1,0).
\]
Determinar dicho punto, la distancia y un vector unitario que indique la dirección desde la cámara hasta la tubería.

\item En un sistema eléctrico tridimensional, dos tramos rectilíneos de cableado idealizados tienen ejes
\[
r:(x,y,z)=(0,0,1)+t(1,2,0),
\]
\[
s:(x,y,z)=(2,0,0)+\lambda(0,1,2).
\]
Estudiar si los ejes se cortan. Si son cruzados, calcular su separación mínima.

\item En una instalación energética, la superficie de un panel tiene normal
\[
n=(1,-2,2)
\]
y la dirección de incidencia de un haz se modeliza mediante
\[
v=(-2,1,-2).
\]
Calcular el ángulo entre el haz y la superficie. Determinar también el ángulo entre el haz y la normal.

\item Una tubería de proceso tiene eje
\[
r:(x,y,z)=(2,1,0)+t(1,-2,2),
\]
y debe atravesar el plano
\[
\pi:x+y+z-6=0.
\]
Determinar el punto de paso y el ángulo de incidencia. Interpretar qué ocurriría si el vector director se modificase a $v=(1,-2,1)$.

\item \M Comprobar con Mathematica los ejercicios 3--9 utilizando únicamente productos escalares, \texttt{Norm} y \texttt{Cross} cuando corresponda. Para el ángulo entre dos vectores puede evaluarse
\[
\operatorname{ArcCos}\left(\frac{u\cdot v}{\|u\|\|v\|}\right).
\]
Comparar siempre el resultado con el cálculo manual.

\item \M En los ejercicios 10--25, utilizar \texttt{Solve} para comprobar pertenencias e intersecciones. Plantear primero manualmente las ecuaciones y utilizar el ordenador únicamente para resolver o verificar el sistema resultante.

\item \M En los ejercicios 26--38, comprobar las distancias y proyecciones con combinaciones de \texttt{Norm}, producto escalar y \texttt{Cross}. No introducir comandos específicos de geometría: el objetivo es reutilizar las mismas operaciones vectoriales estudiadas en la teoría.

\item (Examen) Dados
\[
r:(x,y,z)=(1,-1,2)+t(1,2,-1)
\]
y
\[
\pi:2x-y+z-4=0,
\]
se pide:
\begin{enumerate}[label=\alph*)]
\item estudiar la posición relativa de $r$ y $\pi$;
\item calcular, si existe, su punto de intersección;
\item hallar el ángulo entre la recta y el plano;
\item calcular la distancia del punto $P=(3,0,1)$ al plano y su proyección ortogonal sobre él.
\end{enumerate}

\item (Examen) Sean las rectas
\[
r:(x,y,z)=(1,0,0)+t(1,1,2),
\]
\[
s:(x,y,z)=(0,2,1)+\lambda(2,-1,1).
\]
\begin{enumerate}[label=\alph*)]
\item estudiar su posición relativa;
\item calcular la distancia entre ambas;
\item obtener un vector perpendicular a las dos direcciones;
\item si son cruzadas, explicar geométricamente por qué el vector que realiza la distancia mínima debe ser paralelo a ese vector perpendicular común.
\end{enumerate}

\item (Examen) Una superficie plana queda determinada por
\[
A=(1,0,1),\qquad B=(3,1,0),\qquad C=(0,2,2).
\]
Una trayectoria rectilínea parte de
\[
P=(2,-1,4)
\]
con dirección $v=(1,1,-2)$. Se pide:
\begin{enumerate}[label=\alph*)]
\item obtener la ecuación cartesiana del plano $ABC$;
\item estudiar la posición relativa de la trayectoria y el plano;
\item calcular, si existe, el punto de intersección;
\item calcular el ángulo entre la trayectoria y el plano;
\item obtener la distancia de $P$ al plano y el punto de la superficie más próximo a $P$.
\end{enumerate}

\end{enumerate}

\cleardoublepage

% Glosario computacional final
\chapter*{Glosario de Wolfram Mathematica}
\addcontentsline{toc}{chapter}{Glosario de Wolfram Mathematica}
\chaptermark{Glosario de Wolfram Mathematica}


%\newcommand{\MatCmd}[1]{\NoCaseChange{\MatCmd{#1}}}
\DeclareRobustCommand{\MatCmd}[1]{\texttt{#1}}
\AddToNoCaseChangeList{\MatCmd}

Las prácticas de la asignatura se plantean como una extensión del trabajo realizado en las clases de teoría y problemas. El objetivo no es aprender un lenguaje de programación ni sustituir los razonamientos matemáticos por una salida de ordenador. Se utilizará \textit{Wolfram Mathematica} para \emph{comprobar cálculos}, \emph{explorar ejemplos}, \emph{representar gráficamente resultados} y abordar, una vez comprendido el procedimiento, casos cuyo cálculo manual resulte innecesariamente largo.

Por este motivo, el número de herramientas que se empleará durante el curso es deliberadamente reducido. Este glosario contiene las órdenes y la sintaxis necesarias para realizar los ejercicios marcados con \M a lo largo de los apuntes. No pretende ser un manual de Mathematica.

\section*{Sintaxis básica}

Mathematica distingue entre mayúsculas y minúsculas. Los nombres de las funciones incorporadas comienzan por mayúscula y sus argumentos se escriben entre corchetes. Así, por ejemplo,
\begin{mathematicacode}
Sin[x]
Sqrt[x]
Exp[x]
\end{mathematicacode}
representan, respectivamente, $\sin x$, $\sqrt{x}$ y $e^x$.

Las listas se escriben entre llaves. Esta misma notación permite introducir vectores y matrices:
\begin{mathematicacode}
u = {1, -2, 3}
A = {{1, 2}, {3, 4}}
\end{mathematicacode}
Por tanto, una matriz se introduce como una lista cuyas componentes son sus filas.

La multiplicación puede escribirse mediante un espacio o mediante \MatCmd{*}; las potencias se indican con \MatCmd{\string^}. Por ejemplo,
\begin{mathematicacode}
3 x^2
3*x^2
\end{mathematicacode}
son expresiones equivalentes.

Para plantear ecuaciones se utiliza \MatCmd{==}, no \MatCmd{=}. El signo \MatCmd{=} se reserva para asignar un valor a un símbolo:
\begin{mathematicacode}
a = 3
Solve[x^2 == a, x]
\end{mathematicacode}

Cuando resulte cómodo definir una función utilizaremos el patrón \MatCmd{x\_} y la asignación diferida \MatCmd{:=}:
\begin{mathematicacode}
f[x_] := x^3 - 3 x
f[2]
\end{mathematicacode}
No será necesario utilizar mecanismos de programación más avanzados en esta asignatura.

Algunas constantes y símbolos incorporados que aparecerán con frecuencia son
\begin{center}
\begin{tabular}{ll}
\toprule
\textbf{Mathematica} & \textbf{Significado}\\
\midrule
\MatCmd{Pi} & $\pi$\\
\MatCmd{I} & unidad imaginaria $i$\\
\MatCmd{Infinity} & $\infty$\\
\MatCmd{Sqrt[x]} & $\sqrt{x}$\\
\MatCmd{Exp[x]} & $e^x$\\
\MatCmd{Log[x]} & $\ln x$\\
\MatCmd{Sin[x]}, \MatCmd{Cos[x]}, \MatCmd{Tan[x]} & funciones trigonométricas\\
\MatCmd{ArcCos[x]} & $\arccos x$\\
\bottomrule
\end{tabular}
\end{center}
Los argumentos de las funciones trigonométricas se interpretan en \textbf{radianes}.

\section*{Herramientas de uso general}

\begin{description}

\item[\MatCmd{N}] Obtiene una aproximación decimal de una expresión exacta.
\begin{mathematicacode}
N[Pi/7]
N[Sqrt[2], 12]
\end{mathematicacode}
En el segundo ejemplo se solicita una aproximación con mayor precisión.

\item[\MatCmd{Simplify}] Simplifica una expresión algebraica. Es especialmente útil para comprobar que dos resultados escritos de forma diferente son equivalentes.
\begin{mathematicacode}
Simplify[(x^2 - 1)/(x - 1)]
Simplify[(x + 1)^2 - (x^2 + 2 x + 1)]
\end{mathematicacode}
La simplificación automática no sustituye el estudio del dominio: dos expresiones pueden coincidir donde ambas están definidas y, sin embargo, tener dominios diferentes.

\item[\MatCmd{Solve}] Resuelve simbólicamente ecuaciones o sistemas de ecuaciones.
\begin{mathematicacode}
Solve[x^2 - 5 x + 6 == 0, x]
Solve[{x + y == 3, 2 x - y == 0}, {x, y}]
\end{mathematicacode}
Las ecuaciones se escriben con \MatCmd{==} y un sistema se introduce como una lista de ecuaciones.

\item[\MatCmd{Plot}] Representa una o varias funciones reales en un intervalo.
\begin{mathematicacode}
Plot[x Exp[-x], {x, 0, 8}]
Plot[{Sin[x], Cos[x]}, {x, -2 Pi, 2 Pi}]
\end{mathematicacode}
Una gráfica debe utilizarse para explorar o comprobar el estudio analítico, no para sustituirlo. En particular, la apariencia de una función puede cambiar sustancialmente al modificar el intervalo de representación.

\end{description}

\section*{Números complejos}

\begin{description}

\item[\MatCmd{Re}, \MatCmd{Im}] Devuelven las partes real e imaginaria de un número complejo.
\begin{mathematicacode}
z = 3 - 4 I
Re[z]
Im[z]
\end{mathematicacode}

\item[\MatCmd{Abs}] Para un número complejo devuelve su módulo; para un número real coincide con el valor absoluto.
\begin{mathematicacode}
Abs[3 - 4 I]
\end{mathematicacode}

\item[\MatCmd{Arg}] Devuelve el argumento principal de un número complejo no nulo.
\begin{mathematicacode}
Arg[-1 + Sqrt[3] I]
N[Arg[-1 + Sqrt[3] I]]
\end{mathematicacode}

\end{description}

Para ecuaciones en números complejos se utilizará el mismo comando \MatCmd{Solve} estudiado anteriormente:
\begin{mathematicacode}
Solve[z^3 == -8 I, z]
\end{mathematicacode}

\section*{Funciones y cálculo}

\begin{description}

\item[\MatCmd{Limit}] Calcula simbólicamente un límite.
\begin{mathematicacode}
Limit[(x^2 - 1)/(x - 1), x -> 1]
Limit[(5 x^2 + 3 x - 1)/(-2 x^2 + 7), x -> Infinity]
\end{mathematicacode}
Para los límites laterales puede añadirse la opción \MatCmd{Direction} cuando sea necesaria. En las prácticas ordinarias se priorizarán los límites para los que baste la sintaxis anterior.

\item[\MatCmd{D}] Calcula derivadas simbólicas. El segundo argumento indica la variable respecto de la que se deriva.
\begin{mathematicacode}
D[x^2 Exp[-x], x]
D[Sin[x^2], x]
\end{mathematicacode}
Las derivadas de orden superior pueden obtenerse indicando el orden:
\begin{mathematicacode}
D[Exp[x] Sin[x], {x, 3}]
\end{mathematicacode}
Si se ha definido previamente una función, Mathematica también entiende la notación con primas:
\begin{mathematicacode}
f[x_] := x^3 - 3 x
f'[2]
f''[2]
\end{mathematicacode}

\item[\MatCmd{Series}] Obtiene un desarrollo de Taylor alrededor de un punto hasta el orden indicado.
\begin{mathematicacode}
Series[Sin[x], {x, 0, 5}]
Series[Exp[x], {x, 0, 4}]
\end{mathematicacode}
La salida incluye un término que representa el orden del resto.

\item[\MatCmd{Normal}] Elimina de una serie el término que representa el resto y deja únicamente el polinomio de Taylor. Se utilizará solo cuando sea necesario manipular o representar dicho polinomio.
\begin{mathematicacode}
p = Normal[Series[Exp[x], {x, 0, 4}]]
Plot[{Exp[x], p}, {x, -2, 2}]
\end{mathematicacode}

\item[\MatCmd{Integrate}] Calcula integrales simbólicas, tanto indefinidas como definidas.
\begin{mathematicacode}
Integrate[x Exp[-x], x]
Integrate[x Exp[-x], {x, 0, 4}]
\end{mathematicacode}
En una integral indefinida Mathematica proporciona una primitiva y, por tanto, no añade una constante de integración. Al escribir la solución matemática debe recordarse la constante $C$.

\item[\MatCmd{NIntegrate}] Aproxima numéricamente una integral definida. Se utilizará cuando interese un valor numérico o cuando la primitiva no pueda expresarse mediante funciones elementales sencillas.
\begin{mathematicacode}
NIntegrate[Exp[-x^2], {x, 0, 1}]
\end{mathematicacode}

\end{description}

\section*{Matrices, espacios vectoriales y aplicaciones lineales}

Los vectores y matrices se introducen mediante listas. El operador punto \MatCmd{.} representa el producto matricial y, para dos vectores, el producto escalar:
\begin{mathematicacode}
A = {{1, 2}, {3, 4}}
x = {2, -1}
A.x
x.x
\end{mathematicacode}

\begin{description}

\item[\MatCmd{Det}] Calcula el determinante de una matriz cuadrada.
\begin{mathematicacode}
Det[{{1, 2}, {3, 4}}]
\end{mathematicacode}

\item[\MatCmd{Inverse}] Calcula la inversa de una matriz regular.
\begin{mathematicacode}
Inverse[{{1, 2}, {3, 4}}]
\end{mathematicacode}

\item[\MatCmd{MatrixRank}] Calcula el rango de una matriz.
\begin{mathematicacode}
MatrixRank[{{1, 2, 3}, {2, 4, 6}, {0, 1, 1}}]
\end{mathematicacode}

\item[\MatCmd{RowReduce}] Obtiene la forma escalonada reducida por filas. Es especialmente útil para comprobar una eliminación de Gauss y para estudiar sistemas mediante su matriz ampliada.
\begin{mathematicacode}
RowReduce[{{1, 1, 3}, {2, -1, 0}}]
\end{mathematicacode}
En este ejemplo la última columna puede interpretarse como la columna de términos independientes de un sistema de dos ecuaciones.

\item[\MatCmd{LinearSolve}] Resuelve directamente un sistema matricial $Ax=b$ cuando se proporcionan la matriz de coeficientes y el vector de términos independientes.
\begin{mathematicacode}
A = {{1, 1}, {2, -1}}
b = {3, 0}
LinearSolve[A, b]
\end{mathematicacode}
Para discutir sistemas con parámetros o para obtener una descripción simbólica más flexible se utilizará normalmente \MatCmd{Solve}.

\item[\MatCmd{NullSpace}] Devuelve una base del núcleo de una matriz. Si $A$ representa una aplicación lineal, permite comprobar una base de $\ker A$.
\begin{mathematicacode}
A = {{1, 2, 3}, {2, 4, 6}}
NullSpace[A]
\end{mathematicacode}

\item[\MatCmd{Orthogonalize}] Aplica un proceso de ortonormalización a una familia de vectores linealmente independientes. Se utilizará para comprobar el procedimiento manual de Gram--Schmidt.
\begin{mathematicacode}
Orthogonalize[{{1, 1, 0}, {1, 0, 1}}]
\end{mathematicacode}

\item[\MatCmd{Transpose}] Calcula la matriz traspuesta.
\begin{mathematicacode}
Transpose[{{1, 2, 3}, {4, 5, 6}}]
\end{mathematicacode}
Será útil, entre otros lugares, para comprobar cambios por congruencia del tipo $P^t A P$.

\end{description}

\section*{Diagonalización}

\begin{description}

\item[\MatCmd{Eigenvalues}] Calcula los valores propios de una matriz cuadrada, contando multiplicidades algebraicas.
\begin{mathematicacode}
A = {{2, 1}, {1, 2}}
Eigenvalues[A]
\end{mathematicacode}

\item[\MatCmd{Eigenvectors}] Calcula vectores propios asociados a los valores propios de una matriz.
\begin{mathematicacode}
Eigenvectors[A]
\end{mathematicacode}
Estos comandos se utilizarán para \emph{comprobar} los valores y vectores propios calculados a mano. La decisión de si una matriz es diagonalizable debe justificarse mediante las multiplicidades y las dimensiones de los subespacios propios.

\end{description}

No se introducirá un comando específico que diagonalice automáticamente la matriz. Una vez construidas manualmente las matrices $P$ y $D$, la comprobación se realiza con las herramientas ya conocidas:
\begin{mathematicacode}
Inverse[P].A.P
P.D.Inverse[P]
\end{mathematicacode}
Para formas cuadráticas y diagonalización por congruencia se procede de forma análoga:
\begin{mathematicacode}
Transpose[P].A.P
\end{mathematicacode}

\section*{Geometría vectorial}

No utilizaremos comandos específicos de geometría analítica. Los problemas se resolverán con las mismas operaciones vectoriales estudiadas en el curso.

\begin{description}

\item[\MatCmd{.}] El punto entre dos vectores representa su producto escalar. Entre matrices, o entre una matriz y un vector, representa el producto matricial.
\begin{mathematicacode}
u = {1, 2, -1}
v = {2, 0, 3}
u.v
\end{mathematicacode}

\item[\MatCmd{Norm}] Calcula la norma euclídea de un vector.
\begin{mathematicacode}
Norm[{1, 2, -1}]
\end{mathematicacode}
De esta forma, el ángulo entre dos vectores no nulos puede obtenerse reutilizando el producto escalar:
\begin{mathematicacode}
ArcCos[(u.v)/(Norm[u] Norm[v])]
\end{mathematicacode}

\item[\MatCmd{Cross}] Calcula el producto vectorial de dos vectores de $\R^3$.
\begin{mathematicacode}
Cross[{1, 0, 2}, {0, 1, -1}]
\end{mathematicacode}

\end{description}

Las pertenencias, intersecciones y posiciones relativas que conduzcan a sistemas de ecuaciones se comprobarán con \MatCmd{Solve}, sin introducir órdenes adicionales de geometría.

\section*{Resumen de herramientas}

La siguiente tabla recoge el conjunto de órdenes que se utilizará de forma habitual. Una misma herramienta aparece en varios temas porque precisamente se pretende reutilizar un vocabulario computacional pequeño.

\begin{center}
\small
\begin{tabularx}{\textwidth}{>{\ttfamily}l X}
\toprule
\textbf{Herramienta} & \textbf{Uso principal}\\
\midrule
N & aproximaciones numéricas\\
Simplify & simplificación y comprobación de expresiones\\
Solve & ecuaciones y sistemas simbólicos\\
Plot & representación de funciones\\
Re, Im, Abs, Arg & números complejos\\
Limit & límites\\
D & derivadas\\
Series, Normal & polinomios de Taylor\\
Integrate & integrales simbólicas\\
NIntegrate & integrales numéricas\\
Det, Inverse & determinante e inversa\\
MatrixRank & rango\\
RowReduce & eliminación de Gauss\\
LinearSolve & sistemas lineales en forma matricial\\
NullSpace & núcleo\\
Orthogonalize & Gram--Schmidt\\
Transpose & traspuesta y cambios por congruencia\\
Eigenvalues, Eigenvectors & valores y vectores propios\\
. & producto escalar y producto matricial\\
Norm, Cross & norma y producto vectorial\\
\bottomrule
\end{tabularx}
\end{center}

Esta lista es intencionadamente corta. Siempre que un ejercicio pueda resolverse combinando estas herramientas se evitará introducir un nuevo comando que automatice directamente el procedimiento matemático que se pretende aprender.

\cleardoublepage

% Bibliografía: se mostrará cuando existan referencias citadas.
\include{FrontBackmatter/Bibliography}
\cleardoublepage

\end{document}
